% =====================================================================
%  Chapitre 5 : États quantiques, intrication et états de Bell
%  Fichier importé par main.tex avec % =====================================================================
%  Chapitre 5 : États quantiques, intrication et états de Bell
%  Fichier importé par main.tex avec % =====================================================================
%  Chapitre 5 : États quantiques, intrication et états de Bell
%  Fichier importé par main.tex avec % =====================================================================
%  Chapitre 5 : États quantiques, intrication et états de Bell
%  Fichier importé par main.tex avec \include{chapitre5}.
%  Il ne contient ni préambule ni \begin{document} : les environnements
%  (definition, resultat, remarque, aretenir, savoirfaire, exercice, correction),
%  les macros (\ii, \ee, \dd, \pd, \pdd, \Hop, \ket, \bra, \braket), les symboles
%  de circuits (ctl, cible, croix, mesure, porte) et les couleurs (insarouge,
%  insagris) sont définis dans main.tex. Un chapitre = une \section.
% =====================================================================

% Macros de Dirac (déjà définies dans main.tex ; ce bloc évite l'erreur
% « Undefined control sequence \ket » si main.tex est une ancienne version).
\providecommand{\ket}[1]{|#1\rangle}
\providecommand{\bra}[1]{\langle #1|}
\providecommand{\braket}[2]{\langle #1|#2\rangle}

% =====================================================================
\section{États quantiques, intrication et états de Bell}\label{chap:bell}
% =====================================================================

Les états de deux qubits forment un espace de dimension~$4$, et la plupart d'entre eux ne se
décomposent pas en «~un état pour $A$, un état pour $B$~» : ils sont \emph{intriqués}. Ce chapitre
apprend à mesurer un couple de qubits, à reconnaître l'intrication (critère de séparabilité), à la
produire avec un circuit (états de Bell), à décrire un qubit pris seul (trace partielle), à
quantifier l'information (entropies) et à s'en servir pour téléporter un état.

% ---------------------------------------------------------------------
\subsection{États de deux qubits et mesure}\label{sec:mesure2}
% ---------------------------------------------------------------------

Un état quelconque de $\mathcal{H}^4$ se décompose sur la base de calcul :
\[
  \ket{\Psi}=c_{00}\ket{00}+c_{01}\ket{01}+c_{10}\ket{10}+c_{11}\ket{11}
  =\begin{pmatrix}c_{00}\\ c_{01}\\ c_{10}\\ c_{11}\end{pmatrix},
  \qquad \sum_{i,j}|c_{ij}|^2=1 .
\]

\begin{resultat}{Mesure des deux qubits dans la base de calcul}
\[
  P(ij)=|\braket{ij}{\Psi}|^2=|c_{ij}|^2 .
\]
Pour un seul qubit : $P(A{=}0)=P(00)+P(01)$, $P(A{=}1)=P(10)+P(11)$,
$P(B{=}0)=P(00)+P(10)$, $P(B{=}1)=P(01)+P(11)$.
\end{resultat}

En effet $\braket{ij}{\Psi}=\sum_{k,l}c_{kl}\braket{ij}{kl}=\sum_{k,l}c_{kl}\delta_{ik}\delta_{jl}=c_{ij}$.

\paragraph{Exemple 1.} $\ket{\Psi}=\frac{1}{\sqrt2}\ket{00}+\frac{1}{\sqrt2}\ket{01}
=\big(\tfrac{1}{\sqrt2},\tfrac{1}{\sqrt2},0,0\big)^{\mathsf T}$. Il est normalisé :
$\frac12+\frac12+0+0=1$. Les probabilités jointes sont
$P(00)=P(01)=\frac12$ et $P(10)=P(11)=0$. Pour chaque qubit :
\[
  P(A{=}0)=\tfrac12+\tfrac12=1,\quad P(A{=}1)=0+0=0,\quad
  P(B{=}0)=\tfrac12+0=\tfrac12,\quad P(B{=}1)=\tfrac12+0=\tfrac12 .
\]
$A$ donne toujours $0$ ; $B$ donne $0$ ou $1$ au hasard.

\paragraph{Exemple 2 (mesure de $A$ seul).} Soit
$\ket{\Psi}=\frac12\ket{00}+\frac12\ket{01}+\frac{1}{\sqrt2}\ket{11}$. Normalisation :
$\frac14+\frac14+0+\frac12=1$. Probabilités jointes : $P(00)=\frac14$, $P(01)=\frac14$,
$P(10)=0$, $P(11)=\frac12$. Donc
\[
  P(A{=}0)=\tfrac14+\tfrac14=\tfrac12,\quad P(A{=}1)=0+\tfrac12=\tfrac12,\quad
  P(B{=}0)=\tfrac14+0=\tfrac14,\quad P(B{=}1)=\tfrac14+\tfrac12=\tfrac34 .
\]

\begin{remarque}[Effondrement après la mesure de $A$]
Si $A$ donne le résultat $a$, on ne garde que les termes de $\ket{\Psi}$ dont le premier
chiffre vaut $a$, puis on renormalise :
\[
  \ket{\Psi}\ \to\ \frac{c_{a0}\ket{a0}+c_{a1}\ket{a1}}{\sqrt{|c_{a0}|^2+|c_{a1}|^2}},
  \qquad\text{probabilité }|c_{a0}|^2+|c_{a1}|^2 .
\]
\end{remarque}

Pour l'exemple 2 :
\begin{itemize}[nosep]
  \item $A=0$ (probabilité $\frac12$) : on garde $\frac12\ket{00}+\frac12\ket{01}$, de norme
        au carré $\frac12$ ; renormalisé, $\frac{1}{\sqrt2}\big(\ket{00}+\ket{01}\big)
        =\ket{0}\otimes\ket{+}$. Ensuite $B$ vaut $0$ ou $1$ avec la probabilité $\frac12$.
  \item $A=1$ (probabilité $\frac12$) : on garde $\frac{1}{\sqrt2}\ket{11}$, de norme au carré
        $\frac12$ ; renormalisé, $\ket{11}=\ket{1}\otimes\ket{1}$. Ensuite $B$ vaut $1$
        avec certitude.
\end{itemize}
Contrôle : $P(B{=}0)=\frac12\cdot\frac12+\frac12\cdot0=\frac14$, comme calculé plus haut.
Ici la mesure de $A$ \emph{modifie} l'état de $B$ ($\ket{+}$ ou $\ket{1}$) : c'est la
signature d'un état intriqué (section~\ref{sec:separable}). La figure~\ref{fig:mesure2} résume
les deux mesures successives.

\begin{figure}[!ht]
\centering
\adjustbox{max width=\linewidth}{%
\begin{tikzpicture}[>=Stealth, font=\small,
    noeud/.style={draw=insarouge, fill=insarouge!6, rounded corners=2pt, thick, inner sep=4pt, align=center},
    fin/.style={draw=insagris, fill=insagris!10, rounded corners=2pt, inner sep=4pt, align=center},
    lien/.style={->, thick, insagris},
    proba/.style={font=\footnotesize, fill=white, inner sep=1.5pt, text=black}]
  \node[noeud] (psi) at (0,0) {$\ket{\Psi}$};
  \node[noeud] (a0) at (4.4,1.6) {$A=0$\\ {\footnotesize état $\ket{0}\otimes\ket{+}$}};
  \node[noeud] (a1) at (4.4,-1.6) {$A=1$\\ {\footnotesize état $\ket{1}\otimes\ket{1}$}};
  \node[fin] (b00) at (9.2,2.5) {$B=0$ : $\ket{00}$};
  \node[fin] (b01) at (9.2,0.7) {$B=1$ : $\ket{01}$};
  \node[fin] (b11) at (9.2,-1.6) {$B=1$ : $\ket{11}$};
  \draw[lien] (psi) -- (a0) node[proba, midway] {$\frac12$};
  \draw[lien] (psi) -- (a1) node[proba, midway] {$\frac12$};
  \draw[lien] (a0) -- (b00) node[proba, midway] {$\frac12$};
  \draw[lien] (a0) -- (b01) node[proba, midway] {$\frac12$};
  \draw[lien] (a1) -- (b11) node[proba, midway] {$1$};
  \node[anchor=west] at (11.0,2.5) {$P(00)=\frac12\cdot\frac12=\frac14$};
  \node[anchor=west] at (11.0,0.7) {$P(01)=\frac12\cdot\frac12=\frac14$};
  \node[anchor=west] at (11.0,-1.6) {$P(11)=\frac12\cdot1=\frac12$};
  \node[text=insagris, font=\footnotesize\itshape] at (2.2,3.0) {mesure de $A$};
  \node[text=insagris, font=\footnotesize\itshape] at (6.8,3.6) {mesure de $B$};
  \node[text=insagris, font=\footnotesize\itshape, anchor=west] at (11.0,3.6) {probabilité jointe};
\end{tikzpicture}}
\caption{Mesure de $A$ puis de $B$ pour l'état de l'exemple~2,
$\ket{\Psi}=\frac12\ket{00}+\frac12\ket{01}+\frac{1}{\sqrt2}\ket{11}$. Chaque flèche porte une
probabilité, et le produit des probabilités le long d'un chemin est la probabilité jointe. Après la
mesure de $A$, l'état de $B$ dépend du résultat ($\ket{+}$ ou $\ket{1}$).}
\label{fig:mesure2}
\end{figure}

% ---------------------------------------------------------------------
\subsection{États séparables et états intriqués}\label{sec:separable}
% ---------------------------------------------------------------------

\begin{definition}{Séparable ou intriqué}
Un état $\ket{\Psi}_{AB}$ est \emph{séparable} s'il s'écrit comme un produit
$\ket{\Psi}_{AB}=\ket{\Psi}_A\otimes\ket{\Psi}_B$. Sinon il est \emph{intriqué} : on ne
peut pas attribuer d'état propre à chaque qubit.
\end{definition}

\paragraph{Exemple séparable (exemple 1 ci-dessus).} Par linéarité du produit tensoriel :
\begin{align*}
  \ket{\Psi}&=\tfrac{1}{\sqrt2}\ket{00}_{AB}+\tfrac{1}{\sqrt2}\ket{01}_{AB}
    =\tfrac{1}{\sqrt2}\Big(\ket{0}_A\otimes\ket{0}_B+\ket{0}_A\otimes\ket{1}_B\Big)\\
  &=\ket{0}_A\otimes\Big(\tfrac{1}{\sqrt2}\big(\ket{0}+\ket{1}\big)\Big)_B
    =\ket{0}_A\otimes\ket{+}_B ,
\end{align*}
donc $\ket{\Psi}_A=\ket{0}$ et $\ket{\Psi}_B=\ket{+}$. Contrôle par les composantes :
$\begin{pmatrix}1\\ 0\end{pmatrix}\otimes\tfrac{1}{\sqrt2}\begin{pmatrix}1\\ 1\end{pmatrix}
=\tfrac{1}{\sqrt2}\begin{pmatrix}1\cdot1\\ 1\cdot1\\ 0\cdot1\\ 0\cdot1\end{pmatrix}
=\tfrac{1}{\sqrt2}\begin{pmatrix}1\\ 1\\ 0\\ 0\end{pmatrix}$.

\subsubsection{Un critère de séparabilité}

Un produit quelconque a pour coefficients $c_{00}=ac$, $c_{01}=ad$, $c_{10}=bc$,
$c_{11}=bd$ (formule de la section~\ref{sec:tenseur}, avec $a,b$ pour $A$ et $c,d$ pour $B$), donc
\[
  c_{00}c_{11}=(ac)(bd)=abcd=(ad)(bc)=c_{01}c_{10}.
\]

\begin{resultat}{Test de séparabilité (deux qubits)}
$\ket{\Psi}=\sum c_{ij}\ket{ij}$ est séparable si et seulement si
\[
  D=c_{00}c_{11}-c_{01}c_{10}=0 .
\]
\end{resultat}

\emph{Démonstration.} ($\Rightarrow$) Calcul ci-dessus. ($\Leftarrow$) On regroupe les
termes selon l'état de $A$ :
\[
  \ket{\Psi}=\ket{0}\otimes\big(c_{00}\ket{0}+c_{01}\ket{1}\big)+\ket{1}\otimes\big(c_{10}\ket{0}+c_{11}\ket{1}\big).
\]
Si $D=0$, les lignes $(c_{00},c_{01})$ et $(c_{10},c_{11})$ sont proportionnelles. Si
$(c_{00},c_{01})\neq(0,0)$, il existe $\lambda$ avec $(c_{10},c_{11})=\lambda(c_{00},c_{01})$
et $\ket{\Psi}=\big(\ket{0}+\lambda\ket{1}\big)\otimes\big(c_{00}\ket{0}+c_{01}\ket{1}\big)$.
Si $(c_{00},c_{01})=(0,0)$, alors $\ket{\Psi}=\ket{1}\otimes\big(c_{10}\ket{0}+c_{11}\ket{1}\big)$.
Dans les deux cas $\ket{\Psi}$ est un produit ; on renormalise chaque facteur.
\hfill$\square$

\subsubsection{Application à plusieurs états}

\begin{center}
\renewcommand{\arraystretch}{1.5}
\begin{tabular}{@{}lccl@{}}
\toprule
\textbf{État} & $(c_{00},c_{01},c_{10},c_{11})$ & $D$ & \textbf{Nature} \\
\midrule
$\ket{00}$ & $(1,0,0,0)$ & $1\cdot0-0\cdot0=0$ & séparable \\
$\ket{+}\otimes\ket{+}$ & $\frac12(1,1,1,1)$ & $\frac14-\frac14=0$ & séparable \\
$\ket{+}\otimes\ket{-}$ & $\frac12(1,-1,1,-1)$ & $-\frac14-(-\frac14)=0$ & séparable \\
$\frac{1}{\sqrt2}(\ket{00}+\ket{01})$ & $\frac{1}{\sqrt2}(1,1,0,0)$ & $0-\frac12\cdot0=0$ & séparable \\
$\ket{\Phi^+}=\frac{1}{\sqrt2}(\ket{00}+\ket{11})$ & $\frac{1}{\sqrt2}(1,0,0,1)$ & $\frac12-0=\frac12$ & intriqué \\
$\ket{\Phi^-}=\frac{1}{\sqrt2}(\ket{00}-\ket{11})$ & $\frac{1}{\sqrt2}(1,0,0,-1)$ & $-\frac12-0=-\frac12$ & intriqué \\
$\ket{\Psi^+}=\frac{1}{\sqrt2}(\ket{01}+\ket{10})$ & $\frac{1}{\sqrt2}(0,1,1,0)$ & $0-\frac12=-\frac12$ & intriqué \\
$\ket{\Psi^-}=\frac{1}{\sqrt2}(\ket{01}-\ket{10})$ & $\frac{1}{\sqrt2}(0,1,-1,0)$ & $0-(-\frac12)=\frac12$ & intriqué \\
$\frac12\ket{00}+\frac12\ket{01}+\frac{1}{\sqrt2}\ket{11}$ & $(\frac12,\frac12,0,\frac{1}{\sqrt2})$ & $\frac{1}{2\sqrt2}-0$ & intriqué \\
\bottomrule
\end{tabular}
\end{center}

Les quatre états de Bell $\ket{\Phi^\pm}$ et $\ket{\Psi^\pm}$ de la table sont tous intriqués
($D=\pm\frac12$) ; leurs propriétés font l'objet de la section~\ref{sec:bell}.

\paragraph{Preuve directe que $\ket{\Phi^+}$ est intriqué.} Supposons
$\ket{\Phi^+}=(a\ket{0}+b\ket{1})\otimes(c\ket{0}+d\ket{1})$. En identifiant les
composantes : $ac=\frac{1}{\sqrt2}$, $ad=0$, $bc=0$, $bd=\frac{1}{\sqrt2}$. La condition
$ad=0$ impose $a=0$ (impossible car $ac\neq0$) ou $d=0$ (impossible car $bd\neq0$).
Contradiction : aucun produit ne convient.

\paragraph{Corrélations.} Pour un produit $\ket{\Psi}_A\otimes\ket{\Psi}_B$ avec
$\ket{\Psi}_A=a_0\ket{0}+a_1\ket{1}$ et $\ket{\Psi}_B=b_0\ket{0}+b_1\ket{1}$, on a
$P(ij)=|a_ib_j|^2=|a_i|^2|b_j|^2=P(A{=}i)\,P(B{=}j)$ : les deux qubits sont
\emph{indépendants}. Vérification avec $\ket{+}\otimes\ket{+}$ : $P(00)=\frac14=\frac12\cdot\frac12$.
Pour $\ket{\Phi^+}$ : $P(00)=P(11)=\frac12$ et $P(01)=P(10)=0$, donc
$P(A{=}0)=P(B{=}0)=\frac12$, mais
\[
  P(00)=\tfrac12\neq\tfrac12\cdot\tfrac12=P(A{=}0)\,P(B{=}0) .
\]
Les deux qubits donnent toujours le même résultat, aléatoire : ils sont corrélés.
La figure~\ref{fig:epr} illustre la différence entre les deux cas.

\begin{figure}[!ht]
\centering
\adjustbox{max width=\linewidth}{%
\begin{tikzpicture}[>=Stealth, font=\small,
    labo/.style={draw=insagris, fill=insagris!8, rounded corners=2pt, thick, minimum width=2.6cm,
                 minimum height=1.0cm, align=center},
    source/.style={draw=insarouge, fill=insarouge!10, rounded corners=2pt, thick, minimum width=3.6cm,
                   minimum height=1.0cm, align=center},
    droit/.style={->, thick, insagris},
    onde/.style={->, thick, insarouge, decorate,
                 decoration={snake, amplitude=1.3pt, segment length=7pt, post length=2mm}},
    res/.style={font=\small, align=center}]
  % (a) état séparable
  \node[labo] (Aa) at (0,0) {Alice\\ mesure $A$};
  \node[source] (Sa) at (6,0) {source\\ $\ket{0}_A\otimes\ket{+}_B$};
  \node[labo] (Ba) at (12,0) {Bob\\ mesure $B$};
  \draw[droit] (Sa.west) -- (Aa.east);
  \draw[droit] (Sa.east) -- (Ba.west);
  \node[res, below=3pt of Aa] {$0\ \ 0\ \ 0\ \ 0\ \ 0\ \ 0$};
  \node[res, below=3pt of Ba] {$1\ \ 0\ \ 0\ \ 1\ \ 0\ \ 1$};
  \node[res, text=insagris, font=\footnotesize] at (6,-1.0) {$A$ donne toujours $0$ ; $B$ est aléatoire.\\
        Les deux résultats sont indépendants.};
  \node[anchor=south west, font=\small\bfseries] at (-1.3,0.75) {(a) état séparable};
  % (b) état intriqué
  \begin{scope}[yshift=-3.6cm]
    \node[labo] (Ab) at (0,0) {Alice\\ mesure $A$};
    \node[source] (Sb) at (6,0) {source\\ $\ket{\Phi^+}_{AB}$};
    \node[labo] (Bb) at (12,0) {Bob\\ mesure $B$};
    \draw[onde] (Sb.west) -- (Ab.east);
    \draw[onde] (Sb.east) -- (Bb.west);
    \node[res, below=3pt of Ab] {$1\ \ 0\ \ 0\ \ 1\ \ 1\ \ 0$};
    \node[res, below=3pt of Bb] {$1\ \ 0\ \ 0\ \ 1\ \ 1\ \ 0$};
    \node[res, text=insagris, font=\footnotesize] at (6,-1.0) {$A$ et $B$ sont chacun aléatoires,\\
          mais toujours égaux : ils sont corrélés.};
    \node[anchor=south west, font=\small\bfseries] at (-1.3,0.75) {(b) état intriqué};
  \end{scope}
\end{tikzpicture}}
\caption{Six mesures successives sur des paires préparées dans le même état (les résultats sont
des exemples). (a) Pour un produit, chaque qubit a son propre comportement. (b) Pour
$\ket{\Phi^+}=\frac{1}{\sqrt2}(\ket{00}+\ket{11})$, chaque résultat pris seul est un tirage à pile ou
face, mais les deux résultats coïncident toujours ; la ligne ondulée symbolise l'intrication.}
\label{fig:epr}
\end{figure}

\paragraph{Et un qubit pris seul ?} Dans un état séparable, chaque qubit garde son propre état pur ;
dans un état intriqué, aucun ket ne décrit un qubit isolé. On le décrit alors par une matrice de
densité, obtenue par la trace partielle (section~\ref{sec:trace}).

% ---------------------------------------------------------------------
\subsection{Les états de Bell}\label{sec:bell}
% ---------------------------------------------------------------------

\begin{definition}{Les quatre états de Bell}
\[
  \ket{\Phi^{\pm}}=\frac{1}{\sqrt2}\big(\ket{00}\pm\ket{11}\big),
  \qquad
  \ket{\Psi^{\pm}}=\frac{1}{\sqrt2}\big(\ket{01}\pm\ket{10}\big).
\]
Ils sont intriqués ($|D|=\frac12$, section~\ref{sec:separable}) et le sont au maximum :
$|D|\leq\frac12$ pour tout état normalisé de deux qubits, et chaque qubit d'un état de Bell,
pris seul, est complètement aléatoire (section~\ref{sec:trace}).
\end{definition}

\paragraph{Une base de $\mathcal{H}^4$.} Chaque état est de norme $\frac12+\frac12=1$. Les
$\ket{\Phi^\pm}$ n'ont de composantes que sur $\ket{00}$ et $\ket{11}$, les $\ket{\Psi^\pm}$ que sur
$\ket{01}$ et $\ket{10}$, donc $\braket{\Phi^\pm}{\Psi^\pm}=0$ pour toutes les combinaisons de signes ;
enfin $\braket{\Phi^+}{\Phi^-}=\frac12(1-1)=0$ et $\braket{\Psi^+}{\Psi^-}=\frac12(1-1)=0$. Les quatre
états forment donc une base orthonormée de $\mathcal{H}^4$, la \emph{base de Bell}, constituée
uniquement d'états intriqués. Le circuit de la section~\ref{sec:circuit-bell} fait passer de la
base de calcul à cette base.

\paragraph{Ce qui les distingue : deux parités.} Sur la base de calcul,
$Z\otimes Z\ket{ij}=(-1)^{i+j}\ket{ij}$ (valeur propre $+1$ si les deux chiffres sont égaux, $-1$
sinon) et $X\otimes X\ket{ij}=\ket{\bar i\bar j}$ (les deux chiffres sont inversés, avec
$\bar i=1-i$). On en déduit, par exemple,
\[
  X\otimes X\ket{\Phi^-}=\tfrac{1}{\sqrt2}\big(\ket{11}-\ket{00}\big)=-\ket{\Phi^-},
  \qquad
  Z\otimes Z\ket{\Psi^+}=\tfrac{1}{\sqrt2}\big(-\ket{01}-\ket{10}\big)=-\ket{\Psi^+},
\]
et le tableau complet des valeurs propres :
\begin{center}
\renewcommand{\arraystretch}{1.3}
\begin{tabular}{@{}lcccc@{}}
\toprule
 & $\ket{\Phi^+}$ & $\ket{\Phi^-}$ & $\ket{\Psi^+}$ & $\ket{\Psi^-}$ \\
\midrule
$Z\otimes Z$ & $+1$ & $+1$ & $-1$ & $-1$ \\
$X\otimes X$ & $+1$ & $-1$ & $+1$ & $-1$ \\
\bottomrule
\end{tabular}
\end{center}
Les deux opérateurs commutent, car $ZX=-XZ$ donne
$(Z\otimes Z)(X\otimes X)=ZX\otimes ZX=XZ\otimes XZ=(X\otimes X)(Z\otimes Z)$ ; les quatre états de
Bell sont leurs états propres communs, et chaque couple de valeurs propres apparaît une fois
(figure~\ref{fig:grille-bell}). Le même calcul montre que $X\otimes Z$ et $Z\otimes X$
commutent : $(X\otimes Z)(Z\otimes X)=XZ\otimes ZX=(-ZX)\otimes(-XZ)=ZX\otimes XZ=(Z\otimes X)(X\otimes Z)$.
Plus généralement, $P\otimes Q$ et $P'\otimes Q'$ (matrices de Pauli) commutent si $P,P'$
anticommutent et si $Q,Q'$ anticommutent : les deux signes $-1$ se compensent.

\begin{figure}[!ht]
\centering
\adjustbox{max width=\linewidth}{%
\begin{tikzpicture}[>=Stealth, font=\small,
    cell/.style={draw=insarouge, fill=insarouge!6, rounded corners=2pt, thick, minimum width=5.4cm,
                 minimum height=1.35cm, align=center},
    tete/.style={text=insagris, align=center}]
  \node[tete] at (3.0,2.0) {$X\otimes X=+1$\\ {\footnotesize mêmes résultats en base $X$}};
  \node[tete] at (8.8,2.0) {$X\otimes X=-1$\\ {\footnotesize résultats opposés en base $X$}};
  \node[tete, anchor=east] at (-0.1,0.75) {$Z\otimes Z=+1$\\ {\footnotesize mêmes résultats}\\ {\footnotesize en base $Z$}};
  \node[tete, anchor=east] at (-0.1,-0.85) {$Z\otimes Z=-1$\\ {\footnotesize résultats opposés}\\ {\footnotesize en base $Z$}};
  \node[cell] at (3.0,0.75) {$\ket{\Phi^+}=\frac{1}{\sqrt2}\big(\ket{00}+\ket{11}\big)$};
  \node[cell] at (8.8,0.75) {$\ket{\Phi^-}=\frac{1}{\sqrt2}\big(\ket{00}-\ket{11}\big)$};
  \node[cell] at (3.0,-0.85) {$\ket{\Psi^+}=\frac{1}{\sqrt2}\big(\ket{01}+\ket{10}\big)$};
  \node[cell] at (8.8,-0.85) {$\ket{\Psi^-}=\frac{1}{\sqrt2}\big(\ket{01}-\ket{10}\big)$};
\end{tikzpicture}}
\caption{Les quatre états de Bell, classés par les valeurs propres de $Z\otimes Z$ (lignes) et de
$X\otimes X$ (colonnes). La valeur propre est le produit des deux résultats $\pm1$ d'une mesure des
deux qubits dans la base correspondante : $+1$ si les résultats sont égaux, $-1$ s'ils sont opposés.}
\label{fig:grille-bell}
\end{figure}

Les corrélations ne se limitent donc pas à la base $Z$ : $\ket{\Phi^+}$ s'écrit aussi
$\frac{1}{\sqrt2}\big(\ket{++}+\ket{--}\big)$ et donne des résultats égaux en base $X$ comme en base $Z$
(exercice~\ref{ex:bellX}).

\paragraph{D'un état de Bell à l'autre.} Une porte agissant sur un seul qubit suffit. Avec
$X\ket{0}=\ket{1}$, $X\ket{1}=\ket{0}$ et $Z\ket{1}=-\ket{1}$, en agissant sur le second qubit :
\begin{align*}
  (I\otimes X)\ket{\Phi^+}&=\tfrac{1}{\sqrt2}\big(\ket{01}+\ket{10}\big)=\ket{\Psi^+},\\
  (I\otimes Z)\ket{\Phi^+}&=\tfrac{1}{\sqrt2}\big(\ket{00}-\ket{11}\big)=\ket{\Phi^-},\\
  (I\otimes XZ)\ket{\Phi^+}&=(I\otimes X)\tfrac{1}{\sqrt2}\big(\ket{00}-\ket{11}\big)
    =\tfrac{1}{\sqrt2}\big(\ket{01}-\ket{10}\big)=\ket{\Psi^-}.
\end{align*}
Quatre opérations locales ($I$, $X$, $Z$, $XZ$) sur \emph{un} qubit d'une paire de Bell mènent donc
aux quatre états de la base ; c'est le principe du codage superdense
(exercice~\ref{ex:superdense}).

\begin{remarque}[Le singulet est le même dans toutes les bases]\label{rem:singulet}
Pour toute matrice unitaire $U=\begin{pmatrix}a&b\\c&d\end{pmatrix}$, on a $U\ket{0}=a\ket{0}+c\ket{1}$
et $U\ket{1}=b\ket{0}+d\ket{1}$, donc
\begin{align*}
  (U\otimes U)\big(\ket{01}-\ket{10}\big)
  &=\big(a\ket{0}+c\ket{1}\big)\big(b\ket{0}+d\ket{1}\big)-\big(b\ket{0}+d\ket{1}\big)\big(a\ket{0}+c\ket{1}\big)\\
  &=(ad-bc)\big(\ket{01}-\ket{10}\big),
\end{align*}
c'est-à-dire $(U\otimes U)\ket{\Psi^-}=\det(U)\,\ket{\Psi^-}$. Comme $|\det U|=1$, ce facteur est une
phase globale, sans effet physique : $\ket{\Psi^-}$ s'écrit \emph{de la même façon} dans toutes les
bases orthonormées (avec $U=H$, de déterminant $-1$ : $\ket{\Psi^-}=\frac{1}{\sqrt2}\big(\ket{-+}-\ket{+-}\big)$).
Si Alice et Bob mesurent dans la \emph{même} base, quelle qu'elle soit, leurs résultats sont
toujours opposés : c'est le «~singulet~». Parmi les états de Bell, seul $\ket{\Psi^-}$ a cette propriété.
\end{remarque}

% ---------------------------------------------------------------------
\subsection{Un circuit quantique : le générateur d'états de Bell}\label{sec:circuit-bell}
% ---------------------------------------------------------------------

On assemble une porte de Hadamard et un CNOT. Le système est $(A,B)$ avec le ket
$\ket{a,b}=\ket{a}_A\otimes\ket{b}_B$ ; d'après la convention de la section~\ref{sec:circuits}, $B$ est dessiné
en haut et $A$ en bas (figure~\ref{fig:circuit-bell}). Le circuit a deux étapes :
\begin{enumerate}[nosep]
  \item la porte $I\otimes H$ : identité sur $A$, Hadamard sur $B$ ;
  \item un CNOT de contrôle $B$ et de cible $A$ (celui de la section~\ref{sec:cnot}, dont le
        contrôle est le second chiffre).
\end{enumerate}
Les traits pointillés délimitent les états $\ket{\Pi_0}$, $\ket{\Pi_1}$ et $\ket{\Pi_2}$ :
avant les portes, après $I\otimes H$, après le CNOT. Le circuit transforme la base de calcul
en base de Bell (section~\ref{sec:bell}).

\begin{figure}[!ht]
\centering
\begin{tikzpicture}[>=Stealth, font=\small]
  % fils : B en haut, A en bas
  \draw (0,1.4) -- (7.6,1.4);
  \draw (0,0) -- (7.6,0);
  \node[anchor=east] at (0,1.4) {$\ket{0}_B$};
  \node[anchor=east] at (0,0) {$\ket{0}_A$};
  % première étape : H sur B, identité (pointillée) sur A
  \node[porte] at (2.2,1.4) {$H$};
  \node[draw=insagris, dashed, fill=white, minimum size=7mm] at (2.2,0) {$I$};
  % CNOT : contrôle B (haut), cible A (bas)
  \draw[insarouge, thick] (5,1.4) -- (5,0);
  \pic at (5,1.4) {ctl};
  \pic at (5,0) {cible};
  % accolade de sortie
  \draw[insagris, thick] (7.75,1.6) -- (7.9,1.6) -- (7.9,-0.2) -- (7.75,-0.2);
  \node[anchor=west] at (7.9,0.7) {$\ket{\beta^{00}}_{AB}$};
  % états intermédiaires
  \foreach \x in {0.7,3.6,6.4}
    \draw[dashed, insagris!70] (\x,1.9) -- (\x,-0.7);
  \node[anchor=north, font=\footnotesize] at (0.7,-0.7) {$\ket{\Pi_0}=\ket{00}$};
  \node[anchor=north, font=\footnotesize] at (3.6,-0.7) {$\ket{\Pi_1}=\ket{0}\otimes\ket{+}$};
  \node[anchor=north, font=\footnotesize] at (6.4,-0.7) {$\ket{\Pi_2}=\ket{\beta^{00}}$};
\end{tikzpicture}
\caption{Générateur d'états de Bell : $I\otimes H$, puis un CNOT de contrôle $B$ (fil du haut)
et de cible $A$ (fil du bas). L'entrée représentée est $\ket{00}$.}
\label{fig:circuit-bell}
\end{figure}

\subsubsection{Entrée \texorpdfstring{$\ket{00}$}{|00>} : calcul pas à pas}

\begin{align*}
  \ket{\Pi_0}&=\ket{0}_A\otimes\ket{0}_B=\ket{00},\\[1ex]
  \ket{\Pi_1}&=(I\otimes H)\ket{\Pi_0}=I\ket{0}_A\otimes H\ket{0}_B=\ket{0}\otimes\ket{+}
    =\frac{\ket{00}+\ket{01}}{\sqrt2},\\[1ex]
  \ket{\Pi_2}&=\mathrm{CNOT}\ket{\Pi_1}=\frac{\mathrm{CNOT}\ket{00}+\mathrm{CNOT}\ket{01}}{\sqrt2}
    =\frac{\ket{00}+\ket{11}}{\sqrt2}=\ket{\beta^{00}}=\ket{\Phi^+}.
\end{align*}
Dans $\ket{\Pi_1}$, le premier chiffre est celui de $A$ et le second celui de $B$. Le CNOT a
pour contrôle $B$ : $\ket{00}$ ($B=0$) ne change pas, $\ket{01}$ ($B=1$) devient $\ket{11}$
puisque $A$ est inversé. Le résultat est un \emph{état de Bell}.

Avec les matrices, la règle $A\otimes B=(A_{ij}B)$ donne
\[
  I\otimes H=\begin{pmatrix}1\cdot H&0\\0&1\cdot H\end{pmatrix}
  =\frac{1}{\sqrt2}\begin{pmatrix}1&1&0&0\\1&-1&0&0\\0&0&1&1\\0&0&1&-1\end{pmatrix}.
\]
Sa première colonne est $\ket{\Pi_1}=\frac{1}{\sqrt2}(1,1,0,0)^{\mathsf T}$ (image de $\ket{00}$).
Le CNOT échange les amplitudes de $\ket{01}$ et de $\ket{11}$ (section~\ref{sec:cnot}), donc
\[
  \ket{\Pi_2}=\mathrm{CNOT}\,\frac{1}{\sqrt2}\begin{pmatrix}1\\1\\0\\0\end{pmatrix}
  =\frac{1}{\sqrt2}\begin{pmatrix}1\\0\\0\\1\end{pmatrix}.
\]

\subsubsection{Les quatre entrées}

Prenons $B$ dans l'état $\ket{x}$ et $A$ dans l'état $\ket{y}$, avec $x,y\in\{0,1\}$. Alors
$\ket{\Pi_0}=\ket{y}_A\otimes\ket{x}_B=\ket{y,x}$. Comme
$H\ket{x}=\frac{1}{\sqrt2}\big(\ket{0}+(-1)^x\ket{1}\big)$ (c'est $\ket{+}$ pour $x=0$ et
$\ket{-}$ pour $x=1$),
\[
  \ket{\Pi_1}=\ket{y}\otimes H\ket{x}=\frac{1}{\sqrt2}\Big(\ket{y,0}+(-1)^x\ket{y,1}\Big).
\]
Le CNOT (contrôle $B$) ne change pas $\ket{y,0}$ et transforme $\ket{y,1}$ en
$\ket{\bar y,1}$, avec $\bar y=1-y$ :
\[
  \ket{\Pi_2}=\ket{\beta^{xy}}=\frac{1}{\sqrt2}\Big(\ket{y,0}+(-1)^x\ket{\bar y,1}\Big).
\]
Le premier indice de $\ket{\beta^{xy}}$ est l'entrée de $B$ (le qubit qui reçoit $H$), le
second celle de $A$ (la cible). Les quatre cas :

\begin{center}
\renewcommand{\arraystretch}{1.7}
\begin{tabular}{@{}ccccc@{}}
\toprule
$(B,A)=(x,y)$ & $\ket{\Pi_0}=\ket{y,x}$ & $\ket{\Pi_1}$ & $\ket{\Pi_2}=\ket{\beta^{xy}}$ & Nom \\
\midrule
$(0,0)$ & $\ket{00}$ & $\frac{1}{\sqrt2}(\ket{00}+\ket{01})$ & $\frac{1}{\sqrt2}(\ket{00}+\ket{11})$ & $\ket{\Phi^+}$ \\
$(1,0)$ & $\ket{01}$ & $\frac{1}{\sqrt2}(\ket{00}-\ket{01})$ & $\frac{1}{\sqrt2}(\ket{00}-\ket{11})$ & $\ket{\Phi^-}$ \\
$(0,1)$ & $\ket{10}$ & $\frac{1}{\sqrt2}(\ket{10}+\ket{11})$ & $\frac{1}{\sqrt2}(\ket{10}+\ket{01})$ & $\ket{\Psi^+}$ \\
$(1,1)$ & $\ket{11}$ & $\frac{1}{\sqrt2}(\ket{10}-\ket{11})$ & $\frac{1}{\sqrt2}(\ket{10}-\ket{01})$ & $-\ket{\Psi^-}$ \\
\bottomrule
\end{tabular}
\end{center}
Détail de la dernière ligne : $\ket{\Pi_1}=\ket{1}\otimes\ket{-}=\frac{1}{\sqrt2}(\ket{10}-\ket{11})$ ;
le CNOT laisse $\ket{10}$ et envoie $\ket{11}$ sur $\ket{01}$, d'où
$\frac{1}{\sqrt2}(\ket{10}-\ket{01})=-\frac{1}{\sqrt2}(\ket{01}-\ket{10})=-\ket{\Psi^-}$.

\paragraph{Forme matricielle.} Le CNOT échange les lignes $2$ et $4$ de
$I\otimes H$ (il échange les amplitudes de $\ket{01}$ et $\ket{11}$), donc
\[
  \mathrm{CNOT}\,(I\otimes H)=\frac{1}{\sqrt2}\begin{pmatrix}
    1 & 1 & 0 & 0\\ 0 & 0 & 1 & -1\\ 0 & 0 & 1 & 1\\ 1 & -1 & 0 & 0
  \end{pmatrix}.
\]
Ses colonnes sont les images de $\ket{00},\ket{01},\ket{10},\ket{11}$, c'est-à-dire de
$(B,A)=(0,0),(1,0),(0,1),(1,1)$ : ce sont $\ket{\beta^{00}}$, $\ket{\beta^{10}}$,
$\ket{\beta^{01}}$, $\ket{\beta^{11}}$, soit $\ket{\Phi^+}$, $\ket{\Phi^-}$, $\ket{\Psi^+}$
et $-\ket{\Psi^-}$.

\subsubsection{Matrice de densité de l'état de sortie}

Pour $\ket{\Psi}=\ket{\beta^{00}}=\frac{1}{\sqrt2}\ket{00}+\frac{1}{\sqrt2}\ket{11}$ :
\[
  \rho_{AB}=\ket{\Psi}\bra{\Psi}
  =\begin{pmatrix}\frac{1}{\sqrt2}\\0\\0\\\frac{1}{\sqrt2}\end{pmatrix}
   \begin{pmatrix}\frac{1}{\sqrt2}&0&0&\frac{1}{\sqrt2}\end{pmatrix}
  =\frac12\begin{pmatrix}1&0&0&1\\0&0&0&0\\0&0&0&0\\1&0&0&1\end{pmatrix}.
\]
Sous forme de ket-bra :
$\rho_{AB}=\frac12\big(\ket{00}\bra{00}+\ket{00}\bra{11}+\ket{11}\bra{00}+\ket{11}\bra{11}\big)$.
On vérifie $\operatorname{Tr}\rho_{AB}=\frac12(1+1)=1$ et
\[
  \rho_{AB}^2=\frac14\begin{pmatrix}2&0&0&2\\0&0&0&0\\0&0&0&0\\2&0&0&2\end{pmatrix}=\rho_{AB},
\]
donc $\operatorname{Tr}\rho_{AB}^2=1$ : le couple $(A,B)$ est dans un état \emph{pur}.

\begin{remarque}[Porte intriquante]
Les états d'entrée $\ket{ij}$ sont séparables ($D=0$) ; les sorties sont intriquées
($D=\pm\frac12$). Aucune porte agissant séparément sur $A$ et sur $B$ ($U_A\otimes U_B$,
section~\ref{sec:op2}) ne peut créer de l'intrication à partir d'un produit, puisque
$(U_A\otimes U_B)(\ket{u}\otimes\ket{v})=U_A\ket{u}\otimes U_B\ket{v}$ reste un produit.
Une porte à deux qubits comme le CNOT est indispensable.
\end{remarque}

% ---------------------------------------------------------------------
\subsection{Trace partielle : l'état d'un qubit pris seul}\label{sec:trace}
% ---------------------------------------------------------------------

Un état intriqué n'a pas de ket pour $A$ seul (section~\ref{sec:separable}). On peut pourtant
décrire $A$ par une matrice de densité, obtenue en «~oubliant~» $B$.

\begin{definition}{Trace partielle}
Soit $\rho_{AB}$ la matrice de densité du couple $(A,B)$. L'état du qubit $A$ pris seul est
$\rho_A=\operatorname{Tr}_B(\rho_{AB})$, celui de $B$ est $\rho_B=\operatorname{Tr}_A(\rho_{AB})$. La
trace partielle est linéaire et remplace la partie tracée de chaque terme par sa trace :
\[
  \operatorname{Tr}_B\big(\ket{a}\bra{a'}\otimes\ket{b}\bra{b'}\big)=\ket{a}\bra{a'}\;\braket{b'}{b},
  \qquad
  \operatorname{Tr}_A\big(\ket{a}\bra{a'}\otimes\ket{b}\bra{b'}\big)=\braket{a'}{a}\;\ket{b}\bra{b'}.
\]
\end{definition}

\paragraph{En matrices.} On découpe $\rho_{AB}$ (matrice $4\times4$, lignes et colonnes ordonnées
$00,01,10,11$) en quatre blocs $2\times2$ :
\[
  \rho_{AB}=\begin{pmatrix}R_{00}&R_{01}\\R_{10}&R_{11}\end{pmatrix}
  \quad\Longrightarrow\quad
  \rho_A=\begin{pmatrix}\operatorname{Tr}R_{00}&\operatorname{Tr}R_{01}\\
                        \operatorname{Tr}R_{10}&\operatorname{Tr}R_{11}\end{pmatrix},
  \qquad
  \rho_B=R_{00}+R_{11}.
\]
Autrement dit $(\rho_A)_{ij}=\rho_{(i0),(j0)}+\rho_{(i1),(j1)}$ : $\rho_A$ est la matrice des traces
des blocs, et $\rho_B$ la somme des deux blocs diagonaux.

\begin{remarque}[Pourquoi cette définition]
Toutes les prédictions qui portent sur $A$ seul viennent d'opérateurs de la forme $O_A\otimes I$, et
\[
  \operatorname{Tr}\big(\rho_{AB}\,(O_A\otimes I)\big)=\operatorname{Tr}\big(\rho_A\,O_A\big)
\]
(on le vérifie sur $\ket{a}\bra{a'}\otimes\ket{b}\bra{b'}$, puis par linéarité). La matrice $\rho_A$
contient donc exactement ce que l'on peut prédire sur $A$ sans regarder $B$.
\end{remarque}

\paragraph{Exemple 1 : état séparable $\ket{0}\otimes\ket{+}$.} Ici
$\rho_{AB}=\ket{0}\bra{0}\otimes\ket{+}\bra{+}$, donc
\[
  \rho_A=\ket{0}\bra{0}\,\braket{+}{+}=\ket{0}\bra{0},
  \qquad
  \rho_B=\braket{0}{0}\,\ket{+}\bra{+}=\ket{+}\bra{+}.
\]
Avec les matrices,
\[
  \rho_{AB}=\frac12\begin{pmatrix}1&1&0&0\\1&1&0&0\\0&0&0&0\\0&0&0&0\end{pmatrix},
  \qquad
  R_{00}=\frac12\begin{pmatrix}1&1\\1&1\end{pmatrix},\quad R_{01}=R_{10}=R_{11}=0,
\]
d'où
\[
  \rho_A=\begin{pmatrix}\operatorname{Tr}R_{00}&0\\0&0\end{pmatrix}=\begin{pmatrix}1&0\\0&0\end{pmatrix},
  \qquad
  \rho_B=R_{00}+R_{11}=\frac12\begin{pmatrix}1&1\\1&1\end{pmatrix}.
\]
Un qubit d'un état séparable garde son propre état pur.

\paragraph{Exemple 2 : état de Bell $\ket{\Phi^+}$.} On écrit
$\rho_{AB}=\ket{\Phi^+}\bra{\Phi^+}=\frac12\big(\ket{00}\bra{00}+\ket{00}\bra{11}+\ket{11}\bra{00}
+\ket{11}\bra{11}\big)$. Chaque terme est un produit d'opérateurs à un qubit, par exemple
$\ket{00}\bra{11}=\ket{0}\bra{1}\otimes\ket{0}\bra{1}$, donc
\begin{align*}
  \rho_A=\operatorname{Tr}_B(\rho_{AB})
  &=\tfrac12\Big(\ket{0}\bra{0}\,\braket{0}{0}+\ket{0}\bra{1}\,\braket{1}{0}
      +\ket{1}\bra{0}\,\braket{0}{1}+\ket{1}\bra{1}\,\braket{1}{1}\Big)\\
  &=\tfrac12\Big(\ket{0}\bra{0}\cdot1+\ket{0}\bra{1}\cdot0+\ket{1}\bra{0}\cdot0+\ket{1}\bra{1}\cdot1\Big)
   =\tfrac12\big(\ket{0}\bra{0}+\ket{1}\bra{1}\big)=\tfrac12\,\mathrm{Id}.
\end{align*}
Par les blocs,
\[
  \rho_{AB}=\frac12\begin{pmatrix}1&0&0&1\\0&0&0&0\\0&0&0&0\\1&0&0&1\end{pmatrix},
  \qquad
  R_{00}=\frac12\begin{pmatrix}1&0\\0&0\end{pmatrix},\quad
  R_{11}=\frac12\begin{pmatrix}0&0\\0&1\end{pmatrix},\quad
  \operatorname{Tr}R_{01}=\operatorname{Tr}R_{10}=0,
\]
d'où $\rho_A=\frac12\mathrm{Id}$ et $\rho_B=R_{00}+R_{11}=\frac12\mathrm{Id}$ : le même résultat.

Le qubit $A$ est donc dans l'état \emph{complètement mixte} ($\operatorname{Tr}\rho_A^2=\frac14+\frac14=\frac12$),
au centre de la boule de Bloch (section~\ref{sec:bloch}), alors que le couple est dans un état pur
($\operatorname{Tr}\rho_{AB}^2=1$, section~\ref{sec:circuit-bell}). On connaît parfaitement l'état du
couple mais pas celui d'un qubit pris seul : l'information est entièrement contenue dans les
corrélations entre $A$ et $B$. Mesuré seul, $A$ donne $0$ ou $1$ avec la probabilité $\frac12$, dans
toutes les bases.

\paragraph{Formule pratique pour un état pur.} Pour $\ket{\Psi}=\sum_{ij}c_{ij}\ket{ij}$, notons
$C=\begin{pmatrix}c_{00}&c_{01}\\c_{10}&c_{11}\end{pmatrix}$ la matrice $2\times2$ des coefficients
(ligne $i$ : état de $A$, colonne $j$ : état de $B$). Comme $\rho_{(ij),(i'j')}=c_{ij}\,c_{i'j'}^*$,
\[
  (\rho_A)_{ii'}=\sum_jc_{ij}\,c_{i'j}^*,\qquad\text{soit}\qquad \rho_A=C\,C^\dagger .
\]
\emph{Exemple 3.} Pour $\ket{\Psi}=\frac12\ket{00}+\frac12\ket{01}+\frac{1}{\sqrt2}\ket{11}$ (exemple~2
de la section~\ref{sec:mesure2}), $C=\begin{pmatrix}\frac12&\frac12\\0&\frac{1}{\sqrt2}\end{pmatrix}$ et
\[
  \rho_A=CC^\dagger=\begin{pmatrix}\frac14+\frac14&0+\frac{1}{2\sqrt2}\\[0.4ex]\frac{1}{2\sqrt2}&\frac12\end{pmatrix}
  =\begin{pmatrix}\frac12&\frac{\sqrt2}{4}\\[0.4ex]\frac{\sqrt2}{4}&\frac12\end{pmatrix},
  \qquad
  \operatorname{Tr}\rho_A^2=\tfrac14+\tfrac14+2\cdot\tfrac18=\tfrac34 .
\]
Le qubit $A$ est mixte, mais pas complètement : cet état est intriqué sans être un état de Bell.

\begin{resultat}{Pureté d'un qubit et séparabilité}
Pour un état pur $\ket{\Psi}=\sum c_{ij}\ket{ij}$ de deux qubits, avec $D=c_{00}c_{11}-c_{01}c_{10}$,
\[
  \operatorname{Tr}\rho_A^2=\operatorname{Tr}\rho_B^2=1-2|D|^2 .
\]
Donc $\ket{\Psi}$ est séparable ($D=0$) si et seulement si $\rho_A$ est pur. Plus $|D|$ est grand,
plus $\rho_A$ est mixte, jusqu'à $|D|=\frac12$ et $\rho_A=\mathrm{Id}/2$ pour les états de Bell.
\end{resultat}

\emph{Démonstration.} La matrice $\rho_A=CC^\dagger$ a pour trace $\sum|c_{ij}|^2=1$ et pour
déterminant $\det C\,\overline{\det C}=|D|^2$. Si $\lambda_1,\lambda_2$ sont ses valeurs propres,
$\lambda_1+\lambda_2=1$ et $\operatorname{Tr}\rho_A^2=\lambda_1^2+\lambda_2^2=1-2\lambda_1\lambda_2
=1-2|D|^2$ ; le calcul est le même pour $\rho_B$. \hfill$\square$

Comme $\operatorname{Tr}\rho_A^2\geq\frac12$ pour un qubit, on en déduit $|D|\leq\frac12$. Avec
$\operatorname{Tr}\rho_A^2=\frac{1+|\vec r|^2}{2}$ (section~\ref{sec:bloch}), le vecteur de Bloch de $A$
a pour norme $|\vec r|=\sqrt{1-4|D|^2}$ : il part de la sphère pour un état séparable et se rapproche
du centre quand l'intrication croît (figure~\ref{fig:trace-bloch}).

\begin{center}
\renewcommand{\arraystretch}{1.5}
\begin{tabular}{@{}lcccl@{}}
\toprule
\textbf{État du couple} & $D$ & $\operatorname{Tr}\rho_A^2$ & $|\vec r|$ & \textbf{Qubit $A$} \\
\midrule
$\ket{0}\otimes\ket{+}$ & $0$ & $1$ & $1$ & pur, sur la sphère \\
$\frac12\ket{00}+\frac12\ket{01}+\frac{1}{\sqrt2}\ket{11}$ & $\frac{\sqrt2}{4}$ & $\frac34$
  & $\frac{1}{\sqrt2}$ & mixte, dans la boule \\
$\ket{\Phi^+}$ (comme les trois autres états de Bell) & $\pm\frac12$ & $\frac12$ & $0$
  & complètement mixte, au centre \\
\bottomrule
\end{tabular}
\end{center}

\begin{figure}[!ht]
\centering
\adjustbox{max width=\linewidth}{%
\begin{tikzpicture}[>=Stealth, font=\small]
  \def\R{2.3}
  \fill[insagris!12] (0,0) circle (\R);
  \draw[insagris, dotted] (-\R,0) -- (\R,0);
  \draw[insagris, dotted] (0,-\R) -- (0,\R);
  \draw[insagris, dashed] (0,0) circle ({\R*0.70711});
  \draw[insarouge, very thick] (0,0) circle (\R);
  \node[font=\footnotesize, text=insagris, anchor=south west] at (\R,0) {$x$};
  \node[font=\footnotesize, text=insagris, anchor=south west] at (0.03,\R) {$z$};
  % vecteurs de Bloch de rho_A pour trois états du couple
  \fill[insarouge] (0,\R) circle (2.4pt);
  \fill[insagris!80!black] ({\R*0.70711},0) circle (2.4pt);
  \fill[insagris!80!black] (0,0) circle (2.4pt);
  \draw[insagris] (0,\R) -- (3.1,\R);
  \node[anchor=west, font=\footnotesize, align=left] at (3.1,\R)
    {$\ket{0}\otimes\ket{+}$ : $\rho_A=\ket{0}\bra{0}$\\ $D=0$, \ $|\vec r|=1$ (état pur)};
  \draw[insagris] ({\R*0.70711},0) -- (3.1,0);
  \node[anchor=west, font=\footnotesize, align=left] at (3.1,0)
    {$\frac12\ket{00}+\frac12\ket{01}+\frac{1}{\sqrt2}\ket{11}$\\
     $D=\frac{\sqrt2}{4}$, \ $|\vec r|=\frac{1}{\sqrt2}$ (mixte)};
  \draw[insagris] (0,0) -- (3.1,-\R);
  \node[anchor=west, font=\footnotesize, align=left] at (3.1,-\R)
    {$\ket{\Phi^+}$ : $\rho_A=\mathrm{Id}/2$\\ $D=\frac12$, \ $\vec r=\vec 0$ (centre)};
\end{tikzpicture}}
\caption{Le vecteur de Bloch de $\rho_A$ (boule de Bloch vue en coupe) pour trois états du couple
$(A,B)$. Le cercle pointillé est la sphère de rayon $1/\sqrt2$. Plus l'état est intriqué, plus le
point s'éloigne de la sphère des états purs : $|\vec r|=\sqrt{1-4|D|^2}$.}
\label{fig:trace-bloch}
\end{figure}

% ---------------------------------------------------------------------
\subsection{Information : variables aléatoires et états quantiques}\label{sec:info}
% ---------------------------------------------------------------------

\begin{definition}{Entropie, information mutuelle, entropie conditionnelle}
Soit $A$ une variable aléatoire de loi $p(a)$. Son \emph{entropie de Shannon}, en bits, est
\[
  H(A)=-\sum_a p(a)\log_2p(a)\qquad(\text{avec }0\log_20=0).
\]
Elle mesure l'incertitude moyenne sur $A$ : une pièce équilibrée a $H=1$ bit. Pour un couple
$(A,B)$ de loi $p(a,b)$, $H(A,B)=-\sum_{a,b}p(a,b)\log_2p(a,b)$, et
\begin{align*}
  I(A;B)&=H(A)+H(B)-H(A,B) &&\text{(information mutuelle)},\\
  H(A|B)&=H(A,B)-H(B)=H(A)-I(A;B) &&\text{(entropie conditionnelle)}.
\end{align*}
\end{definition}

Les deux écritures de $H(A|B)$ sont égales : la définition de $I(A;B)$ donne
$H(A,B)=H(A)+H(B)-I(A;B)$, donc $H(A,B)-H(B)=H(A)-I(A;B)$. L'information mutuelle $I(A;B)$
mesure ce que la connaissance de $B$ apprend sur $A$ ; $H(A|B)$ est l'incertitude qui reste sur
$A$ quand on connaît $B$.

\subsubsection{Deux variables aléatoires indépendantes}

Alice et Bob lancent chacun une pièce équilibrée, sans lien entre les deux lancers :
$p(a,b)=\frac14$ pour les quatre couples. Chaque pièce donne $0$ ou $1$ avec la probabilité $\frac12$, donc
\begin{align*}
  H(A)&=-\Big(\tfrac12\log_2\tfrac12+\tfrac12\log_2\tfrac12\Big)=1\text{ bit},\\
  H(B)&=1\text{ bit (même calcul)},\\
  H(A,B)&=-4\cdot\tfrac14\log_2\tfrac14=2\text{ bits},\\
  I(A;B)&=H(A)+H(B)-H(A,B)=1+1-2=0,\\
  H(A|B)&=H(A)-I(A;B)=1-0=H(A)=1\text{ bit}.
\end{align*}
Connaître $B$ n'apprend rien sur $A$.

\subsubsection{Deux variables solidaires}

Les deux pièces tombent toujours du même côté : $p(0,0)=p(1,1)=\frac12$ et
$p(0,1)=p(1,0)=0$. Les lois de chaque pièce sont inchangées ($p(a{=}0)=p(0,0)+p(0,1)=\frac12$,
etc.), donc $H(A)=H(B)=1$ bit, mais
\begin{align*}
  H(A,B)&=-2\cdot\tfrac12\log_2\tfrac12-2\cdot0\log_20=1\text{ bit},\\
  I(A;B)&=1+1-1=1\text{ bit},\qquad H(A|B)=1-1=0.
\end{align*}
Connaître $B$ détermine $A$ : plus aucune incertitude ne reste.

\subsubsection{Deux qubits dans un état de Bell}

Pour un système quantique, on remplace la loi de probabilité par la matrice de densité.

\begin{definition}{Entropie de von Neumann}
L'\emph{entropie de von Neumann} d'un état $\rho$, de valeurs propres $\lambda_i$, est
\[
  S(\rho)=-\operatorname{Tr}(\rho\log_2\rho)=-\sum_i\lambda_i\log_2\lambda_i .
\]
Elle vaut $0$ pour un état pur et $1$ bit pour le qubit complètement mélangé $\mathrm{Id}/2$.
En quantique, $H(A)=S(\rho_A)$, $H(B)=S(\rho_B)$ et $H(A,B)=S(\rho_{AB})$ ; les définitions de
$I(A;B)$ et de $H(A|B)$ sont les mêmes.
\end{definition}

On prend l'état de Bell $\ket{\Phi^+}$ produit par le circuit de la
section~\ref{sec:circuit-bell}. Le couple est dans un état pur,
$\rho_{AB}=\ket{\Phi^+}\bra{\Phi^+}$ ; chaque qubit pris seul est dans l'état
$\rho_A=\rho_B=\frac12\mathrm{Id}$ (trace partielle, section~\ref{sec:trace}).

Calcul des entropies :
\begin{itemize}[nosep]
  \item $\rho_{AB}$ est pur : ses valeurs propres sont $1,0,0,0$, donc
        $H(A,B)=S(\rho_{AB})=-1\cdot\log_21=0$ ;
  \item $\rho_A$ et $\rho_B$ ont pour valeurs propres $\frac12,\frac12$, donc
        $H(A)=H(B)=-\big(\frac12\log_2\frac12+\frac12\log_2\frac12\big)=1$ bit.
\end{itemize}
D'où
\begin{align*}
  I(A;B)&=H(A)+H(B)-H(A,B)=1+1-0=2\text{ bits},\\
  H(A|B)&=H(A,B)-H(B)=0-1=-1\text{ bit}\qquad(=H(A)-I(A;B)=1-2).
\end{align*}

\begin{remarque}[Une entropie conditionnelle négative]
Pour des variables classiques, $p(a,b)\leq p(b)$, donc $-\log_2p(a,b)\geq-\log_2p(b)$ et
$H(A,B)\geq H(B)$ : on a toujours $H(A|B)\geq0$, et $I(A;B)\leq H(A)$. Ici $H(A|B)=-1$ et
$I(A;B)=2>H(A)$ : ces valeurs n'ont pas d'équivalent classique, c'est une signature de
l'intrication. Le résultat est surprenant mais mathématiquement correct.
\end{remarque}

\begin{center}
\renewcommand{\arraystretch}{1.4}
\begin{tabular}{@{}lccc@{}}
\toprule
 & indépendantes & solidaires & état de Bell \\
\midrule
$H(A)$ & $1$ & $1$ & $1$ \\
$H(B)$ & $1$ & $1$ & $1$ \\
$H(A,B)$ & $2$ & $1$ & $0$ \\
$I(A;B)$ & $0$ & $1$ & $2$ \\
$H(A|B)$ & $1$ & $0$ & $-1$ \\
\bottomrule
\end{tabular}
\end{center}
(en bits ; les deux premières colonnes sont classiques, la troisième est quantique).

\begin{figure}[!ht]
\centering
\adjustbox{max width=\linewidth}{%
\begin{tikzpicture}[>=Stealth, font=\small]
  \def\bw{0.8}
  % axe des ordonnées (bits)
  \draw[->] (0,-1.5) -- (0,2.8) node[above] {bits};
  \foreach \v in {-1,0,1,2} {
    \draw (-0.07,\v) -- (0.07,\v);
    \node[left, font=\footnotesize] at (-0.1,\v) {$\v$};
  }
  \draw[insagris!60] (0,0) -- (11.4,0);
  % limite classique de I(A;B) : H(A) = 1
  \draw[dashed, insarouge] (0,1) -- (11.4,1);
  \node[anchor=west, font=\scriptsize, text=insarouge, align=left] at (11.4,1)
    {limite classique\\ $I(A;B)\leq H(A)=1$};
  \node[anchor=west, font=\scriptsize, text=insagris, align=left] at (11.4,0)
    {limite classique\\ $H(A|B)\geq0$};
  % barres : (groupe, H(A,B), I(A;B), H(A|B))
  \foreach \g/\hab/\mi/\hc in {0/2/0/1, 1/1/1/0, 2/0/2/-1} {
    \pgfmathsetmacro{\x}{0.9+3.7*\g}
    \fill[insagris!65] (\x,0) rectangle (\x+\bw,\hab);
    \fill[insarouge] (\x+\bw+0.1,0) rectangle (\x+2*\bw+0.1,\mi);
    \fill[black!80] (\x+2*\bw+0.2,0) rectangle (\x+3*\bw+0.2,\hc);
    \node[above, font=\footnotesize] at (\x+0.5*\bw,{max(\hab,0)}) {$\hab$};
    \node[above, font=\footnotesize] at (\x+1.5*\bw+0.1,{max(\mi,0)}) {$\mi$};
    \ifnum\hc<0
      \node[below, font=\footnotesize] at (\x+2.5*\bw+0.2,\hc) {$\hc$};
    \else
      \node[above, font=\footnotesize] at (\x+2.5*\bw+0.2,\hc) {$\hc$};
    \fi
  }
  \node[align=center, anchor=north] at (2.15,-1.65) {variables\\ indépendantes};
  \node[align=center, anchor=north] at (5.85,-1.65) {variables\\ solidaires};
  \node[align=center, anchor=north] at (9.55,-1.65) {état de Bell\\ $\ket{\Phi^+}$};
  % légende
  \fill[insagris!65] (0.9,3.35) rectangle (1.2,3.65);
  \node[anchor=west, font=\footnotesize] at (1.25,3.5) {$H(A,B)$};
  \fill[insarouge] (3.6,3.35) rectangle (3.9,3.65);
  \node[anchor=west, font=\footnotesize] at (3.95,3.5) {$I(A;B)$};
  \fill[black!80] (6.3,3.35) rectangle (6.6,3.65);
  \node[anchor=west, font=\footnotesize] at (6.65,3.5) {$H(A|B)$};
\end{tikzpicture}}
\caption{Entropies (en bits) pour deux pièces indépendantes, deux pièces solidaires et un état de Bell.
Dans les trois cas $H(A)=H(B)=1$. Les deux premiers groupes respectent les bornes classiques
$I(A;B)\leq H(A)$ et $H(A|B)\geq0$ ; l'état de Bell les dépasse ($I=2$, $H(A|B)=-1$).}
\label{fig:entropies}
\end{figure}

% ---------------------------------------------------------------------
\subsection{Téléportation quantique}\label{sec:teleportation}
% ---------------------------------------------------------------------

\emph{Problème.} Alice détient un qubit $Q$ dans un état qu'elle ne connaît pas,
$\ket{\varphi}=\alpha\ket{0}+\beta\ket{1}$ avec $|\alpha|^2+|\beta|^2=1$, et elle veut que Bob
obtienne cet état sur l'un de ses qubits. Elle ne peut pas mesurer $Q$ pour connaître
$\alpha$ et $\beta$ (une mesure ne donne qu'un bit et détruit l'état), elle ne peut pas le
copier (théorème de non-clonage, remarque en fin de section) et elle n'envoie pas le qubit
lui-même. Le protocole est dû à \textcite{bennett1993}.

\begin{definition}{Protocole de téléportation quantique}
\textbf{Ressources.}
\begin{itemize}[nosep]
  \item Un \emph{e-bit} : Alice et Bob partagent une paire de qubits $(A,B)$ dans l'état de Bell
        $\ket{\Phi^+}_{BA}=\frac{1}{\sqrt2}\big(\ket{00}+\ket{11}\big)$ ; Alice détient $A$,
        Bob détient $B$.
  \item Le qubit $Q$ à téléporter, chez Alice.
  \item Un canal classique (deux bits).
\end{itemize}
Le système est écrit $(B,A,Q)$ : $\ket{b,a,q}=\ket{b}_B\otimes\ket{a}_A\otimes\ket{q}_Q$.

\textbf{Étapes.}
\begin{enumerate}[nosep]
  \item Alice applique un CNOT de contrôle $Q$ et de cible $A$.
  \item Alice applique une porte $H$ sur $Q$.
  \item Alice mesure $A$ et $Q$ dans la base de calcul, obtient deux bits $(a,q)$ et les envoie
        à Bob.
  \item Bob applique $X$ si $a=1$, puis $Z$ si $q=1$, à son qubit $B$.
\end{enumerate}
\end{definition}

La figure~\ref{fig:tele-schema} montre l'organisation du protocole et la
figure~\ref{fig:teleportation} son circuit.

\begin{figure}[!ht]
\centering
\adjustbox{max width=\linewidth}{%
\begin{tikzpicture}[>=Stealth, font=\small,
    qb/.style={circle, draw=insarouge, fill=insarouge!8, thick, minimum size=8mm, inner sep=0pt},
    zone/.style={draw=insagris, rounded corners=4pt, thick, dashed, inner sep=8pt},
    num/.style={circle, fill=insarouge, text=white, font=\footnotesize\bfseries, inner sep=0pt,
                minimum size=4.6mm},
    etape/.style={anchor=west, font=\footnotesize, text=insarouge},
    onde/.style={thick, insarouge, decorate,
                 decoration={snake, amplitude=1.3pt, segment length=7pt}}]
  % Alice : Q (état inconnu) et A
  \node[qb] (Q) at (0,1.3) {$Q$};
  \node[qb] (A) at (0,-0.5) {$A$};
  \node[zone, fit=(Q)(A), label={[text=insagris]above:Alice}] (alice) {};
  \node[anchor=east] at (-0.95,1.3) {$\ket{\varphi}=\alpha\ket{0}+\beta\ket{1}$};
  \draw[insarouge, thick] (Q) -- (A);
  % Bob : B
  \node[qb] (B) at (9,-0.5) {$B$};
  \node[zone, fit=(B), label={[text=insagris]above:Bob}] (bob) {};
  % source de la paire de Bell
  \node[draw=insarouge, fill=insarouge!10, rounded corners=2pt, thick, minimum width=3.2cm,
        minimum height=0.9cm] (src) at (4.5,-2.3) {source $\ket{\Phi^+}_{BA}$};
  \draw[onde] (src.west) -| (A.south);
  \draw[onde] (src.east) -| (B.south);
  % étapes d'Alice : pastilles sur le bord droit de sa zone
  \node[num] (n1) at (alice.east |- Q) {1};
  \node[num] (n2) at (alice.east) {2};
  \node[num] (n3) at (alice.east |- A) {3};
  \node[etape] at (n1.east) {\ CNOT de $Q$ vers $A$};
  \node[etape] at (n2.east) {\ porte $H$ sur $Q$};
  % canal classique
  \draw[double, double distance=1.2pt, ->, insagris] (n3.east) -- (bob.west);
  \node[above, font=\footnotesize] at (4.6,-0.45) {deux bits classiques $(a,q)$};
  \node[below, font=\footnotesize, text=insarouge] at (4.6,-0.55) {mesure de $A$ et $Q$ (étape 3)};
  % correction de Bob
  \node[num] (n4) at (bob.east) {4};
  \node[etape, align=left] at (n4.east) {\ $X$ si $a=1$, puis $Z$ si $q=1$};
  \node[anchor=west, font=\footnotesize, align=left] at ($(n4.east)+(0.05,-0.8)$)
    {\ $B$ est alors dans l'état $\ket{\varphi}$};
\end{tikzpicture}}
\caption{Organisation de la téléportation. Alice possède le qubit $Q$ à téléporter et la moitié $A$
d'une paire de Bell, Bob possède l'autre moitié $B$. (1)~CNOT de $Q$ vers $A$, (2)~porte $H$ sur $Q$,
(3)~mesure de $A$ et $Q$, dont les résultats $(a,q)$ partent par un canal classique, (4)~corrections de
Bob. Aucun qubit ne voyage entre Alice et Bob.}
\label{fig:tele-schema}
\end{figure}

\begin{figure}[!ht]
\centering
\begin{adjustbox}{max width=\linewidth}
\begin{tikzpicture}[>=Stealth, font=\small]
  % fils quantiques : Q en haut, A au milieu, B en bas
  \draw (0,2.4) -- (5.3,2.4);
  \draw (0,1.2) -- (5.3,1.2);
  \draw (0,0) -- (10.6,0);
  \node[anchor=east] at (0,2.4) {$\ket{\varphi}_Q$};
  \node[anchor=south west, text=insagris, font=\footnotesize] at (0.05,1.2) {$A$};
  \node[anchor=south west, text=insagris, font=\footnotesize] at (0.05,0) {$B$};
  % paire de Bell partagée
  \draw[insagris, thick] (-0.1,1.5) -- (-0.25,1.5) -- (-0.25,-0.3) -- (-0.1,-0.3);
  \node[anchor=east] at (-0.35,0.6) {$\ket{\Phi^+}_{BA}$};
  % CNOT : contrôle Q, cible A
  \draw[insarouge, thick] (2.6,2.4) -- (2.6,1.2);
  \pic at (2.6,2.4) {ctl};
  \pic at (2.6,1.2) {cible};
  % Hadamard sur Q
  \node[porte] at (4.2,2.4) {$H$};
  % mesures d'Alice
  \pic at (5.6,2.4) {mesure};
  \pic at (5.6,1.2) {mesure};
  % fils classiques (doubles)
  \draw[double, double distance=1.2pt, insagris] (5.9,1.2) -- (7.6,1.2) -- (7.6,0.35);
  \draw[double, double distance=1.2pt, insagris] (5.9,2.4) -- (9.6,2.4) -- (9.6,0.35);
  \node[above, font=\footnotesize] at (6.5,1.2) {$a$};
  \node[above, font=\footnotesize] at (6.5,2.4) {$q$};
  % corrections de Bob
  \node[porte] at (7.6,0) {$X$};
  \node[porte] at (9.6,0) {$Z$};
  \node[anchor=west] at (10.6,0) {$\ket{\varphi}_B$};
  % états intermédiaires
  \foreach \x in {1.4,3.4,4.9}
    \draw[dashed, insagris!70] (\x,2.8) -- (\x,-0.5);
  \node[anchor=north, font=\footnotesize] at (1.4,-0.5) {$\ket{\Pi_0}$};
  \node[anchor=north, font=\footnotesize] at (3.4,-0.5) {$\ket{\Pi_1}$};
  \node[anchor=north, font=\footnotesize] at (4.9,-0.5) {$\ket{\Pi_2}$};
\end{tikzpicture}
\end{adjustbox}
\caption{Circuit de téléportation. Alice détient $Q$ (dans $\ket{\varphi}=\alpha\ket{0}+\beta\ket{1}$)
et $A$, Bob détient $B$ ; $A$ et $B$ forment la paire $\ket{\Phi^+}_{BA}$. Les traits doubles
sont des fils \emph{classiques} qui portent les résultats $a$ (mesure de $A$) et $q$ (mesure de
$Q$) ; $X$ est appliqué si $a=1$, $Z$ si $q=1$. Le ket s'écrit $\ket{b,a,q}$ : le fil du haut
($Q$) est le dernier chiffre.}
\label{fig:teleportation}
\end{figure}

\subsubsection{Calcul pas à pas}

\paragraph{État initial $\ket{\Pi_0}$.} Alice et Bob partagent la paire, et $Q$ est indépendant
d'elle :
\begin{align*}
  \ket{\Pi_0}&=\ket{\Phi^+}_{BA}\otimes\ket{\varphi}_Q
    =\frac{1}{\sqrt2}\big(\ket{00}+\ket{11}\big)\otimes\big(\alpha\ket{0}+\beta\ket{1}\big)\\
  &=\frac{1}{\sqrt2}\Big(\alpha\ket{00}\ket{0}+\beta\ket{00}\ket{1}+\alpha\ket{11}\ket{0}+\beta\ket{11}\ket{1}\Big)\\
  &=\frac{\alpha\ket{000}+\beta\ket{001}+\alpha\ket{110}+\beta\ket{111}}{\sqrt2}.
\end{align*}

\paragraph{Après le CNOT $Q\to A$ : $\ket{\Pi_1}$.} Le CNOT (contrôle $Q$, dernier chiffre ;
cible $A$, chiffre du milieu) agit ainsi : $\ket{b,a,q}\to\ket{b,\,a\oplus q,\,q}$. Terme par terme :
\begin{center}
\renewcommand{\arraystretch}{1.3}
\begin{tabular}{@{}cccc@{}}
\toprule
Terme & $q$ & Effet sur $A$ & Résultat \\
\midrule
$\alpha\ket{000}$ & $0$ & aucun & $\alpha\ket{000}$ \\
$\beta\ket{001}$ & $1$ & $0\to1$ & $\beta\ket{011}$ \\
$\alpha\ket{110}$ & $0$ & aucun & $\alpha\ket{110}$ \\
$\beta\ket{111}$ & $1$ & $1\to0$ & $\beta\ket{101}$ \\
\bottomrule
\end{tabular}
\end{center}
donc
\[
  \ket{\Pi_1}=\frac{\alpha\ket{000}+\beta\ket{011}+\alpha\ket{110}+\beta\ket{101}}{\sqrt2}.
\]
On isole ensuite le dernier chiffre, celui de $Q$, qui va changer à l'étape suivante
(le premier ket est celui du couple $\ket{b,a}$) :
\[
  \ket{\Pi_1}=\frac{\alpha\ket{00}\otimes\ket{0}+\beta\ket{01}\otimes\ket{1}
    +\alpha\ket{11}\otimes\ket{0}+\beta\ket{10}\otimes\ket{1}}{\sqrt2}.
\]

\paragraph{Après la porte $H$ sur $Q$ : $\ket{\Pi_2}$.} On remplace $\ket{0}$ par
$H\ket{0}=\ket{+}$ et $\ket{1}$ par $H\ket{1}=\ket{-}$ dans le facteur $Q$ :
\[
  \ket{\Pi_2}=\frac{1}{\sqrt2}\Big(\alpha\ket{00}\otimes\ket{+}+\beta\ket{01}\otimes\ket{-}
    +\alpha\ket{11}\otimes\ket{+}+\beta\ket{10}\otimes\ket{-}\Big).
\]
Avec $\ket{\pm}=\frac{1}{\sqrt2}(\ket{0}\pm\ket{1})$, le préfacteur devient
$\frac{1}{\sqrt2}\cdot\frac{1}{\sqrt2}=\frac12$ :
\begin{align*}
  \ket{\Pi_2}&=\frac12\Big(\alpha\ket{00}\big(\ket{0}+\ket{1}\big)+\beta\ket{01}\big(\ket{0}-\ket{1}\big)
     +\alpha\ket{11}\big(\ket{0}+\ket{1}\big)+\beta\ket{10}\big(\ket{0}-\ket{1}\big)\Big)\\
  &=\frac12\Big(\alpha\ket{000}+\alpha\ket{001}+\beta\ket{010}-\beta\ket{011}
     +\alpha\ket{110}+\alpha\ket{111}+\beta\ket{100}-\beta\ket{101}\Big).
\end{align*}

\paragraph{Regroupement selon le résultat d'Alice.} Alice mesurera $A$ et $Q$, donc les deux
derniers chiffres. On regroupe les huit termes qui ont les mêmes chiffres $(a,q)$, en mettant
le ket de Bob en premier :
\begin{center}
\renewcommand{\arraystretch}{1.3}
\begin{tabular}{@{}cll@{}}
\toprule
$(a,q)$ & Termes de $\ket{\Pi_2}$ (facteur $\frac12$ omis) & Ket de Bob \\
\midrule
$(0,0)$ & $\alpha\ket{000}+\beta\ket{100}$ & $(\alpha\ket{0}+\beta\ket{1})\ket{00}$ \\
$(0,1)$ & $\alpha\ket{001}-\beta\ket{101}$ & $(\alpha\ket{0}-\beta\ket{1})\ket{01}$ \\
$(1,0)$ & $\beta\ket{010}+\alpha\ket{110}$ & $(\beta\ket{0}+\alpha\ket{1})\ket{10}$ \\
$(1,1)$ & $-\beta\ket{011}+\alpha\ket{111}$ & $(-\beta\ket{0}+\alpha\ket{1})\ket{11}$ \\
\bottomrule
\end{tabular}
\end{center}
d'où
\[
  \ket{\Pi_2}=\frac12\Big[\big(\alpha\ket{0}+\beta\ket{1}\big)\ket{00}
    +\big(\alpha\ket{0}-\beta\ket{1}\big)\ket{01}
    +\big(\beta\ket{0}+\alpha\ket{1}\big)\ket{10}
    +\big(\alpha\ket{1}-\beta\ket{0}\big)\ket{11}\Big].
\]
Dans chaque terme, le ket de gauche est l'état de $B$ et le ket de droite $\ket{a,q}$ celui du
couple $(A,Q)$ d'Alice.

\subsubsection{Mesure d'Alice et corrections de Bob}

Alice mesure $A$ et $Q$. Comme à la section~\ref{sec:mesure2}, la probabilité d'un résultat
$(a,q)$ est le carré de la norme du terme correspondant, et l'état de Bob devient le ket
associé (déjà normé). Ici
\[
  P(a,q)=\Big\|\tfrac12\ket{\phi_{aq}}\Big\|^2=\tfrac14\big(|\alpha|^2+|\beta|^2\big)=\tfrac14
  \quad\text{pour les quatre résultats,}
\]
qui sont donc équiprobables et \emph{indépendants} de $\alpha$ et $\beta$ : les deux bits
envoyés ne révèlent rien sur $\ket{\varphi}$. Les kets de Bob sont
$\ket{\phi_{00}}=\alpha\ket{0}+\beta\ket{1}$, $\ket{\phi_{01}}=\alpha\ket{0}-\beta\ket{1}$,
$\ket{\phi_{10}}=\beta\ket{0}+\alpha\ket{1}$ et $\ket{\phi_{11}}=-\beta\ket{0}+\alpha\ket{1}$.

\begin{center}
\renewcommand{\arraystretch}{1.5}
\begin{tabular}{@{}cccll@{}}
\toprule
Résultat $(a,q)$ & Probabilité & État de Bob & Lien avec $\ket{\varphi}$ & Correction de Bob \\
\midrule
$(0,0)$ & $\frac14$ & $\alpha\ket{0}+\beta\ket{1}$ & $\ket{\varphi}$ & $I$ (rien) \\
$(0,1)$ & $\frac14$ & $\alpha\ket{0}-\beta\ket{1}$ & $Z\ket{\varphi}$ & $Z$ \\
$(1,0)$ & $\frac14$ & $\beta\ket{0}+\alpha\ket{1}$ & $X\ket{\varphi}$ & $X$ \\
$(1,1)$ & $\frac14$ & $-\beta\ket{0}+\alpha\ket{1}$ & $XZ\ket{\varphi}$ & $X$ puis $Z$ (opérateur $ZX$) \\
\bottomrule
\end{tabular}
\end{center}

\paragraph{Vérification des quatre corrections.} Avec $X\ket{0}=\ket{1}$, $X\ket{1}=\ket{0}$,
$Z\ket{0}=\ket{0}$ et $Z\ket{1}=-\ket{1}$ :
\begin{align*}
  (0,0):\quad & I\big(\alpha\ket{0}+\beta\ket{1}\big)=\alpha\ket{0}+\beta\ket{1},\\
  (0,1):\quad & Z\big(\alpha\ket{0}-\beta\ket{1}\big)=\alpha\ket{0}-\beta Z\ket{1}=\alpha\ket{0}+\beta\ket{1},\\
  (1,0):\quad & X\big(\beta\ket{0}+\alpha\ket{1}\big)=\beta\ket{1}+\alpha\ket{0}=\alpha\ket{0}+\beta\ket{1},\\
  (1,1):\quad & X\big(-\beta\ket{0}+\alpha\ket{1}\big)=-\beta\ket{1}+\alpha\ket{0}=\alpha\ket{0}-\beta\ket{1},
     \ \text{puis}\ Z\big(\alpha\ket{0}-\beta\ket{1}\big)=\alpha\ket{0}+\beta\ket{1}.
\end{align*}
Dans le dernier cas, avec les matrices,
$ZX=\begin{pmatrix}1&0\\0&-1\end{pmatrix}\begin{pmatrix}0&1\\1&0\end{pmatrix}
=\begin{pmatrix}0&1\\-1&0\end{pmatrix}$ et
$ZX\begin{pmatrix}-\beta\\ \alpha\end{pmatrix}=\begin{pmatrix}\alpha\\ \beta\end{pmatrix}$.
L'ordre compte : l'état de Bob est $XZ\ket{\varphi}$ ($Z\ket{\varphi}=\alpha\ket{0}-\beta\ket{1}$,
puis $X$ donne $\alpha\ket{1}-\beta\ket{0}$), et on l'annule avec
$(XZ)^{-1}=Z^{-1}X^{-1}=ZX$ : on applique d'abord $X$, puis $Z$, comme dans le circuit.
Dans les quatre cas, Bob retrouve exactement $\ket{\varphi}=\alpha\ket{0}+\beta\ket{1}$.

\subsubsection{Ce que Bob sait avant de recevoir les bits}

Tant que Bob ignore $(a,q)$, son qubit est dans le mélange des quatre kets $\ket{\phi_{aq}}$, chacun
avec la probabilité $\frac14$ :
\begin{align*}
  \rho_B&=\frac14\sum_{a,q}\ket{\phi_{aq}}\bra{\phi_{aq}}\\
  &=\frac14\Bigg[\begin{pmatrix}|\alpha|^2&\alpha\beta^*\\ \alpha^*\beta&|\beta|^2\end{pmatrix}
    +\begin{pmatrix}|\alpha|^2&-\alpha\beta^*\\ -\alpha^*\beta&|\beta|^2\end{pmatrix}
    +\begin{pmatrix}|\beta|^2&\alpha^*\beta\\ \alpha\beta^*&|\alpha|^2\end{pmatrix}
    +\begin{pmatrix}|\beta|^2&-\alpha^*\beta\\ -\alpha\beta^*&|\alpha|^2\end{pmatrix}\Bigg]\\
  &=\frac14\begin{pmatrix}2(|\alpha|^2+|\beta|^2)&0\\0&2(|\alpha|^2+|\beta|^2)\end{pmatrix}
   =\frac12\begin{pmatrix}1&0\\0&1\end{pmatrix}=\frac12\,\mathrm{Id}.
\end{align*}
Sans les deux bits, Bob a donc un qubit complètement mélangé, qui ne dépend pas de $\alpha$ ni de
$\beta$ : il n'a reçu aucune information sur $\ket{\varphi}$. C'est l'envoi classique de $(a,q)$
qui permet de la retrouver, donc rien ne voyage plus vite que la lumière.

\begin{remarque}[Bilan de la téléportation]
\begin{itemize}[nosep]
  \item Ressources : $1$ e-bit (la paire de Bell) et $2$ bits classiques pour transmettre
        \emph{un qubit} d'état quelconque, sans qu'Alice connaisse $\alpha$ ni $\beta$.
  \item Pas de clonage : après la mesure, $A$ et $Q$ sont dans l'état de base $\ket{a,q}$.
        Alice n'a plus $\ket{\varphi}$ : l'état a été \emph{déplacé}, pas copié. L'intrication
        entre $A$ et $B$ a été consommée.
  \item Les quatre corrections $I,Z,X,ZX$ sont des portes contrôlées par des bits classiques
        (section~\ref{sec:controlees}).
\end{itemize}
\end{remarque}

\begin{remarque}[Les gestes d'Alice forment une mesure de Bell]
Le CNOT de contrôle $Q$ sur $A$ suivi de $H$ sur $Q$ est l'\emph{inverse} du générateur de la
section~\ref{sec:circuit-bell} (les portes $H$ et CNOT sont leurs propres inverses, et l'ordre
s'inverse) : il envoie la base de Bell du couple $(A,Q)$ sur la base de calcul. Mesurer ensuite dans
la base de calcul revient donc à mesurer $(A,Q)$ dans la base de Bell. Les deux bits $(a,q)$
disent \emph{quel} état de Bell Alice a trouvé, ce qui explique les quatre corrections de Bob.
\end{remarque}

\begin{remarque}[Non-clonage]
Alice ne peut pas simplement copier $\ket{\varphi}$ \parencite{wootters1982}. Supposons une porte
unitaire $U$ telle que $U\big(\ket{\psi}\otimes\ket{0}\big)=\ket{\psi}\otimes\ket{\psi}$ pour tout
$\ket{\psi}$. Elle copie $\ket{0}$ et $\ket{1}$, mais par linéarité
\[
  U\big(\ket{+}\otimes\ket{0}\big)=\tfrac{1}{\sqrt2}\big(\ket{00}+\ket{11}\big)
  \neq\ket{+}\otimes\ket{+}=\tfrac12\big(\ket{00}+\ket{01}+\ket{10}+\ket{11}\big),
\]
ce qui contredit l'hypothèse. Le CNOT copie un bit classique, pas une superposition (remarque de la
section~\ref{sec:cnot}). La téléportation respecte cette interdiction : l'état est déplacé, il n'est
jamais présent en deux endroits.
\end{remarque}

% ---------------------------------------------------------------------
\subsection{Corrélations non locales : le jeu CHSH}\label{sec:chsh}
% ---------------------------------------------------------------------

Les états de Bell donnent des résultats corrélés en base $Z$ comme en base $X$
(section~\ref{sec:bell}). Mais deux enveloppes contenant la même carte, préparées à l'avance, donnent
aussi des résultats égaux : des corrélations dans une seule base ne prouvent pas l'intrication. Bell
\parencite{bell1964} a montré qu'on peut trancher en \emph{changeant de base de mesure} ; la version de
Clauser, Horne, Shimony et Holt \parencite{clauser1969} se présente comme un jeu.

\begin{definition}{Le jeu CHSH}
Alice et Bob, éloignés, ne communiquent plus une fois le jeu commencé. Un arbitre tire deux bits $x$ et
$y$ de façon uniforme et indépendante, donne $x$ à Alice et $y$ à Bob. Chacun répond par un bit : $a$ pour
Alice, $b$ pour Bob. Ils \emph{gagnent} si
\[
  a\oplus b=x\cdot y,
\]
c'est-à-dire si $a=b$ pour $(x,y)\in\{(0,0),(0,1),(1,0)\}$ et si $a\neq b$ pour $(x,y)=(1,1)$.
\end{definition}

\paragraph{Stratégie classique : au plus $\frac34$.} Une stratégie déterministe est $a=f(x)$,
$b=g(y)$. Si elle gagnait les quatre parties, on aurait
\[
  f(0)\oplus g(0)=0,\qquad f(0)\oplus g(1)=0,\qquad f(1)\oplus g(0)=0,\qquad f(1)\oplus g(1)=1 .
\]
En additionnant ces quatre égalités modulo $2$, chacune des valeurs $f(0),f(1),g(0),g(1)$ apparaît
deux fois : le membre de gauche vaut $0$, celui de droite vaut $1$. C'est absurde, donc au moins une
partie sur quatre est perdue et $P(\text{gain})\leq\frac34$. La stratégie «~toujours répondre $0$~»
atteint $\frac34$ (elle ne perd que pour $(1,1)$). Un hasard partagé ne change rien : il mélange des
stratégies déterministes, dont aucune ne dépasse $\frac34$.

\paragraph{Stratégie quantique.} Alice et Bob partagent avant le jeu une paire $\ket{\Phi^+}$ (un
e-bit), puis chacun mesure son qubit sans rien dire à l'autre. Pour un angle $\theta$, notons
\[
  \ket{\theta_+}=\cos\tfrac\theta2\,\ket{0}+\sin\tfrac\theta2\,\ket{1},
  \qquad
  \ket{\theta_-}=-\sin\tfrac\theta2\,\ket{0}+\cos\tfrac\theta2\,\ket{1}
\]
la base de mesure d'angle $\theta$ : c'est la direction $\theta$ dans le plan $xz$ de la sphère de
Bloch ($\theta=0$ : base $Z$ ; $\theta=\frac\pi2$ : base $X$). Le résultat vaut $0$ pour
$\ket{\theta_+}$ et $1$ pour $\ket{\theta_-}$. La stratégie est
\begin{itemize}[nosep]
  \item Alice : $\theta=0$ si $x=0$ (base $Z$), $\theta=\frac\pi2$ si $x=1$ (base $X$) ;
  \item Bob : $\theta=\frac\pi4$ si $y=0$, $\theta=-\frac\pi4$ si $y=1$ (bases propres de
        $\frac{Z\pm X}{\sqrt2}$).
\end{itemize}

\begin{figure}[!ht]
\centering
\adjustbox{max width=\linewidth}{%
\begin{tikzpicture}[>=Stealth, font=\footnotesize,
    joueur/.style={draw=insarouge, fill=insarouge!6, rounded corners=2pt, thick,
                   minimum width=2.1cm, minimum height=1.1cm, align=center},
    arb/.style={draw=insagris, fill=insagris!8, rounded corners=2pt, thick,
                minimum width=3.6cm, minimum height=0.8cm, align=center}]
  % --- le jeu ---
  \node[arb] (arb) at (3.2,2.7) {arbitre : $x,y$ au hasard};
  \node[joueur] (al) at (0,0) {Alice};
  \node[joueur] (bo) at (6.4,0) {Bob};
  \draw[->, thick, insagris] (arb.west) -| node[pos=0.2, above] {$x$} (al.north);
  \draw[->, thick, insagris] (arb.east) -| node[pos=0.2, above] {$y$} (bo.north);
  \draw[insarouge, thick, decorate, decoration={snake, amplitude=1.1pt, segment length=6pt}]
    (al.east) -- (bo.west);
  \node[insarouge, above=1pt] at (3.2,0.05) {$\ket{\Phi^+}$};
  \node[text=insagris, below=2pt, font=\scriptsize] at (3.2,-0.05) {paire intriquée};
  \draw[->, thick] (al.south) -- ++(0,-1.0) node[below] {$a$};
  \draw[->, thick] (bo.south) -- ++(0,-1.0) node[below] {$b$};
  \node[draw=insagris, rounded corners=2pt, inner sep=3pt] at (3.2,-1.75) {gain si $a\oplus b=x\cdot y$};
  % --- directions de mesure dans le plan xz de la sphère de Bloch ---
  \begin{scope}[shift={(11.9,0.3)}]
    \draw[insagris!70] (0,0) circle (1.7);
    \draw[insagris!50, dashed] (-1.7,0) -- (1.7,0);
    \draw[insagris!50, dashed] (0,-1.7) -- (0,1.7);
    \node[font=\scriptsize, text=insagris, anchor=north] at (0,-1.72) {$\ket{1}$};
    \node[font=\scriptsize, text=insagris, anchor=east] at (-1.72,-0.02) {$\ket{-}$};
    \draw[->, very thick, insarouge] (0,0) -- (90:1.7);
    \draw[->, very thick, insarouge] (0,0) -- (0:1.7);
    \draw[->, very thick, insagris] (0,0) -- (45:1.7);
    \draw[->, very thick, insagris] (0,0) -- (135:1.7);
    \node[insarouge, anchor=south, inner sep=2pt] at (0,1.72) {$x=0$ : $Z$ $(\ket{0})$};
    \node[insarouge, anchor=west, inner sep=2pt] at (1.72,0) {$x=1$ : $X$ $(\ket{+})$};
    \node[text=insagris, anchor=south west, align=left, inner sep=2pt] at (1.25,1.25) {$y=0$\\ $\frac{Z+X}{\sqrt2}$};
    \node[text=insagris, anchor=south east, align=right, inner sep=2pt] at (-1.25,1.25) {$y=1$\\ $\frac{Z-X}{\sqrt2}$};
    \draw[insagris] (0.85,0) arc (0:45:0.85);
    \node[font=\scriptsize, text=insagris] at (22.5:1.12) {$\frac\pi4$};
    \draw[insagris] (0,0.85) arc (90:135:0.85);
    \node[font=\scriptsize, text=insagris] at (112.5:1.12) {$\frac\pi4$};
    \draw[insagris] (0.6,0.6) arc (45:90:0.85);
    \node[font=\scriptsize, text=insagris] at (67.5:1.12) {$\frac\pi4$};
  \end{scope}
\end{tikzpicture}}
\caption{Le jeu CHSH (à gauche) et les quatre directions de mesure de la stratégie quantique (à
droite), dans le plan $xz$ de la sphère de Bloch : Alice en rouge, Bob en gris. Pour
$(x,y)=(0,0),(0,1),(1,0)$, les directions font un angle $\frac\pi4$ ; pour $(1,1)$, un angle
$\frac{3\pi}4$.}
\label{fig:chsh}
\end{figure}

\begin{resultat}{Corrélations de $\ket{\Phi^+}$ en bases tournées}
Si Alice mesure dans la base d'angle $\alpha$ et Bob dans la base d'angle $\beta$, alors
\[
  P(a=b)=\cos^2\frac{\alpha-\beta}{2}.
\]
\end{resultat}

\noindent\emph{Preuve.} Avec $c=\cos\frac\alpha2$ et $s=\sin\frac\alpha2$,
$\ket{\alpha_+\alpha_+}+\ket{\alpha_-\alpha_-}=(c^2+s^2)\ket{00}+(cs-sc)\big(\ket{01}+\ket{10}\big)+(s^2+c^2)\ket{11}
=\ket{00}+\ket{11}$ : on peut donc écrire $\ket{\Phi^+}=\frac{1}{\sqrt2}\big(\ket{\alpha_+\alpha_+}+\ket{\alpha_-\alpha_-}\big)$.
Si Alice obtient $a=0$ (probabilité $\frac12$), le qubit de Bob est dans $\ket{\alpha_+}$ et Bob obtient $b=0$
avec la probabilité $|\braket{\beta_+}{\alpha_+}|^2$, où $\braket{\beta_+}{\alpha_+}=\cos\frac\beta2\cos\frac\alpha2+\sin\frac\beta2\sin\frac\alpha2=\cos\frac{\alpha-\beta}{2}$.
Si Alice obtient $a=1$, Bob est dans $\ket{\alpha_-}$ et obtient $b=1$ avec la même probabilité, car
$\braket{\beta_-}{\alpha_-}=\braket{\beta_+}{\alpha_+}$. Donc
$P(a=b)=\frac12\cos^2\frac{\alpha-\beta}2+\frac12\cos^2\frac{\alpha-\beta}2$. \hfill$\square$

\medskip
\noindent Pour les quatre couples $(x,y)$ :
\begin{center}
\renewcommand{\arraystretch}{1.3}
\begin{tabular}{@{}c c c c c c@{}}
\toprule
$(x,y)$ & $\alpha$ & $\beta$ & $P(a=b)$ & condition de gain & $P(\text{gain})$\\
\midrule
$(0,0)$ & $0$ & $\frac\pi4$ & $\cos^2\frac\pi8$ & $a=b$ & $\cos^2\frac\pi8$\\
$(0,1)$ & $0$ & $-\frac\pi4$ & $\cos^2\frac\pi8$ & $a=b$ & $\cos^2\frac\pi8$\\
$(1,0)$ & $\frac\pi2$ & $\frac\pi4$ & $\cos^2\frac\pi8$ & $a=b$ & $\cos^2\frac\pi8$\\
$(1,1)$ & $\frac\pi2$ & $-\frac\pi4$ & $\cos^2\frac{3\pi}8=\sin^2\frac\pi8$ & $a\neq b$ &
  $1-\sin^2\frac\pi8=\cos^2\frac\pi8$\\
\bottomrule
\end{tabular}
\end{center}
Alice et Bob gagnent donc chaque partie avec la même probabilité :
\[
  P(\text{gain})=\cos^2\frac\pi8=\frac{2+\sqrt2}{4}\approx0{,}854>\frac34 .
\]

\begin{remarque}[Ce que montre le jeu CHSH]
\begin{itemize}[nosep]
  \item Aucune stratégie classique, même avec du hasard partagé, ne dépasse $\frac34$ : l'intrication
        produit des corrélations qu'aucun modèle «~local~» (résultats fixés à l'avance) ne reproduit.
        Les expériences sur des photons intriqués confirment la valeur quantique.
  \item $\cos^2\frac\pi8$ est le maximum quantique (\emph{borne de Tsirelson}, résultat admis,
        \cite{tsirelson1980}) : on ne peut pas gagner à tous les coups. Avec
        $E_{xy}=P(a=b)-P(a\neq b)=\cos(\alpha-\beta)$ et $S=E_{00}+E_{01}+E_{10}-E_{11}$, on a
        $P(\text{gain})=\frac12+\frac S8$ ; l'inégalité de Bell donne $|S|\leq2$ en classique, et la
        stratégie ci-dessus atteint $S=2\sqrt2$.
  \item Il n'y a pas de communication : la réponse de Bob est uniforme ($P(b=0)=\frac12$) quelle que
        soit la base choisie par Alice. L'intrication ne transmet pas d'information ; elle ne fait
        que corréler les résultats.
\end{itemize}
\end{remarque}

\clearpage
% ---------------------------------------------------------------------
\subsection{Récapitulatif du chapitre}\label{sec:recap5}
% ---------------------------------------------------------------------

\begin{center}
\adjustbox{max width=\linewidth}{%
\begin{tikzpicture}[>=Stealth, font=\footnotesize,
    boite/.style={draw=insarouge, fill=insarouge!6, rounded corners=2pt, thick,
                  minimum width=4.4cm, minimum height=1.45cm, align=center, inner sep=3pt},
    lien/.style={->, very thick, insagris},
    etiq/.style={font=\scriptsize\itshape, text=insagris, inner sep=2pt}]
  \node[boite] (mes) at (0,0)
    {\textbf{Mesure à deux qubits}\\ $P(ij)=|c_{ij}|^2$\\ {\scriptsize\S\,\ref{sec:mesure2}}};
  \node[boite] (sep) at (6.3,0)
    {\textbf{Séparable ou intriqué}\\ $D=c_{00}c_{11}-c_{01}c_{10}$\\ {\scriptsize\S\,\ref{sec:separable}}};
  \node[boite] (bel) at (12.6,0)
    {\textbf{États de Bell}\\ $\ket{\Phi^\pm}$, $\ket{\Psi^\pm}$ ; $\mathrm{CNOT}\,(I\otimes H)$\\
     {\scriptsize\S\,\ref{sec:bell}, \ref{sec:circuit-bell}}};
  \node[boite] (tra) at (12.6,-2.3)
    {\textbf{Trace partielle, entropies}\\ $\rho_A=\operatorname{Tr}_B\rho_{AB}$ ; $S(\rho)$, $I(A;B)$\\
     {\scriptsize\S\,\ref{sec:trace}, \ref{sec:info}}};
  \node[boite] (tel) at (6.3,-2.3)
    {\textbf{Téléportation}\\ e-bit + 2 bits classiques\\ {\scriptsize\S\,\ref{sec:teleportation}}};
  \node[boite] (chs) at (0,-2.3)
    {\textbf{Jeu CHSH}\\ $\frac34$ classique, $\cos^2\frac\pi8$ quantique\\ {\scriptsize\S\,\ref{sec:chsh}}};
  \draw[lien] (mes) -- (sep) node[etiq, midway, above=1pt] {corrélations};
  \draw[lien] (sep) -- (bel) node[etiq, midway, above=1pt] {cas extrême};
  \draw[lien] (bel) -- (tra) node[etiq, midway, right=1pt] {un qubit seul};
  \draw[lien] (tra) -- (tel) node[etiq, midway, above=1pt] {ressource};
  \draw[lien] (tel) -- (chs) node[etiq, midway, above=1pt] {non-localité};
\end{tikzpicture}}
\end{center}

\begin{aretenir}
\begin{itemize}
  \item Deux qubits : quatre amplitudes $c_{ij}$, $P(ij)=|c_{ij}|^2$. Mesurer $A$ supprime les termes
        incompatibles et renormalise : l'état de $B$ peut changer.
  \item Séparable $\Longleftrightarrow D=c_{00}c_{11}-c_{01}c_{10}=0$. Sinon l'état est intriqué et
        aucun qubit n'a d'état propre (pour $\ket{\Phi^+}$, $P(00)\neq P(A{=}0)P(B{=}0)$).
  \item États de Bell : base orthonormée d'états maximalement intriqués, états propres de
        $Z\otimes Z$ et $X\otimes X$ ; $\mathrm{CNOT}\,(I\otimes H)$ les produit à partir de la base de
        calcul. Des portes locales ne créent pas d'intrication.
  \item Qubit seul : $\rho_A=\operatorname{Tr}_B\rho_{AB}$, avec $\operatorname{Tr}\rho_A^2=1-2|D|^2$ pour un état pur.
        Pour $\ket{\Phi^+}$ : $\rho_A=\mathrm{Id}/2$ alors que le couple est pur ; $S(\rho_{AB})=0$,
        $S(\rho_A)=1$ bit, $I(A;B)=2$ bits, $H(A|B)=-1$ (impossible en classique).
  \item Téléportation : un e-bit et deux bits classiques ; $P(a,q)=\frac14$ ; Bob applique $X$ si $a=1$
        puis $Z$ si $q=1$. L'état est déplacé, pas copié ; sans les bits, $\rho_B=\mathrm{Id}/2$.
  \item Jeu CHSH : toute stratégie classique gagne avec une probabilité $\leq\frac34$ ; avec $\ket{\Phi^+}$
        et des bases tournées, $\cos^2\frac\pi8\approx0{,}854$ (borne de Tsirelson), sans aucune
        communication.
\end{itemize}
\end{aretenir}

\begin{center}
\renewcommand{\arraystretch}{1.2}
\begin{tabular}{@{}l l r@{}}
\toprule
\textbf{Notion} & \textbf{Formule} & \textbf{§} \\
\midrule
Mesure de deux qubits & $P(ij)=|c_{ij}|^2$, \quad $P(A{=}a)=P(a0)+P(a1)$ & \ref{sec:mesure2} \\
Séparabilité & $\ket{\Psi}=\ket{\Psi}_A\otimes\ket{\Psi}_B\ \Longleftrightarrow\ D=0$ & \ref{sec:separable} \\
États de Bell & $\ket{\Phi^\pm}=\dfrac{\ket{00}\pm\ket{11}}{\sqrt2}$, \quad
  $\ket{\Psi^\pm}=\dfrac{\ket{01}\pm\ket{10}}{\sqrt2}$ & \ref{sec:bell} \\
Trace partielle & $\operatorname{Tr}_B\big(\ket{a}\bra{a'}\otimes\ket{b}\bra{b'}\big)
  =\ket{a}\bra{a'}\,\braket{b'}{b}$ & \ref{sec:trace} \\
Qubit d'un état pur & $\rho_A=CC^\dagger$, \quad $\operatorname{Tr}\rho_A^2=1-2|D|^2$ & \ref{sec:trace} \\
Entropies & $S(\rho)=-\sum_i\lambda_i\log_2\lambda_i$, \quad $I(A;B)=H(A)+H(B)-H(A,B)$ & \ref{sec:info} \\
Téléportation & $P(a,q)=\frac14$, \quad correction $Z^{q}X^{a}$ (d'abord $X$) & \ref{sec:teleportation} \\
Singulet & $(U\otimes U)\ket{\Psi^-}=\det(U)\,\ket{\Psi^-}$ & \ref{sec:bell} \\
Jeu CHSH & $P_{\mathrm{cl}}\leq\frac34$, \quad $P_{\mathrm{quant}}=\cos^2\frac\pi8=\frac{2+\sqrt2}{4}$ & \ref{sec:chsh} \\
\bottomrule
\end{tabular}
\end{center}

\begin{savoirfaire}
\begin{itemize}
  \item Calculer $P(ij)$, les probabilités d'un seul qubit et l'état après la mesure de l'un d'eux ;
        tester la séparabilité avec $D$ et calculer la pureté de $\rho_A$ (exercice~\ref{ex:separable3}).
  \item Suivre un état dans le générateur de Bell, écrire $\rho_{AB}$ et reconnaître les corrélations
        d'un état de Bell dans deux bases (exercice~\ref{ex:bellX}).
  \item Calculer $H(A)$, $H(A,B)$, $I(A;B)$ et $H(A|B)$ pour une loi ou un état (exercice~\ref{ex:entropies}).
  \item Dérouler la téléportation pour un état donné et utiliser les états de Bell pour le codage
        superdense (exercices~\ref{ex:telenum} et~\ref{ex:superdense}).
  \item Prouver la borne $\frac34$ du jeu CHSH classique et calculer la probabilité de gain d'une
        stratégie quantique (section~\ref{sec:chsh}).
\end{itemize}
\end{savoirfaire}
.
%  Il ne contient ni préambule ni \begin{document} : les environnements
%  (definition, resultat, remarque, aretenir, savoirfaire, exercice, correction),
%  les macros (\ii, \ee, \dd, \pd, \pdd, \Hop, \ket, \bra, \braket), les symboles
%  de circuits (ctl, cible, croix, mesure, porte) et les couleurs (insarouge,
%  insagris) sont définis dans main.tex. Un chapitre = une \section.
% =====================================================================

% Macros de Dirac (déjà définies dans main.tex ; ce bloc évite l'erreur
% « Undefined control sequence \ket » si main.tex est une ancienne version).
\providecommand{\ket}[1]{|#1\rangle}
\providecommand{\bra}[1]{\langle #1|}
\providecommand{\braket}[2]{\langle #1|#2\rangle}

% =====================================================================
\section{États quantiques, intrication et états de Bell}\label{chap:bell}
% =====================================================================

Les états de deux qubits forment un espace de dimension~$4$, et la plupart d'entre eux ne se
décomposent pas en «~un état pour $A$, un état pour $B$~» : ils sont \emph{intriqués}. Ce chapitre
apprend à mesurer un couple de qubits, à reconnaître l'intrication (critère de séparabilité), à la
produire avec un circuit (états de Bell), à décrire un qubit pris seul (trace partielle), à
quantifier l'information (entropies) et à s'en servir pour téléporter un état.

% ---------------------------------------------------------------------
\subsection{États de deux qubits et mesure}\label{sec:mesure2}
% ---------------------------------------------------------------------

Un état quelconque de $\mathcal{H}^4$ se décompose sur la base de calcul :
\[
  \ket{\Psi}=c_{00}\ket{00}+c_{01}\ket{01}+c_{10}\ket{10}+c_{11}\ket{11}
  =\begin{pmatrix}c_{00}\\ c_{01}\\ c_{10}\\ c_{11}\end{pmatrix},
  \qquad \sum_{i,j}|c_{ij}|^2=1 .
\]

\begin{resultat}{Mesure des deux qubits dans la base de calcul}
\[
  P(ij)=|\braket{ij}{\Psi}|^2=|c_{ij}|^2 .
\]
Pour un seul qubit : $P(A{=}0)=P(00)+P(01)$, $P(A{=}1)=P(10)+P(11)$,
$P(B{=}0)=P(00)+P(10)$, $P(B{=}1)=P(01)+P(11)$.
\end{resultat}

En effet $\braket{ij}{\Psi}=\sum_{k,l}c_{kl}\braket{ij}{kl}=\sum_{k,l}c_{kl}\delta_{ik}\delta_{jl}=c_{ij}$.

\paragraph{Exemple 1.} $\ket{\Psi}=\frac{1}{\sqrt2}\ket{00}+\frac{1}{\sqrt2}\ket{01}
=\big(\tfrac{1}{\sqrt2},\tfrac{1}{\sqrt2},0,0\big)^{\mathsf T}$. Il est normalisé :
$\frac12+\frac12+0+0=1$. Les probabilités jointes sont
$P(00)=P(01)=\frac12$ et $P(10)=P(11)=0$. Pour chaque qubit :
\[
  P(A{=}0)=\tfrac12+\tfrac12=1,\quad P(A{=}1)=0+0=0,\quad
  P(B{=}0)=\tfrac12+0=\tfrac12,\quad P(B{=}1)=\tfrac12+0=\tfrac12 .
\]
$A$ donne toujours $0$ ; $B$ donne $0$ ou $1$ au hasard.

\paragraph{Exemple 2 (mesure de $A$ seul).} Soit
$\ket{\Psi}=\frac12\ket{00}+\frac12\ket{01}+\frac{1}{\sqrt2}\ket{11}$. Normalisation :
$\frac14+\frac14+0+\frac12=1$. Probabilités jointes : $P(00)=\frac14$, $P(01)=\frac14$,
$P(10)=0$, $P(11)=\frac12$. Donc
\[
  P(A{=}0)=\tfrac14+\tfrac14=\tfrac12,\quad P(A{=}1)=0+\tfrac12=\tfrac12,\quad
  P(B{=}0)=\tfrac14+0=\tfrac14,\quad P(B{=}1)=\tfrac14+\tfrac12=\tfrac34 .
\]

\begin{remarque}[Effondrement après la mesure de $A$]
Si $A$ donne le résultat $a$, on ne garde que les termes de $\ket{\Psi}$ dont le premier
chiffre vaut $a$, puis on renormalise :
\[
  \ket{\Psi}\ \to\ \frac{c_{a0}\ket{a0}+c_{a1}\ket{a1}}{\sqrt{|c_{a0}|^2+|c_{a1}|^2}},
  \qquad\text{probabilité }|c_{a0}|^2+|c_{a1}|^2 .
\]
\end{remarque}

Pour l'exemple 2 :
\begin{itemize}[nosep]
  \item $A=0$ (probabilité $\frac12$) : on garde $\frac12\ket{00}+\frac12\ket{01}$, de norme
        au carré $\frac12$ ; renormalisé, $\frac{1}{\sqrt2}\big(\ket{00}+\ket{01}\big)
        =\ket{0}\otimes\ket{+}$. Ensuite $B$ vaut $0$ ou $1$ avec la probabilité $\frac12$.
  \item $A=1$ (probabilité $\frac12$) : on garde $\frac{1}{\sqrt2}\ket{11}$, de norme au carré
        $\frac12$ ; renormalisé, $\ket{11}=\ket{1}\otimes\ket{1}$. Ensuite $B$ vaut $1$
        avec certitude.
\end{itemize}
Contrôle : $P(B{=}0)=\frac12\cdot\frac12+\frac12\cdot0=\frac14$, comme calculé plus haut.
Ici la mesure de $A$ \emph{modifie} l'état de $B$ ($\ket{+}$ ou $\ket{1}$) : c'est la
signature d'un état intriqué (section~\ref{sec:separable}). La figure~\ref{fig:mesure2} résume
les deux mesures successives.

\begin{figure}[!ht]
\centering
\adjustbox{max width=\linewidth}{%
\begin{tikzpicture}[>=Stealth, font=\small,
    noeud/.style={draw=insarouge, fill=insarouge!6, rounded corners=2pt, thick, inner sep=4pt, align=center},
    fin/.style={draw=insagris, fill=insagris!10, rounded corners=2pt, inner sep=4pt, align=center},
    lien/.style={->, thick, insagris},
    proba/.style={font=\footnotesize, fill=white, inner sep=1.5pt, text=black}]
  \node[noeud] (psi) at (0,0) {$\ket{\Psi}$};
  \node[noeud] (a0) at (4.4,1.6) {$A=0$\\ {\footnotesize état $\ket{0}\otimes\ket{+}$}};
  \node[noeud] (a1) at (4.4,-1.6) {$A=1$\\ {\footnotesize état $\ket{1}\otimes\ket{1}$}};
  \node[fin] (b00) at (9.2,2.5) {$B=0$ : $\ket{00}$};
  \node[fin] (b01) at (9.2,0.7) {$B=1$ : $\ket{01}$};
  \node[fin] (b11) at (9.2,-1.6) {$B=1$ : $\ket{11}$};
  \draw[lien] (psi) -- (a0) node[proba, midway] {$\frac12$};
  \draw[lien] (psi) -- (a1) node[proba, midway] {$\frac12$};
  \draw[lien] (a0) -- (b00) node[proba, midway] {$\frac12$};
  \draw[lien] (a0) -- (b01) node[proba, midway] {$\frac12$};
  \draw[lien] (a1) -- (b11) node[proba, midway] {$1$};
  \node[anchor=west] at (11.0,2.5) {$P(00)=\frac12\cdot\frac12=\frac14$};
  \node[anchor=west] at (11.0,0.7) {$P(01)=\frac12\cdot\frac12=\frac14$};
  \node[anchor=west] at (11.0,-1.6) {$P(11)=\frac12\cdot1=\frac12$};
  \node[text=insagris, font=\footnotesize\itshape] at (2.2,3.0) {mesure de $A$};
  \node[text=insagris, font=\footnotesize\itshape] at (6.8,3.6) {mesure de $B$};
  \node[text=insagris, font=\footnotesize\itshape, anchor=west] at (11.0,3.6) {probabilité jointe};
\end{tikzpicture}}
\caption{Mesure de $A$ puis de $B$ pour l'état de l'exemple~2,
$\ket{\Psi}=\frac12\ket{00}+\frac12\ket{01}+\frac{1}{\sqrt2}\ket{11}$. Chaque flèche porte une
probabilité, et le produit des probabilités le long d'un chemin est la probabilité jointe. Après la
mesure de $A$, l'état de $B$ dépend du résultat ($\ket{+}$ ou $\ket{1}$).}
\label{fig:mesure2}
\end{figure}

% ---------------------------------------------------------------------
\subsection{États séparables et états intriqués}\label{sec:separable}
% ---------------------------------------------------------------------

\begin{definition}{Séparable ou intriqué}
Un état $\ket{\Psi}_{AB}$ est \emph{séparable} s'il s'écrit comme un produit
$\ket{\Psi}_{AB}=\ket{\Psi}_A\otimes\ket{\Psi}_B$. Sinon il est \emph{intriqué} : on ne
peut pas attribuer d'état propre à chaque qubit.
\end{definition}

\paragraph{Exemple séparable (exemple 1 ci-dessus).} Par linéarité du produit tensoriel :
\begin{align*}
  \ket{\Psi}&=\tfrac{1}{\sqrt2}\ket{00}_{AB}+\tfrac{1}{\sqrt2}\ket{01}_{AB}
    =\tfrac{1}{\sqrt2}\Big(\ket{0}_A\otimes\ket{0}_B+\ket{0}_A\otimes\ket{1}_B\Big)\\
  &=\ket{0}_A\otimes\Big(\tfrac{1}{\sqrt2}\big(\ket{0}+\ket{1}\big)\Big)_B
    =\ket{0}_A\otimes\ket{+}_B ,
\end{align*}
donc $\ket{\Psi}_A=\ket{0}$ et $\ket{\Psi}_B=\ket{+}$. Contrôle par les composantes :
$\begin{pmatrix}1\\ 0\end{pmatrix}\otimes\tfrac{1}{\sqrt2}\begin{pmatrix}1\\ 1\end{pmatrix}
=\tfrac{1}{\sqrt2}\begin{pmatrix}1\cdot1\\ 1\cdot1\\ 0\cdot1\\ 0\cdot1\end{pmatrix}
=\tfrac{1}{\sqrt2}\begin{pmatrix}1\\ 1\\ 0\\ 0\end{pmatrix}$.

\subsubsection{Un critère de séparabilité}

Un produit quelconque a pour coefficients $c_{00}=ac$, $c_{01}=ad$, $c_{10}=bc$,
$c_{11}=bd$ (formule de la section~\ref{sec:tenseur}, avec $a,b$ pour $A$ et $c,d$ pour $B$), donc
\[
  c_{00}c_{11}=(ac)(bd)=abcd=(ad)(bc)=c_{01}c_{10}.
\]

\begin{resultat}{Test de séparabilité (deux qubits)}
$\ket{\Psi}=\sum c_{ij}\ket{ij}$ est séparable si et seulement si
\[
  D=c_{00}c_{11}-c_{01}c_{10}=0 .
\]
\end{resultat}

\emph{Démonstration.} ($\Rightarrow$) Calcul ci-dessus. ($\Leftarrow$) On regroupe les
termes selon l'état de $A$ :
\[
  \ket{\Psi}=\ket{0}\otimes\big(c_{00}\ket{0}+c_{01}\ket{1}\big)+\ket{1}\otimes\big(c_{10}\ket{0}+c_{11}\ket{1}\big).
\]
Si $D=0$, les lignes $(c_{00},c_{01})$ et $(c_{10},c_{11})$ sont proportionnelles. Si
$(c_{00},c_{01})\neq(0,0)$, il existe $\lambda$ avec $(c_{10},c_{11})=\lambda(c_{00},c_{01})$
et $\ket{\Psi}=\big(\ket{0}+\lambda\ket{1}\big)\otimes\big(c_{00}\ket{0}+c_{01}\ket{1}\big)$.
Si $(c_{00},c_{01})=(0,0)$, alors $\ket{\Psi}=\ket{1}\otimes\big(c_{10}\ket{0}+c_{11}\ket{1}\big)$.
Dans les deux cas $\ket{\Psi}$ est un produit ; on renormalise chaque facteur.
\hfill$\square$

\subsubsection{Application à plusieurs états}

\begin{center}
\renewcommand{\arraystretch}{1.5}
\begin{tabular}{@{}lccl@{}}
\toprule
\textbf{État} & $(c_{00},c_{01},c_{10},c_{11})$ & $D$ & \textbf{Nature} \\
\midrule
$\ket{00}$ & $(1,0,0,0)$ & $1\cdot0-0\cdot0=0$ & séparable \\
$\ket{+}\otimes\ket{+}$ & $\frac12(1,1,1,1)$ & $\frac14-\frac14=0$ & séparable \\
$\ket{+}\otimes\ket{-}$ & $\frac12(1,-1,1,-1)$ & $-\frac14-(-\frac14)=0$ & séparable \\
$\frac{1}{\sqrt2}(\ket{00}+\ket{01})$ & $\frac{1}{\sqrt2}(1,1,0,0)$ & $0-\frac12\cdot0=0$ & séparable \\
$\ket{\Phi^+}=\frac{1}{\sqrt2}(\ket{00}+\ket{11})$ & $\frac{1}{\sqrt2}(1,0,0,1)$ & $\frac12-0=\frac12$ & intriqué \\
$\ket{\Phi^-}=\frac{1}{\sqrt2}(\ket{00}-\ket{11})$ & $\frac{1}{\sqrt2}(1,0,0,-1)$ & $-\frac12-0=-\frac12$ & intriqué \\
$\ket{\Psi^+}=\frac{1}{\sqrt2}(\ket{01}+\ket{10})$ & $\frac{1}{\sqrt2}(0,1,1,0)$ & $0-\frac12=-\frac12$ & intriqué \\
$\ket{\Psi^-}=\frac{1}{\sqrt2}(\ket{01}-\ket{10})$ & $\frac{1}{\sqrt2}(0,1,-1,0)$ & $0-(-\frac12)=\frac12$ & intriqué \\
$\frac12\ket{00}+\frac12\ket{01}+\frac{1}{\sqrt2}\ket{11}$ & $(\frac12,\frac12,0,\frac{1}{\sqrt2})$ & $\frac{1}{2\sqrt2}-0$ & intriqué \\
\bottomrule
\end{tabular}
\end{center}

Les quatre états de Bell $\ket{\Phi^\pm}$ et $\ket{\Psi^\pm}$ de la table sont tous intriqués
($D=\pm\frac12$) ; leurs propriétés font l'objet de la section~\ref{sec:bell}.

\paragraph{Preuve directe que $\ket{\Phi^+}$ est intriqué.} Supposons
$\ket{\Phi^+}=(a\ket{0}+b\ket{1})\otimes(c\ket{0}+d\ket{1})$. En identifiant les
composantes : $ac=\frac{1}{\sqrt2}$, $ad=0$, $bc=0$, $bd=\frac{1}{\sqrt2}$. La condition
$ad=0$ impose $a=0$ (impossible car $ac\neq0$) ou $d=0$ (impossible car $bd\neq0$).
Contradiction : aucun produit ne convient.

\paragraph{Corrélations.} Pour un produit $\ket{\Psi}_A\otimes\ket{\Psi}_B$ avec
$\ket{\Psi}_A=a_0\ket{0}+a_1\ket{1}$ et $\ket{\Psi}_B=b_0\ket{0}+b_1\ket{1}$, on a
$P(ij)=|a_ib_j|^2=|a_i|^2|b_j|^2=P(A{=}i)\,P(B{=}j)$ : les deux qubits sont
\emph{indépendants}. Vérification avec $\ket{+}\otimes\ket{+}$ : $P(00)=\frac14=\frac12\cdot\frac12$.
Pour $\ket{\Phi^+}$ : $P(00)=P(11)=\frac12$ et $P(01)=P(10)=0$, donc
$P(A{=}0)=P(B{=}0)=\frac12$, mais
\[
  P(00)=\tfrac12\neq\tfrac12\cdot\tfrac12=P(A{=}0)\,P(B{=}0) .
\]
Les deux qubits donnent toujours le même résultat, aléatoire : ils sont corrélés.
La figure~\ref{fig:epr} illustre la différence entre les deux cas.

\begin{figure}[!ht]
\centering
\adjustbox{max width=\linewidth}{%
\begin{tikzpicture}[>=Stealth, font=\small,
    labo/.style={draw=insagris, fill=insagris!8, rounded corners=2pt, thick, minimum width=2.6cm,
                 minimum height=1.0cm, align=center},
    source/.style={draw=insarouge, fill=insarouge!10, rounded corners=2pt, thick, minimum width=3.6cm,
                   minimum height=1.0cm, align=center},
    droit/.style={->, thick, insagris},
    onde/.style={->, thick, insarouge, decorate,
                 decoration={snake, amplitude=1.3pt, segment length=7pt, post length=2mm}},
    res/.style={font=\small, align=center}]
  % (a) état séparable
  \node[labo] (Aa) at (0,0) {Alice\\ mesure $A$};
  \node[source] (Sa) at (6,0) {source\\ $\ket{0}_A\otimes\ket{+}_B$};
  \node[labo] (Ba) at (12,0) {Bob\\ mesure $B$};
  \draw[droit] (Sa.west) -- (Aa.east);
  \draw[droit] (Sa.east) -- (Ba.west);
  \node[res, below=3pt of Aa] {$0\ \ 0\ \ 0\ \ 0\ \ 0\ \ 0$};
  \node[res, below=3pt of Ba] {$1\ \ 0\ \ 0\ \ 1\ \ 0\ \ 1$};
  \node[res, text=insagris, font=\footnotesize] at (6,-1.0) {$A$ donne toujours $0$ ; $B$ est aléatoire.\\
        Les deux résultats sont indépendants.};
  \node[anchor=south west, font=\small\bfseries] at (-1.3,0.75) {(a) état séparable};
  % (b) état intriqué
  \begin{scope}[yshift=-3.6cm]
    \node[labo] (Ab) at (0,0) {Alice\\ mesure $A$};
    \node[source] (Sb) at (6,0) {source\\ $\ket{\Phi^+}_{AB}$};
    \node[labo] (Bb) at (12,0) {Bob\\ mesure $B$};
    \draw[onde] (Sb.west) -- (Ab.east);
    \draw[onde] (Sb.east) -- (Bb.west);
    \node[res, below=3pt of Ab] {$1\ \ 0\ \ 0\ \ 1\ \ 1\ \ 0$};
    \node[res, below=3pt of Bb] {$1\ \ 0\ \ 0\ \ 1\ \ 1\ \ 0$};
    \node[res, text=insagris, font=\footnotesize] at (6,-1.0) {$A$ et $B$ sont chacun aléatoires,\\
          mais toujours égaux : ils sont corrélés.};
    \node[anchor=south west, font=\small\bfseries] at (-1.3,0.75) {(b) état intriqué};
  \end{scope}
\end{tikzpicture}}
\caption{Six mesures successives sur des paires préparées dans le même état (les résultats sont
des exemples). (a) Pour un produit, chaque qubit a son propre comportement. (b) Pour
$\ket{\Phi^+}=\frac{1}{\sqrt2}(\ket{00}+\ket{11})$, chaque résultat pris seul est un tirage à pile ou
face, mais les deux résultats coïncident toujours ; la ligne ondulée symbolise l'intrication.}
\label{fig:epr}
\end{figure}

\paragraph{Et un qubit pris seul ?} Dans un état séparable, chaque qubit garde son propre état pur ;
dans un état intriqué, aucun ket ne décrit un qubit isolé. On le décrit alors par une matrice de
densité, obtenue par la trace partielle (section~\ref{sec:trace}).

% ---------------------------------------------------------------------
\subsection{Les états de Bell}\label{sec:bell}
% ---------------------------------------------------------------------

\begin{definition}{Les quatre états de Bell}
\[
  \ket{\Phi^{\pm}}=\frac{1}{\sqrt2}\big(\ket{00}\pm\ket{11}\big),
  \qquad
  \ket{\Psi^{\pm}}=\frac{1}{\sqrt2}\big(\ket{01}\pm\ket{10}\big).
\]
Ils sont intriqués ($|D|=\frac12$, section~\ref{sec:separable}) et le sont au maximum :
$|D|\leq\frac12$ pour tout état normalisé de deux qubits, et chaque qubit d'un état de Bell,
pris seul, est complètement aléatoire (section~\ref{sec:trace}).
\end{definition}

\paragraph{Une base de $\mathcal{H}^4$.} Chaque état est de norme $\frac12+\frac12=1$. Les
$\ket{\Phi^\pm}$ n'ont de composantes que sur $\ket{00}$ et $\ket{11}$, les $\ket{\Psi^\pm}$ que sur
$\ket{01}$ et $\ket{10}$, donc $\braket{\Phi^\pm}{\Psi^\pm}=0$ pour toutes les combinaisons de signes ;
enfin $\braket{\Phi^+}{\Phi^-}=\frac12(1-1)=0$ et $\braket{\Psi^+}{\Psi^-}=\frac12(1-1)=0$. Les quatre
états forment donc une base orthonormée de $\mathcal{H}^4$, la \emph{base de Bell}, constituée
uniquement d'états intriqués. Le circuit de la section~\ref{sec:circuit-bell} fait passer de la
base de calcul à cette base.

\paragraph{Ce qui les distingue : deux parités.} Sur la base de calcul,
$Z\otimes Z\ket{ij}=(-1)^{i+j}\ket{ij}$ (valeur propre $+1$ si les deux chiffres sont égaux, $-1$
sinon) et $X\otimes X\ket{ij}=\ket{\bar i\bar j}$ (les deux chiffres sont inversés, avec
$\bar i=1-i$). On en déduit, par exemple,
\[
  X\otimes X\ket{\Phi^-}=\tfrac{1}{\sqrt2}\big(\ket{11}-\ket{00}\big)=-\ket{\Phi^-},
  \qquad
  Z\otimes Z\ket{\Psi^+}=\tfrac{1}{\sqrt2}\big(-\ket{01}-\ket{10}\big)=-\ket{\Psi^+},
\]
et le tableau complet des valeurs propres :
\begin{center}
\renewcommand{\arraystretch}{1.3}
\begin{tabular}{@{}lcccc@{}}
\toprule
 & $\ket{\Phi^+}$ & $\ket{\Phi^-}$ & $\ket{\Psi^+}$ & $\ket{\Psi^-}$ \\
\midrule
$Z\otimes Z$ & $+1$ & $+1$ & $-1$ & $-1$ \\
$X\otimes X$ & $+1$ & $-1$ & $+1$ & $-1$ \\
\bottomrule
\end{tabular}
\end{center}
Les deux opérateurs commutent, car $ZX=-XZ$ donne
$(Z\otimes Z)(X\otimes X)=ZX\otimes ZX=XZ\otimes XZ=(X\otimes X)(Z\otimes Z)$ ; les quatre états de
Bell sont leurs états propres communs, et chaque couple de valeurs propres apparaît une fois
(figure~\ref{fig:grille-bell}). Le même calcul montre que $X\otimes Z$ et $Z\otimes X$
commutent : $(X\otimes Z)(Z\otimes X)=XZ\otimes ZX=(-ZX)\otimes(-XZ)=ZX\otimes XZ=(Z\otimes X)(X\otimes Z)$.
Plus généralement, $P\otimes Q$ et $P'\otimes Q'$ (matrices de Pauli) commutent si $P,P'$
anticommutent et si $Q,Q'$ anticommutent : les deux signes $-1$ se compensent.

\begin{figure}[!ht]
\centering
\adjustbox{max width=\linewidth}{%
\begin{tikzpicture}[>=Stealth, font=\small,
    cell/.style={draw=insarouge, fill=insarouge!6, rounded corners=2pt, thick, minimum width=5.4cm,
                 minimum height=1.35cm, align=center},
    tete/.style={text=insagris, align=center}]
  \node[tete] at (3.0,2.0) {$X\otimes X=+1$\\ {\footnotesize mêmes résultats en base $X$}};
  \node[tete] at (8.8,2.0) {$X\otimes X=-1$\\ {\footnotesize résultats opposés en base $X$}};
  \node[tete, anchor=east] at (-0.1,0.75) {$Z\otimes Z=+1$\\ {\footnotesize mêmes résultats}\\ {\footnotesize en base $Z$}};
  \node[tete, anchor=east] at (-0.1,-0.85) {$Z\otimes Z=-1$\\ {\footnotesize résultats opposés}\\ {\footnotesize en base $Z$}};
  \node[cell] at (3.0,0.75) {$\ket{\Phi^+}=\frac{1}{\sqrt2}\big(\ket{00}+\ket{11}\big)$};
  \node[cell] at (8.8,0.75) {$\ket{\Phi^-}=\frac{1}{\sqrt2}\big(\ket{00}-\ket{11}\big)$};
  \node[cell] at (3.0,-0.85) {$\ket{\Psi^+}=\frac{1}{\sqrt2}\big(\ket{01}+\ket{10}\big)$};
  \node[cell] at (8.8,-0.85) {$\ket{\Psi^-}=\frac{1}{\sqrt2}\big(\ket{01}-\ket{10}\big)$};
\end{tikzpicture}}
\caption{Les quatre états de Bell, classés par les valeurs propres de $Z\otimes Z$ (lignes) et de
$X\otimes X$ (colonnes). La valeur propre est le produit des deux résultats $\pm1$ d'une mesure des
deux qubits dans la base correspondante : $+1$ si les résultats sont égaux, $-1$ s'ils sont opposés.}
\label{fig:grille-bell}
\end{figure}

Les corrélations ne se limitent donc pas à la base $Z$ : $\ket{\Phi^+}$ s'écrit aussi
$\frac{1}{\sqrt2}\big(\ket{++}+\ket{--}\big)$ et donne des résultats égaux en base $X$ comme en base $Z$
(exercice~\ref{ex:bellX}).

\paragraph{D'un état de Bell à l'autre.} Une porte agissant sur un seul qubit suffit. Avec
$X\ket{0}=\ket{1}$, $X\ket{1}=\ket{0}$ et $Z\ket{1}=-\ket{1}$, en agissant sur le second qubit :
\begin{align*}
  (I\otimes X)\ket{\Phi^+}&=\tfrac{1}{\sqrt2}\big(\ket{01}+\ket{10}\big)=\ket{\Psi^+},\\
  (I\otimes Z)\ket{\Phi^+}&=\tfrac{1}{\sqrt2}\big(\ket{00}-\ket{11}\big)=\ket{\Phi^-},\\
  (I\otimes XZ)\ket{\Phi^+}&=(I\otimes X)\tfrac{1}{\sqrt2}\big(\ket{00}-\ket{11}\big)
    =\tfrac{1}{\sqrt2}\big(\ket{01}-\ket{10}\big)=\ket{\Psi^-}.
\end{align*}
Quatre opérations locales ($I$, $X$, $Z$, $XZ$) sur \emph{un} qubit d'une paire de Bell mènent donc
aux quatre états de la base ; c'est le principe du codage superdense
(exercice~\ref{ex:superdense}).

\begin{remarque}[Le singulet est le même dans toutes les bases]\label{rem:singulet}
Pour toute matrice unitaire $U=\begin{pmatrix}a&b\\c&d\end{pmatrix}$, on a $U\ket{0}=a\ket{0}+c\ket{1}$
et $U\ket{1}=b\ket{0}+d\ket{1}$, donc
\begin{align*}
  (U\otimes U)\big(\ket{01}-\ket{10}\big)
  &=\big(a\ket{0}+c\ket{1}\big)\big(b\ket{0}+d\ket{1}\big)-\big(b\ket{0}+d\ket{1}\big)\big(a\ket{0}+c\ket{1}\big)\\
  &=(ad-bc)\big(\ket{01}-\ket{10}\big),
\end{align*}
c'est-à-dire $(U\otimes U)\ket{\Psi^-}=\det(U)\,\ket{\Psi^-}$. Comme $|\det U|=1$, ce facteur est une
phase globale, sans effet physique : $\ket{\Psi^-}$ s'écrit \emph{de la même façon} dans toutes les
bases orthonormées (avec $U=H$, de déterminant $-1$ : $\ket{\Psi^-}=\frac{1}{\sqrt2}\big(\ket{-+}-\ket{+-}\big)$).
Si Alice et Bob mesurent dans la \emph{même} base, quelle qu'elle soit, leurs résultats sont
toujours opposés : c'est le «~singulet~». Parmi les états de Bell, seul $\ket{\Psi^-}$ a cette propriété.
\end{remarque}

% ---------------------------------------------------------------------
\subsection{Un circuit quantique : le générateur d'états de Bell}\label{sec:circuit-bell}
% ---------------------------------------------------------------------

On assemble une porte de Hadamard et un CNOT. Le système est $(A,B)$ avec le ket
$\ket{a,b}=\ket{a}_A\otimes\ket{b}_B$ ; d'après la convention de la section~\ref{sec:circuits}, $B$ est dessiné
en haut et $A$ en bas (figure~\ref{fig:circuit-bell}). Le circuit a deux étapes :
\begin{enumerate}[nosep]
  \item la porte $I\otimes H$ : identité sur $A$, Hadamard sur $B$ ;
  \item un CNOT de contrôle $B$ et de cible $A$ (celui de la section~\ref{sec:cnot}, dont le
        contrôle est le second chiffre).
\end{enumerate}
Les traits pointillés délimitent les états $\ket{\Pi_0}$, $\ket{\Pi_1}$ et $\ket{\Pi_2}$ :
avant les portes, après $I\otimes H$, après le CNOT. Le circuit transforme la base de calcul
en base de Bell (section~\ref{sec:bell}).

\begin{figure}[!ht]
\centering
\begin{tikzpicture}[>=Stealth, font=\small]
  % fils : B en haut, A en bas
  \draw (0,1.4) -- (7.6,1.4);
  \draw (0,0) -- (7.6,0);
  \node[anchor=east] at (0,1.4) {$\ket{0}_B$};
  \node[anchor=east] at (0,0) {$\ket{0}_A$};
  % première étape : H sur B, identité (pointillée) sur A
  \node[porte] at (2.2,1.4) {$H$};
  \node[draw=insagris, dashed, fill=white, minimum size=7mm] at (2.2,0) {$I$};
  % CNOT : contrôle B (haut), cible A (bas)
  \draw[insarouge, thick] (5,1.4) -- (5,0);
  \pic at (5,1.4) {ctl};
  \pic at (5,0) {cible};
  % accolade de sortie
  \draw[insagris, thick] (7.75,1.6) -- (7.9,1.6) -- (7.9,-0.2) -- (7.75,-0.2);
  \node[anchor=west] at (7.9,0.7) {$\ket{\beta^{00}}_{AB}$};
  % états intermédiaires
  \foreach \x in {0.7,3.6,6.4}
    \draw[dashed, insagris!70] (\x,1.9) -- (\x,-0.7);
  \node[anchor=north, font=\footnotesize] at (0.7,-0.7) {$\ket{\Pi_0}=\ket{00}$};
  \node[anchor=north, font=\footnotesize] at (3.6,-0.7) {$\ket{\Pi_1}=\ket{0}\otimes\ket{+}$};
  \node[anchor=north, font=\footnotesize] at (6.4,-0.7) {$\ket{\Pi_2}=\ket{\beta^{00}}$};
\end{tikzpicture}
\caption{Générateur d'états de Bell : $I\otimes H$, puis un CNOT de contrôle $B$ (fil du haut)
et de cible $A$ (fil du bas). L'entrée représentée est $\ket{00}$.}
\label{fig:circuit-bell}
\end{figure}

\subsubsection{Entrée \texorpdfstring{$\ket{00}$}{|00>} : calcul pas à pas}

\begin{align*}
  \ket{\Pi_0}&=\ket{0}_A\otimes\ket{0}_B=\ket{00},\\[1ex]
  \ket{\Pi_1}&=(I\otimes H)\ket{\Pi_0}=I\ket{0}_A\otimes H\ket{0}_B=\ket{0}\otimes\ket{+}
    =\frac{\ket{00}+\ket{01}}{\sqrt2},\\[1ex]
  \ket{\Pi_2}&=\mathrm{CNOT}\ket{\Pi_1}=\frac{\mathrm{CNOT}\ket{00}+\mathrm{CNOT}\ket{01}}{\sqrt2}
    =\frac{\ket{00}+\ket{11}}{\sqrt2}=\ket{\beta^{00}}=\ket{\Phi^+}.
\end{align*}
Dans $\ket{\Pi_1}$, le premier chiffre est celui de $A$ et le second celui de $B$. Le CNOT a
pour contrôle $B$ : $\ket{00}$ ($B=0$) ne change pas, $\ket{01}$ ($B=1$) devient $\ket{11}$
puisque $A$ est inversé. Le résultat est un \emph{état de Bell}.

Avec les matrices, la règle $A\otimes B=(A_{ij}B)$ donne
\[
  I\otimes H=\begin{pmatrix}1\cdot H&0\\0&1\cdot H\end{pmatrix}
  =\frac{1}{\sqrt2}\begin{pmatrix}1&1&0&0\\1&-1&0&0\\0&0&1&1\\0&0&1&-1\end{pmatrix}.
\]
Sa première colonne est $\ket{\Pi_1}=\frac{1}{\sqrt2}(1,1,0,0)^{\mathsf T}$ (image de $\ket{00}$).
Le CNOT échange les amplitudes de $\ket{01}$ et de $\ket{11}$ (section~\ref{sec:cnot}), donc
\[
  \ket{\Pi_2}=\mathrm{CNOT}\,\frac{1}{\sqrt2}\begin{pmatrix}1\\1\\0\\0\end{pmatrix}
  =\frac{1}{\sqrt2}\begin{pmatrix}1\\0\\0\\1\end{pmatrix}.
\]

\subsubsection{Les quatre entrées}

Prenons $B$ dans l'état $\ket{x}$ et $A$ dans l'état $\ket{y}$, avec $x,y\in\{0,1\}$. Alors
$\ket{\Pi_0}=\ket{y}_A\otimes\ket{x}_B=\ket{y,x}$. Comme
$H\ket{x}=\frac{1}{\sqrt2}\big(\ket{0}+(-1)^x\ket{1}\big)$ (c'est $\ket{+}$ pour $x=0$ et
$\ket{-}$ pour $x=1$),
\[
  \ket{\Pi_1}=\ket{y}\otimes H\ket{x}=\frac{1}{\sqrt2}\Big(\ket{y,0}+(-1)^x\ket{y,1}\Big).
\]
Le CNOT (contrôle $B$) ne change pas $\ket{y,0}$ et transforme $\ket{y,1}$ en
$\ket{\bar y,1}$, avec $\bar y=1-y$ :
\[
  \ket{\Pi_2}=\ket{\beta^{xy}}=\frac{1}{\sqrt2}\Big(\ket{y,0}+(-1)^x\ket{\bar y,1}\Big).
\]
Le premier indice de $\ket{\beta^{xy}}$ est l'entrée de $B$ (le qubit qui reçoit $H$), le
second celle de $A$ (la cible). Les quatre cas :

\begin{center}
\renewcommand{\arraystretch}{1.7}
\begin{tabular}{@{}ccccc@{}}
\toprule
$(B,A)=(x,y)$ & $\ket{\Pi_0}=\ket{y,x}$ & $\ket{\Pi_1}$ & $\ket{\Pi_2}=\ket{\beta^{xy}}$ & Nom \\
\midrule
$(0,0)$ & $\ket{00}$ & $\frac{1}{\sqrt2}(\ket{00}+\ket{01})$ & $\frac{1}{\sqrt2}(\ket{00}+\ket{11})$ & $\ket{\Phi^+}$ \\
$(1,0)$ & $\ket{01}$ & $\frac{1}{\sqrt2}(\ket{00}-\ket{01})$ & $\frac{1}{\sqrt2}(\ket{00}-\ket{11})$ & $\ket{\Phi^-}$ \\
$(0,1)$ & $\ket{10}$ & $\frac{1}{\sqrt2}(\ket{10}+\ket{11})$ & $\frac{1}{\sqrt2}(\ket{10}+\ket{01})$ & $\ket{\Psi^+}$ \\
$(1,1)$ & $\ket{11}$ & $\frac{1}{\sqrt2}(\ket{10}-\ket{11})$ & $\frac{1}{\sqrt2}(\ket{10}-\ket{01})$ & $-\ket{\Psi^-}$ \\
\bottomrule
\end{tabular}
\end{center}
Détail de la dernière ligne : $\ket{\Pi_1}=\ket{1}\otimes\ket{-}=\frac{1}{\sqrt2}(\ket{10}-\ket{11})$ ;
le CNOT laisse $\ket{10}$ et envoie $\ket{11}$ sur $\ket{01}$, d'où
$\frac{1}{\sqrt2}(\ket{10}-\ket{01})=-\frac{1}{\sqrt2}(\ket{01}-\ket{10})=-\ket{\Psi^-}$.

\paragraph{Forme matricielle.} Le CNOT échange les lignes $2$ et $4$ de
$I\otimes H$ (il échange les amplitudes de $\ket{01}$ et $\ket{11}$), donc
\[
  \mathrm{CNOT}\,(I\otimes H)=\frac{1}{\sqrt2}\begin{pmatrix}
    1 & 1 & 0 & 0\\ 0 & 0 & 1 & -1\\ 0 & 0 & 1 & 1\\ 1 & -1 & 0 & 0
  \end{pmatrix}.
\]
Ses colonnes sont les images de $\ket{00},\ket{01},\ket{10},\ket{11}$, c'est-à-dire de
$(B,A)=(0,0),(1,0),(0,1),(1,1)$ : ce sont $\ket{\beta^{00}}$, $\ket{\beta^{10}}$,
$\ket{\beta^{01}}$, $\ket{\beta^{11}}$, soit $\ket{\Phi^+}$, $\ket{\Phi^-}$, $\ket{\Psi^+}$
et $-\ket{\Psi^-}$.

\subsubsection{Matrice de densité de l'état de sortie}

Pour $\ket{\Psi}=\ket{\beta^{00}}=\frac{1}{\sqrt2}\ket{00}+\frac{1}{\sqrt2}\ket{11}$ :
\[
  \rho_{AB}=\ket{\Psi}\bra{\Psi}
  =\begin{pmatrix}\frac{1}{\sqrt2}\\0\\0\\\frac{1}{\sqrt2}\end{pmatrix}
   \begin{pmatrix}\frac{1}{\sqrt2}&0&0&\frac{1}{\sqrt2}\end{pmatrix}
  =\frac12\begin{pmatrix}1&0&0&1\\0&0&0&0\\0&0&0&0\\1&0&0&1\end{pmatrix}.
\]
Sous forme de ket-bra :
$\rho_{AB}=\frac12\big(\ket{00}\bra{00}+\ket{00}\bra{11}+\ket{11}\bra{00}+\ket{11}\bra{11}\big)$.
On vérifie $\operatorname{Tr}\rho_{AB}=\frac12(1+1)=1$ et
\[
  \rho_{AB}^2=\frac14\begin{pmatrix}2&0&0&2\\0&0&0&0\\0&0&0&0\\2&0&0&2\end{pmatrix}=\rho_{AB},
\]
donc $\operatorname{Tr}\rho_{AB}^2=1$ : le couple $(A,B)$ est dans un état \emph{pur}.

\begin{remarque}[Porte intriquante]
Les états d'entrée $\ket{ij}$ sont séparables ($D=0$) ; les sorties sont intriquées
($D=\pm\frac12$). Aucune porte agissant séparément sur $A$ et sur $B$ ($U_A\otimes U_B$,
section~\ref{sec:op2}) ne peut créer de l'intrication à partir d'un produit, puisque
$(U_A\otimes U_B)(\ket{u}\otimes\ket{v})=U_A\ket{u}\otimes U_B\ket{v}$ reste un produit.
Une porte à deux qubits comme le CNOT est indispensable.
\end{remarque}

% ---------------------------------------------------------------------
\subsection{Trace partielle : l'état d'un qubit pris seul}\label{sec:trace}
% ---------------------------------------------------------------------

Un état intriqué n'a pas de ket pour $A$ seul (section~\ref{sec:separable}). On peut pourtant
décrire $A$ par une matrice de densité, obtenue en «~oubliant~» $B$.

\begin{definition}{Trace partielle}
Soit $\rho_{AB}$ la matrice de densité du couple $(A,B)$. L'état du qubit $A$ pris seul est
$\rho_A=\operatorname{Tr}_B(\rho_{AB})$, celui de $B$ est $\rho_B=\operatorname{Tr}_A(\rho_{AB})$. La
trace partielle est linéaire et remplace la partie tracée de chaque terme par sa trace :
\[
  \operatorname{Tr}_B\big(\ket{a}\bra{a'}\otimes\ket{b}\bra{b'}\big)=\ket{a}\bra{a'}\;\braket{b'}{b},
  \qquad
  \operatorname{Tr}_A\big(\ket{a}\bra{a'}\otimes\ket{b}\bra{b'}\big)=\braket{a'}{a}\;\ket{b}\bra{b'}.
\]
\end{definition}

\paragraph{En matrices.} On découpe $\rho_{AB}$ (matrice $4\times4$, lignes et colonnes ordonnées
$00,01,10,11$) en quatre blocs $2\times2$ :
\[
  \rho_{AB}=\begin{pmatrix}R_{00}&R_{01}\\R_{10}&R_{11}\end{pmatrix}
  \quad\Longrightarrow\quad
  \rho_A=\begin{pmatrix}\operatorname{Tr}R_{00}&\operatorname{Tr}R_{01}\\
                        \operatorname{Tr}R_{10}&\operatorname{Tr}R_{11}\end{pmatrix},
  \qquad
  \rho_B=R_{00}+R_{11}.
\]
Autrement dit $(\rho_A)_{ij}=\rho_{(i0),(j0)}+\rho_{(i1),(j1)}$ : $\rho_A$ est la matrice des traces
des blocs, et $\rho_B$ la somme des deux blocs diagonaux.

\begin{remarque}[Pourquoi cette définition]
Toutes les prédictions qui portent sur $A$ seul viennent d'opérateurs de la forme $O_A\otimes I$, et
\[
  \operatorname{Tr}\big(\rho_{AB}\,(O_A\otimes I)\big)=\operatorname{Tr}\big(\rho_A\,O_A\big)
\]
(on le vérifie sur $\ket{a}\bra{a'}\otimes\ket{b}\bra{b'}$, puis par linéarité). La matrice $\rho_A$
contient donc exactement ce que l'on peut prédire sur $A$ sans regarder $B$.
\end{remarque}

\paragraph{Exemple 1 : état séparable $\ket{0}\otimes\ket{+}$.} Ici
$\rho_{AB}=\ket{0}\bra{0}\otimes\ket{+}\bra{+}$, donc
\[
  \rho_A=\ket{0}\bra{0}\,\braket{+}{+}=\ket{0}\bra{0},
  \qquad
  \rho_B=\braket{0}{0}\,\ket{+}\bra{+}=\ket{+}\bra{+}.
\]
Avec les matrices,
\[
  \rho_{AB}=\frac12\begin{pmatrix}1&1&0&0\\1&1&0&0\\0&0&0&0\\0&0&0&0\end{pmatrix},
  \qquad
  R_{00}=\frac12\begin{pmatrix}1&1\\1&1\end{pmatrix},\quad R_{01}=R_{10}=R_{11}=0,
\]
d'où
\[
  \rho_A=\begin{pmatrix}\operatorname{Tr}R_{00}&0\\0&0\end{pmatrix}=\begin{pmatrix}1&0\\0&0\end{pmatrix},
  \qquad
  \rho_B=R_{00}+R_{11}=\frac12\begin{pmatrix}1&1\\1&1\end{pmatrix}.
\]
Un qubit d'un état séparable garde son propre état pur.

\paragraph{Exemple 2 : état de Bell $\ket{\Phi^+}$.} On écrit
$\rho_{AB}=\ket{\Phi^+}\bra{\Phi^+}=\frac12\big(\ket{00}\bra{00}+\ket{00}\bra{11}+\ket{11}\bra{00}
+\ket{11}\bra{11}\big)$. Chaque terme est un produit d'opérateurs à un qubit, par exemple
$\ket{00}\bra{11}=\ket{0}\bra{1}\otimes\ket{0}\bra{1}$, donc
\begin{align*}
  \rho_A=\operatorname{Tr}_B(\rho_{AB})
  &=\tfrac12\Big(\ket{0}\bra{0}\,\braket{0}{0}+\ket{0}\bra{1}\,\braket{1}{0}
      +\ket{1}\bra{0}\,\braket{0}{1}+\ket{1}\bra{1}\,\braket{1}{1}\Big)\\
  &=\tfrac12\Big(\ket{0}\bra{0}\cdot1+\ket{0}\bra{1}\cdot0+\ket{1}\bra{0}\cdot0+\ket{1}\bra{1}\cdot1\Big)
   =\tfrac12\big(\ket{0}\bra{0}+\ket{1}\bra{1}\big)=\tfrac12\,\mathrm{Id}.
\end{align*}
Par les blocs,
\[
  \rho_{AB}=\frac12\begin{pmatrix}1&0&0&1\\0&0&0&0\\0&0&0&0\\1&0&0&1\end{pmatrix},
  \qquad
  R_{00}=\frac12\begin{pmatrix}1&0\\0&0\end{pmatrix},\quad
  R_{11}=\frac12\begin{pmatrix}0&0\\0&1\end{pmatrix},\quad
  \operatorname{Tr}R_{01}=\operatorname{Tr}R_{10}=0,
\]
d'où $\rho_A=\frac12\mathrm{Id}$ et $\rho_B=R_{00}+R_{11}=\frac12\mathrm{Id}$ : le même résultat.

Le qubit $A$ est donc dans l'état \emph{complètement mixte} ($\operatorname{Tr}\rho_A^2=\frac14+\frac14=\frac12$),
au centre de la boule de Bloch (section~\ref{sec:bloch}), alors que le couple est dans un état pur
($\operatorname{Tr}\rho_{AB}^2=1$, section~\ref{sec:circuit-bell}). On connaît parfaitement l'état du
couple mais pas celui d'un qubit pris seul : l'information est entièrement contenue dans les
corrélations entre $A$ et $B$. Mesuré seul, $A$ donne $0$ ou $1$ avec la probabilité $\frac12$, dans
toutes les bases.

\paragraph{Formule pratique pour un état pur.} Pour $\ket{\Psi}=\sum_{ij}c_{ij}\ket{ij}$, notons
$C=\begin{pmatrix}c_{00}&c_{01}\\c_{10}&c_{11}\end{pmatrix}$ la matrice $2\times2$ des coefficients
(ligne $i$ : état de $A$, colonne $j$ : état de $B$). Comme $\rho_{(ij),(i'j')}=c_{ij}\,c_{i'j'}^*$,
\[
  (\rho_A)_{ii'}=\sum_jc_{ij}\,c_{i'j}^*,\qquad\text{soit}\qquad \rho_A=C\,C^\dagger .
\]
\emph{Exemple 3.} Pour $\ket{\Psi}=\frac12\ket{00}+\frac12\ket{01}+\frac{1}{\sqrt2}\ket{11}$ (exemple~2
de la section~\ref{sec:mesure2}), $C=\begin{pmatrix}\frac12&\frac12\\0&\frac{1}{\sqrt2}\end{pmatrix}$ et
\[
  \rho_A=CC^\dagger=\begin{pmatrix}\frac14+\frac14&0+\frac{1}{2\sqrt2}\\[0.4ex]\frac{1}{2\sqrt2}&\frac12\end{pmatrix}
  =\begin{pmatrix}\frac12&\frac{\sqrt2}{4}\\[0.4ex]\frac{\sqrt2}{4}&\frac12\end{pmatrix},
  \qquad
  \operatorname{Tr}\rho_A^2=\tfrac14+\tfrac14+2\cdot\tfrac18=\tfrac34 .
\]
Le qubit $A$ est mixte, mais pas complètement : cet état est intriqué sans être un état de Bell.

\begin{resultat}{Pureté d'un qubit et séparabilité}
Pour un état pur $\ket{\Psi}=\sum c_{ij}\ket{ij}$ de deux qubits, avec $D=c_{00}c_{11}-c_{01}c_{10}$,
\[
  \operatorname{Tr}\rho_A^2=\operatorname{Tr}\rho_B^2=1-2|D|^2 .
\]
Donc $\ket{\Psi}$ est séparable ($D=0$) si et seulement si $\rho_A$ est pur. Plus $|D|$ est grand,
plus $\rho_A$ est mixte, jusqu'à $|D|=\frac12$ et $\rho_A=\mathrm{Id}/2$ pour les états de Bell.
\end{resultat}

\emph{Démonstration.} La matrice $\rho_A=CC^\dagger$ a pour trace $\sum|c_{ij}|^2=1$ et pour
déterminant $\det C\,\overline{\det C}=|D|^2$. Si $\lambda_1,\lambda_2$ sont ses valeurs propres,
$\lambda_1+\lambda_2=1$ et $\operatorname{Tr}\rho_A^2=\lambda_1^2+\lambda_2^2=1-2\lambda_1\lambda_2
=1-2|D|^2$ ; le calcul est le même pour $\rho_B$. \hfill$\square$

Comme $\operatorname{Tr}\rho_A^2\geq\frac12$ pour un qubit, on en déduit $|D|\leq\frac12$. Avec
$\operatorname{Tr}\rho_A^2=\frac{1+|\vec r|^2}{2}$ (section~\ref{sec:bloch}), le vecteur de Bloch de $A$
a pour norme $|\vec r|=\sqrt{1-4|D|^2}$ : il part de la sphère pour un état séparable et se rapproche
du centre quand l'intrication croît (figure~\ref{fig:trace-bloch}).

\begin{center}
\renewcommand{\arraystretch}{1.5}
\begin{tabular}{@{}lcccl@{}}
\toprule
\textbf{État du couple} & $D$ & $\operatorname{Tr}\rho_A^2$ & $|\vec r|$ & \textbf{Qubit $A$} \\
\midrule
$\ket{0}\otimes\ket{+}$ & $0$ & $1$ & $1$ & pur, sur la sphère \\
$\frac12\ket{00}+\frac12\ket{01}+\frac{1}{\sqrt2}\ket{11}$ & $\frac{\sqrt2}{4}$ & $\frac34$
  & $\frac{1}{\sqrt2}$ & mixte, dans la boule \\
$\ket{\Phi^+}$ (comme les trois autres états de Bell) & $\pm\frac12$ & $\frac12$ & $0$
  & complètement mixte, au centre \\
\bottomrule
\end{tabular}
\end{center}

\begin{figure}[!ht]
\centering
\adjustbox{max width=\linewidth}{%
\begin{tikzpicture}[>=Stealth, font=\small]
  \def\R{2.3}
  \fill[insagris!12] (0,0) circle (\R);
  \draw[insagris, dotted] (-\R,0) -- (\R,0);
  \draw[insagris, dotted] (0,-\R) -- (0,\R);
  \draw[insagris, dashed] (0,0) circle ({\R*0.70711});
  \draw[insarouge, very thick] (0,0) circle (\R);
  \node[font=\footnotesize, text=insagris, anchor=south west] at (\R,0) {$x$};
  \node[font=\footnotesize, text=insagris, anchor=south west] at (0.03,\R) {$z$};
  % vecteurs de Bloch de rho_A pour trois états du couple
  \fill[insarouge] (0,\R) circle (2.4pt);
  \fill[insagris!80!black] ({\R*0.70711},0) circle (2.4pt);
  \fill[insagris!80!black] (0,0) circle (2.4pt);
  \draw[insagris] (0,\R) -- (3.1,\R);
  \node[anchor=west, font=\footnotesize, align=left] at (3.1,\R)
    {$\ket{0}\otimes\ket{+}$ : $\rho_A=\ket{0}\bra{0}$\\ $D=0$, \ $|\vec r|=1$ (état pur)};
  \draw[insagris] ({\R*0.70711},0) -- (3.1,0);
  \node[anchor=west, font=\footnotesize, align=left] at (3.1,0)
    {$\frac12\ket{00}+\frac12\ket{01}+\frac{1}{\sqrt2}\ket{11}$\\
     $D=\frac{\sqrt2}{4}$, \ $|\vec r|=\frac{1}{\sqrt2}$ (mixte)};
  \draw[insagris] (0,0) -- (3.1,-\R);
  \node[anchor=west, font=\footnotesize, align=left] at (3.1,-\R)
    {$\ket{\Phi^+}$ : $\rho_A=\mathrm{Id}/2$\\ $D=\frac12$, \ $\vec r=\vec 0$ (centre)};
\end{tikzpicture}}
\caption{Le vecteur de Bloch de $\rho_A$ (boule de Bloch vue en coupe) pour trois états du couple
$(A,B)$. Le cercle pointillé est la sphère de rayon $1/\sqrt2$. Plus l'état est intriqué, plus le
point s'éloigne de la sphère des états purs : $|\vec r|=\sqrt{1-4|D|^2}$.}
\label{fig:trace-bloch}
\end{figure}

% ---------------------------------------------------------------------
\subsection{Information : variables aléatoires et états quantiques}\label{sec:info}
% ---------------------------------------------------------------------

\begin{definition}{Entropie, information mutuelle, entropie conditionnelle}
Soit $A$ une variable aléatoire de loi $p(a)$. Son \emph{entropie de Shannon}, en bits, est
\[
  H(A)=-\sum_a p(a)\log_2p(a)\qquad(\text{avec }0\log_20=0).
\]
Elle mesure l'incertitude moyenne sur $A$ : une pièce équilibrée a $H=1$ bit. Pour un couple
$(A,B)$ de loi $p(a,b)$, $H(A,B)=-\sum_{a,b}p(a,b)\log_2p(a,b)$, et
\begin{align*}
  I(A;B)&=H(A)+H(B)-H(A,B) &&\text{(information mutuelle)},\\
  H(A|B)&=H(A,B)-H(B)=H(A)-I(A;B) &&\text{(entropie conditionnelle)}.
\end{align*}
\end{definition}

Les deux écritures de $H(A|B)$ sont égales : la définition de $I(A;B)$ donne
$H(A,B)=H(A)+H(B)-I(A;B)$, donc $H(A,B)-H(B)=H(A)-I(A;B)$. L'information mutuelle $I(A;B)$
mesure ce que la connaissance de $B$ apprend sur $A$ ; $H(A|B)$ est l'incertitude qui reste sur
$A$ quand on connaît $B$.

\subsubsection{Deux variables aléatoires indépendantes}

Alice et Bob lancent chacun une pièce équilibrée, sans lien entre les deux lancers :
$p(a,b)=\frac14$ pour les quatre couples. Chaque pièce donne $0$ ou $1$ avec la probabilité $\frac12$, donc
\begin{align*}
  H(A)&=-\Big(\tfrac12\log_2\tfrac12+\tfrac12\log_2\tfrac12\Big)=1\text{ bit},\\
  H(B)&=1\text{ bit (même calcul)},\\
  H(A,B)&=-4\cdot\tfrac14\log_2\tfrac14=2\text{ bits},\\
  I(A;B)&=H(A)+H(B)-H(A,B)=1+1-2=0,\\
  H(A|B)&=H(A)-I(A;B)=1-0=H(A)=1\text{ bit}.
\end{align*}
Connaître $B$ n'apprend rien sur $A$.

\subsubsection{Deux variables solidaires}

Les deux pièces tombent toujours du même côté : $p(0,0)=p(1,1)=\frac12$ et
$p(0,1)=p(1,0)=0$. Les lois de chaque pièce sont inchangées ($p(a{=}0)=p(0,0)+p(0,1)=\frac12$,
etc.), donc $H(A)=H(B)=1$ bit, mais
\begin{align*}
  H(A,B)&=-2\cdot\tfrac12\log_2\tfrac12-2\cdot0\log_20=1\text{ bit},\\
  I(A;B)&=1+1-1=1\text{ bit},\qquad H(A|B)=1-1=0.
\end{align*}
Connaître $B$ détermine $A$ : plus aucune incertitude ne reste.

\subsubsection{Deux qubits dans un état de Bell}

Pour un système quantique, on remplace la loi de probabilité par la matrice de densité.

\begin{definition}{Entropie de von Neumann}
L'\emph{entropie de von Neumann} d'un état $\rho$, de valeurs propres $\lambda_i$, est
\[
  S(\rho)=-\operatorname{Tr}(\rho\log_2\rho)=-\sum_i\lambda_i\log_2\lambda_i .
\]
Elle vaut $0$ pour un état pur et $1$ bit pour le qubit complètement mélangé $\mathrm{Id}/2$.
En quantique, $H(A)=S(\rho_A)$, $H(B)=S(\rho_B)$ et $H(A,B)=S(\rho_{AB})$ ; les définitions de
$I(A;B)$ et de $H(A|B)$ sont les mêmes.
\end{definition}

On prend l'état de Bell $\ket{\Phi^+}$ produit par le circuit de la
section~\ref{sec:circuit-bell}. Le couple est dans un état pur,
$\rho_{AB}=\ket{\Phi^+}\bra{\Phi^+}$ ; chaque qubit pris seul est dans l'état
$\rho_A=\rho_B=\frac12\mathrm{Id}$ (trace partielle, section~\ref{sec:trace}).

Calcul des entropies :
\begin{itemize}[nosep]
  \item $\rho_{AB}$ est pur : ses valeurs propres sont $1,0,0,0$, donc
        $H(A,B)=S(\rho_{AB})=-1\cdot\log_21=0$ ;
  \item $\rho_A$ et $\rho_B$ ont pour valeurs propres $\frac12,\frac12$, donc
        $H(A)=H(B)=-\big(\frac12\log_2\frac12+\frac12\log_2\frac12\big)=1$ bit.
\end{itemize}
D'où
\begin{align*}
  I(A;B)&=H(A)+H(B)-H(A,B)=1+1-0=2\text{ bits},\\
  H(A|B)&=H(A,B)-H(B)=0-1=-1\text{ bit}\qquad(=H(A)-I(A;B)=1-2).
\end{align*}

\begin{remarque}[Une entropie conditionnelle négative]
Pour des variables classiques, $p(a,b)\leq p(b)$, donc $-\log_2p(a,b)\geq-\log_2p(b)$ et
$H(A,B)\geq H(B)$ : on a toujours $H(A|B)\geq0$, et $I(A;B)\leq H(A)$. Ici $H(A|B)=-1$ et
$I(A;B)=2>H(A)$ : ces valeurs n'ont pas d'équivalent classique, c'est une signature de
l'intrication. Le résultat est surprenant mais mathématiquement correct.
\end{remarque}

\begin{center}
\renewcommand{\arraystretch}{1.4}
\begin{tabular}{@{}lccc@{}}
\toprule
 & indépendantes & solidaires & état de Bell \\
\midrule
$H(A)$ & $1$ & $1$ & $1$ \\
$H(B)$ & $1$ & $1$ & $1$ \\
$H(A,B)$ & $2$ & $1$ & $0$ \\
$I(A;B)$ & $0$ & $1$ & $2$ \\
$H(A|B)$ & $1$ & $0$ & $-1$ \\
\bottomrule
\end{tabular}
\end{center}
(en bits ; les deux premières colonnes sont classiques, la troisième est quantique).

\begin{figure}[!ht]
\centering
\adjustbox{max width=\linewidth}{%
\begin{tikzpicture}[>=Stealth, font=\small]
  \def\bw{0.8}
  % axe des ordonnées (bits)
  \draw[->] (0,-1.5) -- (0,2.8) node[above] {bits};
  \foreach \v in {-1,0,1,2} {
    \draw (-0.07,\v) -- (0.07,\v);
    \node[left, font=\footnotesize] at (-0.1,\v) {$\v$};
  }
  \draw[insagris!60] (0,0) -- (11.4,0);
  % limite classique de I(A;B) : H(A) = 1
  \draw[dashed, insarouge] (0,1) -- (11.4,1);
  \node[anchor=west, font=\scriptsize, text=insarouge, align=left] at (11.4,1)
    {limite classique\\ $I(A;B)\leq H(A)=1$};
  \node[anchor=west, font=\scriptsize, text=insagris, align=left] at (11.4,0)
    {limite classique\\ $H(A|B)\geq0$};
  % barres : (groupe, H(A,B), I(A;B), H(A|B))
  \foreach \g/\hab/\mi/\hc in {0/2/0/1, 1/1/1/0, 2/0/2/-1} {
    \pgfmathsetmacro{\x}{0.9+3.7*\g}
    \fill[insagris!65] (\x,0) rectangle (\x+\bw,\hab);
    \fill[insarouge] (\x+\bw+0.1,0) rectangle (\x+2*\bw+0.1,\mi);
    \fill[black!80] (\x+2*\bw+0.2,0) rectangle (\x+3*\bw+0.2,\hc);
    \node[above, font=\footnotesize] at (\x+0.5*\bw,{max(\hab,0)}) {$\hab$};
    \node[above, font=\footnotesize] at (\x+1.5*\bw+0.1,{max(\mi,0)}) {$\mi$};
    \ifnum\hc<0
      \node[below, font=\footnotesize] at (\x+2.5*\bw+0.2,\hc) {$\hc$};
    \else
      \node[above, font=\footnotesize] at (\x+2.5*\bw+0.2,\hc) {$\hc$};
    \fi
  }
  \node[align=center, anchor=north] at (2.15,-1.65) {variables\\ indépendantes};
  \node[align=center, anchor=north] at (5.85,-1.65) {variables\\ solidaires};
  \node[align=center, anchor=north] at (9.55,-1.65) {état de Bell\\ $\ket{\Phi^+}$};
  % légende
  \fill[insagris!65] (0.9,3.35) rectangle (1.2,3.65);
  \node[anchor=west, font=\footnotesize] at (1.25,3.5) {$H(A,B)$};
  \fill[insarouge] (3.6,3.35) rectangle (3.9,3.65);
  \node[anchor=west, font=\footnotesize] at (3.95,3.5) {$I(A;B)$};
  \fill[black!80] (6.3,3.35) rectangle (6.6,3.65);
  \node[anchor=west, font=\footnotesize] at (6.65,3.5) {$H(A|B)$};
\end{tikzpicture}}
\caption{Entropies (en bits) pour deux pièces indépendantes, deux pièces solidaires et un état de Bell.
Dans les trois cas $H(A)=H(B)=1$. Les deux premiers groupes respectent les bornes classiques
$I(A;B)\leq H(A)$ et $H(A|B)\geq0$ ; l'état de Bell les dépasse ($I=2$, $H(A|B)=-1$).}
\label{fig:entropies}
\end{figure}

% ---------------------------------------------------------------------
\subsection{Téléportation quantique}\label{sec:teleportation}
% ---------------------------------------------------------------------

\emph{Problème.} Alice détient un qubit $Q$ dans un état qu'elle ne connaît pas,
$\ket{\varphi}=\alpha\ket{0}+\beta\ket{1}$ avec $|\alpha|^2+|\beta|^2=1$, et elle veut que Bob
obtienne cet état sur l'un de ses qubits. Elle ne peut pas mesurer $Q$ pour connaître
$\alpha$ et $\beta$ (une mesure ne donne qu'un bit et détruit l'état), elle ne peut pas le
copier (théorème de non-clonage, remarque en fin de section) et elle n'envoie pas le qubit
lui-même. Le protocole est dû à \textcite{bennett1993}.

\begin{definition}{Protocole de téléportation quantique}
\textbf{Ressources.}
\begin{itemize}[nosep]
  \item Un \emph{e-bit} : Alice et Bob partagent une paire de qubits $(A,B)$ dans l'état de Bell
        $\ket{\Phi^+}_{BA}=\frac{1}{\sqrt2}\big(\ket{00}+\ket{11}\big)$ ; Alice détient $A$,
        Bob détient $B$.
  \item Le qubit $Q$ à téléporter, chez Alice.
  \item Un canal classique (deux bits).
\end{itemize}
Le système est écrit $(B,A,Q)$ : $\ket{b,a,q}=\ket{b}_B\otimes\ket{a}_A\otimes\ket{q}_Q$.

\textbf{Étapes.}
\begin{enumerate}[nosep]
  \item Alice applique un CNOT de contrôle $Q$ et de cible $A$.
  \item Alice applique une porte $H$ sur $Q$.
  \item Alice mesure $A$ et $Q$ dans la base de calcul, obtient deux bits $(a,q)$ et les envoie
        à Bob.
  \item Bob applique $X$ si $a=1$, puis $Z$ si $q=1$, à son qubit $B$.
\end{enumerate}
\end{definition}

La figure~\ref{fig:tele-schema} montre l'organisation du protocole et la
figure~\ref{fig:teleportation} son circuit.

\begin{figure}[!ht]
\centering
\adjustbox{max width=\linewidth}{%
\begin{tikzpicture}[>=Stealth, font=\small,
    qb/.style={circle, draw=insarouge, fill=insarouge!8, thick, minimum size=8mm, inner sep=0pt},
    zone/.style={draw=insagris, rounded corners=4pt, thick, dashed, inner sep=8pt},
    num/.style={circle, fill=insarouge, text=white, font=\footnotesize\bfseries, inner sep=0pt,
                minimum size=4.6mm},
    etape/.style={anchor=west, font=\footnotesize, text=insarouge},
    onde/.style={thick, insarouge, decorate,
                 decoration={snake, amplitude=1.3pt, segment length=7pt}}]
  % Alice : Q (état inconnu) et A
  \node[qb] (Q) at (0,1.3) {$Q$};
  \node[qb] (A) at (0,-0.5) {$A$};
  \node[zone, fit=(Q)(A), label={[text=insagris]above:Alice}] (alice) {};
  \node[anchor=east] at (-0.95,1.3) {$\ket{\varphi}=\alpha\ket{0}+\beta\ket{1}$};
  \draw[insarouge, thick] (Q) -- (A);
  % Bob : B
  \node[qb] (B) at (9,-0.5) {$B$};
  \node[zone, fit=(B), label={[text=insagris]above:Bob}] (bob) {};
  % source de la paire de Bell
  \node[draw=insarouge, fill=insarouge!10, rounded corners=2pt, thick, minimum width=3.2cm,
        minimum height=0.9cm] (src) at (4.5,-2.3) {source $\ket{\Phi^+}_{BA}$};
  \draw[onde] (src.west) -| (A.south);
  \draw[onde] (src.east) -| (B.south);
  % étapes d'Alice : pastilles sur le bord droit de sa zone
  \node[num] (n1) at (alice.east |- Q) {1};
  \node[num] (n2) at (alice.east) {2};
  \node[num] (n3) at (alice.east |- A) {3};
  \node[etape] at (n1.east) {\ CNOT de $Q$ vers $A$};
  \node[etape] at (n2.east) {\ porte $H$ sur $Q$};
  % canal classique
  \draw[double, double distance=1.2pt, ->, insagris] (n3.east) -- (bob.west);
  \node[above, font=\footnotesize] at (4.6,-0.45) {deux bits classiques $(a,q)$};
  \node[below, font=\footnotesize, text=insarouge] at (4.6,-0.55) {mesure de $A$ et $Q$ (étape 3)};
  % correction de Bob
  \node[num] (n4) at (bob.east) {4};
  \node[etape, align=left] at (n4.east) {\ $X$ si $a=1$, puis $Z$ si $q=1$};
  \node[anchor=west, font=\footnotesize, align=left] at ($(n4.east)+(0.05,-0.8)$)
    {\ $B$ est alors dans l'état $\ket{\varphi}$};
\end{tikzpicture}}
\caption{Organisation de la téléportation. Alice possède le qubit $Q$ à téléporter et la moitié $A$
d'une paire de Bell, Bob possède l'autre moitié $B$. (1)~CNOT de $Q$ vers $A$, (2)~porte $H$ sur $Q$,
(3)~mesure de $A$ et $Q$, dont les résultats $(a,q)$ partent par un canal classique, (4)~corrections de
Bob. Aucun qubit ne voyage entre Alice et Bob.}
\label{fig:tele-schema}
\end{figure}

\begin{figure}[!ht]
\centering
\begin{adjustbox}{max width=\linewidth}
\begin{tikzpicture}[>=Stealth, font=\small]
  % fils quantiques : Q en haut, A au milieu, B en bas
  \draw (0,2.4) -- (5.3,2.4);
  \draw (0,1.2) -- (5.3,1.2);
  \draw (0,0) -- (10.6,0);
  \node[anchor=east] at (0,2.4) {$\ket{\varphi}_Q$};
  \node[anchor=south west, text=insagris, font=\footnotesize] at (0.05,1.2) {$A$};
  \node[anchor=south west, text=insagris, font=\footnotesize] at (0.05,0) {$B$};
  % paire de Bell partagée
  \draw[insagris, thick] (-0.1,1.5) -- (-0.25,1.5) -- (-0.25,-0.3) -- (-0.1,-0.3);
  \node[anchor=east] at (-0.35,0.6) {$\ket{\Phi^+}_{BA}$};
  % CNOT : contrôle Q, cible A
  \draw[insarouge, thick] (2.6,2.4) -- (2.6,1.2);
  \pic at (2.6,2.4) {ctl};
  \pic at (2.6,1.2) {cible};
  % Hadamard sur Q
  \node[porte] at (4.2,2.4) {$H$};
  % mesures d'Alice
  \pic at (5.6,2.4) {mesure};
  \pic at (5.6,1.2) {mesure};
  % fils classiques (doubles)
  \draw[double, double distance=1.2pt, insagris] (5.9,1.2) -- (7.6,1.2) -- (7.6,0.35);
  \draw[double, double distance=1.2pt, insagris] (5.9,2.4) -- (9.6,2.4) -- (9.6,0.35);
  \node[above, font=\footnotesize] at (6.5,1.2) {$a$};
  \node[above, font=\footnotesize] at (6.5,2.4) {$q$};
  % corrections de Bob
  \node[porte] at (7.6,0) {$X$};
  \node[porte] at (9.6,0) {$Z$};
  \node[anchor=west] at (10.6,0) {$\ket{\varphi}_B$};
  % états intermédiaires
  \foreach \x in {1.4,3.4,4.9}
    \draw[dashed, insagris!70] (\x,2.8) -- (\x,-0.5);
  \node[anchor=north, font=\footnotesize] at (1.4,-0.5) {$\ket{\Pi_0}$};
  \node[anchor=north, font=\footnotesize] at (3.4,-0.5) {$\ket{\Pi_1}$};
  \node[anchor=north, font=\footnotesize] at (4.9,-0.5) {$\ket{\Pi_2}$};
\end{tikzpicture}
\end{adjustbox}
\caption{Circuit de téléportation. Alice détient $Q$ (dans $\ket{\varphi}=\alpha\ket{0}+\beta\ket{1}$)
et $A$, Bob détient $B$ ; $A$ et $B$ forment la paire $\ket{\Phi^+}_{BA}$. Les traits doubles
sont des fils \emph{classiques} qui portent les résultats $a$ (mesure de $A$) et $q$ (mesure de
$Q$) ; $X$ est appliqué si $a=1$, $Z$ si $q=1$. Le ket s'écrit $\ket{b,a,q}$ : le fil du haut
($Q$) est le dernier chiffre.}
\label{fig:teleportation}
\end{figure}

\subsubsection{Calcul pas à pas}

\paragraph{État initial $\ket{\Pi_0}$.} Alice et Bob partagent la paire, et $Q$ est indépendant
d'elle :
\begin{align*}
  \ket{\Pi_0}&=\ket{\Phi^+}_{BA}\otimes\ket{\varphi}_Q
    =\frac{1}{\sqrt2}\big(\ket{00}+\ket{11}\big)\otimes\big(\alpha\ket{0}+\beta\ket{1}\big)\\
  &=\frac{1}{\sqrt2}\Big(\alpha\ket{00}\ket{0}+\beta\ket{00}\ket{1}+\alpha\ket{11}\ket{0}+\beta\ket{11}\ket{1}\Big)\\
  &=\frac{\alpha\ket{000}+\beta\ket{001}+\alpha\ket{110}+\beta\ket{111}}{\sqrt2}.
\end{align*}

\paragraph{Après le CNOT $Q\to A$ : $\ket{\Pi_1}$.} Le CNOT (contrôle $Q$, dernier chiffre ;
cible $A$, chiffre du milieu) agit ainsi : $\ket{b,a,q}\to\ket{b,\,a\oplus q,\,q}$. Terme par terme :
\begin{center}
\renewcommand{\arraystretch}{1.3}
\begin{tabular}{@{}cccc@{}}
\toprule
Terme & $q$ & Effet sur $A$ & Résultat \\
\midrule
$\alpha\ket{000}$ & $0$ & aucun & $\alpha\ket{000}$ \\
$\beta\ket{001}$ & $1$ & $0\to1$ & $\beta\ket{011}$ \\
$\alpha\ket{110}$ & $0$ & aucun & $\alpha\ket{110}$ \\
$\beta\ket{111}$ & $1$ & $1\to0$ & $\beta\ket{101}$ \\
\bottomrule
\end{tabular}
\end{center}
donc
\[
  \ket{\Pi_1}=\frac{\alpha\ket{000}+\beta\ket{011}+\alpha\ket{110}+\beta\ket{101}}{\sqrt2}.
\]
On isole ensuite le dernier chiffre, celui de $Q$, qui va changer à l'étape suivante
(le premier ket est celui du couple $\ket{b,a}$) :
\[
  \ket{\Pi_1}=\frac{\alpha\ket{00}\otimes\ket{0}+\beta\ket{01}\otimes\ket{1}
    +\alpha\ket{11}\otimes\ket{0}+\beta\ket{10}\otimes\ket{1}}{\sqrt2}.
\]

\paragraph{Après la porte $H$ sur $Q$ : $\ket{\Pi_2}$.} On remplace $\ket{0}$ par
$H\ket{0}=\ket{+}$ et $\ket{1}$ par $H\ket{1}=\ket{-}$ dans le facteur $Q$ :
\[
  \ket{\Pi_2}=\frac{1}{\sqrt2}\Big(\alpha\ket{00}\otimes\ket{+}+\beta\ket{01}\otimes\ket{-}
    +\alpha\ket{11}\otimes\ket{+}+\beta\ket{10}\otimes\ket{-}\Big).
\]
Avec $\ket{\pm}=\frac{1}{\sqrt2}(\ket{0}\pm\ket{1})$, le préfacteur devient
$\frac{1}{\sqrt2}\cdot\frac{1}{\sqrt2}=\frac12$ :
\begin{align*}
  \ket{\Pi_2}&=\frac12\Big(\alpha\ket{00}\big(\ket{0}+\ket{1}\big)+\beta\ket{01}\big(\ket{0}-\ket{1}\big)
     +\alpha\ket{11}\big(\ket{0}+\ket{1}\big)+\beta\ket{10}\big(\ket{0}-\ket{1}\big)\Big)\\
  &=\frac12\Big(\alpha\ket{000}+\alpha\ket{001}+\beta\ket{010}-\beta\ket{011}
     +\alpha\ket{110}+\alpha\ket{111}+\beta\ket{100}-\beta\ket{101}\Big).
\end{align*}

\paragraph{Regroupement selon le résultat d'Alice.} Alice mesurera $A$ et $Q$, donc les deux
derniers chiffres. On regroupe les huit termes qui ont les mêmes chiffres $(a,q)$, en mettant
le ket de Bob en premier :
\begin{center}
\renewcommand{\arraystretch}{1.3}
\begin{tabular}{@{}cll@{}}
\toprule
$(a,q)$ & Termes de $\ket{\Pi_2}$ (facteur $\frac12$ omis) & Ket de Bob \\
\midrule
$(0,0)$ & $\alpha\ket{000}+\beta\ket{100}$ & $(\alpha\ket{0}+\beta\ket{1})\ket{00}$ \\
$(0,1)$ & $\alpha\ket{001}-\beta\ket{101}$ & $(\alpha\ket{0}-\beta\ket{1})\ket{01}$ \\
$(1,0)$ & $\beta\ket{010}+\alpha\ket{110}$ & $(\beta\ket{0}+\alpha\ket{1})\ket{10}$ \\
$(1,1)$ & $-\beta\ket{011}+\alpha\ket{111}$ & $(-\beta\ket{0}+\alpha\ket{1})\ket{11}$ \\
\bottomrule
\end{tabular}
\end{center}
d'où
\[
  \ket{\Pi_2}=\frac12\Big[\big(\alpha\ket{0}+\beta\ket{1}\big)\ket{00}
    +\big(\alpha\ket{0}-\beta\ket{1}\big)\ket{01}
    +\big(\beta\ket{0}+\alpha\ket{1}\big)\ket{10}
    +\big(\alpha\ket{1}-\beta\ket{0}\big)\ket{11}\Big].
\]
Dans chaque terme, le ket de gauche est l'état de $B$ et le ket de droite $\ket{a,q}$ celui du
couple $(A,Q)$ d'Alice.

\subsubsection{Mesure d'Alice et corrections de Bob}

Alice mesure $A$ et $Q$. Comme à la section~\ref{sec:mesure2}, la probabilité d'un résultat
$(a,q)$ est le carré de la norme du terme correspondant, et l'état de Bob devient le ket
associé (déjà normé). Ici
\[
  P(a,q)=\Big\|\tfrac12\ket{\phi_{aq}}\Big\|^2=\tfrac14\big(|\alpha|^2+|\beta|^2\big)=\tfrac14
  \quad\text{pour les quatre résultats,}
\]
qui sont donc équiprobables et \emph{indépendants} de $\alpha$ et $\beta$ : les deux bits
envoyés ne révèlent rien sur $\ket{\varphi}$. Les kets de Bob sont
$\ket{\phi_{00}}=\alpha\ket{0}+\beta\ket{1}$, $\ket{\phi_{01}}=\alpha\ket{0}-\beta\ket{1}$,
$\ket{\phi_{10}}=\beta\ket{0}+\alpha\ket{1}$ et $\ket{\phi_{11}}=-\beta\ket{0}+\alpha\ket{1}$.

\begin{center}
\renewcommand{\arraystretch}{1.5}
\begin{tabular}{@{}cccll@{}}
\toprule
Résultat $(a,q)$ & Probabilité & État de Bob & Lien avec $\ket{\varphi}$ & Correction de Bob \\
\midrule
$(0,0)$ & $\frac14$ & $\alpha\ket{0}+\beta\ket{1}$ & $\ket{\varphi}$ & $I$ (rien) \\
$(0,1)$ & $\frac14$ & $\alpha\ket{0}-\beta\ket{1}$ & $Z\ket{\varphi}$ & $Z$ \\
$(1,0)$ & $\frac14$ & $\beta\ket{0}+\alpha\ket{1}$ & $X\ket{\varphi}$ & $X$ \\
$(1,1)$ & $\frac14$ & $-\beta\ket{0}+\alpha\ket{1}$ & $XZ\ket{\varphi}$ & $X$ puis $Z$ (opérateur $ZX$) \\
\bottomrule
\end{tabular}
\end{center}

\paragraph{Vérification des quatre corrections.} Avec $X\ket{0}=\ket{1}$, $X\ket{1}=\ket{0}$,
$Z\ket{0}=\ket{0}$ et $Z\ket{1}=-\ket{1}$ :
\begin{align*}
  (0,0):\quad & I\big(\alpha\ket{0}+\beta\ket{1}\big)=\alpha\ket{0}+\beta\ket{1},\\
  (0,1):\quad & Z\big(\alpha\ket{0}-\beta\ket{1}\big)=\alpha\ket{0}-\beta Z\ket{1}=\alpha\ket{0}+\beta\ket{1},\\
  (1,0):\quad & X\big(\beta\ket{0}+\alpha\ket{1}\big)=\beta\ket{1}+\alpha\ket{0}=\alpha\ket{0}+\beta\ket{1},\\
  (1,1):\quad & X\big(-\beta\ket{0}+\alpha\ket{1}\big)=-\beta\ket{1}+\alpha\ket{0}=\alpha\ket{0}-\beta\ket{1},
     \ \text{puis}\ Z\big(\alpha\ket{0}-\beta\ket{1}\big)=\alpha\ket{0}+\beta\ket{1}.
\end{align*}
Dans le dernier cas, avec les matrices,
$ZX=\begin{pmatrix}1&0\\0&-1\end{pmatrix}\begin{pmatrix}0&1\\1&0\end{pmatrix}
=\begin{pmatrix}0&1\\-1&0\end{pmatrix}$ et
$ZX\begin{pmatrix}-\beta\\ \alpha\end{pmatrix}=\begin{pmatrix}\alpha\\ \beta\end{pmatrix}$.
L'ordre compte : l'état de Bob est $XZ\ket{\varphi}$ ($Z\ket{\varphi}=\alpha\ket{0}-\beta\ket{1}$,
puis $X$ donne $\alpha\ket{1}-\beta\ket{0}$), et on l'annule avec
$(XZ)^{-1}=Z^{-1}X^{-1}=ZX$ : on applique d'abord $X$, puis $Z$, comme dans le circuit.
Dans les quatre cas, Bob retrouve exactement $\ket{\varphi}=\alpha\ket{0}+\beta\ket{1}$.

\subsubsection{Ce que Bob sait avant de recevoir les bits}

Tant que Bob ignore $(a,q)$, son qubit est dans le mélange des quatre kets $\ket{\phi_{aq}}$, chacun
avec la probabilité $\frac14$ :
\begin{align*}
  \rho_B&=\frac14\sum_{a,q}\ket{\phi_{aq}}\bra{\phi_{aq}}\\
  &=\frac14\Bigg[\begin{pmatrix}|\alpha|^2&\alpha\beta^*\\ \alpha^*\beta&|\beta|^2\end{pmatrix}
    +\begin{pmatrix}|\alpha|^2&-\alpha\beta^*\\ -\alpha^*\beta&|\beta|^2\end{pmatrix}
    +\begin{pmatrix}|\beta|^2&\alpha^*\beta\\ \alpha\beta^*&|\alpha|^2\end{pmatrix}
    +\begin{pmatrix}|\beta|^2&-\alpha^*\beta\\ -\alpha\beta^*&|\alpha|^2\end{pmatrix}\Bigg]\\
  &=\frac14\begin{pmatrix}2(|\alpha|^2+|\beta|^2)&0\\0&2(|\alpha|^2+|\beta|^2)\end{pmatrix}
   =\frac12\begin{pmatrix}1&0\\0&1\end{pmatrix}=\frac12\,\mathrm{Id}.
\end{align*}
Sans les deux bits, Bob a donc un qubit complètement mélangé, qui ne dépend pas de $\alpha$ ni de
$\beta$ : il n'a reçu aucune information sur $\ket{\varphi}$. C'est l'envoi classique de $(a,q)$
qui permet de la retrouver, donc rien ne voyage plus vite que la lumière.

\begin{remarque}[Bilan de la téléportation]
\begin{itemize}[nosep]
  \item Ressources : $1$ e-bit (la paire de Bell) et $2$ bits classiques pour transmettre
        \emph{un qubit} d'état quelconque, sans qu'Alice connaisse $\alpha$ ni $\beta$.
  \item Pas de clonage : après la mesure, $A$ et $Q$ sont dans l'état de base $\ket{a,q}$.
        Alice n'a plus $\ket{\varphi}$ : l'état a été \emph{déplacé}, pas copié. L'intrication
        entre $A$ et $B$ a été consommée.
  \item Les quatre corrections $I,Z,X,ZX$ sont des portes contrôlées par des bits classiques
        (section~\ref{sec:controlees}).
\end{itemize}
\end{remarque}

\begin{remarque}[Les gestes d'Alice forment une mesure de Bell]
Le CNOT de contrôle $Q$ sur $A$ suivi de $H$ sur $Q$ est l'\emph{inverse} du générateur de la
section~\ref{sec:circuit-bell} (les portes $H$ et CNOT sont leurs propres inverses, et l'ordre
s'inverse) : il envoie la base de Bell du couple $(A,Q)$ sur la base de calcul. Mesurer ensuite dans
la base de calcul revient donc à mesurer $(A,Q)$ dans la base de Bell. Les deux bits $(a,q)$
disent \emph{quel} état de Bell Alice a trouvé, ce qui explique les quatre corrections de Bob.
\end{remarque}

\begin{remarque}[Non-clonage]
Alice ne peut pas simplement copier $\ket{\varphi}$ \parencite{wootters1982}. Supposons une porte
unitaire $U$ telle que $U\big(\ket{\psi}\otimes\ket{0}\big)=\ket{\psi}\otimes\ket{\psi}$ pour tout
$\ket{\psi}$. Elle copie $\ket{0}$ et $\ket{1}$, mais par linéarité
\[
  U\big(\ket{+}\otimes\ket{0}\big)=\tfrac{1}{\sqrt2}\big(\ket{00}+\ket{11}\big)
  \neq\ket{+}\otimes\ket{+}=\tfrac12\big(\ket{00}+\ket{01}+\ket{10}+\ket{11}\big),
\]
ce qui contredit l'hypothèse. Le CNOT copie un bit classique, pas une superposition (remarque de la
section~\ref{sec:cnot}). La téléportation respecte cette interdiction : l'état est déplacé, il n'est
jamais présent en deux endroits.
\end{remarque}

% ---------------------------------------------------------------------
\subsection{Corrélations non locales : le jeu CHSH}\label{sec:chsh}
% ---------------------------------------------------------------------

Les états de Bell donnent des résultats corrélés en base $Z$ comme en base $X$
(section~\ref{sec:bell}). Mais deux enveloppes contenant la même carte, préparées à l'avance, donnent
aussi des résultats égaux : des corrélations dans une seule base ne prouvent pas l'intrication. Bell
\parencite{bell1964} a montré qu'on peut trancher en \emph{changeant de base de mesure} ; la version de
Clauser, Horne, Shimony et Holt \parencite{clauser1969} se présente comme un jeu.

\begin{definition}{Le jeu CHSH}
Alice et Bob, éloignés, ne communiquent plus une fois le jeu commencé. Un arbitre tire deux bits $x$ et
$y$ de façon uniforme et indépendante, donne $x$ à Alice et $y$ à Bob. Chacun répond par un bit : $a$ pour
Alice, $b$ pour Bob. Ils \emph{gagnent} si
\[
  a\oplus b=x\cdot y,
\]
c'est-à-dire si $a=b$ pour $(x,y)\in\{(0,0),(0,1),(1,0)\}$ et si $a\neq b$ pour $(x,y)=(1,1)$.
\end{definition}

\paragraph{Stratégie classique : au plus $\frac34$.} Une stratégie déterministe est $a=f(x)$,
$b=g(y)$. Si elle gagnait les quatre parties, on aurait
\[
  f(0)\oplus g(0)=0,\qquad f(0)\oplus g(1)=0,\qquad f(1)\oplus g(0)=0,\qquad f(1)\oplus g(1)=1 .
\]
En additionnant ces quatre égalités modulo $2$, chacune des valeurs $f(0),f(1),g(0),g(1)$ apparaît
deux fois : le membre de gauche vaut $0$, celui de droite vaut $1$. C'est absurde, donc au moins une
partie sur quatre est perdue et $P(\text{gain})\leq\frac34$. La stratégie «~toujours répondre $0$~»
atteint $\frac34$ (elle ne perd que pour $(1,1)$). Un hasard partagé ne change rien : il mélange des
stratégies déterministes, dont aucune ne dépasse $\frac34$.

\paragraph{Stratégie quantique.} Alice et Bob partagent avant le jeu une paire $\ket{\Phi^+}$ (un
e-bit), puis chacun mesure son qubit sans rien dire à l'autre. Pour un angle $\theta$, notons
\[
  \ket{\theta_+}=\cos\tfrac\theta2\,\ket{0}+\sin\tfrac\theta2\,\ket{1},
  \qquad
  \ket{\theta_-}=-\sin\tfrac\theta2\,\ket{0}+\cos\tfrac\theta2\,\ket{1}
\]
la base de mesure d'angle $\theta$ : c'est la direction $\theta$ dans le plan $xz$ de la sphère de
Bloch ($\theta=0$ : base $Z$ ; $\theta=\frac\pi2$ : base $X$). Le résultat vaut $0$ pour
$\ket{\theta_+}$ et $1$ pour $\ket{\theta_-}$. La stratégie est
\begin{itemize}[nosep]
  \item Alice : $\theta=0$ si $x=0$ (base $Z$), $\theta=\frac\pi2$ si $x=1$ (base $X$) ;
  \item Bob : $\theta=\frac\pi4$ si $y=0$, $\theta=-\frac\pi4$ si $y=1$ (bases propres de
        $\frac{Z\pm X}{\sqrt2}$).
\end{itemize}

\begin{figure}[!ht]
\centering
\adjustbox{max width=\linewidth}{%
\begin{tikzpicture}[>=Stealth, font=\footnotesize,
    joueur/.style={draw=insarouge, fill=insarouge!6, rounded corners=2pt, thick,
                   minimum width=2.1cm, minimum height=1.1cm, align=center},
    arb/.style={draw=insagris, fill=insagris!8, rounded corners=2pt, thick,
                minimum width=3.6cm, minimum height=0.8cm, align=center}]
  % --- le jeu ---
  \node[arb] (arb) at (3.2,2.7) {arbitre : $x,y$ au hasard};
  \node[joueur] (al) at (0,0) {Alice};
  \node[joueur] (bo) at (6.4,0) {Bob};
  \draw[->, thick, insagris] (arb.west) -| node[pos=0.2, above] {$x$} (al.north);
  \draw[->, thick, insagris] (arb.east) -| node[pos=0.2, above] {$y$} (bo.north);
  \draw[insarouge, thick, decorate, decoration={snake, amplitude=1.1pt, segment length=6pt}]
    (al.east) -- (bo.west);
  \node[insarouge, above=1pt] at (3.2,0.05) {$\ket{\Phi^+}$};
  \node[text=insagris, below=2pt, font=\scriptsize] at (3.2,-0.05) {paire intriquée};
  \draw[->, thick] (al.south) -- ++(0,-1.0) node[below] {$a$};
  \draw[->, thick] (bo.south) -- ++(0,-1.0) node[below] {$b$};
  \node[draw=insagris, rounded corners=2pt, inner sep=3pt] at (3.2,-1.75) {gain si $a\oplus b=x\cdot y$};
  % --- directions de mesure dans le plan xz de la sphère de Bloch ---
  \begin{scope}[shift={(11.9,0.3)}]
    \draw[insagris!70] (0,0) circle (1.7);
    \draw[insagris!50, dashed] (-1.7,0) -- (1.7,0);
    \draw[insagris!50, dashed] (0,-1.7) -- (0,1.7);
    \node[font=\scriptsize, text=insagris, anchor=north] at (0,-1.72) {$\ket{1}$};
    \node[font=\scriptsize, text=insagris, anchor=east] at (-1.72,-0.02) {$\ket{-}$};
    \draw[->, very thick, insarouge] (0,0) -- (90:1.7);
    \draw[->, very thick, insarouge] (0,0) -- (0:1.7);
    \draw[->, very thick, insagris] (0,0) -- (45:1.7);
    \draw[->, very thick, insagris] (0,0) -- (135:1.7);
    \node[insarouge, anchor=south, inner sep=2pt] at (0,1.72) {$x=0$ : $Z$ $(\ket{0})$};
    \node[insarouge, anchor=west, inner sep=2pt] at (1.72,0) {$x=1$ : $X$ $(\ket{+})$};
    \node[text=insagris, anchor=south west, align=left, inner sep=2pt] at (1.25,1.25) {$y=0$\\ $\frac{Z+X}{\sqrt2}$};
    \node[text=insagris, anchor=south east, align=right, inner sep=2pt] at (-1.25,1.25) {$y=1$\\ $\frac{Z-X}{\sqrt2}$};
    \draw[insagris] (0.85,0) arc (0:45:0.85);
    \node[font=\scriptsize, text=insagris] at (22.5:1.12) {$\frac\pi4$};
    \draw[insagris] (0,0.85) arc (90:135:0.85);
    \node[font=\scriptsize, text=insagris] at (112.5:1.12) {$\frac\pi4$};
    \draw[insagris] (0.6,0.6) arc (45:90:0.85);
    \node[font=\scriptsize, text=insagris] at (67.5:1.12) {$\frac\pi4$};
  \end{scope}
\end{tikzpicture}}
\caption{Le jeu CHSH (à gauche) et les quatre directions de mesure de la stratégie quantique (à
droite), dans le plan $xz$ de la sphère de Bloch : Alice en rouge, Bob en gris. Pour
$(x,y)=(0,0),(0,1),(1,0)$, les directions font un angle $\frac\pi4$ ; pour $(1,1)$, un angle
$\frac{3\pi}4$.}
\label{fig:chsh}
\end{figure}

\begin{resultat}{Corrélations de $\ket{\Phi^+}$ en bases tournées}
Si Alice mesure dans la base d'angle $\alpha$ et Bob dans la base d'angle $\beta$, alors
\[
  P(a=b)=\cos^2\frac{\alpha-\beta}{2}.
\]
\end{resultat}

\noindent\emph{Preuve.} Avec $c=\cos\frac\alpha2$ et $s=\sin\frac\alpha2$,
$\ket{\alpha_+\alpha_+}+\ket{\alpha_-\alpha_-}=(c^2+s^2)\ket{00}+(cs-sc)\big(\ket{01}+\ket{10}\big)+(s^2+c^2)\ket{11}
=\ket{00}+\ket{11}$ : on peut donc écrire $\ket{\Phi^+}=\frac{1}{\sqrt2}\big(\ket{\alpha_+\alpha_+}+\ket{\alpha_-\alpha_-}\big)$.
Si Alice obtient $a=0$ (probabilité $\frac12$), le qubit de Bob est dans $\ket{\alpha_+}$ et Bob obtient $b=0$
avec la probabilité $|\braket{\beta_+}{\alpha_+}|^2$, où $\braket{\beta_+}{\alpha_+}=\cos\frac\beta2\cos\frac\alpha2+\sin\frac\beta2\sin\frac\alpha2=\cos\frac{\alpha-\beta}{2}$.
Si Alice obtient $a=1$, Bob est dans $\ket{\alpha_-}$ et obtient $b=1$ avec la même probabilité, car
$\braket{\beta_-}{\alpha_-}=\braket{\beta_+}{\alpha_+}$. Donc
$P(a=b)=\frac12\cos^2\frac{\alpha-\beta}2+\frac12\cos^2\frac{\alpha-\beta}2$. \hfill$\square$

\medskip
\noindent Pour les quatre couples $(x,y)$ :
\begin{center}
\renewcommand{\arraystretch}{1.3}
\begin{tabular}{@{}c c c c c c@{}}
\toprule
$(x,y)$ & $\alpha$ & $\beta$ & $P(a=b)$ & condition de gain & $P(\text{gain})$\\
\midrule
$(0,0)$ & $0$ & $\frac\pi4$ & $\cos^2\frac\pi8$ & $a=b$ & $\cos^2\frac\pi8$\\
$(0,1)$ & $0$ & $-\frac\pi4$ & $\cos^2\frac\pi8$ & $a=b$ & $\cos^2\frac\pi8$\\
$(1,0)$ & $\frac\pi2$ & $\frac\pi4$ & $\cos^2\frac\pi8$ & $a=b$ & $\cos^2\frac\pi8$\\
$(1,1)$ & $\frac\pi2$ & $-\frac\pi4$ & $\cos^2\frac{3\pi}8=\sin^2\frac\pi8$ & $a\neq b$ &
  $1-\sin^2\frac\pi8=\cos^2\frac\pi8$\\
\bottomrule
\end{tabular}
\end{center}
Alice et Bob gagnent donc chaque partie avec la même probabilité :
\[
  P(\text{gain})=\cos^2\frac\pi8=\frac{2+\sqrt2}{4}\approx0{,}854>\frac34 .
\]

\begin{remarque}[Ce que montre le jeu CHSH]
\begin{itemize}[nosep]
  \item Aucune stratégie classique, même avec du hasard partagé, ne dépasse $\frac34$ : l'intrication
        produit des corrélations qu'aucun modèle «~local~» (résultats fixés à l'avance) ne reproduit.
        Les expériences sur des photons intriqués confirment la valeur quantique.
  \item $\cos^2\frac\pi8$ est le maximum quantique (\emph{borne de Tsirelson}, résultat admis,
        \cite{tsirelson1980}) : on ne peut pas gagner à tous les coups. Avec
        $E_{xy}=P(a=b)-P(a\neq b)=\cos(\alpha-\beta)$ et $S=E_{00}+E_{01}+E_{10}-E_{11}$, on a
        $P(\text{gain})=\frac12+\frac S8$ ; l'inégalité de Bell donne $|S|\leq2$ en classique, et la
        stratégie ci-dessus atteint $S=2\sqrt2$.
  \item Il n'y a pas de communication : la réponse de Bob est uniforme ($P(b=0)=\frac12$) quelle que
        soit la base choisie par Alice. L'intrication ne transmet pas d'information ; elle ne fait
        que corréler les résultats.
\end{itemize}
\end{remarque}

\clearpage
% ---------------------------------------------------------------------
\subsection{Récapitulatif du chapitre}\label{sec:recap5}
% ---------------------------------------------------------------------

\begin{center}
\adjustbox{max width=\linewidth}{%
\begin{tikzpicture}[>=Stealth, font=\footnotesize,
    boite/.style={draw=insarouge, fill=insarouge!6, rounded corners=2pt, thick,
                  minimum width=4.4cm, minimum height=1.45cm, align=center, inner sep=3pt},
    lien/.style={->, very thick, insagris},
    etiq/.style={font=\scriptsize\itshape, text=insagris, inner sep=2pt}]
  \node[boite] (mes) at (0,0)
    {\textbf{Mesure à deux qubits}\\ $P(ij)=|c_{ij}|^2$\\ {\scriptsize\S\,\ref{sec:mesure2}}};
  \node[boite] (sep) at (6.3,0)
    {\textbf{Séparable ou intriqué}\\ $D=c_{00}c_{11}-c_{01}c_{10}$\\ {\scriptsize\S\,\ref{sec:separable}}};
  \node[boite] (bel) at (12.6,0)
    {\textbf{États de Bell}\\ $\ket{\Phi^\pm}$, $\ket{\Psi^\pm}$ ; $\mathrm{CNOT}\,(I\otimes H)$\\
     {\scriptsize\S\,\ref{sec:bell}, \ref{sec:circuit-bell}}};
  \node[boite] (tra) at (12.6,-2.3)
    {\textbf{Trace partielle, entropies}\\ $\rho_A=\operatorname{Tr}_B\rho_{AB}$ ; $S(\rho)$, $I(A;B)$\\
     {\scriptsize\S\,\ref{sec:trace}, \ref{sec:info}}};
  \node[boite] (tel) at (6.3,-2.3)
    {\textbf{Téléportation}\\ e-bit + 2 bits classiques\\ {\scriptsize\S\,\ref{sec:teleportation}}};
  \node[boite] (chs) at (0,-2.3)
    {\textbf{Jeu CHSH}\\ $\frac34$ classique, $\cos^2\frac\pi8$ quantique\\ {\scriptsize\S\,\ref{sec:chsh}}};
  \draw[lien] (mes) -- (sep) node[etiq, midway, above=1pt] {corrélations};
  \draw[lien] (sep) -- (bel) node[etiq, midway, above=1pt] {cas extrême};
  \draw[lien] (bel) -- (tra) node[etiq, midway, right=1pt] {un qubit seul};
  \draw[lien] (tra) -- (tel) node[etiq, midway, above=1pt] {ressource};
  \draw[lien] (tel) -- (chs) node[etiq, midway, above=1pt] {non-localité};
\end{tikzpicture}}
\end{center}

\begin{aretenir}
\begin{itemize}
  \item Deux qubits : quatre amplitudes $c_{ij}$, $P(ij)=|c_{ij}|^2$. Mesurer $A$ supprime les termes
        incompatibles et renormalise : l'état de $B$ peut changer.
  \item Séparable $\Longleftrightarrow D=c_{00}c_{11}-c_{01}c_{10}=0$. Sinon l'état est intriqué et
        aucun qubit n'a d'état propre (pour $\ket{\Phi^+}$, $P(00)\neq P(A{=}0)P(B{=}0)$).
  \item États de Bell : base orthonormée d'états maximalement intriqués, états propres de
        $Z\otimes Z$ et $X\otimes X$ ; $\mathrm{CNOT}\,(I\otimes H)$ les produit à partir de la base de
        calcul. Des portes locales ne créent pas d'intrication.
  \item Qubit seul : $\rho_A=\operatorname{Tr}_B\rho_{AB}$, avec $\operatorname{Tr}\rho_A^2=1-2|D|^2$ pour un état pur.
        Pour $\ket{\Phi^+}$ : $\rho_A=\mathrm{Id}/2$ alors que le couple est pur ; $S(\rho_{AB})=0$,
        $S(\rho_A)=1$ bit, $I(A;B)=2$ bits, $H(A|B)=-1$ (impossible en classique).
  \item Téléportation : un e-bit et deux bits classiques ; $P(a,q)=\frac14$ ; Bob applique $X$ si $a=1$
        puis $Z$ si $q=1$. L'état est déplacé, pas copié ; sans les bits, $\rho_B=\mathrm{Id}/2$.
  \item Jeu CHSH : toute stratégie classique gagne avec une probabilité $\leq\frac34$ ; avec $\ket{\Phi^+}$
        et des bases tournées, $\cos^2\frac\pi8\approx0{,}854$ (borne de Tsirelson), sans aucune
        communication.
\end{itemize}
\end{aretenir}

\begin{center}
\renewcommand{\arraystretch}{1.2}
\begin{tabular}{@{}l l r@{}}
\toprule
\textbf{Notion} & \textbf{Formule} & \textbf{§} \\
\midrule
Mesure de deux qubits & $P(ij)=|c_{ij}|^2$, \quad $P(A{=}a)=P(a0)+P(a1)$ & \ref{sec:mesure2} \\
Séparabilité & $\ket{\Psi}=\ket{\Psi}_A\otimes\ket{\Psi}_B\ \Longleftrightarrow\ D=0$ & \ref{sec:separable} \\
États de Bell & $\ket{\Phi^\pm}=\dfrac{\ket{00}\pm\ket{11}}{\sqrt2}$, \quad
  $\ket{\Psi^\pm}=\dfrac{\ket{01}\pm\ket{10}}{\sqrt2}$ & \ref{sec:bell} \\
Trace partielle & $\operatorname{Tr}_B\big(\ket{a}\bra{a'}\otimes\ket{b}\bra{b'}\big)
  =\ket{a}\bra{a'}\,\braket{b'}{b}$ & \ref{sec:trace} \\
Qubit d'un état pur & $\rho_A=CC^\dagger$, \quad $\operatorname{Tr}\rho_A^2=1-2|D|^2$ & \ref{sec:trace} \\
Entropies & $S(\rho)=-\sum_i\lambda_i\log_2\lambda_i$, \quad $I(A;B)=H(A)+H(B)-H(A,B)$ & \ref{sec:info} \\
Téléportation & $P(a,q)=\frac14$, \quad correction $Z^{q}X^{a}$ (d'abord $X$) & \ref{sec:teleportation} \\
Singulet & $(U\otimes U)\ket{\Psi^-}=\det(U)\,\ket{\Psi^-}$ & \ref{sec:bell} \\
Jeu CHSH & $P_{\mathrm{cl}}\leq\frac34$, \quad $P_{\mathrm{quant}}=\cos^2\frac\pi8=\frac{2+\sqrt2}{4}$ & \ref{sec:chsh} \\
\bottomrule
\end{tabular}
\end{center}

\begin{savoirfaire}
\begin{itemize}
  \item Calculer $P(ij)$, les probabilités d'un seul qubit et l'état après la mesure de l'un d'eux ;
        tester la séparabilité avec $D$ et calculer la pureté de $\rho_A$ (exercice~\ref{ex:separable3}).
  \item Suivre un état dans le générateur de Bell, écrire $\rho_{AB}$ et reconnaître les corrélations
        d'un état de Bell dans deux bases (exercice~\ref{ex:bellX}).
  \item Calculer $H(A)$, $H(A,B)$, $I(A;B)$ et $H(A|B)$ pour une loi ou un état (exercice~\ref{ex:entropies}).
  \item Dérouler la téléportation pour un état donné et utiliser les états de Bell pour le codage
        superdense (exercices~\ref{ex:telenum} et~\ref{ex:superdense}).
  \item Prouver la borne $\frac34$ du jeu CHSH classique et calculer la probabilité de gain d'une
        stratégie quantique (section~\ref{sec:chsh}).
\end{itemize}
\end{savoirfaire}
.
%  Il ne contient ni préambule ni \begin{document} : les environnements
%  (definition, resultat, remarque, aretenir, savoirfaire, exercice, correction),
%  les macros (\ii, \ee, \dd, \pd, \pdd, \Hop, \ket, \bra, \braket), les symboles
%  de circuits (ctl, cible, croix, mesure, porte) et les couleurs (insarouge,
%  insagris) sont définis dans main.tex. Un chapitre = une \section.
% =====================================================================

% Macros de Dirac (déjà définies dans main.tex ; ce bloc évite l'erreur
% « Undefined control sequence \ket » si main.tex est une ancienne version).
\providecommand{\ket}[1]{|#1\rangle}
\providecommand{\bra}[1]{\langle #1|}
\providecommand{\braket}[2]{\langle #1|#2\rangle}

% =====================================================================
\section{États quantiques, intrication et états de Bell}\label{chap:bell}
% =====================================================================

Les états de deux qubits forment un espace de dimension~$4$, et la plupart d'entre eux ne se
décomposent pas en «~un état pour $A$, un état pour $B$~» : ils sont \emph{intriqués}. Ce chapitre
apprend à mesurer un couple de qubits, à reconnaître l'intrication (critère de séparabilité), à la
produire avec un circuit (états de Bell), à décrire un qubit pris seul (trace partielle), à
quantifier l'information (entropies) et à s'en servir pour téléporter un état.

% ---------------------------------------------------------------------
\subsection{États de deux qubits et mesure}\label{sec:mesure2}
% ---------------------------------------------------------------------

Un état quelconque de $\mathcal{H}^4$ se décompose sur la base de calcul :
\[
  \ket{\Psi}=c_{00}\ket{00}+c_{01}\ket{01}+c_{10}\ket{10}+c_{11}\ket{11}
  =\begin{pmatrix}c_{00}\\ c_{01}\\ c_{10}\\ c_{11}\end{pmatrix},
  \qquad \sum_{i,j}|c_{ij}|^2=1 .
\]

\begin{resultat}{Mesure des deux qubits dans la base de calcul}
\[
  P(ij)=|\braket{ij}{\Psi}|^2=|c_{ij}|^2 .
\]
Pour un seul qubit : $P(A{=}0)=P(00)+P(01)$, $P(A{=}1)=P(10)+P(11)$,
$P(B{=}0)=P(00)+P(10)$, $P(B{=}1)=P(01)+P(11)$.
\end{resultat}

En effet $\braket{ij}{\Psi}=\sum_{k,l}c_{kl}\braket{ij}{kl}=\sum_{k,l}c_{kl}\delta_{ik}\delta_{jl}=c_{ij}$.

\paragraph{Exemple 1.} $\ket{\Psi}=\frac{1}{\sqrt2}\ket{00}+\frac{1}{\sqrt2}\ket{01}
=\big(\tfrac{1}{\sqrt2},\tfrac{1}{\sqrt2},0,0\big)^{\mathsf T}$. Il est normalisé :
$\frac12+\frac12+0+0=1$. Les probabilités jointes sont
$P(00)=P(01)=\frac12$ et $P(10)=P(11)=0$. Pour chaque qubit :
\[
  P(A{=}0)=\tfrac12+\tfrac12=1,\quad P(A{=}1)=0+0=0,\quad
  P(B{=}0)=\tfrac12+0=\tfrac12,\quad P(B{=}1)=\tfrac12+0=\tfrac12 .
\]
$A$ donne toujours $0$ ; $B$ donne $0$ ou $1$ au hasard.

\paragraph{Exemple 2 (mesure de $A$ seul).} Soit
$\ket{\Psi}=\frac12\ket{00}+\frac12\ket{01}+\frac{1}{\sqrt2}\ket{11}$. Normalisation :
$\frac14+\frac14+0+\frac12=1$. Probabilités jointes : $P(00)=\frac14$, $P(01)=\frac14$,
$P(10)=0$, $P(11)=\frac12$. Donc
\[
  P(A{=}0)=\tfrac14+\tfrac14=\tfrac12,\quad P(A{=}1)=0+\tfrac12=\tfrac12,\quad
  P(B{=}0)=\tfrac14+0=\tfrac14,\quad P(B{=}1)=\tfrac14+\tfrac12=\tfrac34 .
\]

\begin{remarque}[Effondrement après la mesure de $A$]
Si $A$ donne le résultat $a$, on ne garde que les termes de $\ket{\Psi}$ dont le premier
chiffre vaut $a$, puis on renormalise :
\[
  \ket{\Psi}\ \to\ \frac{c_{a0}\ket{a0}+c_{a1}\ket{a1}}{\sqrt{|c_{a0}|^2+|c_{a1}|^2}},
  \qquad\text{probabilité }|c_{a0}|^2+|c_{a1}|^2 .
\]
\end{remarque}

Pour l'exemple 2 :
\begin{itemize}[nosep]
  \item $A=0$ (probabilité $\frac12$) : on garde $\frac12\ket{00}+\frac12\ket{01}$, de norme
        au carré $\frac12$ ; renormalisé, $\frac{1}{\sqrt2}\big(\ket{00}+\ket{01}\big)
        =\ket{0}\otimes\ket{+}$. Ensuite $B$ vaut $0$ ou $1$ avec la probabilité $\frac12$.
  \item $A=1$ (probabilité $\frac12$) : on garde $\frac{1}{\sqrt2}\ket{11}$, de norme au carré
        $\frac12$ ; renormalisé, $\ket{11}=\ket{1}\otimes\ket{1}$. Ensuite $B$ vaut $1$
        avec certitude.
\end{itemize}
Contrôle : $P(B{=}0)=\frac12\cdot\frac12+\frac12\cdot0=\frac14$, comme calculé plus haut.
Ici la mesure de $A$ \emph{modifie} l'état de $B$ ($\ket{+}$ ou $\ket{1}$) : c'est la
signature d'un état intriqué (section~\ref{sec:separable}). La figure~\ref{fig:mesure2} résume
les deux mesures successives.

\begin{figure}[!ht]
\centering
\adjustbox{max width=\linewidth}{%
\begin{tikzpicture}[>=Stealth, font=\small,
    noeud/.style={draw=insarouge, fill=insarouge!6, rounded corners=2pt, thick, inner sep=4pt, align=center},
    fin/.style={draw=insagris, fill=insagris!10, rounded corners=2pt, inner sep=4pt, align=center},
    lien/.style={->, thick, insagris},
    proba/.style={font=\footnotesize, fill=white, inner sep=1.5pt, text=black}]
  \node[noeud] (psi) at (0,0) {$\ket{\Psi}$};
  \node[noeud] (a0) at (4.4,1.6) {$A=0$\\ {\footnotesize état $\ket{0}\otimes\ket{+}$}};
  \node[noeud] (a1) at (4.4,-1.6) {$A=1$\\ {\footnotesize état $\ket{1}\otimes\ket{1}$}};
  \node[fin] (b00) at (9.2,2.5) {$B=0$ : $\ket{00}$};
  \node[fin] (b01) at (9.2,0.7) {$B=1$ : $\ket{01}$};
  \node[fin] (b11) at (9.2,-1.6) {$B=1$ : $\ket{11}$};
  \draw[lien] (psi) -- (a0) node[proba, midway] {$\frac12$};
  \draw[lien] (psi) -- (a1) node[proba, midway] {$\frac12$};
  \draw[lien] (a0) -- (b00) node[proba, midway] {$\frac12$};
  \draw[lien] (a0) -- (b01) node[proba, midway] {$\frac12$};
  \draw[lien] (a1) -- (b11) node[proba, midway] {$1$};
  \node[anchor=west] at (11.0,2.5) {$P(00)=\frac12\cdot\frac12=\frac14$};
  \node[anchor=west] at (11.0,0.7) {$P(01)=\frac12\cdot\frac12=\frac14$};
  \node[anchor=west] at (11.0,-1.6) {$P(11)=\frac12\cdot1=\frac12$};
  \node[text=insagris, font=\footnotesize\itshape] at (2.2,3.0) {mesure de $A$};
  \node[text=insagris, font=\footnotesize\itshape] at (6.8,3.6) {mesure de $B$};
  \node[text=insagris, font=\footnotesize\itshape, anchor=west] at (11.0,3.6) {probabilité jointe};
\end{tikzpicture}}
\caption{Mesure de $A$ puis de $B$ pour l'état de l'exemple~2,
$\ket{\Psi}=\frac12\ket{00}+\frac12\ket{01}+\frac{1}{\sqrt2}\ket{11}$. Chaque flèche porte une
probabilité, et le produit des probabilités le long d'un chemin est la probabilité jointe. Après la
mesure de $A$, l'état de $B$ dépend du résultat ($\ket{+}$ ou $\ket{1}$).}
\label{fig:mesure2}
\end{figure}

% ---------------------------------------------------------------------
\subsection{États séparables et états intriqués}\label{sec:separable}
% ---------------------------------------------------------------------

\begin{definition}{Séparable ou intriqué}
Un état $\ket{\Psi}_{AB}$ est \emph{séparable} s'il s'écrit comme un produit
$\ket{\Psi}_{AB}=\ket{\Psi}_A\otimes\ket{\Psi}_B$. Sinon il est \emph{intriqué} : on ne
peut pas attribuer d'état propre à chaque qubit.
\end{definition}

\paragraph{Exemple séparable (exemple 1 ci-dessus).} Par linéarité du produit tensoriel :
\begin{align*}
  \ket{\Psi}&=\tfrac{1}{\sqrt2}\ket{00}_{AB}+\tfrac{1}{\sqrt2}\ket{01}_{AB}
    =\tfrac{1}{\sqrt2}\Big(\ket{0}_A\otimes\ket{0}_B+\ket{0}_A\otimes\ket{1}_B\Big)\\
  &=\ket{0}_A\otimes\Big(\tfrac{1}{\sqrt2}\big(\ket{0}+\ket{1}\big)\Big)_B
    =\ket{0}_A\otimes\ket{+}_B ,
\end{align*}
donc $\ket{\Psi}_A=\ket{0}$ et $\ket{\Psi}_B=\ket{+}$. Contrôle par les composantes :
$\begin{pmatrix}1\\ 0\end{pmatrix}\otimes\tfrac{1}{\sqrt2}\begin{pmatrix}1\\ 1\end{pmatrix}
=\tfrac{1}{\sqrt2}\begin{pmatrix}1\cdot1\\ 1\cdot1\\ 0\cdot1\\ 0\cdot1\end{pmatrix}
=\tfrac{1}{\sqrt2}\begin{pmatrix}1\\ 1\\ 0\\ 0\end{pmatrix}$.

\subsubsection{Un critère de séparabilité}

Un produit quelconque a pour coefficients $c_{00}=ac$, $c_{01}=ad$, $c_{10}=bc$,
$c_{11}=bd$ (formule de la section~\ref{sec:tenseur}, avec $a,b$ pour $A$ et $c,d$ pour $B$), donc
\[
  c_{00}c_{11}=(ac)(bd)=abcd=(ad)(bc)=c_{01}c_{10}.
\]

\begin{resultat}{Test de séparabilité (deux qubits)}
$\ket{\Psi}=\sum c_{ij}\ket{ij}$ est séparable si et seulement si
\[
  D=c_{00}c_{11}-c_{01}c_{10}=0 .
\]
\end{resultat}

\emph{Démonstration.} ($\Rightarrow$) Calcul ci-dessus. ($\Leftarrow$) On regroupe les
termes selon l'état de $A$ :
\[
  \ket{\Psi}=\ket{0}\otimes\big(c_{00}\ket{0}+c_{01}\ket{1}\big)+\ket{1}\otimes\big(c_{10}\ket{0}+c_{11}\ket{1}\big).
\]
Si $D=0$, les lignes $(c_{00},c_{01})$ et $(c_{10},c_{11})$ sont proportionnelles. Si
$(c_{00},c_{01})\neq(0,0)$, il existe $\lambda$ avec $(c_{10},c_{11})=\lambda(c_{00},c_{01})$
et $\ket{\Psi}=\big(\ket{0}+\lambda\ket{1}\big)\otimes\big(c_{00}\ket{0}+c_{01}\ket{1}\big)$.
Si $(c_{00},c_{01})=(0,0)$, alors $\ket{\Psi}=\ket{1}\otimes\big(c_{10}\ket{0}+c_{11}\ket{1}\big)$.
Dans les deux cas $\ket{\Psi}$ est un produit ; on renormalise chaque facteur.
\hfill$\square$

\subsubsection{Application à plusieurs états}

\begin{center}
\renewcommand{\arraystretch}{1.5}
\begin{tabular}{@{}lccl@{}}
\toprule
\textbf{État} & $(c_{00},c_{01},c_{10},c_{11})$ & $D$ & \textbf{Nature} \\
\midrule
$\ket{00}$ & $(1,0,0,0)$ & $1\cdot0-0\cdot0=0$ & séparable \\
$\ket{+}\otimes\ket{+}$ & $\frac12(1,1,1,1)$ & $\frac14-\frac14=0$ & séparable \\
$\ket{+}\otimes\ket{-}$ & $\frac12(1,-1,1,-1)$ & $-\frac14-(-\frac14)=0$ & séparable \\
$\frac{1}{\sqrt2}(\ket{00}+\ket{01})$ & $\frac{1}{\sqrt2}(1,1,0,0)$ & $0-\frac12\cdot0=0$ & séparable \\
$\ket{\Phi^+}=\frac{1}{\sqrt2}(\ket{00}+\ket{11})$ & $\frac{1}{\sqrt2}(1,0,0,1)$ & $\frac12-0=\frac12$ & intriqué \\
$\ket{\Phi^-}=\frac{1}{\sqrt2}(\ket{00}-\ket{11})$ & $\frac{1}{\sqrt2}(1,0,0,-1)$ & $-\frac12-0=-\frac12$ & intriqué \\
$\ket{\Psi^+}=\frac{1}{\sqrt2}(\ket{01}+\ket{10})$ & $\frac{1}{\sqrt2}(0,1,1,0)$ & $0-\frac12=-\frac12$ & intriqué \\
$\ket{\Psi^-}=\frac{1}{\sqrt2}(\ket{01}-\ket{10})$ & $\frac{1}{\sqrt2}(0,1,-1,0)$ & $0-(-\frac12)=\frac12$ & intriqué \\
$\frac12\ket{00}+\frac12\ket{01}+\frac{1}{\sqrt2}\ket{11}$ & $(\frac12,\frac12,0,\frac{1}{\sqrt2})$ & $\frac{1}{2\sqrt2}-0$ & intriqué \\
\bottomrule
\end{tabular}
\end{center}

Les quatre états de Bell $\ket{\Phi^\pm}$ et $\ket{\Psi^\pm}$ de la table sont tous intriqués
($D=\pm\frac12$) ; leurs propriétés font l'objet de la section~\ref{sec:bell}.

\paragraph{Preuve directe que $\ket{\Phi^+}$ est intriqué.} Supposons
$\ket{\Phi^+}=(a\ket{0}+b\ket{1})\otimes(c\ket{0}+d\ket{1})$. En identifiant les
composantes : $ac=\frac{1}{\sqrt2}$, $ad=0$, $bc=0$, $bd=\frac{1}{\sqrt2}$. La condition
$ad=0$ impose $a=0$ (impossible car $ac\neq0$) ou $d=0$ (impossible car $bd\neq0$).
Contradiction : aucun produit ne convient.

\paragraph{Corrélations.} Pour un produit $\ket{\Psi}_A\otimes\ket{\Psi}_B$ avec
$\ket{\Psi}_A=a_0\ket{0}+a_1\ket{1}$ et $\ket{\Psi}_B=b_0\ket{0}+b_1\ket{1}$, on a
$P(ij)=|a_ib_j|^2=|a_i|^2|b_j|^2=P(A{=}i)\,P(B{=}j)$ : les deux qubits sont
\emph{indépendants}. Vérification avec $\ket{+}\otimes\ket{+}$ : $P(00)=\frac14=\frac12\cdot\frac12$.
Pour $\ket{\Phi^+}$ : $P(00)=P(11)=\frac12$ et $P(01)=P(10)=0$, donc
$P(A{=}0)=P(B{=}0)=\frac12$, mais
\[
  P(00)=\tfrac12\neq\tfrac12\cdot\tfrac12=P(A{=}0)\,P(B{=}0) .
\]
Les deux qubits donnent toujours le même résultat, aléatoire : ils sont corrélés.
La figure~\ref{fig:epr} illustre la différence entre les deux cas.

\begin{figure}[!ht]
\centering
\adjustbox{max width=\linewidth}{%
\begin{tikzpicture}[>=Stealth, font=\small,
    labo/.style={draw=insagris, fill=insagris!8, rounded corners=2pt, thick, minimum width=2.6cm,
                 minimum height=1.0cm, align=center},
    source/.style={draw=insarouge, fill=insarouge!10, rounded corners=2pt, thick, minimum width=3.6cm,
                   minimum height=1.0cm, align=center},
    droit/.style={->, thick, insagris},
    onde/.style={->, thick, insarouge, decorate,
                 decoration={snake, amplitude=1.3pt, segment length=7pt, post length=2mm}},
    res/.style={font=\small, align=center}]
  % (a) état séparable
  \node[labo] (Aa) at (0,0) {Alice\\ mesure $A$};
  \node[source] (Sa) at (6,0) {source\\ $\ket{0}_A\otimes\ket{+}_B$};
  \node[labo] (Ba) at (12,0) {Bob\\ mesure $B$};
  \draw[droit] (Sa.west) -- (Aa.east);
  \draw[droit] (Sa.east) -- (Ba.west);
  \node[res, below=3pt of Aa] {$0\ \ 0\ \ 0\ \ 0\ \ 0\ \ 0$};
  \node[res, below=3pt of Ba] {$1\ \ 0\ \ 0\ \ 1\ \ 0\ \ 1$};
  \node[res, text=insagris, font=\footnotesize] at (6,-1.0) {$A$ donne toujours $0$ ; $B$ est aléatoire.\\
        Les deux résultats sont indépendants.};
  \node[anchor=south west, font=\small\bfseries] at (-1.3,0.75) {(a) état séparable};
  % (b) état intriqué
  \begin{scope}[yshift=-3.6cm]
    \node[labo] (Ab) at (0,0) {Alice\\ mesure $A$};
    \node[source] (Sb) at (6,0) {source\\ $\ket{\Phi^+}_{AB}$};
    \node[labo] (Bb) at (12,0) {Bob\\ mesure $B$};
    \draw[onde] (Sb.west) -- (Ab.east);
    \draw[onde] (Sb.east) -- (Bb.west);
    \node[res, below=3pt of Ab] {$1\ \ 0\ \ 0\ \ 1\ \ 1\ \ 0$};
    \node[res, below=3pt of Bb] {$1\ \ 0\ \ 0\ \ 1\ \ 1\ \ 0$};
    \node[res, text=insagris, font=\footnotesize] at (6,-1.0) {$A$ et $B$ sont chacun aléatoires,\\
          mais toujours égaux : ils sont corrélés.};
    \node[anchor=south west, font=\small\bfseries] at (-1.3,0.75) {(b) état intriqué};
  \end{scope}
\end{tikzpicture}}
\caption{Six mesures successives sur des paires préparées dans le même état (les résultats sont
des exemples). (a) Pour un produit, chaque qubit a son propre comportement. (b) Pour
$\ket{\Phi^+}=\frac{1}{\sqrt2}(\ket{00}+\ket{11})$, chaque résultat pris seul est un tirage à pile ou
face, mais les deux résultats coïncident toujours ; la ligne ondulée symbolise l'intrication.}
\label{fig:epr}
\end{figure}

\paragraph{Et un qubit pris seul ?} Dans un état séparable, chaque qubit garde son propre état pur ;
dans un état intriqué, aucun ket ne décrit un qubit isolé. On le décrit alors par une matrice de
densité, obtenue par la trace partielle (section~\ref{sec:trace}).

% ---------------------------------------------------------------------
\subsection{Les états de Bell}\label{sec:bell}
% ---------------------------------------------------------------------

\begin{definition}{Les quatre états de Bell}
\[
  \ket{\Phi^{\pm}}=\frac{1}{\sqrt2}\big(\ket{00}\pm\ket{11}\big),
  \qquad
  \ket{\Psi^{\pm}}=\frac{1}{\sqrt2}\big(\ket{01}\pm\ket{10}\big).
\]
Ils sont intriqués ($|D|=\frac12$, section~\ref{sec:separable}) et le sont au maximum :
$|D|\leq\frac12$ pour tout état normalisé de deux qubits, et chaque qubit d'un état de Bell,
pris seul, est complètement aléatoire (section~\ref{sec:trace}).
\end{definition}

\paragraph{Une base de $\mathcal{H}^4$.} Chaque état est de norme $\frac12+\frac12=1$. Les
$\ket{\Phi^\pm}$ n'ont de composantes que sur $\ket{00}$ et $\ket{11}$, les $\ket{\Psi^\pm}$ que sur
$\ket{01}$ et $\ket{10}$, donc $\braket{\Phi^\pm}{\Psi^\pm}=0$ pour toutes les combinaisons de signes ;
enfin $\braket{\Phi^+}{\Phi^-}=\frac12(1-1)=0$ et $\braket{\Psi^+}{\Psi^-}=\frac12(1-1)=0$. Les quatre
états forment donc une base orthonormée de $\mathcal{H}^4$, la \emph{base de Bell}, constituée
uniquement d'états intriqués. Le circuit de la section~\ref{sec:circuit-bell} fait passer de la
base de calcul à cette base.

\paragraph{Ce qui les distingue : deux parités.} Sur la base de calcul,
$Z\otimes Z\ket{ij}=(-1)^{i+j}\ket{ij}$ (valeur propre $+1$ si les deux chiffres sont égaux, $-1$
sinon) et $X\otimes X\ket{ij}=\ket{\bar i\bar j}$ (les deux chiffres sont inversés, avec
$\bar i=1-i$). On en déduit, par exemple,
\[
  X\otimes X\ket{\Phi^-}=\tfrac{1}{\sqrt2}\big(\ket{11}-\ket{00}\big)=-\ket{\Phi^-},
  \qquad
  Z\otimes Z\ket{\Psi^+}=\tfrac{1}{\sqrt2}\big(-\ket{01}-\ket{10}\big)=-\ket{\Psi^+},
\]
et le tableau complet des valeurs propres :
\begin{center}
\renewcommand{\arraystretch}{1.3}
\begin{tabular}{@{}lcccc@{}}
\toprule
 & $\ket{\Phi^+}$ & $\ket{\Phi^-}$ & $\ket{\Psi^+}$ & $\ket{\Psi^-}$ \\
\midrule
$Z\otimes Z$ & $+1$ & $+1$ & $-1$ & $-1$ \\
$X\otimes X$ & $+1$ & $-1$ & $+1$ & $-1$ \\
\bottomrule
\end{tabular}
\end{center}
Les deux opérateurs commutent, car $ZX=-XZ$ donne
$(Z\otimes Z)(X\otimes X)=ZX\otimes ZX=XZ\otimes XZ=(X\otimes X)(Z\otimes Z)$ ; les quatre états de
Bell sont leurs états propres communs, et chaque couple de valeurs propres apparaît une fois
(figure~\ref{fig:grille-bell}). Le même calcul montre que $X\otimes Z$ et $Z\otimes X$
commutent : $(X\otimes Z)(Z\otimes X)=XZ\otimes ZX=(-ZX)\otimes(-XZ)=ZX\otimes XZ=(Z\otimes X)(X\otimes Z)$.
Plus généralement, $P\otimes Q$ et $P'\otimes Q'$ (matrices de Pauli) commutent si $P,P'$
anticommutent et si $Q,Q'$ anticommutent : les deux signes $-1$ se compensent.

\begin{figure}[!ht]
\centering
\adjustbox{max width=\linewidth}{%
\begin{tikzpicture}[>=Stealth, font=\small,
    cell/.style={draw=insarouge, fill=insarouge!6, rounded corners=2pt, thick, minimum width=5.4cm,
                 minimum height=1.35cm, align=center},
    tete/.style={text=insagris, align=center}]
  \node[tete] at (3.0,2.0) {$X\otimes X=+1$\\ {\footnotesize mêmes résultats en base $X$}};
  \node[tete] at (8.8,2.0) {$X\otimes X=-1$\\ {\footnotesize résultats opposés en base $X$}};
  \node[tete, anchor=east] at (-0.1,0.75) {$Z\otimes Z=+1$\\ {\footnotesize mêmes résultats}\\ {\footnotesize en base $Z$}};
  \node[tete, anchor=east] at (-0.1,-0.85) {$Z\otimes Z=-1$\\ {\footnotesize résultats opposés}\\ {\footnotesize en base $Z$}};
  \node[cell] at (3.0,0.75) {$\ket{\Phi^+}=\frac{1}{\sqrt2}\big(\ket{00}+\ket{11}\big)$};
  \node[cell] at (8.8,0.75) {$\ket{\Phi^-}=\frac{1}{\sqrt2}\big(\ket{00}-\ket{11}\big)$};
  \node[cell] at (3.0,-0.85) {$\ket{\Psi^+}=\frac{1}{\sqrt2}\big(\ket{01}+\ket{10}\big)$};
  \node[cell] at (8.8,-0.85) {$\ket{\Psi^-}=\frac{1}{\sqrt2}\big(\ket{01}-\ket{10}\big)$};
\end{tikzpicture}}
\caption{Les quatre états de Bell, classés par les valeurs propres de $Z\otimes Z$ (lignes) et de
$X\otimes X$ (colonnes). La valeur propre est le produit des deux résultats $\pm1$ d'une mesure des
deux qubits dans la base correspondante : $+1$ si les résultats sont égaux, $-1$ s'ils sont opposés.}
\label{fig:grille-bell}
\end{figure}

Les corrélations ne se limitent donc pas à la base $Z$ : $\ket{\Phi^+}$ s'écrit aussi
$\frac{1}{\sqrt2}\big(\ket{++}+\ket{--}\big)$ et donne des résultats égaux en base $X$ comme en base $Z$
(exercice~\ref{ex:bellX}).

\paragraph{D'un état de Bell à l'autre.} Une porte agissant sur un seul qubit suffit. Avec
$X\ket{0}=\ket{1}$, $X\ket{1}=\ket{0}$ et $Z\ket{1}=-\ket{1}$, en agissant sur le second qubit :
\begin{align*}
  (I\otimes X)\ket{\Phi^+}&=\tfrac{1}{\sqrt2}\big(\ket{01}+\ket{10}\big)=\ket{\Psi^+},\\
  (I\otimes Z)\ket{\Phi^+}&=\tfrac{1}{\sqrt2}\big(\ket{00}-\ket{11}\big)=\ket{\Phi^-},\\
  (I\otimes XZ)\ket{\Phi^+}&=(I\otimes X)\tfrac{1}{\sqrt2}\big(\ket{00}-\ket{11}\big)
    =\tfrac{1}{\sqrt2}\big(\ket{01}-\ket{10}\big)=\ket{\Psi^-}.
\end{align*}
Quatre opérations locales ($I$, $X$, $Z$, $XZ$) sur \emph{un} qubit d'une paire de Bell mènent donc
aux quatre états de la base ; c'est le principe du codage superdense
(exercice~\ref{ex:superdense}).

\begin{remarque}[Le singulet est le même dans toutes les bases]\label{rem:singulet}
Pour toute matrice unitaire $U=\begin{pmatrix}a&b\\c&d\end{pmatrix}$, on a $U\ket{0}=a\ket{0}+c\ket{1}$
et $U\ket{1}=b\ket{0}+d\ket{1}$, donc
\begin{align*}
  (U\otimes U)\big(\ket{01}-\ket{10}\big)
  &=\big(a\ket{0}+c\ket{1}\big)\big(b\ket{0}+d\ket{1}\big)-\big(b\ket{0}+d\ket{1}\big)\big(a\ket{0}+c\ket{1}\big)\\
  &=(ad-bc)\big(\ket{01}-\ket{10}\big),
\end{align*}
c'est-à-dire $(U\otimes U)\ket{\Psi^-}=\det(U)\,\ket{\Psi^-}$. Comme $|\det U|=1$, ce facteur est une
phase globale, sans effet physique : $\ket{\Psi^-}$ s'écrit \emph{de la même façon} dans toutes les
bases orthonormées (avec $U=H$, de déterminant $-1$ : $\ket{\Psi^-}=\frac{1}{\sqrt2}\big(\ket{-+}-\ket{+-}\big)$).
Si Alice et Bob mesurent dans la \emph{même} base, quelle qu'elle soit, leurs résultats sont
toujours opposés : c'est le «~singulet~». Parmi les états de Bell, seul $\ket{\Psi^-}$ a cette propriété.
\end{remarque}

% ---------------------------------------------------------------------
\subsection{Un circuit quantique : le générateur d'états de Bell}\label{sec:circuit-bell}
% ---------------------------------------------------------------------

On assemble une porte de Hadamard et un CNOT. Le système est $(A,B)$ avec le ket
$\ket{a,b}=\ket{a}_A\otimes\ket{b}_B$ ; d'après la convention de la section~\ref{sec:circuits}, $B$ est dessiné
en haut et $A$ en bas (figure~\ref{fig:circuit-bell}). Le circuit a deux étapes :
\begin{enumerate}[nosep]
  \item la porte $I\otimes H$ : identité sur $A$, Hadamard sur $B$ ;
  \item un CNOT de contrôle $B$ et de cible $A$ (celui de la section~\ref{sec:cnot}, dont le
        contrôle est le second chiffre).
\end{enumerate}
Les traits pointillés délimitent les états $\ket{\Pi_0}$, $\ket{\Pi_1}$ et $\ket{\Pi_2}$ :
avant les portes, après $I\otimes H$, après le CNOT. Le circuit transforme la base de calcul
en base de Bell (section~\ref{sec:bell}).

\begin{figure}[!ht]
\centering
\begin{tikzpicture}[>=Stealth, font=\small]
  % fils : B en haut, A en bas
  \draw (0,1.4) -- (7.6,1.4);
  \draw (0,0) -- (7.6,0);
  \node[anchor=east] at (0,1.4) {$\ket{0}_B$};
  \node[anchor=east] at (0,0) {$\ket{0}_A$};
  % première étape : H sur B, identité (pointillée) sur A
  \node[porte] at (2.2,1.4) {$H$};
  \node[draw=insagris, dashed, fill=white, minimum size=7mm] at (2.2,0) {$I$};
  % CNOT : contrôle B (haut), cible A (bas)
  \draw[insarouge, thick] (5,1.4) -- (5,0);
  \pic at (5,1.4) {ctl};
  \pic at (5,0) {cible};
  % accolade de sortie
  \draw[insagris, thick] (7.75,1.6) -- (7.9,1.6) -- (7.9,-0.2) -- (7.75,-0.2);
  \node[anchor=west] at (7.9,0.7) {$\ket{\beta^{00}}_{AB}$};
  % états intermédiaires
  \foreach \x in {0.7,3.6,6.4}
    \draw[dashed, insagris!70] (\x,1.9) -- (\x,-0.7);
  \node[anchor=north, font=\footnotesize] at (0.7,-0.7) {$\ket{\Pi_0}=\ket{00}$};
  \node[anchor=north, font=\footnotesize] at (3.6,-0.7) {$\ket{\Pi_1}=\ket{0}\otimes\ket{+}$};
  \node[anchor=north, font=\footnotesize] at (6.4,-0.7) {$\ket{\Pi_2}=\ket{\beta^{00}}$};
\end{tikzpicture}
\caption{Générateur d'états de Bell : $I\otimes H$, puis un CNOT de contrôle $B$ (fil du haut)
et de cible $A$ (fil du bas). L'entrée représentée est $\ket{00}$.}
\label{fig:circuit-bell}
\end{figure}

\subsubsection{Entrée \texorpdfstring{$\ket{00}$}{|00>} : calcul pas à pas}

\begin{align*}
  \ket{\Pi_0}&=\ket{0}_A\otimes\ket{0}_B=\ket{00},\\[1ex]
  \ket{\Pi_1}&=(I\otimes H)\ket{\Pi_0}=I\ket{0}_A\otimes H\ket{0}_B=\ket{0}\otimes\ket{+}
    =\frac{\ket{00}+\ket{01}}{\sqrt2},\\[1ex]
  \ket{\Pi_2}&=\mathrm{CNOT}\ket{\Pi_1}=\frac{\mathrm{CNOT}\ket{00}+\mathrm{CNOT}\ket{01}}{\sqrt2}
    =\frac{\ket{00}+\ket{11}}{\sqrt2}=\ket{\beta^{00}}=\ket{\Phi^+}.
\end{align*}
Dans $\ket{\Pi_1}$, le premier chiffre est celui de $A$ et le second celui de $B$. Le CNOT a
pour contrôle $B$ : $\ket{00}$ ($B=0$) ne change pas, $\ket{01}$ ($B=1$) devient $\ket{11}$
puisque $A$ est inversé. Le résultat est un \emph{état de Bell}.

Avec les matrices, la règle $A\otimes B=(A_{ij}B)$ donne
\[
  I\otimes H=\begin{pmatrix}1\cdot H&0\\0&1\cdot H\end{pmatrix}
  =\frac{1}{\sqrt2}\begin{pmatrix}1&1&0&0\\1&-1&0&0\\0&0&1&1\\0&0&1&-1\end{pmatrix}.
\]
Sa première colonne est $\ket{\Pi_1}=\frac{1}{\sqrt2}(1,1,0,0)^{\mathsf T}$ (image de $\ket{00}$).
Le CNOT échange les amplitudes de $\ket{01}$ et de $\ket{11}$ (section~\ref{sec:cnot}), donc
\[
  \ket{\Pi_2}=\mathrm{CNOT}\,\frac{1}{\sqrt2}\begin{pmatrix}1\\1\\0\\0\end{pmatrix}
  =\frac{1}{\sqrt2}\begin{pmatrix}1\\0\\0\\1\end{pmatrix}.
\]

\subsubsection{Les quatre entrées}

Prenons $B$ dans l'état $\ket{x}$ et $A$ dans l'état $\ket{y}$, avec $x,y\in\{0,1\}$. Alors
$\ket{\Pi_0}=\ket{y}_A\otimes\ket{x}_B=\ket{y,x}$. Comme
$H\ket{x}=\frac{1}{\sqrt2}\big(\ket{0}+(-1)^x\ket{1}\big)$ (c'est $\ket{+}$ pour $x=0$ et
$\ket{-}$ pour $x=1$),
\[
  \ket{\Pi_1}=\ket{y}\otimes H\ket{x}=\frac{1}{\sqrt2}\Big(\ket{y,0}+(-1)^x\ket{y,1}\Big).
\]
Le CNOT (contrôle $B$) ne change pas $\ket{y,0}$ et transforme $\ket{y,1}$ en
$\ket{\bar y,1}$, avec $\bar y=1-y$ :
\[
  \ket{\Pi_2}=\ket{\beta^{xy}}=\frac{1}{\sqrt2}\Big(\ket{y,0}+(-1)^x\ket{\bar y,1}\Big).
\]
Le premier indice de $\ket{\beta^{xy}}$ est l'entrée de $B$ (le qubit qui reçoit $H$), le
second celle de $A$ (la cible). Les quatre cas :

\begin{center}
\renewcommand{\arraystretch}{1.7}
\begin{tabular}{@{}ccccc@{}}
\toprule
$(B,A)=(x,y)$ & $\ket{\Pi_0}=\ket{y,x}$ & $\ket{\Pi_1}$ & $\ket{\Pi_2}=\ket{\beta^{xy}}$ & Nom \\
\midrule
$(0,0)$ & $\ket{00}$ & $\frac{1}{\sqrt2}(\ket{00}+\ket{01})$ & $\frac{1}{\sqrt2}(\ket{00}+\ket{11})$ & $\ket{\Phi^+}$ \\
$(1,0)$ & $\ket{01}$ & $\frac{1}{\sqrt2}(\ket{00}-\ket{01})$ & $\frac{1}{\sqrt2}(\ket{00}-\ket{11})$ & $\ket{\Phi^-}$ \\
$(0,1)$ & $\ket{10}$ & $\frac{1}{\sqrt2}(\ket{10}+\ket{11})$ & $\frac{1}{\sqrt2}(\ket{10}+\ket{01})$ & $\ket{\Psi^+}$ \\
$(1,1)$ & $\ket{11}$ & $\frac{1}{\sqrt2}(\ket{10}-\ket{11})$ & $\frac{1}{\sqrt2}(\ket{10}-\ket{01})$ & $-\ket{\Psi^-}$ \\
\bottomrule
\end{tabular}
\end{center}
Détail de la dernière ligne : $\ket{\Pi_1}=\ket{1}\otimes\ket{-}=\frac{1}{\sqrt2}(\ket{10}-\ket{11})$ ;
le CNOT laisse $\ket{10}$ et envoie $\ket{11}$ sur $\ket{01}$, d'où
$\frac{1}{\sqrt2}(\ket{10}-\ket{01})=-\frac{1}{\sqrt2}(\ket{01}-\ket{10})=-\ket{\Psi^-}$.

\paragraph{Forme matricielle.} Le CNOT échange les lignes $2$ et $4$ de
$I\otimes H$ (il échange les amplitudes de $\ket{01}$ et $\ket{11}$), donc
\[
  \mathrm{CNOT}\,(I\otimes H)=\frac{1}{\sqrt2}\begin{pmatrix}
    1 & 1 & 0 & 0\\ 0 & 0 & 1 & -1\\ 0 & 0 & 1 & 1\\ 1 & -1 & 0 & 0
  \end{pmatrix}.
\]
Ses colonnes sont les images de $\ket{00},\ket{01},\ket{10},\ket{11}$, c'est-à-dire de
$(B,A)=(0,0),(1,0),(0,1),(1,1)$ : ce sont $\ket{\beta^{00}}$, $\ket{\beta^{10}}$,
$\ket{\beta^{01}}$, $\ket{\beta^{11}}$, soit $\ket{\Phi^+}$, $\ket{\Phi^-}$, $\ket{\Psi^+}$
et $-\ket{\Psi^-}$.

\subsubsection{Matrice de densité de l'état de sortie}

Pour $\ket{\Psi}=\ket{\beta^{00}}=\frac{1}{\sqrt2}\ket{00}+\frac{1}{\sqrt2}\ket{11}$ :
\[
  \rho_{AB}=\ket{\Psi}\bra{\Psi}
  =\begin{pmatrix}\frac{1}{\sqrt2}\\0\\0\\\frac{1}{\sqrt2}\end{pmatrix}
   \begin{pmatrix}\frac{1}{\sqrt2}&0&0&\frac{1}{\sqrt2}\end{pmatrix}
  =\frac12\begin{pmatrix}1&0&0&1\\0&0&0&0\\0&0&0&0\\1&0&0&1\end{pmatrix}.
\]
Sous forme de ket-bra :
$\rho_{AB}=\frac12\big(\ket{00}\bra{00}+\ket{00}\bra{11}+\ket{11}\bra{00}+\ket{11}\bra{11}\big)$.
On vérifie $\operatorname{Tr}\rho_{AB}=\frac12(1+1)=1$ et
\[
  \rho_{AB}^2=\frac14\begin{pmatrix}2&0&0&2\\0&0&0&0\\0&0&0&0\\2&0&0&2\end{pmatrix}=\rho_{AB},
\]
donc $\operatorname{Tr}\rho_{AB}^2=1$ : le couple $(A,B)$ est dans un état \emph{pur}.

\begin{remarque}[Porte intriquante]
Les états d'entrée $\ket{ij}$ sont séparables ($D=0$) ; les sorties sont intriquées
($D=\pm\frac12$). Aucune porte agissant séparément sur $A$ et sur $B$ ($U_A\otimes U_B$,
section~\ref{sec:op2}) ne peut créer de l'intrication à partir d'un produit, puisque
$(U_A\otimes U_B)(\ket{u}\otimes\ket{v})=U_A\ket{u}\otimes U_B\ket{v}$ reste un produit.
Une porte à deux qubits comme le CNOT est indispensable.
\end{remarque}

% ---------------------------------------------------------------------
\subsection{Trace partielle : l'état d'un qubit pris seul}\label{sec:trace}
% ---------------------------------------------------------------------

Un état intriqué n'a pas de ket pour $A$ seul (section~\ref{sec:separable}). On peut pourtant
décrire $A$ par une matrice de densité, obtenue en «~oubliant~» $B$.

\begin{definition}{Trace partielle}
Soit $\rho_{AB}$ la matrice de densité du couple $(A,B)$. L'état du qubit $A$ pris seul est
$\rho_A=\operatorname{Tr}_B(\rho_{AB})$, celui de $B$ est $\rho_B=\operatorname{Tr}_A(\rho_{AB})$. La
trace partielle est linéaire et remplace la partie tracée de chaque terme par sa trace :
\[
  \operatorname{Tr}_B\big(\ket{a}\bra{a'}\otimes\ket{b}\bra{b'}\big)=\ket{a}\bra{a'}\;\braket{b'}{b},
  \qquad
  \operatorname{Tr}_A\big(\ket{a}\bra{a'}\otimes\ket{b}\bra{b'}\big)=\braket{a'}{a}\;\ket{b}\bra{b'}.
\]
\end{definition}

\paragraph{En matrices.} On découpe $\rho_{AB}$ (matrice $4\times4$, lignes et colonnes ordonnées
$00,01,10,11$) en quatre blocs $2\times2$ :
\[
  \rho_{AB}=\begin{pmatrix}R_{00}&R_{01}\\R_{10}&R_{11}\end{pmatrix}
  \quad\Longrightarrow\quad
  \rho_A=\begin{pmatrix}\operatorname{Tr}R_{00}&\operatorname{Tr}R_{01}\\
                        \operatorname{Tr}R_{10}&\operatorname{Tr}R_{11}\end{pmatrix},
  \qquad
  \rho_B=R_{00}+R_{11}.
\]
Autrement dit $(\rho_A)_{ij}=\rho_{(i0),(j0)}+\rho_{(i1),(j1)}$ : $\rho_A$ est la matrice des traces
des blocs, et $\rho_B$ la somme des deux blocs diagonaux.

\begin{remarque}[Pourquoi cette définition]
Toutes les prédictions qui portent sur $A$ seul viennent d'opérateurs de la forme $O_A\otimes I$, et
\[
  \operatorname{Tr}\big(\rho_{AB}\,(O_A\otimes I)\big)=\operatorname{Tr}\big(\rho_A\,O_A\big)
\]
(on le vérifie sur $\ket{a}\bra{a'}\otimes\ket{b}\bra{b'}$, puis par linéarité). La matrice $\rho_A$
contient donc exactement ce que l'on peut prédire sur $A$ sans regarder $B$.
\end{remarque}

\paragraph{Exemple 1 : état séparable $\ket{0}\otimes\ket{+}$.} Ici
$\rho_{AB}=\ket{0}\bra{0}\otimes\ket{+}\bra{+}$, donc
\[
  \rho_A=\ket{0}\bra{0}\,\braket{+}{+}=\ket{0}\bra{0},
  \qquad
  \rho_B=\braket{0}{0}\,\ket{+}\bra{+}=\ket{+}\bra{+}.
\]
Avec les matrices,
\[
  \rho_{AB}=\frac12\begin{pmatrix}1&1&0&0\\1&1&0&0\\0&0&0&0\\0&0&0&0\end{pmatrix},
  \qquad
  R_{00}=\frac12\begin{pmatrix}1&1\\1&1\end{pmatrix},\quad R_{01}=R_{10}=R_{11}=0,
\]
d'où
\[
  \rho_A=\begin{pmatrix}\operatorname{Tr}R_{00}&0\\0&0\end{pmatrix}=\begin{pmatrix}1&0\\0&0\end{pmatrix},
  \qquad
  \rho_B=R_{00}+R_{11}=\frac12\begin{pmatrix}1&1\\1&1\end{pmatrix}.
\]
Un qubit d'un état séparable garde son propre état pur.

\paragraph{Exemple 2 : état de Bell $\ket{\Phi^+}$.} On écrit
$\rho_{AB}=\ket{\Phi^+}\bra{\Phi^+}=\frac12\big(\ket{00}\bra{00}+\ket{00}\bra{11}+\ket{11}\bra{00}
+\ket{11}\bra{11}\big)$. Chaque terme est un produit d'opérateurs à un qubit, par exemple
$\ket{00}\bra{11}=\ket{0}\bra{1}\otimes\ket{0}\bra{1}$, donc
\begin{align*}
  \rho_A=\operatorname{Tr}_B(\rho_{AB})
  &=\tfrac12\Big(\ket{0}\bra{0}\,\braket{0}{0}+\ket{0}\bra{1}\,\braket{1}{0}
      +\ket{1}\bra{0}\,\braket{0}{1}+\ket{1}\bra{1}\,\braket{1}{1}\Big)\\
  &=\tfrac12\Big(\ket{0}\bra{0}\cdot1+\ket{0}\bra{1}\cdot0+\ket{1}\bra{0}\cdot0+\ket{1}\bra{1}\cdot1\Big)
   =\tfrac12\big(\ket{0}\bra{0}+\ket{1}\bra{1}\big)=\tfrac12\,\mathrm{Id}.
\end{align*}
Par les blocs,
\[
  \rho_{AB}=\frac12\begin{pmatrix}1&0&0&1\\0&0&0&0\\0&0&0&0\\1&0&0&1\end{pmatrix},
  \qquad
  R_{00}=\frac12\begin{pmatrix}1&0\\0&0\end{pmatrix},\quad
  R_{11}=\frac12\begin{pmatrix}0&0\\0&1\end{pmatrix},\quad
  \operatorname{Tr}R_{01}=\operatorname{Tr}R_{10}=0,
\]
d'où $\rho_A=\frac12\mathrm{Id}$ et $\rho_B=R_{00}+R_{11}=\frac12\mathrm{Id}$ : le même résultat.

Le qubit $A$ est donc dans l'état \emph{complètement mixte} ($\operatorname{Tr}\rho_A^2=\frac14+\frac14=\frac12$),
au centre de la boule de Bloch (section~\ref{sec:bloch}), alors que le couple est dans un état pur
($\operatorname{Tr}\rho_{AB}^2=1$, section~\ref{sec:circuit-bell}). On connaît parfaitement l'état du
couple mais pas celui d'un qubit pris seul : l'information est entièrement contenue dans les
corrélations entre $A$ et $B$. Mesuré seul, $A$ donne $0$ ou $1$ avec la probabilité $\frac12$, dans
toutes les bases.

\paragraph{Formule pratique pour un état pur.} Pour $\ket{\Psi}=\sum_{ij}c_{ij}\ket{ij}$, notons
$C=\begin{pmatrix}c_{00}&c_{01}\\c_{10}&c_{11}\end{pmatrix}$ la matrice $2\times2$ des coefficients
(ligne $i$ : état de $A$, colonne $j$ : état de $B$). Comme $\rho_{(ij),(i'j')}=c_{ij}\,c_{i'j'}^*$,
\[
  (\rho_A)_{ii'}=\sum_jc_{ij}\,c_{i'j}^*,\qquad\text{soit}\qquad \rho_A=C\,C^\dagger .
\]
\emph{Exemple 3.} Pour $\ket{\Psi}=\frac12\ket{00}+\frac12\ket{01}+\frac{1}{\sqrt2}\ket{11}$ (exemple~2
de la section~\ref{sec:mesure2}), $C=\begin{pmatrix}\frac12&\frac12\\0&\frac{1}{\sqrt2}\end{pmatrix}$ et
\[
  \rho_A=CC^\dagger=\begin{pmatrix}\frac14+\frac14&0+\frac{1}{2\sqrt2}\\[0.4ex]\frac{1}{2\sqrt2}&\frac12\end{pmatrix}
  =\begin{pmatrix}\frac12&\frac{\sqrt2}{4}\\[0.4ex]\frac{\sqrt2}{4}&\frac12\end{pmatrix},
  \qquad
  \operatorname{Tr}\rho_A^2=\tfrac14+\tfrac14+2\cdot\tfrac18=\tfrac34 .
\]
Le qubit $A$ est mixte, mais pas complètement : cet état est intriqué sans être un état de Bell.

\begin{resultat}{Pureté d'un qubit et séparabilité}
Pour un état pur $\ket{\Psi}=\sum c_{ij}\ket{ij}$ de deux qubits, avec $D=c_{00}c_{11}-c_{01}c_{10}$,
\[
  \operatorname{Tr}\rho_A^2=\operatorname{Tr}\rho_B^2=1-2|D|^2 .
\]
Donc $\ket{\Psi}$ est séparable ($D=0$) si et seulement si $\rho_A$ est pur. Plus $|D|$ est grand,
plus $\rho_A$ est mixte, jusqu'à $|D|=\frac12$ et $\rho_A=\mathrm{Id}/2$ pour les états de Bell.
\end{resultat}

\emph{Démonstration.} La matrice $\rho_A=CC^\dagger$ a pour trace $\sum|c_{ij}|^2=1$ et pour
déterminant $\det C\,\overline{\det C}=|D|^2$. Si $\lambda_1,\lambda_2$ sont ses valeurs propres,
$\lambda_1+\lambda_2=1$ et $\operatorname{Tr}\rho_A^2=\lambda_1^2+\lambda_2^2=1-2\lambda_1\lambda_2
=1-2|D|^2$ ; le calcul est le même pour $\rho_B$. \hfill$\square$

Comme $\operatorname{Tr}\rho_A^2\geq\frac12$ pour un qubit, on en déduit $|D|\leq\frac12$. Avec
$\operatorname{Tr}\rho_A^2=\frac{1+|\vec r|^2}{2}$ (section~\ref{sec:bloch}), le vecteur de Bloch de $A$
a pour norme $|\vec r|=\sqrt{1-4|D|^2}$ : il part de la sphère pour un état séparable et se rapproche
du centre quand l'intrication croît (figure~\ref{fig:trace-bloch}).

\begin{center}
\renewcommand{\arraystretch}{1.5}
\begin{tabular}{@{}lcccl@{}}
\toprule
\textbf{État du couple} & $D$ & $\operatorname{Tr}\rho_A^2$ & $|\vec r|$ & \textbf{Qubit $A$} \\
\midrule
$\ket{0}\otimes\ket{+}$ & $0$ & $1$ & $1$ & pur, sur la sphère \\
$\frac12\ket{00}+\frac12\ket{01}+\frac{1}{\sqrt2}\ket{11}$ & $\frac{\sqrt2}{4}$ & $\frac34$
  & $\frac{1}{\sqrt2}$ & mixte, dans la boule \\
$\ket{\Phi^+}$ (comme les trois autres états de Bell) & $\pm\frac12$ & $\frac12$ & $0$
  & complètement mixte, au centre \\
\bottomrule
\end{tabular}
\end{center}

\begin{figure}[!ht]
\centering
\adjustbox{max width=\linewidth}{%
\begin{tikzpicture}[>=Stealth, font=\small]
  \def\R{2.3}
  \fill[insagris!12] (0,0) circle (\R);
  \draw[insagris, dotted] (-\R,0) -- (\R,0);
  \draw[insagris, dotted] (0,-\R) -- (0,\R);
  \draw[insagris, dashed] (0,0) circle ({\R*0.70711});
  \draw[insarouge, very thick] (0,0) circle (\R);
  \node[font=\footnotesize, text=insagris, anchor=south west] at (\R,0) {$x$};
  \node[font=\footnotesize, text=insagris, anchor=south west] at (0.03,\R) {$z$};
  % vecteurs de Bloch de rho_A pour trois états du couple
  \fill[insarouge] (0,\R) circle (2.4pt);
  \fill[insagris!80!black] ({\R*0.70711},0) circle (2.4pt);
  \fill[insagris!80!black] (0,0) circle (2.4pt);
  \draw[insagris] (0,\R) -- (3.1,\R);
  \node[anchor=west, font=\footnotesize, align=left] at (3.1,\R)
    {$\ket{0}\otimes\ket{+}$ : $\rho_A=\ket{0}\bra{0}$\\ $D=0$, \ $|\vec r|=1$ (état pur)};
  \draw[insagris] ({\R*0.70711},0) -- (3.1,0);
  \node[anchor=west, font=\footnotesize, align=left] at (3.1,0)
    {$\frac12\ket{00}+\frac12\ket{01}+\frac{1}{\sqrt2}\ket{11}$\\
     $D=\frac{\sqrt2}{4}$, \ $|\vec r|=\frac{1}{\sqrt2}$ (mixte)};
  \draw[insagris] (0,0) -- (3.1,-\R);
  \node[anchor=west, font=\footnotesize, align=left] at (3.1,-\R)
    {$\ket{\Phi^+}$ : $\rho_A=\mathrm{Id}/2$\\ $D=\frac12$, \ $\vec r=\vec 0$ (centre)};
\end{tikzpicture}}
\caption{Le vecteur de Bloch de $\rho_A$ (boule de Bloch vue en coupe) pour trois états du couple
$(A,B)$. Le cercle pointillé est la sphère de rayon $1/\sqrt2$. Plus l'état est intriqué, plus le
point s'éloigne de la sphère des états purs : $|\vec r|=\sqrt{1-4|D|^2}$.}
\label{fig:trace-bloch}
\end{figure}

% ---------------------------------------------------------------------
\subsection{Information : variables aléatoires et états quantiques}\label{sec:info}
% ---------------------------------------------------------------------

\begin{definition}{Entropie, information mutuelle, entropie conditionnelle}
Soit $A$ une variable aléatoire de loi $p(a)$. Son \emph{entropie de Shannon}, en bits, est
\[
  H(A)=-\sum_a p(a)\log_2p(a)\qquad(\text{avec }0\log_20=0).
\]
Elle mesure l'incertitude moyenne sur $A$ : une pièce équilibrée a $H=1$ bit. Pour un couple
$(A,B)$ de loi $p(a,b)$, $H(A,B)=-\sum_{a,b}p(a,b)\log_2p(a,b)$, et
\begin{align*}
  I(A;B)&=H(A)+H(B)-H(A,B) &&\text{(information mutuelle)},\\
  H(A|B)&=H(A,B)-H(B)=H(A)-I(A;B) &&\text{(entropie conditionnelle)}.
\end{align*}
\end{definition}

Les deux écritures de $H(A|B)$ sont égales : la définition de $I(A;B)$ donne
$H(A,B)=H(A)+H(B)-I(A;B)$, donc $H(A,B)-H(B)=H(A)-I(A;B)$. L'information mutuelle $I(A;B)$
mesure ce que la connaissance de $B$ apprend sur $A$ ; $H(A|B)$ est l'incertitude qui reste sur
$A$ quand on connaît $B$.

\subsubsection{Deux variables aléatoires indépendantes}

Alice et Bob lancent chacun une pièce équilibrée, sans lien entre les deux lancers :
$p(a,b)=\frac14$ pour les quatre couples. Chaque pièce donne $0$ ou $1$ avec la probabilité $\frac12$, donc
\begin{align*}
  H(A)&=-\Big(\tfrac12\log_2\tfrac12+\tfrac12\log_2\tfrac12\Big)=1\text{ bit},\\
  H(B)&=1\text{ bit (même calcul)},\\
  H(A,B)&=-4\cdot\tfrac14\log_2\tfrac14=2\text{ bits},\\
  I(A;B)&=H(A)+H(B)-H(A,B)=1+1-2=0,\\
  H(A|B)&=H(A)-I(A;B)=1-0=H(A)=1\text{ bit}.
\end{align*}
Connaître $B$ n'apprend rien sur $A$.

\subsubsection{Deux variables solidaires}

Les deux pièces tombent toujours du même côté : $p(0,0)=p(1,1)=\frac12$ et
$p(0,1)=p(1,0)=0$. Les lois de chaque pièce sont inchangées ($p(a{=}0)=p(0,0)+p(0,1)=\frac12$,
etc.), donc $H(A)=H(B)=1$ bit, mais
\begin{align*}
  H(A,B)&=-2\cdot\tfrac12\log_2\tfrac12-2\cdot0\log_20=1\text{ bit},\\
  I(A;B)&=1+1-1=1\text{ bit},\qquad H(A|B)=1-1=0.
\end{align*}
Connaître $B$ détermine $A$ : plus aucune incertitude ne reste.

\subsubsection{Deux qubits dans un état de Bell}

Pour un système quantique, on remplace la loi de probabilité par la matrice de densité.

\begin{definition}{Entropie de von Neumann}
L'\emph{entropie de von Neumann} d'un état $\rho$, de valeurs propres $\lambda_i$, est
\[
  S(\rho)=-\operatorname{Tr}(\rho\log_2\rho)=-\sum_i\lambda_i\log_2\lambda_i .
\]
Elle vaut $0$ pour un état pur et $1$ bit pour le qubit complètement mélangé $\mathrm{Id}/2$.
En quantique, $H(A)=S(\rho_A)$, $H(B)=S(\rho_B)$ et $H(A,B)=S(\rho_{AB})$ ; les définitions de
$I(A;B)$ et de $H(A|B)$ sont les mêmes.
\end{definition}

On prend l'état de Bell $\ket{\Phi^+}$ produit par le circuit de la
section~\ref{sec:circuit-bell}. Le couple est dans un état pur,
$\rho_{AB}=\ket{\Phi^+}\bra{\Phi^+}$ ; chaque qubit pris seul est dans l'état
$\rho_A=\rho_B=\frac12\mathrm{Id}$ (trace partielle, section~\ref{sec:trace}).

Calcul des entropies :
\begin{itemize}[nosep]
  \item $\rho_{AB}$ est pur : ses valeurs propres sont $1,0,0,0$, donc
        $H(A,B)=S(\rho_{AB})=-1\cdot\log_21=0$ ;
  \item $\rho_A$ et $\rho_B$ ont pour valeurs propres $\frac12,\frac12$, donc
        $H(A)=H(B)=-\big(\frac12\log_2\frac12+\frac12\log_2\frac12\big)=1$ bit.
\end{itemize}
D'où
\begin{align*}
  I(A;B)&=H(A)+H(B)-H(A,B)=1+1-0=2\text{ bits},\\
  H(A|B)&=H(A,B)-H(B)=0-1=-1\text{ bit}\qquad(=H(A)-I(A;B)=1-2).
\end{align*}

\begin{remarque}[Une entropie conditionnelle négative]
Pour des variables classiques, $p(a,b)\leq p(b)$, donc $-\log_2p(a,b)\geq-\log_2p(b)$ et
$H(A,B)\geq H(B)$ : on a toujours $H(A|B)\geq0$, et $I(A;B)\leq H(A)$. Ici $H(A|B)=-1$ et
$I(A;B)=2>H(A)$ : ces valeurs n'ont pas d'équivalent classique, c'est une signature de
l'intrication. Le résultat est surprenant mais mathématiquement correct.
\end{remarque}

\begin{center}
\renewcommand{\arraystretch}{1.4}
\begin{tabular}{@{}lccc@{}}
\toprule
 & indépendantes & solidaires & état de Bell \\
\midrule
$H(A)$ & $1$ & $1$ & $1$ \\
$H(B)$ & $1$ & $1$ & $1$ \\
$H(A,B)$ & $2$ & $1$ & $0$ \\
$I(A;B)$ & $0$ & $1$ & $2$ \\
$H(A|B)$ & $1$ & $0$ & $-1$ \\
\bottomrule
\end{tabular}
\end{center}
(en bits ; les deux premières colonnes sont classiques, la troisième est quantique).

\begin{figure}[!ht]
\centering
\adjustbox{max width=\linewidth}{%
\begin{tikzpicture}[>=Stealth, font=\small]
  \def\bw{0.8}
  % axe des ordonnées (bits)
  \draw[->] (0,-1.5) -- (0,2.8) node[above] {bits};
  \foreach \v in {-1,0,1,2} {
    \draw (-0.07,\v) -- (0.07,\v);
    \node[left, font=\footnotesize] at (-0.1,\v) {$\v$};
  }
  \draw[insagris!60] (0,0) -- (11.4,0);
  % limite classique de I(A;B) : H(A) = 1
  \draw[dashed, insarouge] (0,1) -- (11.4,1);
  \node[anchor=west, font=\scriptsize, text=insarouge, align=left] at (11.4,1)
    {limite classique\\ $I(A;B)\leq H(A)=1$};
  \node[anchor=west, font=\scriptsize, text=insagris, align=left] at (11.4,0)
    {limite classique\\ $H(A|B)\geq0$};
  % barres : (groupe, H(A,B), I(A;B), H(A|B))
  \foreach \g/\hab/\mi/\hc in {0/2/0/1, 1/1/1/0, 2/0/2/-1} {
    \pgfmathsetmacro{\x}{0.9+3.7*\g}
    \fill[insagris!65] (\x,0) rectangle (\x+\bw,\hab);
    \fill[insarouge] (\x+\bw+0.1,0) rectangle (\x+2*\bw+0.1,\mi);
    \fill[black!80] (\x+2*\bw+0.2,0) rectangle (\x+3*\bw+0.2,\hc);
    \node[above, font=\footnotesize] at (\x+0.5*\bw,{max(\hab,0)}) {$\hab$};
    \node[above, font=\footnotesize] at (\x+1.5*\bw+0.1,{max(\mi,0)}) {$\mi$};
    \ifnum\hc<0
      \node[below, font=\footnotesize] at (\x+2.5*\bw+0.2,\hc) {$\hc$};
    \else
      \node[above, font=\footnotesize] at (\x+2.5*\bw+0.2,\hc) {$\hc$};
    \fi
  }
  \node[align=center, anchor=north] at (2.15,-1.65) {variables\\ indépendantes};
  \node[align=center, anchor=north] at (5.85,-1.65) {variables\\ solidaires};
  \node[align=center, anchor=north] at (9.55,-1.65) {état de Bell\\ $\ket{\Phi^+}$};
  % légende
  \fill[insagris!65] (0.9,3.35) rectangle (1.2,3.65);
  \node[anchor=west, font=\footnotesize] at (1.25,3.5) {$H(A,B)$};
  \fill[insarouge] (3.6,3.35) rectangle (3.9,3.65);
  \node[anchor=west, font=\footnotesize] at (3.95,3.5) {$I(A;B)$};
  \fill[black!80] (6.3,3.35) rectangle (6.6,3.65);
  \node[anchor=west, font=\footnotesize] at (6.65,3.5) {$H(A|B)$};
\end{tikzpicture}}
\caption{Entropies (en bits) pour deux pièces indépendantes, deux pièces solidaires et un état de Bell.
Dans les trois cas $H(A)=H(B)=1$. Les deux premiers groupes respectent les bornes classiques
$I(A;B)\leq H(A)$ et $H(A|B)\geq0$ ; l'état de Bell les dépasse ($I=2$, $H(A|B)=-1$).}
\label{fig:entropies}
\end{figure}

% ---------------------------------------------------------------------
\subsection{Téléportation quantique}\label{sec:teleportation}
% ---------------------------------------------------------------------

\emph{Problème.} Alice détient un qubit $Q$ dans un état qu'elle ne connaît pas,
$\ket{\varphi}=\alpha\ket{0}+\beta\ket{1}$ avec $|\alpha|^2+|\beta|^2=1$, et elle veut que Bob
obtienne cet état sur l'un de ses qubits. Elle ne peut pas mesurer $Q$ pour connaître
$\alpha$ et $\beta$ (une mesure ne donne qu'un bit et détruit l'état), elle ne peut pas le
copier (théorème de non-clonage, remarque en fin de section) et elle n'envoie pas le qubit
lui-même. Le protocole est dû à \textcite{bennett1993}.

\begin{definition}{Protocole de téléportation quantique}
\textbf{Ressources.}
\begin{itemize}[nosep]
  \item Un \emph{e-bit} : Alice et Bob partagent une paire de qubits $(A,B)$ dans l'état de Bell
        $\ket{\Phi^+}_{BA}=\frac{1}{\sqrt2}\big(\ket{00}+\ket{11}\big)$ ; Alice détient $A$,
        Bob détient $B$.
  \item Le qubit $Q$ à téléporter, chez Alice.
  \item Un canal classique (deux bits).
\end{itemize}
Le système est écrit $(B,A,Q)$ : $\ket{b,a,q}=\ket{b}_B\otimes\ket{a}_A\otimes\ket{q}_Q$.

\textbf{Étapes.}
\begin{enumerate}[nosep]
  \item Alice applique un CNOT de contrôle $Q$ et de cible $A$.
  \item Alice applique une porte $H$ sur $Q$.
  \item Alice mesure $A$ et $Q$ dans la base de calcul, obtient deux bits $(a,q)$ et les envoie
        à Bob.
  \item Bob applique $X$ si $a=1$, puis $Z$ si $q=1$, à son qubit $B$.
\end{enumerate}
\end{definition}

La figure~\ref{fig:tele-schema} montre l'organisation du protocole et la
figure~\ref{fig:teleportation} son circuit.

\begin{figure}[!ht]
\centering
\adjustbox{max width=\linewidth}{%
\begin{tikzpicture}[>=Stealth, font=\small,
    qb/.style={circle, draw=insarouge, fill=insarouge!8, thick, minimum size=8mm, inner sep=0pt},
    zone/.style={draw=insagris, rounded corners=4pt, thick, dashed, inner sep=8pt},
    num/.style={circle, fill=insarouge, text=white, font=\footnotesize\bfseries, inner sep=0pt,
                minimum size=4.6mm},
    etape/.style={anchor=west, font=\footnotesize, text=insarouge},
    onde/.style={thick, insarouge, decorate,
                 decoration={snake, amplitude=1.3pt, segment length=7pt}}]
  % Alice : Q (état inconnu) et A
  \node[qb] (Q) at (0,1.3) {$Q$};
  \node[qb] (A) at (0,-0.5) {$A$};
  \node[zone, fit=(Q)(A), label={[text=insagris]above:Alice}] (alice) {};
  \node[anchor=east] at (-0.95,1.3) {$\ket{\varphi}=\alpha\ket{0}+\beta\ket{1}$};
  \draw[insarouge, thick] (Q) -- (A);
  % Bob : B
  \node[qb] (B) at (9,-0.5) {$B$};
  \node[zone, fit=(B), label={[text=insagris]above:Bob}] (bob) {};
  % source de la paire de Bell
  \node[draw=insarouge, fill=insarouge!10, rounded corners=2pt, thick, minimum width=3.2cm,
        minimum height=0.9cm] (src) at (4.5,-2.3) {source $\ket{\Phi^+}_{BA}$};
  \draw[onde] (src.west) -| (A.south);
  \draw[onde] (src.east) -| (B.south);
  % étapes d'Alice : pastilles sur le bord droit de sa zone
  \node[num] (n1) at (alice.east |- Q) {1};
  \node[num] (n2) at (alice.east) {2};
  \node[num] (n3) at (alice.east |- A) {3};
  \node[etape] at (n1.east) {\ CNOT de $Q$ vers $A$};
  \node[etape] at (n2.east) {\ porte $H$ sur $Q$};
  % canal classique
  \draw[double, double distance=1.2pt, ->, insagris] (n3.east) -- (bob.west);
  \node[above, font=\footnotesize] at (4.6,-0.45) {deux bits classiques $(a,q)$};
  \node[below, font=\footnotesize, text=insarouge] at (4.6,-0.55) {mesure de $A$ et $Q$ (étape 3)};
  % correction de Bob
  \node[num] (n4) at (bob.east) {4};
  \node[etape, align=left] at (n4.east) {\ $X$ si $a=1$, puis $Z$ si $q=1$};
  \node[anchor=west, font=\footnotesize, align=left] at ($(n4.east)+(0.05,-0.8)$)
    {\ $B$ est alors dans l'état $\ket{\varphi}$};
\end{tikzpicture}}
\caption{Organisation de la téléportation. Alice possède le qubit $Q$ à téléporter et la moitié $A$
d'une paire de Bell, Bob possède l'autre moitié $B$. (1)~CNOT de $Q$ vers $A$, (2)~porte $H$ sur $Q$,
(3)~mesure de $A$ et $Q$, dont les résultats $(a,q)$ partent par un canal classique, (4)~corrections de
Bob. Aucun qubit ne voyage entre Alice et Bob.}
\label{fig:tele-schema}
\end{figure}

\begin{figure}[!ht]
\centering
\begin{adjustbox}{max width=\linewidth}
\begin{tikzpicture}[>=Stealth, font=\small]
  % fils quantiques : Q en haut, A au milieu, B en bas
  \draw (0,2.4) -- (5.3,2.4);
  \draw (0,1.2) -- (5.3,1.2);
  \draw (0,0) -- (10.6,0);
  \node[anchor=east] at (0,2.4) {$\ket{\varphi}_Q$};
  \node[anchor=south west, text=insagris, font=\footnotesize] at (0.05,1.2) {$A$};
  \node[anchor=south west, text=insagris, font=\footnotesize] at (0.05,0) {$B$};
  % paire de Bell partagée
  \draw[insagris, thick] (-0.1,1.5) -- (-0.25,1.5) -- (-0.25,-0.3) -- (-0.1,-0.3);
  \node[anchor=east] at (-0.35,0.6) {$\ket{\Phi^+}_{BA}$};
  % CNOT : contrôle Q, cible A
  \draw[insarouge, thick] (2.6,2.4) -- (2.6,1.2);
  \pic at (2.6,2.4) {ctl};
  \pic at (2.6,1.2) {cible};
  % Hadamard sur Q
  \node[porte] at (4.2,2.4) {$H$};
  % mesures d'Alice
  \pic at (5.6,2.4) {mesure};
  \pic at (5.6,1.2) {mesure};
  % fils classiques (doubles)
  \draw[double, double distance=1.2pt, insagris] (5.9,1.2) -- (7.6,1.2) -- (7.6,0.35);
  \draw[double, double distance=1.2pt, insagris] (5.9,2.4) -- (9.6,2.4) -- (9.6,0.35);
  \node[above, font=\footnotesize] at (6.5,1.2) {$a$};
  \node[above, font=\footnotesize] at (6.5,2.4) {$q$};
  % corrections de Bob
  \node[porte] at (7.6,0) {$X$};
  \node[porte] at (9.6,0) {$Z$};
  \node[anchor=west] at (10.6,0) {$\ket{\varphi}_B$};
  % états intermédiaires
  \foreach \x in {1.4,3.4,4.9}
    \draw[dashed, insagris!70] (\x,2.8) -- (\x,-0.5);
  \node[anchor=north, font=\footnotesize] at (1.4,-0.5) {$\ket{\Pi_0}$};
  \node[anchor=north, font=\footnotesize] at (3.4,-0.5) {$\ket{\Pi_1}$};
  \node[anchor=north, font=\footnotesize] at (4.9,-0.5) {$\ket{\Pi_2}$};
\end{tikzpicture}
\end{adjustbox}
\caption{Circuit de téléportation. Alice détient $Q$ (dans $\ket{\varphi}=\alpha\ket{0}+\beta\ket{1}$)
et $A$, Bob détient $B$ ; $A$ et $B$ forment la paire $\ket{\Phi^+}_{BA}$. Les traits doubles
sont des fils \emph{classiques} qui portent les résultats $a$ (mesure de $A$) et $q$ (mesure de
$Q$) ; $X$ est appliqué si $a=1$, $Z$ si $q=1$. Le ket s'écrit $\ket{b,a,q}$ : le fil du haut
($Q$) est le dernier chiffre.}
\label{fig:teleportation}
\end{figure}

\subsubsection{Calcul pas à pas}

\paragraph{État initial $\ket{\Pi_0}$.} Alice et Bob partagent la paire, et $Q$ est indépendant
d'elle :
\begin{align*}
  \ket{\Pi_0}&=\ket{\Phi^+}_{BA}\otimes\ket{\varphi}_Q
    =\frac{1}{\sqrt2}\big(\ket{00}+\ket{11}\big)\otimes\big(\alpha\ket{0}+\beta\ket{1}\big)\\
  &=\frac{1}{\sqrt2}\Big(\alpha\ket{00}\ket{0}+\beta\ket{00}\ket{1}+\alpha\ket{11}\ket{0}+\beta\ket{11}\ket{1}\Big)\\
  &=\frac{\alpha\ket{000}+\beta\ket{001}+\alpha\ket{110}+\beta\ket{111}}{\sqrt2}.
\end{align*}

\paragraph{Après le CNOT $Q\to A$ : $\ket{\Pi_1}$.} Le CNOT (contrôle $Q$, dernier chiffre ;
cible $A$, chiffre du milieu) agit ainsi : $\ket{b,a,q}\to\ket{b,\,a\oplus q,\,q}$. Terme par terme :
\begin{center}
\renewcommand{\arraystretch}{1.3}
\begin{tabular}{@{}cccc@{}}
\toprule
Terme & $q$ & Effet sur $A$ & Résultat \\
\midrule
$\alpha\ket{000}$ & $0$ & aucun & $\alpha\ket{000}$ \\
$\beta\ket{001}$ & $1$ & $0\to1$ & $\beta\ket{011}$ \\
$\alpha\ket{110}$ & $0$ & aucun & $\alpha\ket{110}$ \\
$\beta\ket{111}$ & $1$ & $1\to0$ & $\beta\ket{101}$ \\
\bottomrule
\end{tabular}
\end{center}
donc
\[
  \ket{\Pi_1}=\frac{\alpha\ket{000}+\beta\ket{011}+\alpha\ket{110}+\beta\ket{101}}{\sqrt2}.
\]
On isole ensuite le dernier chiffre, celui de $Q$, qui va changer à l'étape suivante
(le premier ket est celui du couple $\ket{b,a}$) :
\[
  \ket{\Pi_1}=\frac{\alpha\ket{00}\otimes\ket{0}+\beta\ket{01}\otimes\ket{1}
    +\alpha\ket{11}\otimes\ket{0}+\beta\ket{10}\otimes\ket{1}}{\sqrt2}.
\]

\paragraph{Après la porte $H$ sur $Q$ : $\ket{\Pi_2}$.} On remplace $\ket{0}$ par
$H\ket{0}=\ket{+}$ et $\ket{1}$ par $H\ket{1}=\ket{-}$ dans le facteur $Q$ :
\[
  \ket{\Pi_2}=\frac{1}{\sqrt2}\Big(\alpha\ket{00}\otimes\ket{+}+\beta\ket{01}\otimes\ket{-}
    +\alpha\ket{11}\otimes\ket{+}+\beta\ket{10}\otimes\ket{-}\Big).
\]
Avec $\ket{\pm}=\frac{1}{\sqrt2}(\ket{0}\pm\ket{1})$, le préfacteur devient
$\frac{1}{\sqrt2}\cdot\frac{1}{\sqrt2}=\frac12$ :
\begin{align*}
  \ket{\Pi_2}&=\frac12\Big(\alpha\ket{00}\big(\ket{0}+\ket{1}\big)+\beta\ket{01}\big(\ket{0}-\ket{1}\big)
     +\alpha\ket{11}\big(\ket{0}+\ket{1}\big)+\beta\ket{10}\big(\ket{0}-\ket{1}\big)\Big)\\
  &=\frac12\Big(\alpha\ket{000}+\alpha\ket{001}+\beta\ket{010}-\beta\ket{011}
     +\alpha\ket{110}+\alpha\ket{111}+\beta\ket{100}-\beta\ket{101}\Big).
\end{align*}

\paragraph{Regroupement selon le résultat d'Alice.} Alice mesurera $A$ et $Q$, donc les deux
derniers chiffres. On regroupe les huit termes qui ont les mêmes chiffres $(a,q)$, en mettant
le ket de Bob en premier :
\begin{center}
\renewcommand{\arraystretch}{1.3}
\begin{tabular}{@{}cll@{}}
\toprule
$(a,q)$ & Termes de $\ket{\Pi_2}$ (facteur $\frac12$ omis) & Ket de Bob \\
\midrule
$(0,0)$ & $\alpha\ket{000}+\beta\ket{100}$ & $(\alpha\ket{0}+\beta\ket{1})\ket{00}$ \\
$(0,1)$ & $\alpha\ket{001}-\beta\ket{101}$ & $(\alpha\ket{0}-\beta\ket{1})\ket{01}$ \\
$(1,0)$ & $\beta\ket{010}+\alpha\ket{110}$ & $(\beta\ket{0}+\alpha\ket{1})\ket{10}$ \\
$(1,1)$ & $-\beta\ket{011}+\alpha\ket{111}$ & $(-\beta\ket{0}+\alpha\ket{1})\ket{11}$ \\
\bottomrule
\end{tabular}
\end{center}
d'où
\[
  \ket{\Pi_2}=\frac12\Big[\big(\alpha\ket{0}+\beta\ket{1}\big)\ket{00}
    +\big(\alpha\ket{0}-\beta\ket{1}\big)\ket{01}
    +\big(\beta\ket{0}+\alpha\ket{1}\big)\ket{10}
    +\big(\alpha\ket{1}-\beta\ket{0}\big)\ket{11}\Big].
\]
Dans chaque terme, le ket de gauche est l'état de $B$ et le ket de droite $\ket{a,q}$ celui du
couple $(A,Q)$ d'Alice.

\subsubsection{Mesure d'Alice et corrections de Bob}

Alice mesure $A$ et $Q$. Comme à la section~\ref{sec:mesure2}, la probabilité d'un résultat
$(a,q)$ est le carré de la norme du terme correspondant, et l'état de Bob devient le ket
associé (déjà normé). Ici
\[
  P(a,q)=\Big\|\tfrac12\ket{\phi_{aq}}\Big\|^2=\tfrac14\big(|\alpha|^2+|\beta|^2\big)=\tfrac14
  \quad\text{pour les quatre résultats,}
\]
qui sont donc équiprobables et \emph{indépendants} de $\alpha$ et $\beta$ : les deux bits
envoyés ne révèlent rien sur $\ket{\varphi}$. Les kets de Bob sont
$\ket{\phi_{00}}=\alpha\ket{0}+\beta\ket{1}$, $\ket{\phi_{01}}=\alpha\ket{0}-\beta\ket{1}$,
$\ket{\phi_{10}}=\beta\ket{0}+\alpha\ket{1}$ et $\ket{\phi_{11}}=-\beta\ket{0}+\alpha\ket{1}$.

\begin{center}
\renewcommand{\arraystretch}{1.5}
\begin{tabular}{@{}cccll@{}}
\toprule
Résultat $(a,q)$ & Probabilité & État de Bob & Lien avec $\ket{\varphi}$ & Correction de Bob \\
\midrule
$(0,0)$ & $\frac14$ & $\alpha\ket{0}+\beta\ket{1}$ & $\ket{\varphi}$ & $I$ (rien) \\
$(0,1)$ & $\frac14$ & $\alpha\ket{0}-\beta\ket{1}$ & $Z\ket{\varphi}$ & $Z$ \\
$(1,0)$ & $\frac14$ & $\beta\ket{0}+\alpha\ket{1}$ & $X\ket{\varphi}$ & $X$ \\
$(1,1)$ & $\frac14$ & $-\beta\ket{0}+\alpha\ket{1}$ & $XZ\ket{\varphi}$ & $X$ puis $Z$ (opérateur $ZX$) \\
\bottomrule
\end{tabular}
\end{center}

\paragraph{Vérification des quatre corrections.} Avec $X\ket{0}=\ket{1}$, $X\ket{1}=\ket{0}$,
$Z\ket{0}=\ket{0}$ et $Z\ket{1}=-\ket{1}$ :
\begin{align*}
  (0,0):\quad & I\big(\alpha\ket{0}+\beta\ket{1}\big)=\alpha\ket{0}+\beta\ket{1},\\
  (0,1):\quad & Z\big(\alpha\ket{0}-\beta\ket{1}\big)=\alpha\ket{0}-\beta Z\ket{1}=\alpha\ket{0}+\beta\ket{1},\\
  (1,0):\quad & X\big(\beta\ket{0}+\alpha\ket{1}\big)=\beta\ket{1}+\alpha\ket{0}=\alpha\ket{0}+\beta\ket{1},\\
  (1,1):\quad & X\big(-\beta\ket{0}+\alpha\ket{1}\big)=-\beta\ket{1}+\alpha\ket{0}=\alpha\ket{0}-\beta\ket{1},
     \ \text{puis}\ Z\big(\alpha\ket{0}-\beta\ket{1}\big)=\alpha\ket{0}+\beta\ket{1}.
\end{align*}
Dans le dernier cas, avec les matrices,
$ZX=\begin{pmatrix}1&0\\0&-1\end{pmatrix}\begin{pmatrix}0&1\\1&0\end{pmatrix}
=\begin{pmatrix}0&1\\-1&0\end{pmatrix}$ et
$ZX\begin{pmatrix}-\beta\\ \alpha\end{pmatrix}=\begin{pmatrix}\alpha\\ \beta\end{pmatrix}$.
L'ordre compte : l'état de Bob est $XZ\ket{\varphi}$ ($Z\ket{\varphi}=\alpha\ket{0}-\beta\ket{1}$,
puis $X$ donne $\alpha\ket{1}-\beta\ket{0}$), et on l'annule avec
$(XZ)^{-1}=Z^{-1}X^{-1}=ZX$ : on applique d'abord $X$, puis $Z$, comme dans le circuit.
Dans les quatre cas, Bob retrouve exactement $\ket{\varphi}=\alpha\ket{0}+\beta\ket{1}$.

\subsubsection{Ce que Bob sait avant de recevoir les bits}

Tant que Bob ignore $(a,q)$, son qubit est dans le mélange des quatre kets $\ket{\phi_{aq}}$, chacun
avec la probabilité $\frac14$ :
\begin{align*}
  \rho_B&=\frac14\sum_{a,q}\ket{\phi_{aq}}\bra{\phi_{aq}}\\
  &=\frac14\Bigg[\begin{pmatrix}|\alpha|^2&\alpha\beta^*\\ \alpha^*\beta&|\beta|^2\end{pmatrix}
    +\begin{pmatrix}|\alpha|^2&-\alpha\beta^*\\ -\alpha^*\beta&|\beta|^2\end{pmatrix}
    +\begin{pmatrix}|\beta|^2&\alpha^*\beta\\ \alpha\beta^*&|\alpha|^2\end{pmatrix}
    +\begin{pmatrix}|\beta|^2&-\alpha^*\beta\\ -\alpha\beta^*&|\alpha|^2\end{pmatrix}\Bigg]\\
  &=\frac14\begin{pmatrix}2(|\alpha|^2+|\beta|^2)&0\\0&2(|\alpha|^2+|\beta|^2)\end{pmatrix}
   =\frac12\begin{pmatrix}1&0\\0&1\end{pmatrix}=\frac12\,\mathrm{Id}.
\end{align*}
Sans les deux bits, Bob a donc un qubit complètement mélangé, qui ne dépend pas de $\alpha$ ni de
$\beta$ : il n'a reçu aucune information sur $\ket{\varphi}$. C'est l'envoi classique de $(a,q)$
qui permet de la retrouver, donc rien ne voyage plus vite que la lumière.

\begin{remarque}[Bilan de la téléportation]
\begin{itemize}[nosep]
  \item Ressources : $1$ e-bit (la paire de Bell) et $2$ bits classiques pour transmettre
        \emph{un qubit} d'état quelconque, sans qu'Alice connaisse $\alpha$ ni $\beta$.
  \item Pas de clonage : après la mesure, $A$ et $Q$ sont dans l'état de base $\ket{a,q}$.
        Alice n'a plus $\ket{\varphi}$ : l'état a été \emph{déplacé}, pas copié. L'intrication
        entre $A$ et $B$ a été consommée.
  \item Les quatre corrections $I,Z,X,ZX$ sont des portes contrôlées par des bits classiques
        (section~\ref{sec:controlees}).
\end{itemize}
\end{remarque}

\begin{remarque}[Les gestes d'Alice forment une mesure de Bell]
Le CNOT de contrôle $Q$ sur $A$ suivi de $H$ sur $Q$ est l'\emph{inverse} du générateur de la
section~\ref{sec:circuit-bell} (les portes $H$ et CNOT sont leurs propres inverses, et l'ordre
s'inverse) : il envoie la base de Bell du couple $(A,Q)$ sur la base de calcul. Mesurer ensuite dans
la base de calcul revient donc à mesurer $(A,Q)$ dans la base de Bell. Les deux bits $(a,q)$
disent \emph{quel} état de Bell Alice a trouvé, ce qui explique les quatre corrections de Bob.
\end{remarque}

\begin{remarque}[Non-clonage]
Alice ne peut pas simplement copier $\ket{\varphi}$ \parencite{wootters1982}. Supposons une porte
unitaire $U$ telle que $U\big(\ket{\psi}\otimes\ket{0}\big)=\ket{\psi}\otimes\ket{\psi}$ pour tout
$\ket{\psi}$. Elle copie $\ket{0}$ et $\ket{1}$, mais par linéarité
\[
  U\big(\ket{+}\otimes\ket{0}\big)=\tfrac{1}{\sqrt2}\big(\ket{00}+\ket{11}\big)
  \neq\ket{+}\otimes\ket{+}=\tfrac12\big(\ket{00}+\ket{01}+\ket{10}+\ket{11}\big),
\]
ce qui contredit l'hypothèse. Le CNOT copie un bit classique, pas une superposition (remarque de la
section~\ref{sec:cnot}). La téléportation respecte cette interdiction : l'état est déplacé, il n'est
jamais présent en deux endroits.
\end{remarque}

% ---------------------------------------------------------------------
\subsection{Corrélations non locales : le jeu CHSH}\label{sec:chsh}
% ---------------------------------------------------------------------

Les états de Bell donnent des résultats corrélés en base $Z$ comme en base $X$
(section~\ref{sec:bell}). Mais deux enveloppes contenant la même carte, préparées à l'avance, donnent
aussi des résultats égaux : des corrélations dans une seule base ne prouvent pas l'intrication. Bell
\parencite{bell1964} a montré qu'on peut trancher en \emph{changeant de base de mesure} ; la version de
Clauser, Horne, Shimony et Holt \parencite{clauser1969} se présente comme un jeu.

\begin{definition}{Le jeu CHSH}
Alice et Bob, éloignés, ne communiquent plus une fois le jeu commencé. Un arbitre tire deux bits $x$ et
$y$ de façon uniforme et indépendante, donne $x$ à Alice et $y$ à Bob. Chacun répond par un bit : $a$ pour
Alice, $b$ pour Bob. Ils \emph{gagnent} si
\[
  a\oplus b=x\cdot y,
\]
c'est-à-dire si $a=b$ pour $(x,y)\in\{(0,0),(0,1),(1,0)\}$ et si $a\neq b$ pour $(x,y)=(1,1)$.
\end{definition}

\paragraph{Stratégie classique : au plus $\frac34$.} Une stratégie déterministe est $a=f(x)$,
$b=g(y)$. Si elle gagnait les quatre parties, on aurait
\[
  f(0)\oplus g(0)=0,\qquad f(0)\oplus g(1)=0,\qquad f(1)\oplus g(0)=0,\qquad f(1)\oplus g(1)=1 .
\]
En additionnant ces quatre égalités modulo $2$, chacune des valeurs $f(0),f(1),g(0),g(1)$ apparaît
deux fois : le membre de gauche vaut $0$, celui de droite vaut $1$. C'est absurde, donc au moins une
partie sur quatre est perdue et $P(\text{gain})\leq\frac34$. La stratégie «~toujours répondre $0$~»
atteint $\frac34$ (elle ne perd que pour $(1,1)$). Un hasard partagé ne change rien : il mélange des
stratégies déterministes, dont aucune ne dépasse $\frac34$.

\paragraph{Stratégie quantique.} Alice et Bob partagent avant le jeu une paire $\ket{\Phi^+}$ (un
e-bit), puis chacun mesure son qubit sans rien dire à l'autre. Pour un angle $\theta$, notons
\[
  \ket{\theta_+}=\cos\tfrac\theta2\,\ket{0}+\sin\tfrac\theta2\,\ket{1},
  \qquad
  \ket{\theta_-}=-\sin\tfrac\theta2\,\ket{0}+\cos\tfrac\theta2\,\ket{1}
\]
la base de mesure d'angle $\theta$ : c'est la direction $\theta$ dans le plan $xz$ de la sphère de
Bloch ($\theta=0$ : base $Z$ ; $\theta=\frac\pi2$ : base $X$). Le résultat vaut $0$ pour
$\ket{\theta_+}$ et $1$ pour $\ket{\theta_-}$. La stratégie est
\begin{itemize}[nosep]
  \item Alice : $\theta=0$ si $x=0$ (base $Z$), $\theta=\frac\pi2$ si $x=1$ (base $X$) ;
  \item Bob : $\theta=\frac\pi4$ si $y=0$, $\theta=-\frac\pi4$ si $y=1$ (bases propres de
        $\frac{Z\pm X}{\sqrt2}$).
\end{itemize}

\begin{figure}[!ht]
\centering
\adjustbox{max width=\linewidth}{%
\begin{tikzpicture}[>=Stealth, font=\footnotesize,
    joueur/.style={draw=insarouge, fill=insarouge!6, rounded corners=2pt, thick,
                   minimum width=2.1cm, minimum height=1.1cm, align=center},
    arb/.style={draw=insagris, fill=insagris!8, rounded corners=2pt, thick,
                minimum width=3.6cm, minimum height=0.8cm, align=center}]
  % --- le jeu ---
  \node[arb] (arb) at (3.2,2.7) {arbitre : $x,y$ au hasard};
  \node[joueur] (al) at (0,0) {Alice};
  \node[joueur] (bo) at (6.4,0) {Bob};
  \draw[->, thick, insagris] (arb.west) -| node[pos=0.2, above] {$x$} (al.north);
  \draw[->, thick, insagris] (arb.east) -| node[pos=0.2, above] {$y$} (bo.north);
  \draw[insarouge, thick, decorate, decoration={snake, amplitude=1.1pt, segment length=6pt}]
    (al.east) -- (bo.west);
  \node[insarouge, above=1pt] at (3.2,0.05) {$\ket{\Phi^+}$};
  \node[text=insagris, below=2pt, font=\scriptsize] at (3.2,-0.05) {paire intriquée};
  \draw[->, thick] (al.south) -- ++(0,-1.0) node[below] {$a$};
  \draw[->, thick] (bo.south) -- ++(0,-1.0) node[below] {$b$};
  \node[draw=insagris, rounded corners=2pt, inner sep=3pt] at (3.2,-1.75) {gain si $a\oplus b=x\cdot y$};
  % --- directions de mesure dans le plan xz de la sphère de Bloch ---
  \begin{scope}[shift={(11.9,0.3)}]
    \draw[insagris!70] (0,0) circle (1.7);
    \draw[insagris!50, dashed] (-1.7,0) -- (1.7,0);
    \draw[insagris!50, dashed] (0,-1.7) -- (0,1.7);
    \node[font=\scriptsize, text=insagris, anchor=north] at (0,-1.72) {$\ket{1}$};
    \node[font=\scriptsize, text=insagris, anchor=east] at (-1.72,-0.02) {$\ket{-}$};
    \draw[->, very thick, insarouge] (0,0) -- (90:1.7);
    \draw[->, very thick, insarouge] (0,0) -- (0:1.7);
    \draw[->, very thick, insagris] (0,0) -- (45:1.7);
    \draw[->, very thick, insagris] (0,0) -- (135:1.7);
    \node[insarouge, anchor=south, inner sep=2pt] at (0,1.72) {$x=0$ : $Z$ $(\ket{0})$};
    \node[insarouge, anchor=west, inner sep=2pt] at (1.72,0) {$x=1$ : $X$ $(\ket{+})$};
    \node[text=insagris, anchor=south west, align=left, inner sep=2pt] at (1.25,1.25) {$y=0$\\ $\frac{Z+X}{\sqrt2}$};
    \node[text=insagris, anchor=south east, align=right, inner sep=2pt] at (-1.25,1.25) {$y=1$\\ $\frac{Z-X}{\sqrt2}$};
    \draw[insagris] (0.85,0) arc (0:45:0.85);
    \node[font=\scriptsize, text=insagris] at (22.5:1.12) {$\frac\pi4$};
    \draw[insagris] (0,0.85) arc (90:135:0.85);
    \node[font=\scriptsize, text=insagris] at (112.5:1.12) {$\frac\pi4$};
    \draw[insagris] (0.6,0.6) arc (45:90:0.85);
    \node[font=\scriptsize, text=insagris] at (67.5:1.12) {$\frac\pi4$};
  \end{scope}
\end{tikzpicture}}
\caption{Le jeu CHSH (à gauche) et les quatre directions de mesure de la stratégie quantique (à
droite), dans le plan $xz$ de la sphère de Bloch : Alice en rouge, Bob en gris. Pour
$(x,y)=(0,0),(0,1),(1,0)$, les directions font un angle $\frac\pi4$ ; pour $(1,1)$, un angle
$\frac{3\pi}4$.}
\label{fig:chsh}
\end{figure}

\begin{resultat}{Corrélations de $\ket{\Phi^+}$ en bases tournées}
Si Alice mesure dans la base d'angle $\alpha$ et Bob dans la base d'angle $\beta$, alors
\[
  P(a=b)=\cos^2\frac{\alpha-\beta}{2}.
\]
\end{resultat}

\noindent\emph{Preuve.} Avec $c=\cos\frac\alpha2$ et $s=\sin\frac\alpha2$,
$\ket{\alpha_+\alpha_+}+\ket{\alpha_-\alpha_-}=(c^2+s^2)\ket{00}+(cs-sc)\big(\ket{01}+\ket{10}\big)+(s^2+c^2)\ket{11}
=\ket{00}+\ket{11}$ : on peut donc écrire $\ket{\Phi^+}=\frac{1}{\sqrt2}\big(\ket{\alpha_+\alpha_+}+\ket{\alpha_-\alpha_-}\big)$.
Si Alice obtient $a=0$ (probabilité $\frac12$), le qubit de Bob est dans $\ket{\alpha_+}$ et Bob obtient $b=0$
avec la probabilité $|\braket{\beta_+}{\alpha_+}|^2$, où $\braket{\beta_+}{\alpha_+}=\cos\frac\beta2\cos\frac\alpha2+\sin\frac\beta2\sin\frac\alpha2=\cos\frac{\alpha-\beta}{2}$.
Si Alice obtient $a=1$, Bob est dans $\ket{\alpha_-}$ et obtient $b=1$ avec la même probabilité, car
$\braket{\beta_-}{\alpha_-}=\braket{\beta_+}{\alpha_+}$. Donc
$P(a=b)=\frac12\cos^2\frac{\alpha-\beta}2+\frac12\cos^2\frac{\alpha-\beta}2$. \hfill$\square$

\medskip
\noindent Pour les quatre couples $(x,y)$ :
\begin{center}
\renewcommand{\arraystretch}{1.3}
\begin{tabular}{@{}c c c c c c@{}}
\toprule
$(x,y)$ & $\alpha$ & $\beta$ & $P(a=b)$ & condition de gain & $P(\text{gain})$\\
\midrule
$(0,0)$ & $0$ & $\frac\pi4$ & $\cos^2\frac\pi8$ & $a=b$ & $\cos^2\frac\pi8$\\
$(0,1)$ & $0$ & $-\frac\pi4$ & $\cos^2\frac\pi8$ & $a=b$ & $\cos^2\frac\pi8$\\
$(1,0)$ & $\frac\pi2$ & $\frac\pi4$ & $\cos^2\frac\pi8$ & $a=b$ & $\cos^2\frac\pi8$\\
$(1,1)$ & $\frac\pi2$ & $-\frac\pi4$ & $\cos^2\frac{3\pi}8=\sin^2\frac\pi8$ & $a\neq b$ &
  $1-\sin^2\frac\pi8=\cos^2\frac\pi8$\\
\bottomrule
\end{tabular}
\end{center}
Alice et Bob gagnent donc chaque partie avec la même probabilité :
\[
  P(\text{gain})=\cos^2\frac\pi8=\frac{2+\sqrt2}{4}\approx0{,}854>\frac34 .
\]

\begin{remarque}[Ce que montre le jeu CHSH]
\begin{itemize}[nosep]
  \item Aucune stratégie classique, même avec du hasard partagé, ne dépasse $\frac34$ : l'intrication
        produit des corrélations qu'aucun modèle «~local~» (résultats fixés à l'avance) ne reproduit.
        Les expériences sur des photons intriqués confirment la valeur quantique.
  \item $\cos^2\frac\pi8$ est le maximum quantique (\emph{borne de Tsirelson}, résultat admis,
        \cite{tsirelson1980}) : on ne peut pas gagner à tous les coups. Avec
        $E_{xy}=P(a=b)-P(a\neq b)=\cos(\alpha-\beta)$ et $S=E_{00}+E_{01}+E_{10}-E_{11}$, on a
        $P(\text{gain})=\frac12+\frac S8$ ; l'inégalité de Bell donne $|S|\leq2$ en classique, et la
        stratégie ci-dessus atteint $S=2\sqrt2$.
  \item Il n'y a pas de communication : la réponse de Bob est uniforme ($P(b=0)=\frac12$) quelle que
        soit la base choisie par Alice. L'intrication ne transmet pas d'information ; elle ne fait
        que corréler les résultats.
\end{itemize}
\end{remarque}

\clearpage
% ---------------------------------------------------------------------
\subsection{Récapitulatif du chapitre}\label{sec:recap5}
% ---------------------------------------------------------------------

\begin{center}
\adjustbox{max width=\linewidth}{%
\begin{tikzpicture}[>=Stealth, font=\footnotesize,
    boite/.style={draw=insarouge, fill=insarouge!6, rounded corners=2pt, thick,
                  minimum width=4.4cm, minimum height=1.45cm, align=center, inner sep=3pt},
    lien/.style={->, very thick, insagris},
    etiq/.style={font=\scriptsize\itshape, text=insagris, inner sep=2pt}]
  \node[boite] (mes) at (0,0)
    {\textbf{Mesure à deux qubits}\\ $P(ij)=|c_{ij}|^2$\\ {\scriptsize\S\,\ref{sec:mesure2}}};
  \node[boite] (sep) at (6.3,0)
    {\textbf{Séparable ou intriqué}\\ $D=c_{00}c_{11}-c_{01}c_{10}$\\ {\scriptsize\S\,\ref{sec:separable}}};
  \node[boite] (bel) at (12.6,0)
    {\textbf{États de Bell}\\ $\ket{\Phi^\pm}$, $\ket{\Psi^\pm}$ ; $\mathrm{CNOT}\,(I\otimes H)$\\
     {\scriptsize\S\,\ref{sec:bell}, \ref{sec:circuit-bell}}};
  \node[boite] (tra) at (12.6,-2.3)
    {\textbf{Trace partielle, entropies}\\ $\rho_A=\operatorname{Tr}_B\rho_{AB}$ ; $S(\rho)$, $I(A;B)$\\
     {\scriptsize\S\,\ref{sec:trace}, \ref{sec:info}}};
  \node[boite] (tel) at (6.3,-2.3)
    {\textbf{Téléportation}\\ e-bit + 2 bits classiques\\ {\scriptsize\S\,\ref{sec:teleportation}}};
  \node[boite] (chs) at (0,-2.3)
    {\textbf{Jeu CHSH}\\ $\frac34$ classique, $\cos^2\frac\pi8$ quantique\\ {\scriptsize\S\,\ref{sec:chsh}}};
  \draw[lien] (mes) -- (sep) node[etiq, midway, above=1pt] {corrélations};
  \draw[lien] (sep) -- (bel) node[etiq, midway, above=1pt] {cas extrême};
  \draw[lien] (bel) -- (tra) node[etiq, midway, right=1pt] {un qubit seul};
  \draw[lien] (tra) -- (tel) node[etiq, midway, above=1pt] {ressource};
  \draw[lien] (tel) -- (chs) node[etiq, midway, above=1pt] {non-localité};
\end{tikzpicture}}
\end{center}

\begin{aretenir}
\begin{itemize}
  \item Deux qubits : quatre amplitudes $c_{ij}$, $P(ij)=|c_{ij}|^2$. Mesurer $A$ supprime les termes
        incompatibles et renormalise : l'état de $B$ peut changer.
  \item Séparable $\Longleftrightarrow D=c_{00}c_{11}-c_{01}c_{10}=0$. Sinon l'état est intriqué et
        aucun qubit n'a d'état propre (pour $\ket{\Phi^+}$, $P(00)\neq P(A{=}0)P(B{=}0)$).
  \item États de Bell : base orthonormée d'états maximalement intriqués, états propres de
        $Z\otimes Z$ et $X\otimes X$ ; $\mathrm{CNOT}\,(I\otimes H)$ les produit à partir de la base de
        calcul. Des portes locales ne créent pas d'intrication.
  \item Qubit seul : $\rho_A=\operatorname{Tr}_B\rho_{AB}$, avec $\operatorname{Tr}\rho_A^2=1-2|D|^2$ pour un état pur.
        Pour $\ket{\Phi^+}$ : $\rho_A=\mathrm{Id}/2$ alors que le couple est pur ; $S(\rho_{AB})=0$,
        $S(\rho_A)=1$ bit, $I(A;B)=2$ bits, $H(A|B)=-1$ (impossible en classique).
  \item Téléportation : un e-bit et deux bits classiques ; $P(a,q)=\frac14$ ; Bob applique $X$ si $a=1$
        puis $Z$ si $q=1$. L'état est déplacé, pas copié ; sans les bits, $\rho_B=\mathrm{Id}/2$.
  \item Jeu CHSH : toute stratégie classique gagne avec une probabilité $\leq\frac34$ ; avec $\ket{\Phi^+}$
        et des bases tournées, $\cos^2\frac\pi8\approx0{,}854$ (borne de Tsirelson), sans aucune
        communication.
\end{itemize}
\end{aretenir}

\begin{center}
\renewcommand{\arraystretch}{1.2}
\begin{tabular}{@{}l l r@{}}
\toprule
\textbf{Notion} & \textbf{Formule} & \textbf{§} \\
\midrule
Mesure de deux qubits & $P(ij)=|c_{ij}|^2$, \quad $P(A{=}a)=P(a0)+P(a1)$ & \ref{sec:mesure2} \\
Séparabilité & $\ket{\Psi}=\ket{\Psi}_A\otimes\ket{\Psi}_B\ \Longleftrightarrow\ D=0$ & \ref{sec:separable} \\
États de Bell & $\ket{\Phi^\pm}=\dfrac{\ket{00}\pm\ket{11}}{\sqrt2}$, \quad
  $\ket{\Psi^\pm}=\dfrac{\ket{01}\pm\ket{10}}{\sqrt2}$ & \ref{sec:bell} \\
Trace partielle & $\operatorname{Tr}_B\big(\ket{a}\bra{a'}\otimes\ket{b}\bra{b'}\big)
  =\ket{a}\bra{a'}\,\braket{b'}{b}$ & \ref{sec:trace} \\
Qubit d'un état pur & $\rho_A=CC^\dagger$, \quad $\operatorname{Tr}\rho_A^2=1-2|D|^2$ & \ref{sec:trace} \\
Entropies & $S(\rho)=-\sum_i\lambda_i\log_2\lambda_i$, \quad $I(A;B)=H(A)+H(B)-H(A,B)$ & \ref{sec:info} \\
Téléportation & $P(a,q)=\frac14$, \quad correction $Z^{q}X^{a}$ (d'abord $X$) & \ref{sec:teleportation} \\
Singulet & $(U\otimes U)\ket{\Psi^-}=\det(U)\,\ket{\Psi^-}$ & \ref{sec:bell} \\
Jeu CHSH & $P_{\mathrm{cl}}\leq\frac34$, \quad $P_{\mathrm{quant}}=\cos^2\frac\pi8=\frac{2+\sqrt2}{4}$ & \ref{sec:chsh} \\
\bottomrule
\end{tabular}
\end{center}

\begin{savoirfaire}
\begin{itemize}
  \item Calculer $P(ij)$, les probabilités d'un seul qubit et l'état après la mesure de l'un d'eux ;
        tester la séparabilité avec $D$ et calculer la pureté de $\rho_A$ (exercice~\ref{ex:separable3}).
  \item Suivre un état dans le générateur de Bell, écrire $\rho_{AB}$ et reconnaître les corrélations
        d'un état de Bell dans deux bases (exercice~\ref{ex:bellX}).
  \item Calculer $H(A)$, $H(A,B)$, $I(A;B)$ et $H(A|B)$ pour une loi ou un état (exercice~\ref{ex:entropies}).
  \item Dérouler la téléportation pour un état donné et utiliser les états de Bell pour le codage
        superdense (exercices~\ref{ex:telenum} et~\ref{ex:superdense}).
  \item Prouver la borne $\frac34$ du jeu CHSH classique et calculer la probabilité de gain d'une
        stratégie quantique (section~\ref{sec:chsh}).
\end{itemize}
\end{savoirfaire}
.
%  Il ne contient ni préambule ni \begin{document} : les environnements
%  (definition, resultat, remarque, aretenir, savoirfaire, exercice, correction),
%  les macros (\ii, \ee, \dd, \pd, \pdd, \Hop, \ket, \bra, \braket), les symboles
%  de circuits (ctl, cible, croix, mesure, porte) et les couleurs (insarouge,
%  insagris) sont définis dans main.tex. Un chapitre = une \section.
% =====================================================================

% Macros de Dirac (déjà définies dans main.tex ; ce bloc évite l'erreur
% « Undefined control sequence \ket » si main.tex est une ancienne version).
\providecommand{\ket}[1]{|#1\rangle}
\providecommand{\bra}[1]{\langle #1|}
\providecommand{\braket}[2]{\langle #1|#2\rangle}

% =====================================================================
\section{États quantiques, intrication et états de Bell}\label{chap:bell}
% =====================================================================

Les états de deux qubits forment un espace de dimension~$4$, et la plupart d'entre eux ne se
décomposent pas en «~un état pour $A$, un état pour $B$~» : ils sont \emph{intriqués}. Ce chapitre
apprend à mesurer un couple de qubits, à reconnaître l'intrication (critère de séparabilité), à la
produire avec un circuit (états de Bell), à décrire un qubit pris seul (trace partielle), à
quantifier l'information (entropies) et à s'en servir pour téléporter un état.

% ---------------------------------------------------------------------
\subsection{États de deux qubits et mesure}\label{sec:mesure2}
% ---------------------------------------------------------------------

Un état quelconque de $\mathcal{H}^4$ se décompose sur la base de calcul :
\[
  \ket{\Psi}=c_{00}\ket{00}+c_{01}\ket{01}+c_{10}\ket{10}+c_{11}\ket{11}
  =\begin{pmatrix}c_{00}\\ c_{01}\\ c_{10}\\ c_{11}\end{pmatrix},
  \qquad \sum_{i,j}|c_{ij}|^2=1 .
\]

\begin{resultat}{Mesure des deux qubits dans la base de calcul}
\[
  P(ij)=|\braket{ij}{\Psi}|^2=|c_{ij}|^2 .
\]
Pour un seul qubit : $P(A{=}0)=P(00)+P(01)$, $P(A{=}1)=P(10)+P(11)$,
$P(B{=}0)=P(00)+P(10)$, $P(B{=}1)=P(01)+P(11)$.
\end{resultat}

En effet $\braket{ij}{\Psi}=\sum_{k,l}c_{kl}\braket{ij}{kl}=\sum_{k,l}c_{kl}\delta_{ik}\delta_{jl}=c_{ij}$.

\paragraph{Exemple 1.} $\ket{\Psi}=\frac{1}{\sqrt2}\ket{00}+\frac{1}{\sqrt2}\ket{01}
=\big(\tfrac{1}{\sqrt2},\tfrac{1}{\sqrt2},0,0\big)^{\mathsf T}$. Il est normalisé :
$\frac12+\frac12+0+0=1$. Les probabilités jointes sont
$P(00)=P(01)=\frac12$ et $P(10)=P(11)=0$. Pour chaque qubit :
\[
  P(A{=}0)=\tfrac12+\tfrac12=1,\quad P(A{=}1)=0+0=0,\quad
  P(B{=}0)=\tfrac12+0=\tfrac12,\quad P(B{=}1)=\tfrac12+0=\tfrac12 .
\]
$A$ donne toujours $0$ ; $B$ donne $0$ ou $1$ au hasard.

\paragraph{Exemple 2 (mesure de $A$ seul).} Soit
$\ket{\Psi}=\frac12\ket{00}+\frac12\ket{01}+\frac{1}{\sqrt2}\ket{11}$. Normalisation :
$\frac14+\frac14+0+\frac12=1$. Probabilités jointes : $P(00)=\frac14$, $P(01)=\frac14$,
$P(10)=0$, $P(11)=\frac12$. Donc
\[
  P(A{=}0)=\tfrac14+\tfrac14=\tfrac12,\quad P(A{=}1)=0+\tfrac12=\tfrac12,\quad
  P(B{=}0)=\tfrac14+0=\tfrac14,\quad P(B{=}1)=\tfrac14+\tfrac12=\tfrac34 .
\]

\begin{remarque}[Effondrement après la mesure de $A$]
Si $A$ donne le résultat $a$, on ne garde que les termes de $\ket{\Psi}$ dont le premier
chiffre vaut $a$, puis on renormalise :
\[
  \ket{\Psi}\ \to\ \frac{c_{a0}\ket{a0}+c_{a1}\ket{a1}}{\sqrt{|c_{a0}|^2+|c_{a1}|^2}},
  \qquad\text{probabilité }|c_{a0}|^2+|c_{a1}|^2 .
\]
\end{remarque}

Pour l'exemple 2 :
\begin{itemize}[nosep]
  \item $A=0$ (probabilité $\frac12$) : on garde $\frac12\ket{00}+\frac12\ket{01}$, de norme
        au carré $\frac12$ ; renormalisé, $\frac{1}{\sqrt2}\big(\ket{00}+\ket{01}\big)
        =\ket{0}\otimes\ket{+}$. Ensuite $B$ vaut $0$ ou $1$ avec la probabilité $\frac12$.
  \item $A=1$ (probabilité $\frac12$) : on garde $\frac{1}{\sqrt2}\ket{11}$, de norme au carré
        $\frac12$ ; renormalisé, $\ket{11}=\ket{1}\otimes\ket{1}$. Ensuite $B$ vaut $1$
        avec certitude.
\end{itemize}
Contrôle : $P(B{=}0)=\frac12\cdot\frac12+\frac12\cdot0=\frac14$, comme calculé plus haut.
Ici la mesure de $A$ \emph{modifie} l'état de $B$ ($\ket{+}$ ou $\ket{1}$) : c'est la
signature d'un état intriqué (section~\ref{sec:separable}). La figure~\ref{fig:mesure2} résume
les deux mesures successives.

\begin{figure}[!ht]
\centering
\adjustbox{max width=\linewidth}{%
\begin{tikzpicture}[>=Stealth, font=\small,
    noeud/.style={draw=insarouge, fill=insarouge!6, rounded corners=2pt, thick, inner sep=4pt, align=center},
    fin/.style={draw=insagris, fill=insagris!10, rounded corners=2pt, inner sep=4pt, align=center},
    lien/.style={->, thick, insagris},
    proba/.style={font=\footnotesize, fill=white, inner sep=1.5pt, text=black}]
  \node[noeud] (psi) at (0,0) {$\ket{\Psi}$};
  \node[noeud] (a0) at (4.4,1.6) {$A=0$\\ {\footnotesize état $\ket{0}\otimes\ket{+}$}};
  \node[noeud] (a1) at (4.4,-1.6) {$A=1$\\ {\footnotesize état $\ket{1}\otimes\ket{1}$}};
  \node[fin] (b00) at (9.2,2.5) {$B=0$ : $\ket{00}$};
  \node[fin] (b01) at (9.2,0.7) {$B=1$ : $\ket{01}$};
  \node[fin] (b11) at (9.2,-1.6) {$B=1$ : $\ket{11}$};
  \draw[lien] (psi) -- (a0) node[proba, midway] {$\frac12$};
  \draw[lien] (psi) -- (a1) node[proba, midway] {$\frac12$};
  \draw[lien] (a0) -- (b00) node[proba, midway] {$\frac12$};
  \draw[lien] (a0) -- (b01) node[proba, midway] {$\frac12$};
  \draw[lien] (a1) -- (b11) node[proba, midway] {$1$};
  \node[anchor=west] at (11.0,2.5) {$P(00)=\frac12\cdot\frac12=\frac14$};
  \node[anchor=west] at (11.0,0.7) {$P(01)=\frac12\cdot\frac12=\frac14$};
  \node[anchor=west] at (11.0,-1.6) {$P(11)=\frac12\cdot1=\frac12$};
  \node[text=insagris, font=\footnotesize\itshape] at (2.2,3.0) {mesure de $A$};
  \node[text=insagris, font=\footnotesize\itshape] at (6.8,3.6) {mesure de $B$};
  \node[text=insagris, font=\footnotesize\itshape, anchor=west] at (11.0,3.6) {probabilité jointe};
\end{tikzpicture}}
\caption{Mesure de $A$ puis de $B$ pour l'état de l'exemple~2,
$\ket{\Psi}=\frac12\ket{00}+\frac12\ket{01}+\frac{1}{\sqrt2}\ket{11}$. Chaque flèche porte une
probabilité, et le produit des probabilités le long d'un chemin est la probabilité jointe. Après la
mesure de $A$, l'état de $B$ dépend du résultat ($\ket{+}$ ou $\ket{1}$).}
\label{fig:mesure2}
\end{figure}

% ---------------------------------------------------------------------
\subsection{États séparables et états intriqués}\label{sec:separable}
% ---------------------------------------------------------------------

\begin{definition}{Séparable ou intriqué}
Un état $\ket{\Psi}_{AB}$ est \emph{séparable} s'il s'écrit comme un produit
$\ket{\Psi}_{AB}=\ket{\Psi}_A\otimes\ket{\Psi}_B$. Sinon il est \emph{intriqué} : on ne
peut pas attribuer d'état propre à chaque qubit.
\end{definition}

\paragraph{Exemple séparable (exemple 1 ci-dessus).} Par linéarité du produit tensoriel :
\begin{align*}
  \ket{\Psi}&=\tfrac{1}{\sqrt2}\ket{00}_{AB}+\tfrac{1}{\sqrt2}\ket{01}_{AB}
    =\tfrac{1}{\sqrt2}\Big(\ket{0}_A\otimes\ket{0}_B+\ket{0}_A\otimes\ket{1}_B\Big)\\
  &=\ket{0}_A\otimes\Big(\tfrac{1}{\sqrt2}\big(\ket{0}+\ket{1}\big)\Big)_B
    =\ket{0}_A\otimes\ket{+}_B ,
\end{align*}
donc $\ket{\Psi}_A=\ket{0}$ et $\ket{\Psi}_B=\ket{+}$. Contrôle par les composantes :
$\begin{pmatrix}1\\ 0\end{pmatrix}\otimes\tfrac{1}{\sqrt2}\begin{pmatrix}1\\ 1\end{pmatrix}
=\tfrac{1}{\sqrt2}\begin{pmatrix}1\cdot1\\ 1\cdot1\\ 0\cdot1\\ 0\cdot1\end{pmatrix}
=\tfrac{1}{\sqrt2}\begin{pmatrix}1\\ 1\\ 0\\ 0\end{pmatrix}$.

\subsubsection{Un critère de séparabilité}

Un produit quelconque a pour coefficients $c_{00}=ac$, $c_{01}=ad$, $c_{10}=bc$,
$c_{11}=bd$ (formule de la section~\ref{sec:tenseur}, avec $a,b$ pour $A$ et $c,d$ pour $B$), donc
\[
  c_{00}c_{11}=(ac)(bd)=abcd=(ad)(bc)=c_{01}c_{10}.
\]

\begin{resultat}{Test de séparabilité (deux qubits)}
$\ket{\Psi}=\sum c_{ij}\ket{ij}$ est séparable si et seulement si
\[
  D=c_{00}c_{11}-c_{01}c_{10}=0 .
\]
\end{resultat}

\emph{Démonstration.} ($\Rightarrow$) Calcul ci-dessus. ($\Leftarrow$) On regroupe les
termes selon l'état de $A$ :
\[
  \ket{\Psi}=\ket{0}\otimes\big(c_{00}\ket{0}+c_{01}\ket{1}\big)+\ket{1}\otimes\big(c_{10}\ket{0}+c_{11}\ket{1}\big).
\]
Si $D=0$, les lignes $(c_{00},c_{01})$ et $(c_{10},c_{11})$ sont proportionnelles. Si
$(c_{00},c_{01})\neq(0,0)$, il existe $\lambda$ avec $(c_{10},c_{11})=\lambda(c_{00},c_{01})$
et $\ket{\Psi}=\big(\ket{0}+\lambda\ket{1}\big)\otimes\big(c_{00}\ket{0}+c_{01}\ket{1}\big)$.
Si $(c_{00},c_{01})=(0,0)$, alors $\ket{\Psi}=\ket{1}\otimes\big(c_{10}\ket{0}+c_{11}\ket{1}\big)$.
Dans les deux cas $\ket{\Psi}$ est un produit ; on renormalise chaque facteur.
\hfill$\square$

\subsubsection{Application à plusieurs états}

\begin{center}
\renewcommand{\arraystretch}{1.5}
\begin{tabular}{@{}lccl@{}}
\toprule
\textbf{État} & $(c_{00},c_{01},c_{10},c_{11})$ & $D$ & \textbf{Nature} \\
\midrule
$\ket{00}$ & $(1,0,0,0)$ & $1\cdot0-0\cdot0=0$ & séparable \\
$\ket{+}\otimes\ket{+}$ & $\frac12(1,1,1,1)$ & $\frac14-\frac14=0$ & séparable \\
$\ket{+}\otimes\ket{-}$ & $\frac12(1,-1,1,-1)$ & $-\frac14-(-\frac14)=0$ & séparable \\
$\frac{1}{\sqrt2}(\ket{00}+\ket{01})$ & $\frac{1}{\sqrt2}(1,1,0,0)$ & $0-\frac12\cdot0=0$ & séparable \\
$\ket{\Phi^+}=\frac{1}{\sqrt2}(\ket{00}+\ket{11})$ & $\frac{1}{\sqrt2}(1,0,0,1)$ & $\frac12-0=\frac12$ & intriqué \\
$\ket{\Phi^-}=\frac{1}{\sqrt2}(\ket{00}-\ket{11})$ & $\frac{1}{\sqrt2}(1,0,0,-1)$ & $-\frac12-0=-\frac12$ & intriqué \\
$\ket{\Psi^+}=\frac{1}{\sqrt2}(\ket{01}+\ket{10})$ & $\frac{1}{\sqrt2}(0,1,1,0)$ & $0-\frac12=-\frac12$ & intriqué \\
$\ket{\Psi^-}=\frac{1}{\sqrt2}(\ket{01}-\ket{10})$ & $\frac{1}{\sqrt2}(0,1,-1,0)$ & $0-(-\frac12)=\frac12$ & intriqué \\
$\frac12\ket{00}+\frac12\ket{01}+\frac{1}{\sqrt2}\ket{11}$ & $(\frac12,\frac12,0,\frac{1}{\sqrt2})$ & $\frac{1}{2\sqrt2}-0$ & intriqué \\
\bottomrule
\end{tabular}
\end{center}

Les quatre états de Bell $\ket{\Phi^\pm}$ et $\ket{\Psi^\pm}$ de la table sont tous intriqués
($D=\pm\frac12$) ; leurs propriétés font l'objet de la section~\ref{sec:bell}.

\paragraph{Preuve directe que $\ket{\Phi^+}$ est intriqué.} Supposons
$\ket{\Phi^+}=(a\ket{0}+b\ket{1})\otimes(c\ket{0}+d\ket{1})$. En identifiant les
composantes : $ac=\frac{1}{\sqrt2}$, $ad=0$, $bc=0$, $bd=\frac{1}{\sqrt2}$. La condition
$ad=0$ impose $a=0$ (impossible car $ac\neq0$) ou $d=0$ (impossible car $bd\neq0$).
Contradiction : aucun produit ne convient.

\paragraph{Corrélations.} Pour un produit $\ket{\Psi}_A\otimes\ket{\Psi}_B$ avec
$\ket{\Psi}_A=a_0\ket{0}+a_1\ket{1}$ et $\ket{\Psi}_B=b_0\ket{0}+b_1\ket{1}$, on a
$P(ij)=|a_ib_j|^2=|a_i|^2|b_j|^2=P(A{=}i)\,P(B{=}j)$ : les deux qubits sont
\emph{indépendants}. Vérification avec $\ket{+}\otimes\ket{+}$ : $P(00)=\frac14=\frac12\cdot\frac12$.
Pour $\ket{\Phi^+}$ : $P(00)=P(11)=\frac12$ et $P(01)=P(10)=0$, donc
$P(A{=}0)=P(B{=}0)=\frac12$, mais
\[
  P(00)=\tfrac12\neq\tfrac12\cdot\tfrac12=P(A{=}0)\,P(B{=}0) .
\]
Les deux qubits donnent toujours le même résultat, aléatoire : ils sont corrélés.
La figure~\ref{fig:epr} illustre la différence entre les deux cas.

\begin{figure}[!ht]
\centering
\adjustbox{max width=\linewidth}{%
\begin{tikzpicture}[>=Stealth, font=\small,
    labo/.style={draw=insagris, fill=insagris!8, rounded corners=2pt, thick, minimum width=2.6cm,
                 minimum height=1.0cm, align=center},
    source/.style={draw=insarouge, fill=insarouge!10, rounded corners=2pt, thick, minimum width=3.6cm,
                   minimum height=1.0cm, align=center},
    droit/.style={->, thick, insagris},
    onde/.style={->, thick, insarouge, decorate,
                 decoration={snake, amplitude=1.3pt, segment length=7pt, post length=2mm}},
    res/.style={font=\small, align=center}]
  % (a) état séparable
  \node[labo] (Aa) at (0,0) {Alice\\ mesure $A$};
  \node[source] (Sa) at (6,0) {source\\ $\ket{0}_A\otimes\ket{+}_B$};
  \node[labo] (Ba) at (12,0) {Bob\\ mesure $B$};
  \draw[droit] (Sa.west) -- (Aa.east);
  \draw[droit] (Sa.east) -- (Ba.west);
  \node[res, below=3pt of Aa] {$0\ \ 0\ \ 0\ \ 0\ \ 0\ \ 0$};
  \node[res, below=3pt of Ba] {$1\ \ 0\ \ 0\ \ 1\ \ 0\ \ 1$};
  \node[res, text=insagris, font=\footnotesize] at (6,-1.0) {$A$ donne toujours $0$ ; $B$ est aléatoire.\\
        Les deux résultats sont indépendants.};
  \node[anchor=south west, font=\small\bfseries] at (-1.3,0.75) {(a) état séparable};
  % (b) état intriqué
  \begin{scope}[yshift=-3.6cm]
    \node[labo] (Ab) at (0,0) {Alice\\ mesure $A$};
    \node[source] (Sb) at (6,0) {source\\ $\ket{\Phi^+}_{AB}$};
    \node[labo] (Bb) at (12,0) {Bob\\ mesure $B$};
    \draw[onde] (Sb.west) -- (Ab.east);
    \draw[onde] (Sb.east) -- (Bb.west);
    \node[res, below=3pt of Ab] {$1\ \ 0\ \ 0\ \ 1\ \ 1\ \ 0$};
    \node[res, below=3pt of Bb] {$1\ \ 0\ \ 0\ \ 1\ \ 1\ \ 0$};
    \node[res, text=insagris, font=\footnotesize] at (6,-1.0) {$A$ et $B$ sont chacun aléatoires,\\
          mais toujours égaux : ils sont corrélés.};
    \node[anchor=south west, font=\small\bfseries] at (-1.3,0.75) {(b) état intriqué};
  \end{scope}
\end{tikzpicture}}
\caption{Six mesures successives sur des paires préparées dans le même état (les résultats sont
des exemples). (a) Pour un produit, chaque qubit a son propre comportement. (b) Pour
$\ket{\Phi^+}=\frac{1}{\sqrt2}(\ket{00}+\ket{11})$, chaque résultat pris seul est un tirage à pile ou
face, mais les deux résultats coïncident toujours ; la ligne ondulée symbolise l'intrication.}
\label{fig:epr}
\end{figure}

\paragraph{Et un qubit pris seul ?} Dans un état séparable, chaque qubit garde son propre état pur ;
dans un état intriqué, aucun ket ne décrit un qubit isolé. On le décrit alors par une matrice de
densité, obtenue par la trace partielle (section~\ref{sec:trace}).

% ---------------------------------------------------------------------
\subsection{Les états de Bell}\label{sec:bell}
% ---------------------------------------------------------------------

\begin{definition}{Les quatre états de Bell}
\[
  \ket{\Phi^{\pm}}=\frac{1}{\sqrt2}\big(\ket{00}\pm\ket{11}\big),
  \qquad
  \ket{\Psi^{\pm}}=\frac{1}{\sqrt2}\big(\ket{01}\pm\ket{10}\big).
\]
Ils sont intriqués ($|D|=\frac12$, section~\ref{sec:separable}) et le sont au maximum :
$|D|\leq\frac12$ pour tout état normalisé de deux qubits, et chaque qubit d'un état de Bell,
pris seul, est complètement aléatoire (section~\ref{sec:trace}).
\end{definition}

\paragraph{Une base de $\mathcal{H}^4$.} Chaque état est de norme $\frac12+\frac12=1$. Les
$\ket{\Phi^\pm}$ n'ont de composantes que sur $\ket{00}$ et $\ket{11}$, les $\ket{\Psi^\pm}$ que sur
$\ket{01}$ et $\ket{10}$, donc $\braket{\Phi^\pm}{\Psi^\pm}=0$ pour toutes les combinaisons de signes ;
enfin $\braket{\Phi^+}{\Phi^-}=\frac12(1-1)=0$ et $\braket{\Psi^+}{\Psi^-}=\frac12(1-1)=0$. Les quatre
états forment donc une base orthonormée de $\mathcal{H}^4$, la \emph{base de Bell}, constituée
uniquement d'états intriqués. Le circuit de la section~\ref{sec:circuit-bell} fait passer de la
base de calcul à cette base.

\paragraph{Ce qui les distingue : deux parités.} Sur la base de calcul,
$Z\otimes Z\ket{ij}=(-1)^{i+j}\ket{ij}$ (valeur propre $+1$ si les deux chiffres sont égaux, $-1$
sinon) et $X\otimes X\ket{ij}=\ket{\bar i\bar j}$ (les deux chiffres sont inversés, avec
$\bar i=1-i$). On en déduit, par exemple,
\[
  X\otimes X\ket{\Phi^-}=\tfrac{1}{\sqrt2}\big(\ket{11}-\ket{00}\big)=-\ket{\Phi^-},
  \qquad
  Z\otimes Z\ket{\Psi^+}=\tfrac{1}{\sqrt2}\big(-\ket{01}-\ket{10}\big)=-\ket{\Psi^+},
\]
et le tableau complet des valeurs propres :
\begin{center}
\renewcommand{\arraystretch}{1.3}
\begin{tabular}{@{}lcccc@{}}
\toprule
 & $\ket{\Phi^+}$ & $\ket{\Phi^-}$ & $\ket{\Psi^+}$ & $\ket{\Psi^-}$ \\
\midrule
$Z\otimes Z$ & $+1$ & $+1$ & $-1$ & $-1$ \\
$X\otimes X$ & $+1$ & $-1$ & $+1$ & $-1$ \\
\bottomrule
\end{tabular}
\end{center}
Les deux opérateurs commutent, car $ZX=-XZ$ donne
$(Z\otimes Z)(X\otimes X)=ZX\otimes ZX=XZ\otimes XZ=(X\otimes X)(Z\otimes Z)$ ; les quatre états de
Bell sont leurs états propres communs, et chaque couple de valeurs propres apparaît une fois
(figure~\ref{fig:grille-bell}). Le même calcul montre que $X\otimes Z$ et $Z\otimes X$
commutent : $(X\otimes Z)(Z\otimes X)=XZ\otimes ZX=(-ZX)\otimes(-XZ)=ZX\otimes XZ=(Z\otimes X)(X\otimes Z)$.
Plus généralement, $P\otimes Q$ et $P'\otimes Q'$ (matrices de Pauli) commutent si $P,P'$
anticommutent et si $Q,Q'$ anticommutent : les deux signes $-1$ se compensent.

\begin{figure}[!ht]
\centering
\adjustbox{max width=\linewidth}{%
\begin{tikzpicture}[>=Stealth, font=\small,
    cell/.style={draw=insarouge, fill=insarouge!6, rounded corners=2pt, thick, minimum width=5.4cm,
                 minimum height=1.35cm, align=center},
    tete/.style={text=insagris, align=center}]
  \node[tete] at (3.0,2.0) {$X\otimes X=+1$\\ {\footnotesize mêmes résultats en base $X$}};
  \node[tete] at (8.8,2.0) {$X\otimes X=-1$\\ {\footnotesize résultats opposés en base $X$}};
  \node[tete, anchor=east] at (-0.1,0.75) {$Z\otimes Z=+1$\\ {\footnotesize mêmes résultats}\\ {\footnotesize en base $Z$}};
  \node[tete, anchor=east] at (-0.1,-0.85) {$Z\otimes Z=-1$\\ {\footnotesize résultats opposés}\\ {\footnotesize en base $Z$}};
  \node[cell] at (3.0,0.75) {$\ket{\Phi^+}=\frac{1}{\sqrt2}\big(\ket{00}+\ket{11}\big)$};
  \node[cell] at (8.8,0.75) {$\ket{\Phi^-}=\frac{1}{\sqrt2}\big(\ket{00}-\ket{11}\big)$};
  \node[cell] at (3.0,-0.85) {$\ket{\Psi^+}=\frac{1}{\sqrt2}\big(\ket{01}+\ket{10}\big)$};
  \node[cell] at (8.8,-0.85) {$\ket{\Psi^-}=\frac{1}{\sqrt2}\big(\ket{01}-\ket{10}\big)$};
\end{tikzpicture}}
\caption{Les quatre états de Bell, classés par les valeurs propres de $Z\otimes Z$ (lignes) et de
$X\otimes X$ (colonnes). La valeur propre est le produit des deux résultats $\pm1$ d'une mesure des
deux qubits dans la base correspondante : $+1$ si les résultats sont égaux, $-1$ s'ils sont opposés.}
\label{fig:grille-bell}
\end{figure}

Les corrélations ne se limitent donc pas à la base $Z$ : $\ket{\Phi^+}$ s'écrit aussi
$\frac{1}{\sqrt2}\big(\ket{++}+\ket{--}\big)$ et donne des résultats égaux en base $X$ comme en base $Z$
(exercice~\ref{ex:bellX}).

\paragraph{D'un état de Bell à l'autre.} Une porte agissant sur un seul qubit suffit. Avec
$X\ket{0}=\ket{1}$, $X\ket{1}=\ket{0}$ et $Z\ket{1}=-\ket{1}$, en agissant sur le second qubit :
\begin{align*}
  (I\otimes X)\ket{\Phi^+}&=\tfrac{1}{\sqrt2}\big(\ket{01}+\ket{10}\big)=\ket{\Psi^+},\\
  (I\otimes Z)\ket{\Phi^+}&=\tfrac{1}{\sqrt2}\big(\ket{00}-\ket{11}\big)=\ket{\Phi^-},\\
  (I\otimes XZ)\ket{\Phi^+}&=(I\otimes X)\tfrac{1}{\sqrt2}\big(\ket{00}-\ket{11}\big)
    =\tfrac{1}{\sqrt2}\big(\ket{01}-\ket{10}\big)=\ket{\Psi^-}.
\end{align*}
Quatre opérations locales ($I$, $X$, $Z$, $XZ$) sur \emph{un} qubit d'une paire de Bell mènent donc
aux quatre états de la base ; c'est le principe du codage superdense
(exercice~\ref{ex:superdense}).

\begin{remarque}[Le singulet est le même dans toutes les bases]\label{rem:singulet}
Pour toute matrice unitaire $U=\begin{pmatrix}a&b\\c&d\end{pmatrix}$, on a $U\ket{0}=a\ket{0}+c\ket{1}$
et $U\ket{1}=b\ket{0}+d\ket{1}$, donc
\begin{align*}
  (U\otimes U)\big(\ket{01}-\ket{10}\big)
  &=\big(a\ket{0}+c\ket{1}\big)\big(b\ket{0}+d\ket{1}\big)-\big(b\ket{0}+d\ket{1}\big)\big(a\ket{0}+c\ket{1}\big)\\
  &=(ad-bc)\big(\ket{01}-\ket{10}\big),
\end{align*}
c'est-à-dire $(U\otimes U)\ket{\Psi^-}=\det(U)\,\ket{\Psi^-}$. Comme $|\det U|=1$, ce facteur est une
phase globale, sans effet physique : $\ket{\Psi^-}$ s'écrit \emph{de la même façon} dans toutes les
bases orthonormées (avec $U=H$, de déterminant $-1$ : $\ket{\Psi^-}=\frac{1}{\sqrt2}\big(\ket{-+}-\ket{+-}\big)$).
Si Alice et Bob mesurent dans la \emph{même} base, quelle qu'elle soit, leurs résultats sont
toujours opposés : c'est le «~singulet~». Parmi les états de Bell, seul $\ket{\Psi^-}$ a cette propriété.
\end{remarque}

% ---------------------------------------------------------------------
\subsection{Un circuit quantique : le générateur d'états de Bell}\label{sec:circuit-bell}
% ---------------------------------------------------------------------

On assemble une porte de Hadamard et un CNOT. Le système est $(A,B)$ avec le ket
$\ket{a,b}=\ket{a}_A\otimes\ket{b}_B$ ; d'après la convention de la section~\ref{sec:circuits}, $B$ est dessiné
en haut et $A$ en bas (figure~\ref{fig:circuit-bell}). Le circuit a deux étapes :
\begin{enumerate}[nosep]
  \item la porte $I\otimes H$ : identité sur $A$, Hadamard sur $B$ ;
  \item un CNOT de contrôle $B$ et de cible $A$ (celui de la section~\ref{sec:cnot}, dont le
        contrôle est le second chiffre).
\end{enumerate}
Les traits pointillés délimitent les états $\ket{\Pi_0}$, $\ket{\Pi_1}$ et $\ket{\Pi_2}$ :
avant les portes, après $I\otimes H$, après le CNOT. Le circuit transforme la base de calcul
en base de Bell (section~\ref{sec:bell}).

\begin{figure}[!ht]
\centering
\begin{tikzpicture}[>=Stealth, font=\small]
  % fils : B en haut, A en bas
  \draw (0,1.4) -- (7.6,1.4);
  \draw (0,0) -- (7.6,0);
  \node[anchor=east] at (0,1.4) {$\ket{0}_B$};
  \node[anchor=east] at (0,0) {$\ket{0}_A$};
  % première étape : H sur B, identité (pointillée) sur A
  \node[porte] at (2.2,1.4) {$H$};
  \node[draw=insagris, dashed, fill=white, minimum size=7mm] at (2.2,0) {$I$};
  % CNOT : contrôle B (haut), cible A (bas)
  \draw[insarouge, thick] (5,1.4) -- (5,0);
  \pic at (5,1.4) {ctl};
  \pic at (5,0) {cible};
  % accolade de sortie
  \draw[insagris, thick] (7.75,1.6) -- (7.9,1.6) -- (7.9,-0.2) -- (7.75,-0.2);
  \node[anchor=west] at (7.9,0.7) {$\ket{\beta^{00}}_{AB}$};
  % états intermédiaires
  \foreach \x in {0.7,3.6,6.4}
    \draw[dashed, insagris!70] (\x,1.9) -- (\x,-0.7);
  \node[anchor=north, font=\footnotesize] at (0.7,-0.7) {$\ket{\Pi_0}=\ket{00}$};
  \node[anchor=north, font=\footnotesize] at (3.6,-0.7) {$\ket{\Pi_1}=\ket{0}\otimes\ket{+}$};
  \node[anchor=north, font=\footnotesize] at (6.4,-0.7) {$\ket{\Pi_2}=\ket{\beta^{00}}$};
\end{tikzpicture}
\caption{Générateur d'états de Bell : $I\otimes H$, puis un CNOT de contrôle $B$ (fil du haut)
et de cible $A$ (fil du bas). L'entrée représentée est $\ket{00}$.}
\label{fig:circuit-bell}
\end{figure}

\subsubsection{Entrée \texorpdfstring{$\ket{00}$}{|00>} : calcul pas à pas}

\begin{align*}
  \ket{\Pi_0}&=\ket{0}_A\otimes\ket{0}_B=\ket{00},\\[1ex]
  \ket{\Pi_1}&=(I\otimes H)\ket{\Pi_0}=I\ket{0}_A\otimes H\ket{0}_B=\ket{0}\otimes\ket{+}
    =\frac{\ket{00}+\ket{01}}{\sqrt2},\\[1ex]
  \ket{\Pi_2}&=\mathrm{CNOT}\ket{\Pi_1}=\frac{\mathrm{CNOT}\ket{00}+\mathrm{CNOT}\ket{01}}{\sqrt2}
    =\frac{\ket{00}+\ket{11}}{\sqrt2}=\ket{\beta^{00}}=\ket{\Phi^+}.
\end{align*}
Dans $\ket{\Pi_1}$, le premier chiffre est celui de $A$ et le second celui de $B$. Le CNOT a
pour contrôle $B$ : $\ket{00}$ ($B=0$) ne change pas, $\ket{01}$ ($B=1$) devient $\ket{11}$
puisque $A$ est inversé. Le résultat est un \emph{état de Bell}.

Avec les matrices, la règle $A\otimes B=(A_{ij}B)$ donne
\[
  I\otimes H=\begin{pmatrix}1\cdot H&0\\0&1\cdot H\end{pmatrix}
  =\frac{1}{\sqrt2}\begin{pmatrix}1&1&0&0\\1&-1&0&0\\0&0&1&1\\0&0&1&-1\end{pmatrix}.
\]
Sa première colonne est $\ket{\Pi_1}=\frac{1}{\sqrt2}(1,1,0,0)^{\mathsf T}$ (image de $\ket{00}$).
Le CNOT échange les amplitudes de $\ket{01}$ et de $\ket{11}$ (section~\ref{sec:cnot}), donc
\[
  \ket{\Pi_2}=\mathrm{CNOT}\,\frac{1}{\sqrt2}\begin{pmatrix}1\\1\\0\\0\end{pmatrix}
  =\frac{1}{\sqrt2}\begin{pmatrix}1\\0\\0\\1\end{pmatrix}.
\]

\subsubsection{Les quatre entrées}

Prenons $B$ dans l'état $\ket{x}$ et $A$ dans l'état $\ket{y}$, avec $x,y\in\{0,1\}$. Alors
$\ket{\Pi_0}=\ket{y}_A\otimes\ket{x}_B=\ket{y,x}$. Comme
$H\ket{x}=\frac{1}{\sqrt2}\big(\ket{0}+(-1)^x\ket{1}\big)$ (c'est $\ket{+}$ pour $x=0$ et
$\ket{-}$ pour $x=1$),
\[
  \ket{\Pi_1}=\ket{y}\otimes H\ket{x}=\frac{1}{\sqrt2}\Big(\ket{y,0}+(-1)^x\ket{y,1}\Big).
\]
Le CNOT (contrôle $B$) ne change pas $\ket{y,0}$ et transforme $\ket{y,1}$ en
$\ket{\bar y,1}$, avec $\bar y=1-y$ :
\[
  \ket{\Pi_2}=\ket{\beta^{xy}}=\frac{1}{\sqrt2}\Big(\ket{y,0}+(-1)^x\ket{\bar y,1}\Big).
\]
Le premier indice de $\ket{\beta^{xy}}$ est l'entrée de $B$ (le qubit qui reçoit $H$), le
second celle de $A$ (la cible). Les quatre cas :

\begin{center}
\renewcommand{\arraystretch}{1.7}
\begin{tabular}{@{}ccccc@{}}
\toprule
$(B,A)=(x,y)$ & $\ket{\Pi_0}=\ket{y,x}$ & $\ket{\Pi_1}$ & $\ket{\Pi_2}=\ket{\beta^{xy}}$ & Nom \\
\midrule
$(0,0)$ & $\ket{00}$ & $\frac{1}{\sqrt2}(\ket{00}+\ket{01})$ & $\frac{1}{\sqrt2}(\ket{00}+\ket{11})$ & $\ket{\Phi^+}$ \\
$(1,0)$ & $\ket{01}$ & $\frac{1}{\sqrt2}(\ket{00}-\ket{01})$ & $\frac{1}{\sqrt2}(\ket{00}-\ket{11})$ & $\ket{\Phi^-}$ \\
$(0,1)$ & $\ket{10}$ & $\frac{1}{\sqrt2}(\ket{10}+\ket{11})$ & $\frac{1}{\sqrt2}(\ket{10}+\ket{01})$ & $\ket{\Psi^+}$ \\
$(1,1)$ & $\ket{11}$ & $\frac{1}{\sqrt2}(\ket{10}-\ket{11})$ & $\frac{1}{\sqrt2}(\ket{10}-\ket{01})$ & $-\ket{\Psi^-}$ \\
\bottomrule
\end{tabular}
\end{center}
Détail de la dernière ligne : $\ket{\Pi_1}=\ket{1}\otimes\ket{-}=\frac{1}{\sqrt2}(\ket{10}-\ket{11})$ ;
le CNOT laisse $\ket{10}$ et envoie $\ket{11}$ sur $\ket{01}$, d'où
$\frac{1}{\sqrt2}(\ket{10}-\ket{01})=-\frac{1}{\sqrt2}(\ket{01}-\ket{10})=-\ket{\Psi^-}$.

\paragraph{Forme matricielle.} Le CNOT échange les lignes $2$ et $4$ de
$I\otimes H$ (il échange les amplitudes de $\ket{01}$ et $\ket{11}$), donc
\[
  \mathrm{CNOT}\,(I\otimes H)=\frac{1}{\sqrt2}\begin{pmatrix}
    1 & 1 & 0 & 0\\ 0 & 0 & 1 & -1\\ 0 & 0 & 1 & 1\\ 1 & -1 & 0 & 0
  \end{pmatrix}.
\]
Ses colonnes sont les images de $\ket{00},\ket{01},\ket{10},\ket{11}$, c'est-à-dire de
$(B,A)=(0,0),(1,0),(0,1),(1,1)$ : ce sont $\ket{\beta^{00}}$, $\ket{\beta^{10}}$,
$\ket{\beta^{01}}$, $\ket{\beta^{11}}$, soit $\ket{\Phi^+}$, $\ket{\Phi^-}$, $\ket{\Psi^+}$
et $-\ket{\Psi^-}$.

\subsubsection{Matrice de densité de l'état de sortie}

Pour $\ket{\Psi}=\ket{\beta^{00}}=\frac{1}{\sqrt2}\ket{00}+\frac{1}{\sqrt2}\ket{11}$ :
\[
  \rho_{AB}=\ket{\Psi}\bra{\Psi}
  =\begin{pmatrix}\frac{1}{\sqrt2}\\0\\0\\\frac{1}{\sqrt2}\end{pmatrix}
   \begin{pmatrix}\frac{1}{\sqrt2}&0&0&\frac{1}{\sqrt2}\end{pmatrix}
  =\frac12\begin{pmatrix}1&0&0&1\\0&0&0&0\\0&0&0&0\\1&0&0&1\end{pmatrix}.
\]
Sous forme de ket-bra :
$\rho_{AB}=\frac12\big(\ket{00}\bra{00}+\ket{00}\bra{11}+\ket{11}\bra{00}+\ket{11}\bra{11}\big)$.
On vérifie $\operatorname{Tr}\rho_{AB}=\frac12(1+1)=1$ et
\[
  \rho_{AB}^2=\frac14\begin{pmatrix}2&0&0&2\\0&0&0&0\\0&0&0&0\\2&0&0&2\end{pmatrix}=\rho_{AB},
\]
donc $\operatorname{Tr}\rho_{AB}^2=1$ : le couple $(A,B)$ est dans un état \emph{pur}.

\begin{remarque}[Porte intriquante]
Les états d'entrée $\ket{ij}$ sont séparables ($D=0$) ; les sorties sont intriquées
($D=\pm\frac12$). Aucune porte agissant séparément sur $A$ et sur $B$ ($U_A\otimes U_B$,
section~\ref{sec:op2}) ne peut créer de l'intrication à partir d'un produit, puisque
$(U_A\otimes U_B)(\ket{u}\otimes\ket{v})=U_A\ket{u}\otimes U_B\ket{v}$ reste un produit.
Une porte à deux qubits comme le CNOT est indispensable.
\end{remarque}

% ---------------------------------------------------------------------
\subsection{Trace partielle : l'état d'un qubit pris seul}\label{sec:trace}
% ---------------------------------------------------------------------

Un état intriqué n'a pas de ket pour $A$ seul (section~\ref{sec:separable}). On peut pourtant
décrire $A$ par une matrice de densité, obtenue en «~oubliant~» $B$.

\begin{definition}{Trace partielle}
Soit $\rho_{AB}$ la matrice de densité du couple $(A,B)$. L'état du qubit $A$ pris seul est
$\rho_A=\operatorname{Tr}_B(\rho_{AB})$, celui de $B$ est $\rho_B=\operatorname{Tr}_A(\rho_{AB})$. La
trace partielle est linéaire et remplace la partie tracée de chaque terme par sa trace :
\[
  \operatorname{Tr}_B\big(\ket{a}\bra{a'}\otimes\ket{b}\bra{b'}\big)=\ket{a}\bra{a'}\;\braket{b'}{b},
  \qquad
  \operatorname{Tr}_A\big(\ket{a}\bra{a'}\otimes\ket{b}\bra{b'}\big)=\braket{a'}{a}\;\ket{b}\bra{b'}.
\]
\end{definition}

\paragraph{En matrices.} On découpe $\rho_{AB}$ (matrice $4\times4$, lignes et colonnes ordonnées
$00,01,10,11$) en quatre blocs $2\times2$ :
\[
  \rho_{AB}=\begin{pmatrix}R_{00}&R_{01}\\R_{10}&R_{11}\end{pmatrix}
  \quad\Longrightarrow\quad
  \rho_A=\begin{pmatrix}\operatorname{Tr}R_{00}&\operatorname{Tr}R_{01}\\
                        \operatorname{Tr}R_{10}&\operatorname{Tr}R_{11}\end{pmatrix},
  \qquad
  \rho_B=R_{00}+R_{11}.
\]
Autrement dit $(\rho_A)_{ij}=\rho_{(i0),(j0)}+\rho_{(i1),(j1)}$ : $\rho_A$ est la matrice des traces
des blocs, et $\rho_B$ la somme des deux blocs diagonaux.

\begin{remarque}[Pourquoi cette définition]
Toutes les prédictions qui portent sur $A$ seul viennent d'opérateurs de la forme $O_A\otimes I$, et
\[
  \operatorname{Tr}\big(\rho_{AB}\,(O_A\otimes I)\big)=\operatorname{Tr}\big(\rho_A\,O_A\big)
\]
(on le vérifie sur $\ket{a}\bra{a'}\otimes\ket{b}\bra{b'}$, puis par linéarité). La matrice $\rho_A$
contient donc exactement ce que l'on peut prédire sur $A$ sans regarder $B$.
\end{remarque}

\paragraph{Exemple 1 : état séparable $\ket{0}\otimes\ket{+}$.} Ici
$\rho_{AB}=\ket{0}\bra{0}\otimes\ket{+}\bra{+}$, donc
\[
  \rho_A=\ket{0}\bra{0}\,\braket{+}{+}=\ket{0}\bra{0},
  \qquad
  \rho_B=\braket{0}{0}\,\ket{+}\bra{+}=\ket{+}\bra{+}.
\]
Avec les matrices,
\[
  \rho_{AB}=\frac12\begin{pmatrix}1&1&0&0\\1&1&0&0\\0&0&0&0\\0&0&0&0\end{pmatrix},
  \qquad
  R_{00}=\frac12\begin{pmatrix}1&1\\1&1\end{pmatrix},\quad R_{01}=R_{10}=R_{11}=0,
\]
d'où
\[
  \rho_A=\begin{pmatrix}\operatorname{Tr}R_{00}&0\\0&0\end{pmatrix}=\begin{pmatrix}1&0\\0&0\end{pmatrix},
  \qquad
  \rho_B=R_{00}+R_{11}=\frac12\begin{pmatrix}1&1\\1&1\end{pmatrix}.
\]
Un qubit d'un état séparable garde son propre état pur.

\paragraph{Exemple 2 : état de Bell $\ket{\Phi^+}$.} On écrit
$\rho_{AB}=\ket{\Phi^+}\bra{\Phi^+}=\frac12\big(\ket{00}\bra{00}+\ket{00}\bra{11}+\ket{11}\bra{00}
+\ket{11}\bra{11}\big)$. Chaque terme est un produit d'opérateurs à un qubit, par exemple
$\ket{00}\bra{11}=\ket{0}\bra{1}\otimes\ket{0}\bra{1}$, donc
\begin{align*}
  \rho_A=\operatorname{Tr}_B(\rho_{AB})
  &=\tfrac12\Big(\ket{0}\bra{0}\,\braket{0}{0}+\ket{0}\bra{1}\,\braket{1}{0}
      +\ket{1}\bra{0}\,\braket{0}{1}+\ket{1}\bra{1}\,\braket{1}{1}\Big)\\
  &=\tfrac12\Big(\ket{0}\bra{0}\cdot1+\ket{0}\bra{1}\cdot0+\ket{1}\bra{0}\cdot0+\ket{1}\bra{1}\cdot1\Big)
   =\tfrac12\big(\ket{0}\bra{0}+\ket{1}\bra{1}\big)=\tfrac12\,\mathrm{Id}.
\end{align*}
Par les blocs,
\[
  \rho_{AB}=\frac12\begin{pmatrix}1&0&0&1\\0&0&0&0\\0&0&0&0\\1&0&0&1\end{pmatrix},
  \qquad
  R_{00}=\frac12\begin{pmatrix}1&0\\0&0\end{pmatrix},\quad
  R_{11}=\frac12\begin{pmatrix}0&0\\0&1\end{pmatrix},\quad
  \operatorname{Tr}R_{01}=\operatorname{Tr}R_{10}=0,
\]
d'où $\rho_A=\frac12\mathrm{Id}$ et $\rho_B=R_{00}+R_{11}=\frac12\mathrm{Id}$ : le même résultat.

Le qubit $A$ est donc dans l'état \emph{complètement mixte} ($\operatorname{Tr}\rho_A^2=\frac14+\frac14=\frac12$),
au centre de la boule de Bloch (section~\ref{sec:bloch}), alors que le couple est dans un état pur
($\operatorname{Tr}\rho_{AB}^2=1$, section~\ref{sec:circuit-bell}). On connaît parfaitement l'état du
couple mais pas celui d'un qubit pris seul : l'information est entièrement contenue dans les
corrélations entre $A$ et $B$. Mesuré seul, $A$ donne $0$ ou $1$ avec la probabilité $\frac12$, dans
toutes les bases.

\paragraph{Formule pratique pour un état pur.} Pour $\ket{\Psi}=\sum_{ij}c_{ij}\ket{ij}$, notons
$C=\begin{pmatrix}c_{00}&c_{01}\\c_{10}&c_{11}\end{pmatrix}$ la matrice $2\times2$ des coefficients
(ligne $i$ : état de $A$, colonne $j$ : état de $B$). Comme $\rho_{(ij),(i'j')}=c_{ij}\,c_{i'j'}^*$,
\[
  (\rho_A)_{ii'}=\sum_jc_{ij}\,c_{i'j}^*,\qquad\text{soit}\qquad \rho_A=C\,C^\dagger .
\]
\emph{Exemple 3.} Pour $\ket{\Psi}=\frac12\ket{00}+\frac12\ket{01}+\frac{1}{\sqrt2}\ket{11}$ (exemple~2
de la section~\ref{sec:mesure2}), $C=\begin{pmatrix}\frac12&\frac12\\0&\frac{1}{\sqrt2}\end{pmatrix}$ et
\[
  \rho_A=CC^\dagger=\begin{pmatrix}\frac14+\frac14&0+\frac{1}{2\sqrt2}\\[0.4ex]\frac{1}{2\sqrt2}&\frac12\end{pmatrix}
  =\begin{pmatrix}\frac12&\frac{\sqrt2}{4}\\[0.4ex]\frac{\sqrt2}{4}&\frac12\end{pmatrix},
  \qquad
  \operatorname{Tr}\rho_A^2=\tfrac14+\tfrac14+2\cdot\tfrac18=\tfrac34 .
\]
Le qubit $A$ est mixte, mais pas complètement : cet état est intriqué sans être un état de Bell.

\begin{resultat}{Pureté d'un qubit et séparabilité}
Pour un état pur $\ket{\Psi}=\sum c_{ij}\ket{ij}$ de deux qubits, avec $D=c_{00}c_{11}-c_{01}c_{10}$,
\[
  \operatorname{Tr}\rho_A^2=\operatorname{Tr}\rho_B^2=1-2|D|^2 .
\]
Donc $\ket{\Psi}$ est séparable ($D=0$) si et seulement si $\rho_A$ est pur. Plus $|D|$ est grand,
plus $\rho_A$ est mixte, jusqu'à $|D|=\frac12$ et $\rho_A=\mathrm{Id}/2$ pour les états de Bell.
\end{resultat}

\emph{Démonstration.} La matrice $\rho_A=CC^\dagger$ a pour trace $\sum|c_{ij}|^2=1$ et pour
déterminant $\det C\,\overline{\det C}=|D|^2$. Si $\lambda_1,\lambda_2$ sont ses valeurs propres,
$\lambda_1+\lambda_2=1$ et $\operatorname{Tr}\rho_A^2=\lambda_1^2+\lambda_2^2=1-2\lambda_1\lambda_2
=1-2|D|^2$ ; le calcul est le même pour $\rho_B$. \hfill$\square$

Comme $\operatorname{Tr}\rho_A^2\geq\frac12$ pour un qubit, on en déduit $|D|\leq\frac12$. Avec
$\operatorname{Tr}\rho_A^2=\frac{1+|\vec r|^2}{2}$ (section~\ref{sec:bloch}), le vecteur de Bloch de $A$
a pour norme $|\vec r|=\sqrt{1-4|D|^2}$ : il part de la sphère pour un état séparable et se rapproche
du centre quand l'intrication croît (figure~\ref{fig:trace-bloch}).

\begin{center}
\renewcommand{\arraystretch}{1.5}
\begin{tabular}{@{}lcccl@{}}
\toprule
\textbf{État du couple} & $D$ & $\operatorname{Tr}\rho_A^2$ & $|\vec r|$ & \textbf{Qubit $A$} \\
\midrule
$\ket{0}\otimes\ket{+}$ & $0$ & $1$ & $1$ & pur, sur la sphère \\
$\frac12\ket{00}+\frac12\ket{01}+\frac{1}{\sqrt2}\ket{11}$ & $\frac{\sqrt2}{4}$ & $\frac34$
  & $\frac{1}{\sqrt2}$ & mixte, dans la boule \\
$\ket{\Phi^+}$ (comme les trois autres états de Bell) & $\pm\frac12$ & $\frac12$ & $0$
  & complètement mixte, au centre \\
\bottomrule
\end{tabular}
\end{center}

\begin{figure}[!ht]
\centering
\adjustbox{max width=\linewidth}{%
\begin{tikzpicture}[>=Stealth, font=\small]
  \def\R{2.3}
  \fill[insagris!12] (0,0) circle (\R);
  \draw[insagris, dotted] (-\R,0) -- (\R,0);
  \draw[insagris, dotted] (0,-\R) -- (0,\R);
  \draw[insagris, dashed] (0,0) circle ({\R*0.70711});
  \draw[insarouge, very thick] (0,0) circle (\R);
  \node[font=\footnotesize, text=insagris, anchor=south west] at (\R,0) {$x$};
  \node[font=\footnotesize, text=insagris, anchor=south west] at (0.03,\R) {$z$};
  % vecteurs de Bloch de rho_A pour trois états du couple
  \fill[insarouge] (0,\R) circle (2.4pt);
  \fill[insagris!80!black] ({\R*0.70711},0) circle (2.4pt);
  \fill[insagris!80!black] (0,0) circle (2.4pt);
  \draw[insagris] (0,\R) -- (3.1,\R);
  \node[anchor=west, font=\footnotesize, align=left] at (3.1,\R)
    {$\ket{0}\otimes\ket{+}$ : $\rho_A=\ket{0}\bra{0}$\\ $D=0$, \ $|\vec r|=1$ (état pur)};
  \draw[insagris] ({\R*0.70711},0) -- (3.1,0);
  \node[anchor=west, font=\footnotesize, align=left] at (3.1,0)
    {$\frac12\ket{00}+\frac12\ket{01}+\frac{1}{\sqrt2}\ket{11}$\\
     $D=\frac{\sqrt2}{4}$, \ $|\vec r|=\frac{1}{\sqrt2}$ (mixte)};
  \draw[insagris] (0,0) -- (3.1,-\R);
  \node[anchor=west, font=\footnotesize, align=left] at (3.1,-\R)
    {$\ket{\Phi^+}$ : $\rho_A=\mathrm{Id}/2$\\ $D=\frac12$, \ $\vec r=\vec 0$ (centre)};
\end{tikzpicture}}
\caption{Le vecteur de Bloch de $\rho_A$ (boule de Bloch vue en coupe) pour trois états du couple
$(A,B)$. Le cercle pointillé est la sphère de rayon $1/\sqrt2$. Plus l'état est intriqué, plus le
point s'éloigne de la sphère des états purs : $|\vec r|=\sqrt{1-4|D|^2}$.}
\label{fig:trace-bloch}
\end{figure}

% ---------------------------------------------------------------------
\subsection{Information : variables aléatoires et états quantiques}\label{sec:info}
% ---------------------------------------------------------------------

\begin{definition}{Entropie, information mutuelle, entropie conditionnelle}
Soit $A$ une variable aléatoire de loi $p(a)$. Son \emph{entropie de Shannon}, en bits, est
\[
  H(A)=-\sum_a p(a)\log_2p(a)\qquad(\text{avec }0\log_20=0).
\]
Elle mesure l'incertitude moyenne sur $A$ : une pièce équilibrée a $H=1$ bit. Pour un couple
$(A,B)$ de loi $p(a,b)$, $H(A,B)=-\sum_{a,b}p(a,b)\log_2p(a,b)$, et
\begin{align*}
  I(A;B)&=H(A)+H(B)-H(A,B) &&\text{(information mutuelle)},\\
  H(A|B)&=H(A,B)-H(B)=H(A)-I(A;B) &&\text{(entropie conditionnelle)}.
\end{align*}
\end{definition}

Les deux écritures de $H(A|B)$ sont égales : la définition de $I(A;B)$ donne
$H(A,B)=H(A)+H(B)-I(A;B)$, donc $H(A,B)-H(B)=H(A)-I(A;B)$. L'information mutuelle $I(A;B)$
mesure ce que la connaissance de $B$ apprend sur $A$ ; $H(A|B)$ est l'incertitude qui reste sur
$A$ quand on connaît $B$.

\subsubsection{Deux variables aléatoires indépendantes}

Alice et Bob lancent chacun une pièce équilibrée, sans lien entre les deux lancers :
$p(a,b)=\frac14$ pour les quatre couples. Chaque pièce donne $0$ ou $1$ avec la probabilité $\frac12$, donc
\begin{align*}
  H(A)&=-\Big(\tfrac12\log_2\tfrac12+\tfrac12\log_2\tfrac12\Big)=1\text{ bit},\\
  H(B)&=1\text{ bit (même calcul)},\\
  H(A,B)&=-4\cdot\tfrac14\log_2\tfrac14=2\text{ bits},\\
  I(A;B)&=H(A)+H(B)-H(A,B)=1+1-2=0,\\
  H(A|B)&=H(A)-I(A;B)=1-0=H(A)=1\text{ bit}.
\end{align*}
Connaître $B$ n'apprend rien sur $A$.

\subsubsection{Deux variables solidaires}

Les deux pièces tombent toujours du même côté : $p(0,0)=p(1,1)=\frac12$ et
$p(0,1)=p(1,0)=0$. Les lois de chaque pièce sont inchangées ($p(a{=}0)=p(0,0)+p(0,1)=\frac12$,
etc.), donc $H(A)=H(B)=1$ bit, mais
\begin{align*}
  H(A,B)&=-2\cdot\tfrac12\log_2\tfrac12-2\cdot0\log_20=1\text{ bit},\\
  I(A;B)&=1+1-1=1\text{ bit},\qquad H(A|B)=1-1=0.
\end{align*}
Connaître $B$ détermine $A$ : plus aucune incertitude ne reste.

\subsubsection{Deux qubits dans un état de Bell}

Pour un système quantique, on remplace la loi de probabilité par la matrice de densité.

\begin{definition}{Entropie de von Neumann}
L'\emph{entropie de von Neumann} d'un état $\rho$, de valeurs propres $\lambda_i$, est
\[
  S(\rho)=-\operatorname{Tr}(\rho\log_2\rho)=-\sum_i\lambda_i\log_2\lambda_i .
\]
Elle vaut $0$ pour un état pur et $1$ bit pour le qubit complètement mélangé $\mathrm{Id}/2$.
En quantique, $H(A)=S(\rho_A)$, $H(B)=S(\rho_B)$ et $H(A,B)=S(\rho_{AB})$ ; les définitions de
$I(A;B)$ et de $H(A|B)$ sont les mêmes.
\end{definition}

On prend l'état de Bell $\ket{\Phi^+}$ produit par le circuit de la
section~\ref{sec:circuit-bell}. Le couple est dans un état pur,
$\rho_{AB}=\ket{\Phi^+}\bra{\Phi^+}$ ; chaque qubit pris seul est dans l'état
$\rho_A=\rho_B=\frac12\mathrm{Id}$ (trace partielle, section~\ref{sec:trace}).

Calcul des entropies :
\begin{itemize}[nosep]
  \item $\rho_{AB}$ est pur : ses valeurs propres sont $1,0,0,0$, donc
        $H(A,B)=S(\rho_{AB})=-1\cdot\log_21=0$ ;
  \item $\rho_A$ et $\rho_B$ ont pour valeurs propres $\frac12,\frac12$, donc
        $H(A)=H(B)=-\big(\frac12\log_2\frac12+\frac12\log_2\frac12\big)=1$ bit.
\end{itemize}
D'où
\begin{align*}
  I(A;B)&=H(A)+H(B)-H(A,B)=1+1-0=2\text{ bits},\\
  H(A|B)&=H(A,B)-H(B)=0-1=-1\text{ bit}\qquad(=H(A)-I(A;B)=1-2).
\end{align*}

\begin{remarque}[Une entropie conditionnelle négative]
Pour des variables classiques, $p(a,b)\leq p(b)$, donc $-\log_2p(a,b)\geq-\log_2p(b)$ et
$H(A,B)\geq H(B)$ : on a toujours $H(A|B)\geq0$, et $I(A;B)\leq H(A)$. Ici $H(A|B)=-1$ et
$I(A;B)=2>H(A)$ : ces valeurs n'ont pas d'équivalent classique, c'est une signature de
l'intrication. Le résultat est surprenant mais mathématiquement correct.
\end{remarque}

\begin{center}
\renewcommand{\arraystretch}{1.4}
\begin{tabular}{@{}lccc@{}}
\toprule
 & indépendantes & solidaires & état de Bell \\
\midrule
$H(A)$ & $1$ & $1$ & $1$ \\
$H(B)$ & $1$ & $1$ & $1$ \\
$H(A,B)$ & $2$ & $1$ & $0$ \\
$I(A;B)$ & $0$ & $1$ & $2$ \\
$H(A|B)$ & $1$ & $0$ & $-1$ \\
\bottomrule
\end{tabular}
\end{center}
(en bits ; les deux premières colonnes sont classiques, la troisième est quantique).

\begin{figure}[!ht]
\centering
\adjustbox{max width=\linewidth}{%
\begin{tikzpicture}[>=Stealth, font=\small]
  \def\bw{0.8}
  % axe des ordonnées (bits)
  \draw[->] (0,-1.5) -- (0,2.8) node[above] {bits};
  \foreach \v in {-1,0,1,2} {
    \draw (-0.07,\v) -- (0.07,\v);
    \node[left, font=\footnotesize] at (-0.1,\v) {$\v$};
  }
  \draw[insagris!60] (0,0) -- (11.4,0);
  % limite classique de I(A;B) : H(A) = 1
  \draw[dashed, insarouge] (0,1) -- (11.4,1);
  \node[anchor=west, font=\scriptsize, text=insarouge, align=left] at (11.4,1)
    {limite classique\\ $I(A;B)\leq H(A)=1$};
  \node[anchor=west, font=\scriptsize, text=insagris, align=left] at (11.4,0)
    {limite classique\\ $H(A|B)\geq0$};
  % barres : (groupe, H(A,B), I(A;B), H(A|B))
  \foreach \g/\hab/\mi/\hc in {0/2/0/1, 1/1/1/0, 2/0/2/-1} {
    \pgfmathsetmacro{\x}{0.9+3.7*\g}
    \fill[insagris!65] (\x,0) rectangle (\x+\bw,\hab);
    \fill[insarouge] (\x+\bw+0.1,0) rectangle (\x+2*\bw+0.1,\mi);
    \fill[black!80] (\x+2*\bw+0.2,0) rectangle (\x+3*\bw+0.2,\hc);
    \node[above, font=\footnotesize] at (\x+0.5*\bw,{max(\hab,0)}) {$\hab$};
    \node[above, font=\footnotesize] at (\x+1.5*\bw+0.1,{max(\mi,0)}) {$\mi$};
    \ifnum\hc<0
      \node[below, font=\footnotesize] at (\x+2.5*\bw+0.2,\hc) {$\hc$};
    \else
      \node[above, font=\footnotesize] at (\x+2.5*\bw+0.2,\hc) {$\hc$};
    \fi
  }
  \node[align=center, anchor=north] at (2.15,-1.65) {variables\\ indépendantes};
  \node[align=center, anchor=north] at (5.85,-1.65) {variables\\ solidaires};
  \node[align=center, anchor=north] at (9.55,-1.65) {état de Bell\\ $\ket{\Phi^+}$};
  % légende
  \fill[insagris!65] (0.9,3.35) rectangle (1.2,3.65);
  \node[anchor=west, font=\footnotesize] at (1.25,3.5) {$H(A,B)$};
  \fill[insarouge] (3.6,3.35) rectangle (3.9,3.65);
  \node[anchor=west, font=\footnotesize] at (3.95,3.5) {$I(A;B)$};
  \fill[black!80] (6.3,3.35) rectangle (6.6,3.65);
  \node[anchor=west, font=\footnotesize] at (6.65,3.5) {$H(A|B)$};
\end{tikzpicture}}
\caption{Entropies (en bits) pour deux pièces indépendantes, deux pièces solidaires et un état de Bell.
Dans les trois cas $H(A)=H(B)=1$. Les deux premiers groupes respectent les bornes classiques
$I(A;B)\leq H(A)$ et $H(A|B)\geq0$ ; l'état de Bell les dépasse ($I=2$, $H(A|B)=-1$).}
\label{fig:entropies}
\end{figure}

% ---------------------------------------------------------------------
\subsection{Téléportation quantique}\label{sec:teleportation}
% ---------------------------------------------------------------------

\emph{Problème.} Alice détient un qubit $Q$ dans un état qu'elle ne connaît pas,
$\ket{\varphi}=\alpha\ket{0}+\beta\ket{1}$ avec $|\alpha|^2+|\beta|^2=1$, et elle veut que Bob
obtienne cet état sur l'un de ses qubits. Elle ne peut pas mesurer $Q$ pour connaître
$\alpha$ et $\beta$ (une mesure ne donne qu'un bit et détruit l'état), elle ne peut pas le
copier (théorème de non-clonage, remarque en fin de section) et elle n'envoie pas le qubit
lui-même. Le protocole est dû à \textcite{bennett1993}.

\begin{definition}{Protocole de téléportation quantique}
\textbf{Ressources.}
\begin{itemize}[nosep]
  \item Un \emph{e-bit} : Alice et Bob partagent une paire de qubits $(A,B)$ dans l'état de Bell
        $\ket{\Phi^+}_{BA}=\frac{1}{\sqrt2}\big(\ket{00}+\ket{11}\big)$ ; Alice détient $A$,
        Bob détient $B$.
  \item Le qubit $Q$ à téléporter, chez Alice.
  \item Un canal classique (deux bits).
\end{itemize}
Le système est écrit $(B,A,Q)$ : $\ket{b,a,q}=\ket{b}_B\otimes\ket{a}_A\otimes\ket{q}_Q$.

\textbf{Étapes.}
\begin{enumerate}[nosep]
  \item Alice applique un CNOT de contrôle $Q$ et de cible $A$.
  \item Alice applique une porte $H$ sur $Q$.
  \item Alice mesure $A$ et $Q$ dans la base de calcul, obtient deux bits $(a,q)$ et les envoie
        à Bob.
  \item Bob applique $X$ si $a=1$, puis $Z$ si $q=1$, à son qubit $B$.
\end{enumerate}
\end{definition}

La figure~\ref{fig:tele-schema} montre l'organisation du protocole et la
figure~\ref{fig:teleportation} son circuit.

\begin{figure}[!ht]
\centering
\adjustbox{max width=\linewidth}{%
\begin{tikzpicture}[>=Stealth, font=\small,
    qb/.style={circle, draw=insarouge, fill=insarouge!8, thick, minimum size=8mm, inner sep=0pt},
    zone/.style={draw=insagris, rounded corners=4pt, thick, dashed, inner sep=8pt},
    num/.style={circle, fill=insarouge, text=white, font=\footnotesize\bfseries, inner sep=0pt,
                minimum size=4.6mm},
    etape/.style={anchor=west, font=\footnotesize, text=insarouge},
    onde/.style={thick, insarouge, decorate,
                 decoration={snake, amplitude=1.3pt, segment length=7pt}}]
  % Alice : Q (état inconnu) et A
  \node[qb] (Q) at (0,1.3) {$Q$};
  \node[qb] (A) at (0,-0.5) {$A$};
  \node[zone, fit=(Q)(A), label={[text=insagris]above:Alice}] (alice) {};
  \node[anchor=east] at (-0.95,1.3) {$\ket{\varphi}=\alpha\ket{0}+\beta\ket{1}$};
  \draw[insarouge, thick] (Q) -- (A);
  % Bob : B
  \node[qb] (B) at (9,-0.5) {$B$};
  \node[zone, fit=(B), label={[text=insagris]above:Bob}] (bob) {};
  % source de la paire de Bell
  \node[draw=insarouge, fill=insarouge!10, rounded corners=2pt, thick, minimum width=3.2cm,
        minimum height=0.9cm] (src) at (4.5,-2.3) {source $\ket{\Phi^+}_{BA}$};
  \draw[onde] (src.west) -| (A.south);
  \draw[onde] (src.east) -| (B.south);
  % étapes d'Alice : pastilles sur le bord droit de sa zone
  \node[num] (n1) at (alice.east |- Q) {1};
  \node[num] (n2) at (alice.east) {2};
  \node[num] (n3) at (alice.east |- A) {3};
  \node[etape] at (n1.east) {\ CNOT de $Q$ vers $A$};
  \node[etape] at (n2.east) {\ porte $H$ sur $Q$};
  % canal classique
  \draw[double, double distance=1.2pt, ->, insagris] (n3.east) -- (bob.west);
  \node[above, font=\footnotesize] at (4.6,-0.45) {deux bits classiques $(a,q)$};
  \node[below, font=\footnotesize, text=insarouge] at (4.6,-0.55) {mesure de $A$ et $Q$ (étape 3)};
  % correction de Bob
  \node[num] (n4) at (bob.east) {4};
  \node[etape, align=left] at (n4.east) {\ $X$ si $a=1$, puis $Z$ si $q=1$};
  \node[anchor=west, font=\footnotesize, align=left] at ($(n4.east)+(0.05,-0.8)$)
    {\ $B$ est alors dans l'état $\ket{\varphi}$};
\end{tikzpicture}}
\caption{Organisation de la téléportation. Alice possède le qubit $Q$ à téléporter et la moitié $A$
d'une paire de Bell, Bob possède l'autre moitié $B$. (1)~CNOT de $Q$ vers $A$, (2)~porte $H$ sur $Q$,
(3)~mesure de $A$ et $Q$, dont les résultats $(a,q)$ partent par un canal classique, (4)~corrections de
Bob. Aucun qubit ne voyage entre Alice et Bob.}
\label{fig:tele-schema}
\end{figure}

\begin{figure}[!ht]
\centering
\begin{adjustbox}{max width=\linewidth}
\begin{tikzpicture}[>=Stealth, font=\small]
  % fils quantiques : Q en haut, A au milieu, B en bas
  \draw (0,2.4) -- (5.3,2.4);
  \draw (0,1.2) -- (5.3,1.2);
  \draw (0,0) -- (10.6,0);
  \node[anchor=east] at (0,2.4) {$\ket{\varphi}_Q$};
  \node[anchor=south west, text=insagris, font=\footnotesize] at (0.05,1.2) {$A$};
  \node[anchor=south west, text=insagris, font=\footnotesize] at (0.05,0) {$B$};
  % paire de Bell partagée
  \draw[insagris, thick] (-0.1,1.5) -- (-0.25,1.5) -- (-0.25,-0.3) -- (-0.1,-0.3);
  \node[anchor=east] at (-0.35,0.6) {$\ket{\Phi^+}_{BA}$};
  % CNOT : contrôle Q, cible A
  \draw[insarouge, thick] (2.6,2.4) -- (2.6,1.2);
  \pic at (2.6,2.4) {ctl};
  \pic at (2.6,1.2) {cible};
  % Hadamard sur Q
  \node[porte] at (4.2,2.4) {$H$};
  % mesures d'Alice
  \pic at (5.6,2.4) {mesure};
  \pic at (5.6,1.2) {mesure};
  % fils classiques (doubles)
  \draw[double, double distance=1.2pt, insagris] (5.9,1.2) -- (7.6,1.2) -- (7.6,0.35);
  \draw[double, double distance=1.2pt, insagris] (5.9,2.4) -- (9.6,2.4) -- (9.6,0.35);
  \node[above, font=\footnotesize] at (6.5,1.2) {$a$};
  \node[above, font=\footnotesize] at (6.5,2.4) {$q$};
  % corrections de Bob
  \node[porte] at (7.6,0) {$X$};
  \node[porte] at (9.6,0) {$Z$};
  \node[anchor=west] at (10.6,0) {$\ket{\varphi}_B$};
  % états intermédiaires
  \foreach \x in {1.4,3.4,4.9}
    \draw[dashed, insagris!70] (\x,2.8) -- (\x,-0.5);
  \node[anchor=north, font=\footnotesize] at (1.4,-0.5) {$\ket{\Pi_0}$};
  \node[anchor=north, font=\footnotesize] at (3.4,-0.5) {$\ket{\Pi_1}$};
  \node[anchor=north, font=\footnotesize] at (4.9,-0.5) {$\ket{\Pi_2}$};
\end{tikzpicture}
\end{adjustbox}
\caption{Circuit de téléportation. Alice détient $Q$ (dans $\ket{\varphi}=\alpha\ket{0}+\beta\ket{1}$)
et $A$, Bob détient $B$ ; $A$ et $B$ forment la paire $\ket{\Phi^+}_{BA}$. Les traits doubles
sont des fils \emph{classiques} qui portent les résultats $a$ (mesure de $A$) et $q$ (mesure de
$Q$) ; $X$ est appliqué si $a=1$, $Z$ si $q=1$. Le ket s'écrit $\ket{b,a,q}$ : le fil du haut
($Q$) est le dernier chiffre.}
\label{fig:teleportation}
\end{figure}

\subsubsection{Calcul pas à pas}

\paragraph{État initial $\ket{\Pi_0}$.} Alice et Bob partagent la paire, et $Q$ est indépendant
d'elle :
\begin{align*}
  \ket{\Pi_0}&=\ket{\Phi^+}_{BA}\otimes\ket{\varphi}_Q
    =\frac{1}{\sqrt2}\big(\ket{00}+\ket{11}\big)\otimes\big(\alpha\ket{0}+\beta\ket{1}\big)\\
  &=\frac{1}{\sqrt2}\Big(\alpha\ket{00}\ket{0}+\beta\ket{00}\ket{1}+\alpha\ket{11}\ket{0}+\beta\ket{11}\ket{1}\Big)\\
  &=\frac{\alpha\ket{000}+\beta\ket{001}+\alpha\ket{110}+\beta\ket{111}}{\sqrt2}.
\end{align*}

\paragraph{Après le CNOT $Q\to A$ : $\ket{\Pi_1}$.} Le CNOT (contrôle $Q$, dernier chiffre ;
cible $A$, chiffre du milieu) agit ainsi : $\ket{b,a,q}\to\ket{b,\,a\oplus q,\,q}$. Terme par terme :
\begin{center}
\renewcommand{\arraystretch}{1.3}
\begin{tabular}{@{}cccc@{}}
\toprule
Terme & $q$ & Effet sur $A$ & Résultat \\
\midrule
$\alpha\ket{000}$ & $0$ & aucun & $\alpha\ket{000}$ \\
$\beta\ket{001}$ & $1$ & $0\to1$ & $\beta\ket{011}$ \\
$\alpha\ket{110}$ & $0$ & aucun & $\alpha\ket{110}$ \\
$\beta\ket{111}$ & $1$ & $1\to0$ & $\beta\ket{101}$ \\
\bottomrule
\end{tabular}
\end{center}
donc
\[
  \ket{\Pi_1}=\frac{\alpha\ket{000}+\beta\ket{011}+\alpha\ket{110}+\beta\ket{101}}{\sqrt2}.
\]
On isole ensuite le dernier chiffre, celui de $Q$, qui va changer à l'étape suivante
(le premier ket est celui du couple $\ket{b,a}$) :
\[
  \ket{\Pi_1}=\frac{\alpha\ket{00}\otimes\ket{0}+\beta\ket{01}\otimes\ket{1}
    +\alpha\ket{11}\otimes\ket{0}+\beta\ket{10}\otimes\ket{1}}{\sqrt2}.
\]

\paragraph{Après la porte $H$ sur $Q$ : $\ket{\Pi_2}$.} On remplace $\ket{0}$ par
$H\ket{0}=\ket{+}$ et $\ket{1}$ par $H\ket{1}=\ket{-}$ dans le facteur $Q$ :
\[
  \ket{\Pi_2}=\frac{1}{\sqrt2}\Big(\alpha\ket{00}\otimes\ket{+}+\beta\ket{01}\otimes\ket{-}
    +\alpha\ket{11}\otimes\ket{+}+\beta\ket{10}\otimes\ket{-}\Big).
\]
Avec $\ket{\pm}=\frac{1}{\sqrt2}(\ket{0}\pm\ket{1})$, le préfacteur devient
$\frac{1}{\sqrt2}\cdot\frac{1}{\sqrt2}=\frac12$ :
\begin{align*}
  \ket{\Pi_2}&=\frac12\Big(\alpha\ket{00}\big(\ket{0}+\ket{1}\big)+\beta\ket{01}\big(\ket{0}-\ket{1}\big)
     +\alpha\ket{11}\big(\ket{0}+\ket{1}\big)+\beta\ket{10}\big(\ket{0}-\ket{1}\big)\Big)\\
  &=\frac12\Big(\alpha\ket{000}+\alpha\ket{001}+\beta\ket{010}-\beta\ket{011}
     +\alpha\ket{110}+\alpha\ket{111}+\beta\ket{100}-\beta\ket{101}\Big).
\end{align*}

\paragraph{Regroupement selon le résultat d'Alice.} Alice mesurera $A$ et $Q$, donc les deux
derniers chiffres. On regroupe les huit termes qui ont les mêmes chiffres $(a,q)$, en mettant
le ket de Bob en premier :
\begin{center}
\renewcommand{\arraystretch}{1.3}
\begin{tabular}{@{}cll@{}}
\toprule
$(a,q)$ & Termes de $\ket{\Pi_2}$ (facteur $\frac12$ omis) & Ket de Bob \\
\midrule
$(0,0)$ & $\alpha\ket{000}+\beta\ket{100}$ & $(\alpha\ket{0}+\beta\ket{1})\ket{00}$ \\
$(0,1)$ & $\alpha\ket{001}-\beta\ket{101}$ & $(\alpha\ket{0}-\beta\ket{1})\ket{01}$ \\
$(1,0)$ & $\beta\ket{010}+\alpha\ket{110}$ & $(\beta\ket{0}+\alpha\ket{1})\ket{10}$ \\
$(1,1)$ & $-\beta\ket{011}+\alpha\ket{111}$ & $(-\beta\ket{0}+\alpha\ket{1})\ket{11}$ \\
\bottomrule
\end{tabular}
\end{center}
d'où
\[
  \ket{\Pi_2}=\frac12\Big[\big(\alpha\ket{0}+\beta\ket{1}\big)\ket{00}
    +\big(\alpha\ket{0}-\beta\ket{1}\big)\ket{01}
    +\big(\beta\ket{0}+\alpha\ket{1}\big)\ket{10}
    +\big(\alpha\ket{1}-\beta\ket{0}\big)\ket{11}\Big].
\]
Dans chaque terme, le ket de gauche est l'état de $B$ et le ket de droite $\ket{a,q}$ celui du
couple $(A,Q)$ d'Alice.

\subsubsection{Mesure d'Alice et corrections de Bob}

Alice mesure $A$ et $Q$. Comme à la section~\ref{sec:mesure2}, la probabilité d'un résultat
$(a,q)$ est le carré de la norme du terme correspondant, et l'état de Bob devient le ket
associé (déjà normé). Ici
\[
  P(a,q)=\Big\|\tfrac12\ket{\phi_{aq}}\Big\|^2=\tfrac14\big(|\alpha|^2+|\beta|^2\big)=\tfrac14
  \quad\text{pour les quatre résultats,}
\]
qui sont donc équiprobables et \emph{indépendants} de $\alpha$ et $\beta$ : les deux bits
envoyés ne révèlent rien sur $\ket{\varphi}$. Les kets de Bob sont
$\ket{\phi_{00}}=\alpha\ket{0}+\beta\ket{1}$, $\ket{\phi_{01}}=\alpha\ket{0}-\beta\ket{1}$,
$\ket{\phi_{10}}=\beta\ket{0}+\alpha\ket{1}$ et $\ket{\phi_{11}}=-\beta\ket{0}+\alpha\ket{1}$.

\begin{center}
\renewcommand{\arraystretch}{1.5}
\begin{tabular}{@{}cccll@{}}
\toprule
Résultat $(a,q)$ & Probabilité & État de Bob & Lien avec $\ket{\varphi}$ & Correction de Bob \\
\midrule
$(0,0)$ & $\frac14$ & $\alpha\ket{0}+\beta\ket{1}$ & $\ket{\varphi}$ & $I$ (rien) \\
$(0,1)$ & $\frac14$ & $\alpha\ket{0}-\beta\ket{1}$ & $Z\ket{\varphi}$ & $Z$ \\
$(1,0)$ & $\frac14$ & $\beta\ket{0}+\alpha\ket{1}$ & $X\ket{\varphi}$ & $X$ \\
$(1,1)$ & $\frac14$ & $-\beta\ket{0}+\alpha\ket{1}$ & $XZ\ket{\varphi}$ & $X$ puis $Z$ (opérateur $ZX$) \\
\bottomrule
\end{tabular}
\end{center}

\paragraph{Vérification des quatre corrections.} Avec $X\ket{0}=\ket{1}$, $X\ket{1}=\ket{0}$,
$Z\ket{0}=\ket{0}$ et $Z\ket{1}=-\ket{1}$ :
\begin{align*}
  (0,0):\quad & I\big(\alpha\ket{0}+\beta\ket{1}\big)=\alpha\ket{0}+\beta\ket{1},\\
  (0,1):\quad & Z\big(\alpha\ket{0}-\beta\ket{1}\big)=\alpha\ket{0}-\beta Z\ket{1}=\alpha\ket{0}+\beta\ket{1},\\
  (1,0):\quad & X\big(\beta\ket{0}+\alpha\ket{1}\big)=\beta\ket{1}+\alpha\ket{0}=\alpha\ket{0}+\beta\ket{1},\\
  (1,1):\quad & X\big(-\beta\ket{0}+\alpha\ket{1}\big)=-\beta\ket{1}+\alpha\ket{0}=\alpha\ket{0}-\beta\ket{1},
     \ \text{puis}\ Z\big(\alpha\ket{0}-\beta\ket{1}\big)=\alpha\ket{0}+\beta\ket{1}.
\end{align*}
Dans le dernier cas, avec les matrices,
$ZX=\begin{pmatrix}1&0\\0&-1\end{pmatrix}\begin{pmatrix}0&1\\1&0\end{pmatrix}
=\begin{pmatrix}0&1\\-1&0\end{pmatrix}$ et
$ZX\begin{pmatrix}-\beta\\ \alpha\end{pmatrix}=\begin{pmatrix}\alpha\\ \beta\end{pmatrix}$.
L'ordre compte : l'état de Bob est $XZ\ket{\varphi}$ ($Z\ket{\varphi}=\alpha\ket{0}-\beta\ket{1}$,
puis $X$ donne $\alpha\ket{1}-\beta\ket{0}$), et on l'annule avec
$(XZ)^{-1}=Z^{-1}X^{-1}=ZX$ : on applique d'abord $X$, puis $Z$, comme dans le circuit.
Dans les quatre cas, Bob retrouve exactement $\ket{\varphi}=\alpha\ket{0}+\beta\ket{1}$.

\subsubsection{Ce que Bob sait avant de recevoir les bits}

Tant que Bob ignore $(a,q)$, son qubit est dans le mélange des quatre kets $\ket{\phi_{aq}}$, chacun
avec la probabilité $\frac14$ :
\begin{align*}
  \rho_B&=\frac14\sum_{a,q}\ket{\phi_{aq}}\bra{\phi_{aq}}\\
  &=\frac14\Bigg[\begin{pmatrix}|\alpha|^2&\alpha\beta^*\\ \alpha^*\beta&|\beta|^2\end{pmatrix}
    +\begin{pmatrix}|\alpha|^2&-\alpha\beta^*\\ -\alpha^*\beta&|\beta|^2\end{pmatrix}
    +\begin{pmatrix}|\beta|^2&\alpha^*\beta\\ \alpha\beta^*&|\alpha|^2\end{pmatrix}
    +\begin{pmatrix}|\beta|^2&-\alpha^*\beta\\ -\alpha\beta^*&|\alpha|^2\end{pmatrix}\Bigg]\\
  &=\frac14\begin{pmatrix}2(|\alpha|^2+|\beta|^2)&0\\0&2(|\alpha|^2+|\beta|^2)\end{pmatrix}
   =\frac12\begin{pmatrix}1&0\\0&1\end{pmatrix}=\frac12\,\mathrm{Id}.
\end{align*}
Sans les deux bits, Bob a donc un qubit complètement mélangé, qui ne dépend pas de $\alpha$ ni de
$\beta$ : il n'a reçu aucune information sur $\ket{\varphi}$. C'est l'envoi classique de $(a,q)$
qui permet de la retrouver, donc rien ne voyage plus vite que la lumière.

\begin{remarque}[Bilan de la téléportation]
\begin{itemize}[nosep]
  \item Ressources : $1$ e-bit (la paire de Bell) et $2$ bits classiques pour transmettre
        \emph{un qubit} d'état quelconque, sans qu'Alice connaisse $\alpha$ ni $\beta$.
  \item Pas de clonage : après la mesure, $A$ et $Q$ sont dans l'état de base $\ket{a,q}$.
        Alice n'a plus $\ket{\varphi}$ : l'état a été \emph{déplacé}, pas copié. L'intrication
        entre $A$ et $B$ a été consommée.
  \item Les quatre corrections $I,Z,X,ZX$ sont des portes contrôlées par des bits classiques
        (section~\ref{sec:controlees}).
\end{itemize}
\end{remarque}

\begin{remarque}[Les gestes d'Alice forment une mesure de Bell]
Le CNOT de contrôle $Q$ sur $A$ suivi de $H$ sur $Q$ est l'\emph{inverse} du générateur de la
section~\ref{sec:circuit-bell} (les portes $H$ et CNOT sont leurs propres inverses, et l'ordre
s'inverse) : il envoie la base de Bell du couple $(A,Q)$ sur la base de calcul. Mesurer ensuite dans
la base de calcul revient donc à mesurer $(A,Q)$ dans la base de Bell. Les deux bits $(a,q)$
disent \emph{quel} état de Bell Alice a trouvé, ce qui explique les quatre corrections de Bob.
\end{remarque}

\begin{remarque}[Non-clonage]
Alice ne peut pas simplement copier $\ket{\varphi}$ \parencite{wootters1982}. Supposons une porte
unitaire $U$ telle que $U\big(\ket{\psi}\otimes\ket{0}\big)=\ket{\psi}\otimes\ket{\psi}$ pour tout
$\ket{\psi}$. Elle copie $\ket{0}$ et $\ket{1}$, mais par linéarité
\[
  U\big(\ket{+}\otimes\ket{0}\big)=\tfrac{1}{\sqrt2}\big(\ket{00}+\ket{11}\big)
  \neq\ket{+}\otimes\ket{+}=\tfrac12\big(\ket{00}+\ket{01}+\ket{10}+\ket{11}\big),
\]
ce qui contredit l'hypothèse. Le CNOT copie un bit classique, pas une superposition (remarque de la
section~\ref{sec:cnot}). La téléportation respecte cette interdiction : l'état est déplacé, il n'est
jamais présent en deux endroits.
\end{remarque}

% ---------------------------------------------------------------------
\subsection{Corrélations non locales : le jeu CHSH}\label{sec:chsh}
% ---------------------------------------------------------------------

Les états de Bell donnent des résultats corrélés en base $Z$ comme en base $X$
(section~\ref{sec:bell}). Mais deux enveloppes contenant la même carte, préparées à l'avance, donnent
aussi des résultats égaux : des corrélations dans une seule base ne prouvent pas l'intrication. Bell
\parencite{bell1964} a montré qu'on peut trancher en \emph{changeant de base de mesure} ; la version de
Clauser, Horne, Shimony et Holt \parencite{clauser1969} se présente comme un jeu.

\begin{definition}{Le jeu CHSH}
Alice et Bob, éloignés, ne communiquent plus une fois le jeu commencé. Un arbitre tire deux bits $x$ et
$y$ de façon uniforme et indépendante, donne $x$ à Alice et $y$ à Bob. Chacun répond par un bit : $a$ pour
Alice, $b$ pour Bob. Ils \emph{gagnent} si
\[
  a\oplus b=x\cdot y,
\]
c'est-à-dire si $a=b$ pour $(x,y)\in\{(0,0),(0,1),(1,0)\}$ et si $a\neq b$ pour $(x,y)=(1,1)$.
\end{definition}

\paragraph{Stratégie classique : au plus $\frac34$.} Une stratégie déterministe est $a=f(x)$,
$b=g(y)$. Si elle gagnait les quatre parties, on aurait
\[
  f(0)\oplus g(0)=0,\qquad f(0)\oplus g(1)=0,\qquad f(1)\oplus g(0)=0,\qquad f(1)\oplus g(1)=1 .
\]
En additionnant ces quatre égalités modulo $2$, chacune des valeurs $f(0),f(1),g(0),g(1)$ apparaît
deux fois : le membre de gauche vaut $0$, celui de droite vaut $1$. C'est absurde, donc au moins une
partie sur quatre est perdue et $P(\text{gain})\leq\frac34$. La stratégie «~toujours répondre $0$~»
atteint $\frac34$ (elle ne perd que pour $(1,1)$). Un hasard partagé ne change rien : il mélange des
stratégies déterministes, dont aucune ne dépasse $\frac34$.

\paragraph{Stratégie quantique.} Alice et Bob partagent avant le jeu une paire $\ket{\Phi^+}$ (un
e-bit), puis chacun mesure son qubit sans rien dire à l'autre. Pour un angle $\theta$, notons
\[
  \ket{\theta_+}=\cos\tfrac\theta2\,\ket{0}+\sin\tfrac\theta2\,\ket{1},
  \qquad
  \ket{\theta_-}=-\sin\tfrac\theta2\,\ket{0}+\cos\tfrac\theta2\,\ket{1}
\]
la base de mesure d'angle $\theta$ : c'est la direction $\theta$ dans le plan $xz$ de la sphère de
Bloch ($\theta=0$ : base $Z$ ; $\theta=\frac\pi2$ : base $X$). Le résultat vaut $0$ pour
$\ket{\theta_+}$ et $1$ pour $\ket{\theta_-}$. La stratégie est
\begin{itemize}[nosep]
  \item Alice : $\theta=0$ si $x=0$ (base $Z$), $\theta=\frac\pi2$ si $x=1$ (base $X$) ;
  \item Bob : $\theta=\frac\pi4$ si $y=0$, $\theta=-\frac\pi4$ si $y=1$ (bases propres de
        $\frac{Z\pm X}{\sqrt2}$).
\end{itemize}

\begin{figure}[!ht]
\centering
\adjustbox{max width=\linewidth}{%
\begin{tikzpicture}[>=Stealth, font=\footnotesize,
    joueur/.style={draw=insarouge, fill=insarouge!6, rounded corners=2pt, thick,
                   minimum width=2.1cm, minimum height=1.1cm, align=center},
    arb/.style={draw=insagris, fill=insagris!8, rounded corners=2pt, thick,
                minimum width=3.6cm, minimum height=0.8cm, align=center}]
  % --- le jeu ---
  \node[arb] (arb) at (3.2,2.7) {arbitre : $x,y$ au hasard};
  \node[joueur] (al) at (0,0) {Alice};
  \node[joueur] (bo) at (6.4,0) {Bob};
  \draw[->, thick, insagris] (arb.west) -| node[pos=0.2, above] {$x$} (al.north);
  \draw[->, thick, insagris] (arb.east) -| node[pos=0.2, above] {$y$} (bo.north);
  \draw[insarouge, thick, decorate, decoration={snake, amplitude=1.1pt, segment length=6pt}]
    (al.east) -- (bo.west);
  \node[insarouge, above=1pt] at (3.2,0.05) {$\ket{\Phi^+}$};
  \node[text=insagris, below=2pt, font=\scriptsize] at (3.2,-0.05) {paire intriquée};
  \draw[->, thick] (al.south) -- ++(0,-1.0) node[below] {$a$};
  \draw[->, thick] (bo.south) -- ++(0,-1.0) node[below] {$b$};
  \node[draw=insagris, rounded corners=2pt, inner sep=3pt] at (3.2,-1.75) {gain si $a\oplus b=x\cdot y$};
  % --- directions de mesure dans le plan xz de la sphère de Bloch ---
  \begin{scope}[shift={(11.9,0.3)}]
    \draw[insagris!70] (0,0) circle (1.7);
    \draw[insagris!50, dashed] (-1.7,0) -- (1.7,0);
    \draw[insagris!50, dashed] (0,-1.7) -- (0,1.7);
    \node[font=\scriptsize, text=insagris, anchor=north] at (0,-1.72) {$\ket{1}$};
    \node[font=\scriptsize, text=insagris, anchor=east] at (-1.72,-0.02) {$\ket{-}$};
    \draw[->, very thick, insarouge] (0,0) -- (90:1.7);
    \draw[->, very thick, insarouge] (0,0) -- (0:1.7);
    \draw[->, very thick, insagris] (0,0) -- (45:1.7);
    \draw[->, very thick, insagris] (0,0) -- (135:1.7);
    \node[insarouge, anchor=south, inner sep=2pt] at (0,1.72) {$x=0$ : $Z$ $(\ket{0})$};
    \node[insarouge, anchor=west, inner sep=2pt] at (1.72,0) {$x=1$ : $X$ $(\ket{+})$};
    \node[text=insagris, anchor=south west, align=left, inner sep=2pt] at (1.25,1.25) {$y=0$\\ $\frac{Z+X}{\sqrt2}$};
    \node[text=insagris, anchor=south east, align=right, inner sep=2pt] at (-1.25,1.25) {$y=1$\\ $\frac{Z-X}{\sqrt2}$};
    \draw[insagris] (0.85,0) arc (0:45:0.85);
    \node[font=\scriptsize, text=insagris] at (22.5:1.12) {$\frac\pi4$};
    \draw[insagris] (0,0.85) arc (90:135:0.85);
    \node[font=\scriptsize, text=insagris] at (112.5:1.12) {$\frac\pi4$};
    \draw[insagris] (0.6,0.6) arc (45:90:0.85);
    \node[font=\scriptsize, text=insagris] at (67.5:1.12) {$\frac\pi4$};
  \end{scope}
\end{tikzpicture}}
\caption{Le jeu CHSH (à gauche) et les quatre directions de mesure de la stratégie quantique (à
droite), dans le plan $xz$ de la sphère de Bloch : Alice en rouge, Bob en gris. Pour
$(x,y)=(0,0),(0,1),(1,0)$, les directions font un angle $\frac\pi4$ ; pour $(1,1)$, un angle
$\frac{3\pi}4$.}
\label{fig:chsh}
\end{figure}

\begin{resultat}{Corrélations de $\ket{\Phi^+}$ en bases tournées}
Si Alice mesure dans la base d'angle $\alpha$ et Bob dans la base d'angle $\beta$, alors
\[
  P(a=b)=\cos^2\frac{\alpha-\beta}{2}.
\]
\end{resultat}

\noindent\emph{Preuve.} Avec $c=\cos\frac\alpha2$ et $s=\sin\frac\alpha2$,
$\ket{\alpha_+\alpha_+}+\ket{\alpha_-\alpha_-}=(c^2+s^2)\ket{00}+(cs-sc)\big(\ket{01}+\ket{10}\big)+(s^2+c^2)\ket{11}
=\ket{00}+\ket{11}$ : on peut donc écrire $\ket{\Phi^+}=\frac{1}{\sqrt2}\big(\ket{\alpha_+\alpha_+}+\ket{\alpha_-\alpha_-}\big)$.
Si Alice obtient $a=0$ (probabilité $\frac12$), le qubit de Bob est dans $\ket{\alpha_+}$ et Bob obtient $b=0$
avec la probabilité $|\braket{\beta_+}{\alpha_+}|^2$, où $\braket{\beta_+}{\alpha_+}=\cos\frac\beta2\cos\frac\alpha2+\sin\frac\beta2\sin\frac\alpha2=\cos\frac{\alpha-\beta}{2}$.
Si Alice obtient $a=1$, Bob est dans $\ket{\alpha_-}$ et obtient $b=1$ avec la même probabilité, car
$\braket{\beta_-}{\alpha_-}=\braket{\beta_+}{\alpha_+}$. Donc
$P(a=b)=\frac12\cos^2\frac{\alpha-\beta}2+\frac12\cos^2\frac{\alpha-\beta}2$. \hfill$\square$

\medskip
\noindent Pour les quatre couples $(x,y)$ :
\begin{center}
\renewcommand{\arraystretch}{1.3}
\begin{tabular}{@{}c c c c c c@{}}
\toprule
$(x,y)$ & $\alpha$ & $\beta$ & $P(a=b)$ & condition de gain & $P(\text{gain})$\\
\midrule
$(0,0)$ & $0$ & $\frac\pi4$ & $\cos^2\frac\pi8$ & $a=b$ & $\cos^2\frac\pi8$\\
$(0,1)$ & $0$ & $-\frac\pi4$ & $\cos^2\frac\pi8$ & $a=b$ & $\cos^2\frac\pi8$\\
$(1,0)$ & $\frac\pi2$ & $\frac\pi4$ & $\cos^2\frac\pi8$ & $a=b$ & $\cos^2\frac\pi8$\\
$(1,1)$ & $\frac\pi2$ & $-\frac\pi4$ & $\cos^2\frac{3\pi}8=\sin^2\frac\pi8$ & $a\neq b$ &
  $1-\sin^2\frac\pi8=\cos^2\frac\pi8$\\
\bottomrule
\end{tabular}
\end{center}
Alice et Bob gagnent donc chaque partie avec la même probabilité :
\[
  P(\text{gain})=\cos^2\frac\pi8=\frac{2+\sqrt2}{4}\approx0{,}854>\frac34 .
\]

\begin{remarque}[Ce que montre le jeu CHSH]
\begin{itemize}[nosep]
  \item Aucune stratégie classique, même avec du hasard partagé, ne dépasse $\frac34$ : l'intrication
        produit des corrélations qu'aucun modèle «~local~» (résultats fixés à l'avance) ne reproduit.
        Les expériences sur des photons intriqués confirment la valeur quantique.
  \item $\cos^2\frac\pi8$ est le maximum quantique (\emph{borne de Tsirelson}, résultat admis,
        \cite{tsirelson1980}) : on ne peut pas gagner à tous les coups. Avec
        $E_{xy}=P(a=b)-P(a\neq b)=\cos(\alpha-\beta)$ et $S=E_{00}+E_{01}+E_{10}-E_{11}$, on a
        $P(\text{gain})=\frac12+\frac S8$ ; l'inégalité de Bell donne $|S|\leq2$ en classique, et la
        stratégie ci-dessus atteint $S=2\sqrt2$.
  \item Il n'y a pas de communication : la réponse de Bob est uniforme ($P(b=0)=\frac12$) quelle que
        soit la base choisie par Alice. L'intrication ne transmet pas d'information ; elle ne fait
        que corréler les résultats.
\end{itemize}
\end{remarque}

\clearpage
% ---------------------------------------------------------------------
\subsection{Récapitulatif du chapitre}\label{sec:recap5}
% ---------------------------------------------------------------------

\begin{center}
\adjustbox{max width=\linewidth}{%
\begin{tikzpicture}[>=Stealth, font=\footnotesize,
    boite/.style={draw=insarouge, fill=insarouge!6, rounded corners=2pt, thick,
                  minimum width=4.4cm, minimum height=1.45cm, align=center, inner sep=3pt},
    lien/.style={->, very thick, insagris},
    etiq/.style={font=\scriptsize\itshape, text=insagris, inner sep=2pt}]
  \node[boite] (mes) at (0,0)
    {\textbf{Mesure à deux qubits}\\ $P(ij)=|c_{ij}|^2$\\ {\scriptsize\S\,\ref{sec:mesure2}}};
  \node[boite] (sep) at (6.3,0)
    {\textbf{Séparable ou intriqué}\\ $D=c_{00}c_{11}-c_{01}c_{10}$\\ {\scriptsize\S\,\ref{sec:separable}}};
  \node[boite] (bel) at (12.6,0)
    {\textbf{États de Bell}\\ $\ket{\Phi^\pm}$, $\ket{\Psi^\pm}$ ; $\mathrm{CNOT}\,(I\otimes H)$\\
     {\scriptsize\S\,\ref{sec:bell}, \ref{sec:circuit-bell}}};
  \node[boite] (tra) at (12.6,-2.3)
    {\textbf{Trace partielle, entropies}\\ $\rho_A=\operatorname{Tr}_B\rho_{AB}$ ; $S(\rho)$, $I(A;B)$\\
     {\scriptsize\S\,\ref{sec:trace}, \ref{sec:info}}};
  \node[boite] (tel) at (6.3,-2.3)
    {\textbf{Téléportation}\\ e-bit + 2 bits classiques\\ {\scriptsize\S\,\ref{sec:teleportation}}};
  \node[boite] (chs) at (0,-2.3)
    {\textbf{Jeu CHSH}\\ $\frac34$ classique, $\cos^2\frac\pi8$ quantique\\ {\scriptsize\S\,\ref{sec:chsh}}};
  \draw[lien] (mes) -- (sep) node[etiq, midway, above=1pt] {corrélations};
  \draw[lien] (sep) -- (bel) node[etiq, midway, above=1pt] {cas extrême};
  \draw[lien] (bel) -- (tra) node[etiq, midway, right=1pt] {un qubit seul};
  \draw[lien] (tra) -- (tel) node[etiq, midway, above=1pt] {ressource};
  \draw[lien] (tel) -- (chs) node[etiq, midway, above=1pt] {non-localité};
\end{tikzpicture}}
\end{center}

\begin{aretenir}
\begin{itemize}
  \item Deux qubits : quatre amplitudes $c_{ij}$, $P(ij)=|c_{ij}|^2$. Mesurer $A$ supprime les termes
        incompatibles et renormalise : l'état de $B$ peut changer.
  \item Séparable $\Longleftrightarrow D=c_{00}c_{11}-c_{01}c_{10}=0$. Sinon l'état est intriqué et
        aucun qubit n'a d'état propre (pour $\ket{\Phi^+}$, $P(00)\neq P(A{=}0)P(B{=}0)$).
  \item États de Bell : base orthonormée d'états maximalement intriqués, états propres de
        $Z\otimes Z$ et $X\otimes X$ ; $\mathrm{CNOT}\,(I\otimes H)$ les produit à partir de la base de
        calcul. Des portes locales ne créent pas d'intrication.
  \item Qubit seul : $\rho_A=\operatorname{Tr}_B\rho_{AB}$, avec $\operatorname{Tr}\rho_A^2=1-2|D|^2$ pour un état pur.
        Pour $\ket{\Phi^+}$ : $\rho_A=\mathrm{Id}/2$ alors que le couple est pur ; $S(\rho_{AB})=0$,
        $S(\rho_A)=1$ bit, $I(A;B)=2$ bits, $H(A|B)=-1$ (impossible en classique).
  \item Téléportation : un e-bit et deux bits classiques ; $P(a,q)=\frac14$ ; Bob applique $X$ si $a=1$
        puis $Z$ si $q=1$. L'état est déplacé, pas copié ; sans les bits, $\rho_B=\mathrm{Id}/2$.
  \item Jeu CHSH : toute stratégie classique gagne avec une probabilité $\leq\frac34$ ; avec $\ket{\Phi^+}$
        et des bases tournées, $\cos^2\frac\pi8\approx0{,}854$ (borne de Tsirelson), sans aucune
        communication.
\end{itemize}
\end{aretenir}

\begin{center}
\renewcommand{\arraystretch}{1.2}
\begin{tabular}{@{}l l r@{}}
\toprule
\textbf{Notion} & \textbf{Formule} & \textbf{§} \\
\midrule
Mesure de deux qubits & $P(ij)=|c_{ij}|^2$, \quad $P(A{=}a)=P(a0)+P(a1)$ & \ref{sec:mesure2} \\
Séparabilité & $\ket{\Psi}=\ket{\Psi}_A\otimes\ket{\Psi}_B\ \Longleftrightarrow\ D=0$ & \ref{sec:separable} \\
États de Bell & $\ket{\Phi^\pm}=\dfrac{\ket{00}\pm\ket{11}}{\sqrt2}$, \quad
  $\ket{\Psi^\pm}=\dfrac{\ket{01}\pm\ket{10}}{\sqrt2}$ & \ref{sec:bell} \\
Trace partielle & $\operatorname{Tr}_B\big(\ket{a}\bra{a'}\otimes\ket{b}\bra{b'}\big)
  =\ket{a}\bra{a'}\,\braket{b'}{b}$ & \ref{sec:trace} \\
Qubit d'un état pur & $\rho_A=CC^\dagger$, \quad $\operatorname{Tr}\rho_A^2=1-2|D|^2$ & \ref{sec:trace} \\
Entropies & $S(\rho)=-\sum_i\lambda_i\log_2\lambda_i$, \quad $I(A;B)=H(A)+H(B)-H(A,B)$ & \ref{sec:info} \\
Téléportation & $P(a,q)=\frac14$, \quad correction $Z^{q}X^{a}$ (d'abord $X$) & \ref{sec:teleportation} \\
Singulet & $(U\otimes U)\ket{\Psi^-}=\det(U)\,\ket{\Psi^-}$ & \ref{sec:bell} \\
Jeu CHSH & $P_{\mathrm{cl}}\leq\frac34$, \quad $P_{\mathrm{quant}}=\cos^2\frac\pi8=\frac{2+\sqrt2}{4}$ & \ref{sec:chsh} \\
\bottomrule
\end{tabular}
\end{center}

\begin{savoirfaire}
\begin{itemize}
  \item Calculer $P(ij)$, les probabilités d'un seul qubit et l'état après la mesure de l'un d'eux ;
        tester la séparabilité avec $D$ et calculer la pureté de $\rho_A$ (exercice~\ref{ex:separable3}).
  \item Suivre un état dans le générateur de Bell, écrire $\rho_{AB}$ et reconnaître les corrélations
        d'un état de Bell dans deux bases (exercice~\ref{ex:bellX}).
  \item Calculer $H(A)$, $H(A,B)$, $I(A;B)$ et $H(A|B)$ pour une loi ou un état (exercice~\ref{ex:entropies}).
  \item Dérouler la téléportation pour un état donné et utiliser les états de Bell pour le codage
        superdense (exercices~\ref{ex:telenum} et~\ref{ex:superdense}).
  \item Prouver la borne $\frac34$ du jeu CHSH classique et calculer la probabilité de gain d'une
        stratégie quantique (section~\ref{sec:chsh}).
\end{itemize}
\end{savoirfaire}
